% !TEX root = print.tex

\part{Besov spaces}
\label{part:besov}

This part describes the Besov spaces $\Besov{s}{p}{q}$ in each setting by the characterisations
available there. Its first chapter, \Cref{chap:representations}, is the overview: it lists the
descriptions, the settings and parameters in which each is proved, and how they differ. The
chapters after it prove them: Fourier coefficients, moduli of smoothness and Littlewood--Paley
blocks on the torus (\Cref{ch:torus} and \Cref{ch:classical}); martingale differences of a
filtration on a finite measure space (\Cref{ch:filtrations}); the Haar system, piecewise
polynomials, splines and multiwavelets on intervals and on the cube (\Cref{ch:interval} and
\Cref{ch:cube}); band-limited approximation, Littlewood--Paley blocks, differences, local means
and atoms on $\R^d$, with the invariance under smooth changes of variables (\Cref{ch:rd}); the
coefficients of the Daubechies wavelets on $\R$ (\Cref{ch:daubechies}), with their tensor products
on $\R^d$ (\Cref{ch:daubechies-rd}); and, for dominating mixed smoothness on $\R^d$
(\Cref{ch:mixed-smoothness}), mixed differences, mixed local means and tensor Littlewood--Paley
blocks.

\chapter{The representations of Besov spaces}
\label{chap:representations}

The chapters of this part, the abstract chapters on approximation spaces
(\Cref{ch:approximation-spaces}, \Cref{ch:interpolation}, \Cref{ch:decompositions} and
\Cref{ch:spectral}), the chapter on Triebel--Lizorkin spaces (\Cref{ch:triebel-lizorkin}) and
the chapters on domains (\Cref{ch:lipschitz-domains}, \Cref{ch:traces-lipschitz},
\Cref{ch:laplacian-eigenbases}, \Cref{ch:manifolds} and \Cref{ch:manifold-eigenbases})
describe Besov spaces in many ways and prove that the descriptions agree. This chapter collects
them: what each description measures
(\Cref{sec:descriptions}), in which setting and for which parameters it is proved to give the
Besov space (\Cref{sec:where}), and how the descriptions differ (\Cref{sec:differences}).
Throughout $1 \leq q \leq \infty$, and $1 \leq p < \infty$ unless stated otherwise. Two
descriptions agree when the same functions have a finite quasi-norm; most theorems moreover bound
each quasi-norm by a constant times the other. The dimension of the underlying space is $n$.

\section{The descriptions}
\label{sec:descriptions}

\begin{enumerate}
  \item \emph{Approximation} (\Cref{def:approx-space} and \Cref{def:dyadic-norm}).
    $\|f\| + \|(2^{ks} E_{2^k}(f))_k\|_{\ell_q}$, where $E_m(f)$ is the error of best
    approximation of $f$ by the $m$-th set of a scheme: trigonometric polynomials, band-limited
    functions, piecewise polynomials, splines, step functions, the spans of the translates of a
    Daubechies scaling function, or combinations of $m$ elements of a dictionary.
  \item \emph{Differences} (\Cref{def:rd-modulus} and \Cref{def:rd-modulus-on};
    \Cref{def:modulus}, \Cref{def:interval-modulus}, \Cref{def:iso-modulus} and
    \Cref{def:cube-modulus}).
    $\|f\|_p + \|(2^{ks} \omega_r(f, 2^{-k})_p)_k\|_{\ell_q}$ with a modulus of smoothness of
    order $r$ on a group such as $\R^n$ or $\T^n$, or on a set $\Omega$, where only the
    differences with all their points in $\Omega$ count. The moduli on the torus and on an interval
    and the isotropic modulus on the cube are of this form; the moduli along the coordinates of the
    cube take differences in one coordinate at a time. For $r = 1$, $q = p$ and $0 < s < 1$ the
    Gagliardo form $\int\!\int |f(y) - f(x)|^p / |y - x|^{n + sp} \, dy \, dx$ of the
    Sobolev--Slobodeckij space $W^{s,p}(\R^n)$, added to $\|f\|_p^p$, is equivalent to the
    $p$-th power of this norm (\Cref{thm:sobolev-slobodeckij-besov}).
  \item \emph{Littlewood--Paley blocks} (\Cref{def:lp-blocks}).
    $\|(2^{js} \|\psi_j * f\|_p)_j\|_{\ell_q}$, where the Fourier transforms of the kernels
    $\psi_j$ form a smooth dyadic decomposition of unity. For $p = 2$ on the torus the sharp
    dyadic blocks of the Fourier coefficients serve as well:
    \[
      |\hat f(0)| + \Bigl\|\Bigl(2^{ks} \Bigl(\sum_{2^k \leq \deg m < 2^{k+1}} |\hat f(m)|^2
        \Bigr)^{1/2}\Bigr)_k\Bigr\|_{\ell_q},
      \qquad \deg m = \max_i |m_i| .
    \]
  \item \emph{Local means} (\Cref{def:local-means}).
    $\|(2^{js} \|k_j * f\|_p)_j\|_{\ell_q}$ with kernels $k_j$ supported in balls of radius
    $C 2^{-j}$, all built from one bump.
  \item \emph{Atoms} (\Cref{def:atomic-system} and \Cref{def:regular-atoms}).
    Decompositions into atoms $a_{j,m}$ of level $j$, supported near a cube of side $2^{-j}$,
    $f = \sum_j \sum_m \lambda_{j,m} a_{j,m}$, measured by
    \[
      \|\lambda\|_{b^s_{p,q}} = \bigl\|\bigl(2^{j(s - n/p)} \|\lambda_{j,\cdot}\|_{\ell_p}\bigr)_j
        \bigr\|_{\ell_q} .
    \]
  \item \emph{Coefficients} (\Cref{def:haar}, \Cref{def:multiwavelets}, \Cref{def:daub-wavelet} and
    \Cref{def:spectral}).
    Weighted $\ell_q(\ell_p)$ norms of the coefficients of $f$ in a basis, level by level: Haar
    functions, multiwavelets, Daubechies wavelets, the Fourier basis, a Hilbert basis of a
    spectral scheme, among them the Neumann eigenbasis of the Laplacian on a bounded Lipschitz
    domain and the eigenbasis of a chart form on a compact manifold.
  \item \emph{Blocks} (\Cref{def:block-decomposition} and \Cref{def:blocks}).
    Decompositions $f = \sum_k g_k$ with $g_k$ in the $2^k$-th set of the scheme, measured by
    $\|(2^{ks} \|g_k\|)_k\|_{\ell_q}$; for a scheme with uniformly bounded projections $P_m$, the
    dyadic blocks $P_{2^{k+1}} f - P_{2^k} f$, among them martingale differences.
  \item \emph{Interpolation} (\Cref{def:k-functional} and \Cref{def:interpolation}).
    The real interpolation space $(X, Y)_{s/r, q}$ between $X$ and a space $Y$ of smoothness $r$,
    defined through the K-functional on the grid $t = 2^{-kr}$ (the discrete $K$-method), or the
    real interpolation space $(N_0, N_1)_{\eta, q}$ of the continuous $K$-method between two
    approximation, Besov or classical spaces (\Cref{def:approximation-space-couples},
    \Cref{def:besov-couples} and \Cref{def:modulus-couples}).
\end{enumerate}

On a domain $\Omega$ two further quasi-norms enter: the \emph{restriction} quasi-norm, the
infimum of a quasi-norm on $\R^n$ over all extensions of $f$, and the \emph{intrinsic local
means} $\|(2^{js} \|k_j * f\|_{L_p(\Omega)})_j\|_{\ell_q}$ (\Cref{def:intrinsic-local-means}).
Spaces of \emph{dominating mixed smoothness} replace the level $k \in \N$ by $k \in \N^n$, the
weight $2^{ks}$ by $2^{s|k|_1}$, and every kernel by a tensor product
(\Cref{def:mixed-moduli}, \Cref{def:mixed-local-means} and \Cref{def:mixed-lp}).

\section{Where each description is proved}
\label{sec:where}

Each table lists the descriptions that are shown to agree in one setting, the smoothness range of
the statement, and the result. A range $0 \leq s < \theta$ or $0 < s < \theta$ is the range of
the cited statement; outside it the descriptions may differ (\Cref{sec:differences}).

\paragraph{Abstract schemes.} Here $A$ is an approximation scheme in a normed space $X$.

\begin{center}
\small
\begin{tabular}{p{0.4\textwidth}p{0.25\textwidth}p{0.22\textwidth}}
\hline
Descriptions & Range & Result \\
\hline
\raggedright Pietsch's quasi-norm and the dyadic quasi-norm & $s \geq 0$ & \Cref{thm:dyadic-norm} \\
\raggedright Approximation and interpolation, under Jackson's and Bernstein's inequalities of order $r$
  & $0 < s < r$ & \Cref{thm:representation}\newline \Cref{cor:equal-spaces} \\
\raggedright Approximation and moduli, under Jackson's and the inverse inequality of order $\theta$
  & $0 \leq s < \theta$ & \Cref{thm:modulus-coincidence} \\
\raggedright Approximation and real interpolation of two approximation spaces, also of two
  Besov spaces of a scheme of functions
  & $0 < s_0 < s_1$ & \Cref{thm:approximation-spaces-reiteration}\newline
  \Cref{thm:besov-interpolate} \\
\raggedright Moduli and real interpolation of two classical spaces, under Jackson's and the
  inverse inequality of order $\theta$
  & $0 < s_0 < s_1 < \theta$ & \Cref{cor:modulus-spaces-interpolate} \\
\raggedright Approximation and blocks & $s > 0$ & \Cref{thm:block-decomposition} \\
\raggedright Approximation and dyadic blocks of projections & $s > 0$ & \Cref{thm:blocks} \\
\raggedright Approximation and martingale differences in $L_p$ & $s > 0$ & \Cref{thm:martingale} \\
\raggedright Approximation and atoms of decay $\sigma$ & $0 < s < \sigma$
  & \Cref{thm:atomic-synthesis}\newline \Cref{thm:atomic-analysis} \\
\raggedright Approximation and coefficients in a Hilbert space & $s > 0$
  & \Cref{thm:spectral-coefficients} \\
\raggedright $n$-term approximation from a Hilbert basis and coefficients in $\ell_\tau$,
  $q = \tau$
  & $s = 1/\tau - 1/2$,\newline $0 < \tau < 2$ & \Cref{thm:stechkin} \\
\raggedright Intrinsic approximation and restrictions, given an extension & $s \geq 0$
  & \Cref{thm:retract} \\
\hline
\end{tabular}
\end{center}

\paragraph{The torus $\T^n$.} The scheme consists of the trigonometric polynomials.
\begin{center}
\small
\begin{tabular}{p{0.4\textwidth}p{0.25\textwidth}p{0.22\textwidth}}
\hline
Descriptions & Range & Result \\
\hline
\raggedright Trigonometric approximation and Fourier coefficients, $p = 2$; $\Besov{s}{2}{2} = H^s$
  & $s > 0$ & \Cref{thm:torus-coefficients} \\
\raggedright Trigonometric approximation and differences of order $r$, $1 \leq p < \infty$
  & $0 < s < r$ & \Cref{thm:torus-lp-moduli} \\
\raggedright Trigonometric approximation and smooth Littlewood--Paley blocks, $1 \leq p < \infty$
  & $s > 0$ & \Cref{thm:torus-lp-besov} \\
\raggedright Differences of order $r$ and real interpolation of two such spaces,
  $1 \leq p < \infty$
  & $0 < s_0 < s_1 < r$ & \Cref{thm:torus-besov-interpolate} \\
\raggedright Weak derivatives, Fourier coefficients and approximation: $\Sob{m}{2} = H^m = \Besov{m}{2}{2}$
  & $m \geq 1$ & \Cref{thm:weak-sobolev} \\
\raggedright $n$-term trigonometric approximation in $L_2$ and Fourier coefficients in
  $\ell_\tau$, $q = \tau$
  & $s = 1/\tau - 1/2$,\newline $0 < \tau < 2$ & \Cref{cor:torus-nterm} \\
\hline
\end{tabular}
\end{center}

\paragraph{Spaces with a continuous spectral basis.} The scheme consists of the spans of the
continuous basis functions (\Cref{def:continuous-spectral-basis}); the basis is a hypothesis, and
the torus satisfies it (\Cref{prop:fourier-eigenbasis}).
\begin{center}
\small
\begin{tabular}{p{0.4\textwidth}p{0.25\textwidth}p{0.22\textwidth}}
\hline
Descriptions & Range & Result \\
\hline
\raggedright Approximation and coefficients, $p = 2$; $\Besov{s}{2}{2} = H^s$ & $s > 0$
  & \Cref{thm:continuous-spectral-basis-besov} \\
\raggedright $H^s$ and continuous functions, under the local Weyl bound of dimension $D$
  & $s > D/2$ & \Cref{thm:local-weyl-continuity} \\
\hline
\end{tabular}
\end{center}

\paragraph{The interval $[0, 1]$ and the circle.} The schemes are the dyadic piecewise polynomials of degree $< r$ and the dyadic splines of order
$r$; the differences are of order $r$.

\begin{center}
\small
\begin{tabular}{p{0.4\textwidth}p{0.25\textwidth}p{0.22\textwidth}}
\hline
Descriptions & Range & Result \\
\hline
\raggedright Piecewise polynomials and differences on $[0, 1]$ & $0 \leq s < 1/p$ & \Cref{thm:pp-besov} \\
\raggedright Piecewise polynomials and periodic differences; $\Besov{s}{p}{q}[0, 1] = \Besov{s}{p}{q}(\T)$
  & $0 \leq s < 1/p$ & \Cref{thm:pp-besov-per}\newline \Cref{cor:interval-circle} \\
\raggedright Splines and differences on $[0, 1]$ & $0 \leq s < r - 1 + 1/p$ & \Cref{thm:spline-besov} \\
\raggedright Periodic splines and periodic differences & $0 \leq s < r - 1 + 1/p$
  & \Cref{thm:per-spline-besov} \\
\raggedright Piecewise polynomials and multiscale details & $s > 0$ & \Cref{thm:pp-multiscale} \\
\raggedright Piecewise polynomials and multiwavelet coefficients & $s > 0$
  & \Cref{thm:multiwavelet-besov} \\
\raggedright Step functions and Haar coefficients & $s > 0$ & \Cref{thm:haar} \\
\raggedright Haar coefficients and differences of order $1$ & $0 < s < 1/p$ & \Cref{cor:haar-besov} \\
\raggedright Step functions in $L_p$, $0 < p \leq 1$, Haar coefficients and differences of order $1$
  & $s > 1/p - 1$\newline (differences: $s < 1/p$) & \Cref{thm:haar-below-one}\newline
  \Cref{cor:haar-besov-below-one} \\
\raggedright $n$-term Haar approximation in $L_2$ and differences of order $1$ in $L_\tau$,
  $p = q = \tau$
  & $s = 1/\tau - 1/2$,\newline $0 < \tau < 2$ & \Cref{cor:haar-nterm-besov}\newline
  \Cref{thm:haar-nterm-besov-all} \\
\hline
\end{tabular}
\end{center}

\paragraph{The cube $[0, 1]^n$.} The schemes are the piecewise polynomials of degree $< r$ in
each variable and the tensor splines of order $r$; the differences are of order $r$.
\begin{center}
\small
\begin{tabular}{p{0.4\textwidth}p{0.25\textwidth}p{0.22\textwidth}}
\hline
Descriptions & Range & Result \\
\hline
\raggedright Piecewise polynomials and coordinate differences & $0 \leq s < 1/p$ & \Cref{thm:cube-besov} \\
\raggedright Tensor splines and coordinate differences & $0 \leq s < r - 1 + 1/p$
  & \Cref{thm:cube-spline-besov} \\
\raggedright Piecewise polynomials and isotropic differences & $0 \leq s < 1/p$ & \Cref{thm:iso-besov} \\
\raggedright Tensor splines and isotropic differences; coordinate and isotropic differences
  & $0 \leq s < r - 1 + 1/p$ & \Cref{thm:iso-besov} \\
\raggedright Tensor splines and differences on $[0, 1]^n \subseteq \R^n$ for the Euclidean norm
  & $0 \leq s < r - 1 + 1/p$ & \Cref{thm:cube-splines-euclidean} \\
\hline
\end{tabular}
\end{center}

\paragraph{The space $\R^n$.} The scheme consists of the band-limited functions.
\begin{center}
\small
\begin{tabular}{p{0.4\textwidth}p{0.25\textwidth}p{0.22\textwidth}}
\hline
Descriptions & Range & Result \\
\hline
\raggedright Band-limited approximation and Littlewood--Paley blocks & $s > 0$ & \Cref{thm:lp-band} \\
\raggedright Band-limited approximation, Littlewood--Paley blocks and differences of order $r$
  & $0 < s < r$ & \Cref{thm:rd-besov} \\
\raggedright Local means and differences of order $L$ & $0 < s < L$ & \Cref{thm:local-means} \\
\raggedright Atoms of order $L$ and differences of order $L$ & $0 < s < L$
  & \Cref{thm:atomic-characterisation} \\
\raggedright Differences of order $r$ and real interpolation of two such spaces
  & $0 < s_0 < s_1 < r$ & \Cref{thm:rd-besov-interpolate} \\
\raggedright Band-limited approximation and the Gagliardo form, $q = p$ & $0 < s < 1$
  & \Cref{thm:sobolev-slobodeckij-besov} \\
\hline
\end{tabular}
\end{center}

\paragraph{The line $\R$ with Daubechies wavelets.} The scheme consists of the spans $V_j$ of the
functions $\varphi(2^j \cdot - k)$, $k \in \Z$, for the scaling function $\varphi$ of a Daubechies
mask of order $N \geq 2$; the differences are of order $N$. The order $\theta$ is that of a bound
$\omega_N(\varphi, t)_p \leq K t^\theta$ (\Cref{def:daub-bernstein-bound}), and $\lambda$ is the
rate of a growth bound $T^n 1 \leq C \lambda^n$ for the transfer operator.
\begin{center}
\small
\begin{tabular}{p{0.4\textwidth}p{0.25\textwidth}p{0.22\textwidth}}
\hline
Descriptions & Range & Result \\
\hline
\raggedright Daubechies approximation and differences of order $N$ & $0 \leq s < \theta$
  & \Cref{thm:daub-approx-modulus} \\
\raggedright Daubechies approximation and wavelet coefficients & $s > 0$
  & \Cref{thm:daub-approx-coeff} \\
\raggedright Wavelet coefficients and differences of order $N$; Littlewood--Paley blocks if
  $s < N$ & $0 < s < \theta$ & \Cref{thm:daub-besov} \\
\raggedright Wavelet coefficients and differences, $1 \leq p \leq 2$
  & $0 < s < N - \log_4 \lambda$ & \Cref{cor:daub-besov-le-two} \\
\raggedright Wavelet coefficients and differences, $2 \leq p < \infty$, with
  $s' < N - \log_4 \lambda$ and a H\"older exponent $a$ of $\varphi$
  & $0 < s < \frac{2}{p} s' + (1 - \frac{2}{p}) a$ & \Cref{cor:daub-besov-two-le} \\
\raggedright D4 coefficients and differences of order $2$, $1 \leq p \leq 2$; $2 \leq p < \infty$;
  every $p$, with $\alpha = 2 - \log_2(1 + \sqrt{3})$
  & $0 < s < 1$; $0 < s < \frac{2}{p} + (1 - \frac{2}{p}) \alpha$; $0 < s < \alpha$
  & \Cref{thm:daub-besov-d4} \\
\raggedright Coefficients of order three and differences of order $3$, $1 \leq p \leq 2$;
  $2 \leq p < \infty$; every $p$
  & $0 < s < 3 - \log_2 3$; $0 < s < \frac{5}{2} - \log_2 3 + \frac{1}{p}$;
  $0 < s < \frac{5}{2} - \log_2 3$ & \Cref{thm:daub-besov-order-three} \\
\hline
\end{tabular}
\end{center}

\paragraph{The space $\R^n$ with tensor Daubechies wavelets.} The scheme consists of the spans
$V_j$ of the functions $\Phi(2^j \cdot - k)$, $k \in \Z^n$, for the tensor product $\Phi$ of $n$
copies of the scaling function $\varphi$ of a Daubechies mask of order $N \geq 2$; the
coefficients are those of $\Phi$ and of the $2^n - 1$ tensor wavelets. Here $\R^n$ carries the
maximum norm, the differences are of order $N$, and $\theta$ is the order of a bound
$\omega_N(\varphi, t)_p \leq K t^\theta$.
\begin{center}
\small
\begin{tabular}{p{0.4\textwidth}p{0.25\textwidth}p{0.22\textwidth}}
\hline
Descriptions & Range & Result \\
\hline
\raggedright Tensor Daubechies approximation and differences of order $N$
  & $0 \leq s < \min(\theta, N)$ & \Cref{thm:daub-rd-approx-modulus} \\
\raggedright Tensor Daubechies approximation and wavelet coefficients & $s > 0$
  & \Cref{thm:daub-rd-approx-coeff} \\
\raggedright Wavelet coefficients and differences of order $N$ & $0 < s < \min(\theta, N)$
  & \Cref{thm:daub-rd-besov} \\
\raggedright Wavelet coefficients in the coordinates of the dyadic grid and differences, on
  every inner product space & $0 < s < \min(\theta, N)$ & \Cref{thm:daub-rd-besov-grid} \\
\raggedright Wavelet coefficients and differences, masks of order three; D4
  & $0 < s < \frac{5}{2} - \log_2 3$; $0 < s < 2 - \log_2(1 + \sqrt{3})$
  & \Cref{cor:daub-rd-besov-orders} \\
\raggedright $H^s$ and wavelet coefficients, $p = q = 2$; masks of order three
  & $0 < s < \min(\theta, N)$; $0 < s < 3 - \log_2 3$ & \Cref{cor:daub-rd-sobolev} \\
\hline
\end{tabular}
\end{center}

\paragraph{Domains in $\R^n$.} Restrictions are taken from the functions on $\R^n$ with a
finite quasi-norm of the same kind.
\begin{center}
\small
\begin{tabular}{p{0.4\textwidth}p{0.25\textwidth}p{0.22\textwidth}}
\hline
Descriptions & Range & Result \\
\hline
\raggedright Lipschitz graph domains: intrinsic local means and restrictions & $0 < s < L$
  & \Cref{thm:restriction-local-means} \\
\raggedright Lipschitz graph domains: intrinsic differences, restrictions of differences and of
  band-limited approximation & $0 < s < r$ & \Cref{thm:graph-besov} \\
\raggedright Bounded Lipschitz domains: the same three & $0 < s < r$
  & \Cref{thm:bounded-lipschitz-besov} \\
\raggedright The cube $[0, 1]^n$: tensor splines and restrictions of the functions on $\R^n$
  & $0 < s < r - 1 + 1/p$ & \Cref{thm:cube-splines-restriction} \\
\raggedright Bounded Lipschitz domains in a cube: intrinsic differences, restrictions of tensor
  splines of the cube and of spline blocks & $0 < s < r - 1 + 1/p$
  & \Cref{thm:domain-splines}\newline \Cref{cor:domain-spline-blocks} \\
\raggedright Smooth bounded domains: charts of the closure, intrinsic differences and
  restrictions & $0 < s < r$ & \Cref{thm:smooth-domain-besov} \\
\raggedright Bounded Lipschitz domains, $p = 2$: Neumann eigenfunction expansions and
  intrinsic differences & $0 < s < 1$ & \Cref{thm:spectral-besov-neumann}\newline
  \Cref{cor:spectral-besov-neumann-blocks} \\
\hline
\end{tabular}
\end{center}

\paragraph{Compact manifolds.} The pieces of a function in the charts of a Besov atlas are
measured by the intrinsic differences on the model space.
\begin{center}
\small
\begin{tabular}{p{0.4\textwidth}p{0.25\textwidth}p{0.22\textwidth}}
\hline
Descriptions & Range & Result \\
\hline
\raggedright Charts of any Besov atlas, differences of any order $r > s$ & $s > 0$
  & \Cref{def:manifold-besov} \\
\raggedright The boundary: charts of the boundary hyperplane & $s > 0$
  & \Cref{thm:boundary-atlas-independence} \\
\raggedright Charts and eigenfunction expansions of a chart form, $p = 2$ & $0 < s < 1$
  & \Cref{thm:spectral-besov-manifold}\newline \Cref{cor:spectral-besov-manifold-blocks} \\
\hline
\end{tabular}
\end{center}

\paragraph{Traces.} The trace maps $\Besov{s+1/p}{p}{q}$ onto $\Besov{s}{p}{q}$ of the boundary; for abstract schemes
the loss $1/p$ is an exponent $\beta \geq 0$ of the trace bound.

\begin{center}
\small
\begin{tabular}{p{0.4\textwidth}p{0.25\textwidth}p{0.22\textwidth}}
\hline
Setting & Range & Result \\
\hline
\raggedright Abstract schemes with a trace bound & $s > 0$
  & \Cref{thm:trace-transfer}\newline \Cref{thm:trace-surjective} \\
\raggedright From $\T^{n+1}$ to a face $\T^n$ & $s > 0$ & \Cref{thm:trace-torus} \\
\raggedright From a Lipschitz graph domain to its boundary graph & $0 < s < 1$
  & \Cref{thm:graph-trace-space} \\
\raggedright From a bounded Lipschitz domain to its boundary & $0 < s < 1$
  & \Cref{thm:boundary-trace-space} \\
\raggedright From the half-space to its boundary hyperplane & $s > 0$
  & \Cref{thm:half-space-trace-space} \\
\raggedright From a compact manifold with boundary to its boundary & $s > 0$
  & \Cref{thm:manifold-trace} \\
\hline
\end{tabular}
\end{center}

\paragraph{Dominating mixed smoothness on $\R^n$.} The kernels are tensor products of
one-dimensional kernels.
\begin{center}
\small
\begin{tabular}{p{0.4\textwidth}p{0.25\textwidth}p{0.22\textwidth}}
\hline
Descriptions & Range & Result \\
\hline
\raggedright Mixed local means and mixed differences of order $L$ & $0 < s < L$
  & \Cref{thm:mixed-local-means} \\
\raggedright Tensor Littlewood--Paley blocks and mixed differences & $0 < s < L$ & \Cref{thm:mixed-lp} \\
\raggedright Approximation by Littlewood--Paley sums over hyperbolic crosses & $s > 0$
  & \Cref{thm:hyperbolic-cross} \\
\hline
\end{tabular}
\end{center}

\paragraph{Triebel--Lizorkin spaces.} The space $\TL{s}{p}{q}$ takes the $\ell_q$ norm of the
Littlewood--Paley blocks at every point before the $L_p$ norm
(\Cref{def:triebel-lizorkin} and \Cref{def:torus-triebel-lizorkin}); each row compares a description
of $\TL{s}{p}{q}$ with this definition, on $\R^n$ unless $\T^n$ is named. The descriptions also take
the $\ell_q$ norm at every point first: the local means of
\Cref{def:tl-local-means} and \Cref{def:torus-tl-local-means}, the ball means of differences of
\Cref{def:ball-means-differences} and the coefficient space $f^s_{p,q}$ of
\Cref{def:tl-sequence-space}. The first two rows hold for $1 \leq p < \infty$, the others for
$1 < p < \infty$ and $1 < q \leq \infty$.
\begin{center}
\small
\begin{tabular}{p{0.4\textwidth}p{0.25\textwidth}p{0.22\textwidth}}
\hline
Descriptions & Range & Result \\
\hline
\raggedright Littlewood--Paley blocks: $\TL{s}{p}{p} = \Besov{s}{p}{p}$ and
  $\Besov{s}{p}{\min(p,q)} \subseteq \TL{s}{p}{q} \subseteq \Besov{s}{p}{\max(p,q)}$
  & $s > 0$ & \Cref{thm:tl-diagonal}\newline \Cref{thm:tl-besov-sandwich} \\
\raggedright The same on $\T^n$ & $s > 0$ & \Cref{thm:torus-tl-besov-sandwich} \\
\raggedright Local means of order $L$ and Littlewood--Paley blocks & $0 < s < L$
  & \Cref{thm:tl-local-means}\newline \Cref{cor:tl-local-means-independence} \\
\raggedright The same on $\T^n$ & $0 < s < L$ & \Cref{thm:torus-tl-local-means} \\
\raggedright Ball means of differences of order $L$ & $0 < s < L$ & \Cref{cor:tl-ball-means} \\
\raggedright Atoms of order $L$ normalised in $L_\infty$, coefficients in $f^s_{p,q}$
  & $0 < s < L$ & \Cref{cor:tl-atomic} \\
\hline
\end{tabular}
\end{center}

\section{How the descriptions differ}
\label{sec:differences}

\paragraph{Linear and nonlinear.} The approximation quasi-norm measures best approximation, and a
best approximation need not depend linearly on $f$, although $A^s_q$ is a linear space and $A_m$
may be a nonlinear set, as for combinations of $m$ dictionary elements. Littlewood--Paley blocks,
local means, projections, martingale differences and coefficients in a basis are linear in $f$,
and the coefficients in a basis are unique. Decompositions into blocks or atoms are not unique,
and their quasi-norm is an infimum over all decompositions. The decompositions constructed here
come from near-best approximations in the abstract setting (\Cref{thm:atomic-analysis}) and from
the local means on $\R^n$ (\Cref{thm:atomic-decomposition}), where every term
$\lambda_{j,m} a_{j,m} = \eta_j * (1_{P_{j,m}} (k_j * f))$ is linear in $f$.

\paragraph{The range of smoothness.} Every description detects smoothness only up to a bound.
Differences of order $r$ detect $s < r$, and local means and atoms of order $L$ detect $s < L$.
Littlewood--Paley blocks and Fourier coefficients have no such bound. The approximation space of a scheme is defined for every $s$, and it is the Besov
space below the inverse order $\theta$ of the scheme: $\theta = 1/p$ for piecewise polynomials and
step functions, $\theta = r - 1 + 1/p$ for splines of order $r$, and every $\theta$ for
band-limited functions (\Cref{thm:rd-besov}). For the spaces of a Daubechies scaling function
$\varphi$ it is the order of a bound $\omega_N(\varphi, t)_p \leq K t^\theta$, so the range of
the wavelet coefficients is set by the regularity of $\varphi$ (\Cref{thm:daub-besov}). A scheme
of piecewise polynomials contains
functions with jumps, and these lie in its approximation spaces for every $s$ but in no Besov
space of smoothness $s > 1/p$.

\paragraph{The order.} For $0 < s < r$ the differences of order $r$ give the band-limited
approximation space on $\R^n$ (\Cref{thm:rd-besov}), so the spaces defined by differences of any
order $r > s$ coincide. The same holds for local means of any order $L > s$
(\Cref{thm:local-means}) and their mixed analogues (\Cref{thm:mixed-lp}).

\paragraph{Fourier analysis and $p$.} Fourier coefficients describe the spaces only for $p = 2$,
through Parseval's identity (\Cref{thm:torus-coefficients} and \Cref{thm:spectral-coefficients}).
For other $p$ the sharp dyadic blocks are replaced by the smooth kernels $\psi_j$. The proofs here
need only integrable kernels, Young's inequality and the kernel identities
(\Cref{thm:young} and \Cref{thm:kernel-identities}). On the torus the kernels of $\R^n$ act by
convolution with their periodisations, and their Fourier transforms at the lattice points are
Fourier multipliers; this transfers Jackson's and Bernstein's inequalities to every $L_p(\T^n)$
(\Cref{thm:torus-econv-fourier} and \Cref{thm:torus-lp-moduli}), and the smooth blocks $\psi_j * f$
describe $\Besov{s}{p}{q}(\T^n)$ for every $1 \leq p < \infty$ (\Cref{thm:torus-lp-besov}).

\paragraph{Besov and Triebel--Lizorkin spaces.} A Besov quasi-norm takes the $L_p$ norm of each
level first, so a bound of one kind of block by convolutions of the other kind is integrated
level by level with Young's inequality. A Triebel--Lizorkin quasi-norm takes the $\ell_q$ norm
at every point first, so the bounds must stay pointwise: a convolution $K * h$ with
$|K(y)| \leq C 2^{\nu n} (1 + 2^\nu \|y\|)^{-N}$, $N > n$, is at most $C_{n,N} C \, M|h|$
(\Cref{thm:decay-kernels-maximal}), and the Fefferman--Stein inequality
(\Cref{thm:fefferman-stein}) integrates the maximal functions in $L_p(\ell_q)$. It holds for
$1 < p < \infty$ and $1 < q \leq \infty$, which is the range of the descriptions of
$\TL{s}{p}{q}$ beyond the Littlewood--Paley blocks.

\paragraph{Coordinates, directions and mixed smoothness.} On the cube the coordinate moduli, with
differences along one coordinate at a time, and the isotropic modulus, with differences in every
direction, give the same spaces (\Cref{thm:iso-besov}). Dominating mixed smoothness measures
differences in several coordinates at once, $\Delta^{L,e}_h$, with the weight $2^{s|k|_1}$ for
$k \in \N^n$; its Littlewood--Paley blocks are tensor products, and its natural approximation is by
sums over the hyperbolic crosses $\{|k|_1 \leq N\}$, with the rate
$2^{-sN} (N + 1)^{(n-1)(1-1/q)}$ (\Cref{thm:hyperbolic-cross}).

\paragraph{Domains.} On a domain $\Omega$ the intrinsic differences use only the values of $f$ in
$\Omega$, the restriction quasi-norm uses extensions to $\R^n$, and the intrinsic local means use
the values in $\Omega$ through kernels supported in a cone. On Lipschitz domains they agree, by
Rychkov's extension operator (\Cref{thm:graph-extension}), which is linear and is one operator for
all $0 < s < L$, $p$ and $q$; on bounded Lipschitz domains the graph domains are glued by a
partition of unity and pointwise multipliers (\Cref{thm:multipliers} and
\Cref{thm:local-extension}).

\paragraph{Traces.} Restriction to a hypersurface costs the smoothness $1/p$. On the torus every
$s > 0$ is allowed (\Cref{thm:trace-torus}). On the boundary of a Lipschitz domain the Besov space
is defined in charts by differences of order $1$, the order that a Lipschitz change of variables
preserves, and the trace theorem holds for $0 < s < 1$
(\Cref{thm:graph-trace-space} and \Cref{thm:boundary-trace-space}). Smooth changes of charts keep
differences of every order, so on the half-space and on compact manifolds with boundary the trace
theorem holds for every $s > 0$ (\Cref{thm:half-space-trace-space} and \Cref{thm:manifold-trace}).

\chapter{The torus}
\label{ch:torus}

Let $\T = \R / \Z$ with its normalised Haar measure, let $d$ be a finite index set, and let
$\T^d$ carry the product measure, a probability measure charging every nonempty open set.
For $k \in \Z^d$ the trigonometric monomial is $e_k(x) = \exp(2\pi i \sum_j k_j x_j)$, and
$\hat f(k) = \int_{\T^d} f \, \overline{e_k}$ is the Fourier coefficient of $f \in
L_1(\T^d)$. The norm of $h \in \T^d$ is $\|h\| = \max_j \dist(h_j, \Z)$.

In Mathlib the Haar measure on $\T$, over which the Fourier basis of $L_2(\T^d)$ is stated, is
a local instance of one file; the project re-declares it as a scoped instance.

\begin{definition}[Trigonometric polynomials]
\label{def:trig-poly}
\lean{Torus.trigDeg, Torus.boxLt, Torus.card_boxLt_le, Torus.trigMonomial, Torus.trigPoly,
  Torus.trigScheme, TorusHaar.volume_eq_haarAddCircle, TorusHaar.isAddHaarMeasure,
  TorusHaar.isProbabilityMeasure, Torus.natAbs_le_trigDeg, Torus.trigDeg_le_iff,
  Torus.trigDeg_add_le, Torus.trigDeg_zero, Torus.mem_boxLt, Torus.boxLt_zero, Torus.boxLt_mono,
  Torus.trigMonomial_apply, Torus.trigPoly_zero, Torus.trigPoly_mono,
  Torus.trigMonomial_mem_trigPoly, Torus.mem_trigScheme, Torus.map_mem_trigPoly}
\leanok
\uses{def:function-scheme}
The degree of $k \in \Z^d$ is $\deg k = \max_j |k_j|$, with $\deg k = 0$ for $d = \emptyset$.
The space $\mathcal{T}_n$ of trigonometric polynomials of degree $< n$ is the span of
$\{e_k : \deg k < n\}$ in $C(\T^d)$, a space of dimension at most $(2n)^d$. The spaces
$\mathcal{T}_n$ form a scheme of functions on $\T^d$, the \emph{trigonometric scheme}. For all
finite index sets $\iota, \kappa$, every $m \in \N$ and every linear map $L$ from the bounded
continuous functions on $\T^\iota$ to those on $\T^\kappa$ with $L e_k \in \mathcal{T}_m$ for all
$k$ with $\deg k < m$, $L$ maps $\mathcal{T}_m$ on $\T^\iota$ into $\mathcal{T}_m$ on $\T^\kappa$,
since $\mathcal{T}_m$ is the span of these monomials.
\end{definition}

\begin{lemma}[Products of trigonometric polynomials]
\label{lem:trig-products}
\lean{Torus.trigPoly_mul_le, Torus.mul_mem_trigPoly, Torus.pow_mem_trigPoly,
  Torus.trigMonomial_add, Torus.trigMonomial_zero, Torus.norm_trigMonomial_apply}
\leanok
\uses{def:trig-poly}
For all $m, n \in \N$, $S \in \mathcal{T}_m$ and $T \in \mathcal{T}_n$: $S T \in
\mathcal{T}_{m+n}$ and $T^{m+1} \in \mathcal{T}_{(m+1)n}$.
\end{lemma}

\begin{proof}
\leanok
$e_k e_l = e_{k+l}$ and $\deg(k + l) \leq \deg k + \deg l$.
\end{proof}

\begin{lemma}[Supremum norm and $L_2$ norm]
\label{lem:trig-sup-l2}
\lean{Torus.norm_sq_le_card_mul_norm_toLp_sq, Torus.exists_sum_of_mem_trigPoly,
  Torus.periodic_norm_toLp_sum_sq, Torus.toLp_trigMonomial}
\leanok
\uses{def:trig-poly}
For every $n \in \N$ and $T \in \mathcal{T}_n$, $\sup |T|^2 \leq (2n)^d \, \|T\|_{L_2}^2$.
\end{lemma}

\begin{proof}
\leanok
Write $T = \sum_{\deg k < n} c_k e_k$. Then $|T(x)| \leq \sum |c_k|$, the Cauchy--Schwarz
inequality bounds this by $(\#\{k : \deg k < n\})^{1/2} (\sum |c_k|^2)^{1/2}$, and
$\sum |c_k|^2 = \|T\|_2^2$ by orthonormality.
\end{proof}

The Nikolskii inequality for trigonometric polynomials~\cite{Nikolskii, Temlyakov} is
usually proved with the Dirichlet or de la Vall\'ee Poussin kernel. The proof below avoids
kernels: it passes to a power of the polynomial and uses \Cref{lem:trig-sup-l2} for it.

\begin{theorem}[Nikolskii inequality on the torus]
\label{thm:torus-nikolskii}
\lean{Torus.norm_le_rpow_mul_eLpNorm, Torus.norm_rpow_le_card_mul_eLpNorm_rpow,
  Torus.norm_pow_le_norm_pow, Torus.norm_toLp_two_pow, Torus.trigScheme_hasNikolskii}
\leanok
\uses{lem:trig-products, lem:trig-sup-l2, def:nikolskii-property, lem:bc-linfty}
For all $n \geq 1$, $T \in \mathcal{T}_n$ and $1 \leq p < \infty$,
\[
  \sup |T| \leq \bigl(2 \lceil p/2 \rceil\bigr)^{d/p} \, n^{d/p} \, \|T\|_{L_p} .
\]
The trigonometric scheme has the Nikolskii property of dimension $d$.
\end{theorem}

\begin{proof}
\leanok
Let $m = \lceil p/2 \rceil$, so $2m \geq p$, and $T \neq 0$. By \Cref{lem:trig-products},
$T^m \in \mathcal{T}_{mn}$, and \Cref{lem:trig-sup-l2} gives
\[
  (\sup|T|)^{2m} \leq (\sup|T^m|)^2 \leq (2mn)^d \int |T|^{2m} ,
\]
and $\int |T|^{2m} \leq (\sup|T|)^{2m-p} \int |T|^p$.
Divide by $(\sup|T|)^{2m-p}$ and take $p$-th roots. For $p = \infty$ the constant is $1$ by
\Cref{lem:bc-linfty}.
\end{proof}

\begin{definition}[Sobolev space $H^s(\T^d)$]
\label{def:torus-sobolev}
\lean{Torus.fourierScheme, Torus.toLp_trigScheme_eq, Torus.periodicSobolevSpace,
  Torus.periodic_mem_range_sobolev_iff, Torus.periodicSobolevSpace_order,
  Torus.periodicSobolevSpace_incl}
\leanok
\uses{def:spectral, def:sobolev-embed, prop:spectral-jackson-bernstein, def:sobolev-scheme}
The Fourier basis $(e_k)_{k \in \Z^d}$ of $L_2(\T^d)$ with the level $\deg$ is a spectral
scheme whose approximants are the classes of $\mathcal{T}_n$. For $s > 0$, the Sobolev space
$H^s(\T^d)$ consists of the $f \in L_2(\T^d)$ with
\[
  \sum_{k \in \Z^d} (1 + \deg k)^{2s} |\hat f(k)|^2 < \infty ,
\]
the range of $J_s$. With $J_s$ it is a Sobolev space of order $s$ for the trigonometric
scheme in $L_2$, with constants $1$.
\end{definition}

The weight $(1 + \max_j |k_j|)^{2s}$ is equivalent to the more common $(1 + |k|^2)^s$, up to
constants depending on $d$ and $s$; the degree is used because it indexes the trigonometric
polynomials.

The monomials $e_k$ are the eigenfunctions of the Laplacian on $\T^d$,
$\Delta e_k = -4\pi^2 |k|^2 e_k$, and the degree groups them into the levels of a continuous
spectral basis.

\begin{proposition}[The Fourier basis as a continuous spectral basis]
\label{prop:fourier-eigenbasis}
\lean{Torus.fourierEigenbasis, Torus.fourierEigenbasis_level, Torus.fourierEigenbasis_basisFun,
  Torus.basisSpan_fourierEigenbasis, Torus.basisScheme_fourierEigenbasis,
  Torus.spectralLevelSet_fourierEigenbasis, Torus.localWeylBound_fourierEigenbasis,
  Torus.hasNikolskii_fourierEigenbasis}
\leanok
\uses{def:continuous-spectral-basis, def:local-weyl, def:torus-sobolev, thm:torus-nikolskii}
The trigonometric monomials $e_k$ with the level $\deg$ form a continuous spectral basis of
$L_2(\T^d)$ whose scheme is the trigonometric scheme. It satisfies the local Weyl bound of
dimension $d$ with constant $2^d$, and its scheme has the Nikolskii property of dimension $d$.
\end{proposition}

\begin{proof}
\leanok
The $e_k$ are bounded continuous representatives of the Fourier basis, and their spans are the
spaces $\mathcal{T}_n$. Since $|e_k| = 1$ and at most $(2n)^d$ frequencies have degree $< n$,
$\sum_{\deg k < n} |e_k(x)|^2 \leq 2^d n^d$. The Nikolskii property is \Cref{thm:torus-nikolskii}.
\end{proof}

\begin{theorem}[Embeddings on the torus]
\label{thm:torus-embeddings}
\lean{Torus.memBesov_exists_inclusion_eq, Torus.memBesov_exists_continuous,
  Torus.periodic_sobolev_memBesov_top, Torus.periodic_exists_sobolev_of_memBesov,
  Torus.sobolev_exists_inclusion_eq, Torus.sobolev_exists_continuous, Torus.toLp_toContinuousMap}
\leanok
\uses{thm:torus-nikolskii, thm:besov-embedding, def:torus-sobolev, prop:smoothness}
Let $1 \leq p \leq r \leq \infty$ and $1 \leq q \leq \infty$.
\begin{enumerate}
  \item If $s > d(1/p - 1/r)$, then every $f \in \Besov{s}{p}{q}(\T^d)$ is $\iota_{r,p} g$ for
    some $g \in \Besov{s - d(1/p - 1/r)}{r}{q}(\T^d)$.
  \item If $s > d/p$, then every $f \in \Besov{s}{p}{q}(\T^d)$ has a continuous
    representative.
  \item If $0 < s < t$, then $H^s(\T^d) \subseteq \Besov{s}{2}{\infty}(\T^d)$ and
    $\Besov{t}{2}{q}(\T^d) \subseteq H^s(\T^d)$.
  \item If $2 \leq r$ and $s > d(1/2 - 1/r)$, then every $f \in H^s(\T^d)$ is
    $\iota_{r,2} g$ for some $g \in \Besov{s - d(1/2 - 1/r)}{r}{\infty}(\T^d)$.
  \item If $s > d/2$, then every $f \in H^s(\T^d)$ has a continuous representative.
\end{enumerate}
\end{theorem}

\begin{proof}
\leanok
The first two items are \Cref{thm:besov-embedding} with \Cref{thm:torus-nikolskii}, the
third is \Cref{prop:smoothness} for $H^s(\T^d)$, and the last two combine both.
\end{proof}

The next theorem is the Littlewood--Paley characterisation of $\Besov{s}{2}{q}(\T^d)$ with
sharp dyadic cut-offs. For $p \neq 2$ it needs smooth cut-offs and a Fourier multiplier
theorem, which Mathlib does not have.

\begin{theorem}[Fourier coefficients]
\label{thm:torus-coefficients}
\lean{Torus.periodic_memBesov_two_two_iff, Torus.fourierBlock, Torus.enorm_fourierBlock_sq,
  Torus.periodic_memBesov_two_iff_blockNorm_lt_top}
\leanok
\uses{thm:spectral-coefficients, def:torus-sobolev}
Let $s > 0$, $1 \leq q \leq \infty$ and $f \in L_2(\T^d)$. Then $f \in
\Besov{s}{2}{q}(\T^d)$ if and only if
\[
  |\hat f(0)| + \Bigl\| \Bigl(2^{ks} \Bigl(\sum_{2^k \leq \deg m < 2^{k+1}} |\hat f(m)|^2
    \Bigr)^{1/2}\Bigr)_k \Bigr\|_{\ell_q} < \infty ,
\]
and $\Besov{s}{2}{2}(\T^d) = H^s(\T^d)$.
\end{theorem}

\begin{proof}
\leanok
\Cref{thm:spectral-coefficients} for the Fourier basis; the only index of level $0$ is
$k = 0$.
\end{proof}

\section{\texorpdfstring{$n$}{n}-term trigonometric approximation}
\label{sec:torus-nterm}

Best $n$-term approximation from the Fourier basis approximates $f \in L_2(\T^d)$ by
combinations $\sum_{k \in \Lambda} c_k e_k$ of at most $n$ monomials, with the set
$\Lambda \subseteq \Z^d$ of frequencies chosen for each $f$.

\begin{corollary}[$n$-term trigonometric approximation]
\label{cor:torus-nterm}
\lean{Torus.periodic_memApprox_nTerm_iff}
\leanok
\uses{thm:stechkin}
Let $0 < \tau < 2$ and $\theta = 1/\tau - 1/2$. Then $f \in L_2(\T^d)$ lies in the
approximation space $A^\theta_\tau$ of $n$-term approximation from the dictionary
$\{c \, e_k : k \in \Z^d, \ c \in \C\}$ if and only if
\[
  \sum_{k \in \Z^d} |\hat f(k)|^\tau < \infty .
\]
\end{corollary}

\begin{proof}
\leanok
\Cref{thm:stechkin} for Mathlib's Fourier basis of $L_2(\T^d)$, whose coefficients are
$\langle e_k, f \rangle = \hat f(k)$.
\end{proof}

For $\tau = 1$ the condition says that the Fourier series of $f$ converges absolutely, the
setting of Stechkin~\cite{Stechkin1955}; for $\tau < 1$ it is stronger. Linear approximation by
the trigonometric polynomials of degree $< n$ leads instead to conditions on the $\ell_2$ norms
of the dyadic blocks of coefficients ordered by degree, such as $\Besov{s}{2}{q}(\T^d)$ in
\Cref{thm:torus-coefficients}, and to the weighted $\ell_2$ condition of $H^s(\T^d)$ in
\Cref{def:torus-sobolev}.

\section{Traces on the torus}
\label{sec:torus-traces}

The trace of a function on $\T^{n+1}$ is its restriction to the face $\{x_0 = 0\} \cong \T^n$,
$f \mapsto f(0, \cdot)$. It is defined on continuous functions and not on $L_p$; the abstract
trace theorems of \Cref{ch:traces-abstract} extend it to Besov spaces.

\begin{definition}[Restriction to the face]
\label{def:face-restriction}
\lean{Torus.faceEmbedding, Torus.restrictFace, Torus.compLinear, Torus.lineThrough,
  Torus.restrictFace_trigMonomial, Torus.restrictFace_mem_trigPoly,
  Torus.compLinear_lineThrough_mem_trigPoly, Torus.faceEmbedding_apply, Torus.lineThrough_apply,
  Torus.compLinear_apply, Torus.trigDeg_tail_le, Torus.trigDeg_const_le,
  Torus.compLinear_lineThrough_trigMonomial}
\leanok
\uses{def:trig-poly}
For a bounded continuous $a$ on $\T^{n+1}$, its restriction to the face is
$a(0, \cdot) : x' \mapsto a(0, x')$ on $\T^n$. It maps $\mathcal{T}_m$ on $\T^{n+1}$ into
$\mathcal{T}_m$ on $\T^n$, since $e_k(0, x') = e_{(k_1, \dots, k_n)}(x')$; so does the
restriction to every line $t \mapsto a(t, x')$, onto $\mathcal{T}_m$ on $\T^1$; both follow
from the statement on linear maps in \Cref{def:trig-poly}.
\end{definition}

\begin{lemma}[Fubini on $\T^{n+1} = \T \times \T^n$]
\label{lem:fubini-cons}
\lean{Torus.lintegral_eq_lintegral_cons}
\leanok
For every measurable $F : \T^{n+1} \to [0, \infty]$,
$\int_{\T^{n+1}} F = \int_{\T^n} \int_{\T} F(t, x') \, dt \, dx'$.
\end{lemma}

\begin{proof}
\leanok
The map $(t, x') \mapsto (t, x')$ is measure preserving from the product measure, and Tonelli's
theorem applies.
\end{proof}

\begin{theorem}[Trace inequality on the torus]
\label{thm:trace-inequality-torus}
\lean{Torus.eLpNorm_restrictFace_le, Torus.enorm_rpow_le_lintegral}
\leanok
\uses{def:face-restriction, lem:fubini-cons, thm:torus-nikolskii}
For all $m \in \N$, $a \in \mathcal{T}_m$ on $\T^{n+1}$ and $1 \leq p < \infty$,
\[
  \|a(0, \cdot)\|_{L_p(\T^n)} \leq \bigl(2 \lceil p/2 \rceil m\bigr)^{1/p} \,
  \|a\|_{L_p(\T^{n+1})} .
\]
\end{theorem}

\begin{proof}
\leanok
For fixed $x'$ the line $t \mapsto a(t, x')$ is a trigonometric polynomial of degree $< m$ in
one variable, so \Cref{thm:torus-nikolskii} with $d = 1$ gives
$|a(0, x')|^p \leq 2 \lceil p/2 \rceil m \int_\T |a(t, x')|^p \, dt$. Integrate over $x'$ and
apply \Cref{lem:fubini-cons}.
\end{proof}

The opposite direction needs a polynomial in $x_0$ that equals $1$ at $0$ and is small in
$L_p$. The normalised Fejér kernel serves, and Parseval's identity replaces the usual pointwise
kernel estimates.

\begin{lemma}[Extension by the Fejér kernel]
\label{lem:fejer-extension}
\lean{Torus.dirichletSum, Torus.fejerFace, Torus.extendFace, Torus.tailMap,
  Torus.fejerFace_cons_zero, Torus.enorm_fejerFace_le, Torus.lintegral_enorm_fejerFace,
  Torus.lintegral_enorm_dirichletSum_sq, Torus.extendFace_mem_trigPoly,
  Torus.restrictFace_extendFace, Torus.eLpNorm_extendFace_le, Torus.tailMap_apply,
  Torus.trigDeg_cons_zero_le, Torus.compLinear_tailMap_trigMonomial,
  Torus.compLinear_tailMap_mem_trigPoly, Torus.trigMonomial_single_cons, Torus.trigDeg_single_le,
  Torus.star_trigMonomial, Torus.star_mem_trigPoly, Torus.dirichletSum_cons,
  Torus.dirichletSum_mem_trigPoly, Torus.norm_dirichletSum_le, Torus.fejerFace_apply,
  Torus.fejerFace_mem_trigPoly, Torus.toLpR, Torus.toLpR_injective, Torus.enorm_toLpR,
  Torus.iUnionSubmodule_le_range}
\leanok
\uses{def:face-restriction, lem:fubini-cons, lem:trig-products}
Let $m \geq 1$ and $\varphi_m(x) = m^{-2} \bigl|\sum_{j<m} e^{2\pi i j x_0}\bigr|^2$ on
$\T^{n+1}$. For every $b \in \mathcal{T}_m$ on $\T^n$, the function
$E b : x \mapsto \varphi_m(x) \, b(x_1, \dots, x_n)$ lies in $\mathcal{T}_{3m}$, satisfies
$(E b)(0, \cdot) = b$, and for every $1 \leq p < \infty$,
\[
  \|E b\|_{L_p(\T^{n+1})} \leq m^{-1/p} \, \|b\|_{L_p(\T^n)} .
\]
\end{lemma}

\begin{proof}
\leanok
The Dirichlet sum has degree $< m$ and modulus at most $m$, so $\varphi_m$ has degree $< 2m$,
$0 \leq \varphi_m \leq 1$ and $\varphi_m(0, \cdot) = 1$. By Parseval,
$\int \varphi_m = m^{-2} \cdot m = 1/m$, hence $\int \varphi_m^p \leq \int \varphi_m = 1/m$,
and \Cref{lem:fubini-cons} gives $\int |E b|^p = \int \varphi_m^p \cdot \int |b|^p$.
\end{proof}

\begin{theorem}[Traces of Besov spaces on the torus]
\label{thm:trace-torus}
\lean{Torus.periodicTraceBound, Torus.periodicTraceOnApprox,
  Torus.periodicTraceBound_limitTrace_toLp, Torus.periodicTraceBound_rightInverse,
  Torus.trace_besov, Torus.periodicTraceOnApprox_apply, Torus.periodicTraceConst,
  Torus.periodicTraceBound_T_toLp}
\leanok
\uses{thm:trace-inequality-torus, lem:fejer-extension, thm:trace-transfer,
  thm:trace-surjective}
Let $1 \leq p < \infty$, $1 \leq q \leq \infty$ and $s > 0$. The restriction to the face
extends from trigonometric polynomials to a map $\operatorname{Tr}$ from
$\Besov{s+1/p}{p}{q}(\T^{n+1})$ into $\Besov{s}{p}{q}(\T^n)$, and every
$g \in \Besov{s}{p}{q}(\T^n)$ is $\operatorname{Tr} f$ for some
$f \in \Besov{s+1/p}{p}{q}(\T^{n+1})$.
\end{theorem}

\begin{proof}
\leanok
By \Cref{thm:trace-inequality-torus} the restriction is a trace bound of exponent $1/p$ from
the trigonometric scheme of $\T^{n+1}$ to that of $\T^n$ in $L_p$, and by
\Cref{lem:fejer-extension} with $m = 2^k$ it has a right inverse bound with $\gamma = 2$ and
$C' = 1$. Apply \Cref{thm:trace-transfer} with $s + 1/p$ in place of $s$, and
\Cref{thm:trace-surjective}.
\end{proof}

The map $\operatorname{Tr}$ is $f \mapsto f(0, \cdot)$ on trigonometric polynomials, and on
every continuous function of $\Besov{1/p}{p}{1}(\T^{n+1})$ it is the restriction as well
(\Cref{thm:trace-torus-continuous}). The proof uses approximants that are near-best both in
$L_p$ and in the sup norm: smooth de la Vallée Poussin operators, built from the band-limiting
kernel $\Theta_J$ of \Cref{def:lp-kernels} acting on $\T^d$ by the convolution of
\Cref{def:torus-econv} in \Cref{ch:classical}, on bounded continuous functions as well
(\Cref{lem:torus-uniform-density} and \Cref{lem:torus-econv-continuous}). Their multiplier is $1$
on frequencies of degree $< 2^k$ and $0$ from degree $2^{k+\#d+2}$ on.

\begin{lemma}[De la Vallée Poussin operators on the torus]
\label{lem:torus-vallee-poussin}
\lean{Torus.fourier_dilatedBandKernel_latticeVec_eq_zero, Torus.valleePoussin,
  Torus.valleePoussinBCF, Torus.valleePoussin_mem_trigPolyLp, Torus.valleePoussinBCF_mem_trigPoly,
  Torus.valleePoussin_eq_self, Torus.valleePoussinBCF_eq_self, Torus.toLp_valleePoussinBCF,
  Torus.valleePoussinConst, Torus.valleePoussinConst_ne_top, Torus.toReal_valleePoussinConst,
  Torus.enorm_sub_valleePoussin_le, Torus.norm_sub_valleePoussinBCF_le,
  Torus.tendsto_valleePoussinBCF, Torus.tendsto_dyadicWeight_mul_enorm_sub_valleePoussin}
\leanok
\uses{lem:torus-uniform-density, lem:torus-econv-continuous, thm:torus-econv-fourier,
  thm:torus-lp-bernstein, lem:lp-kernels-basic, lem:lp-kernels-frequency, lem:torus-proj,
  lem:lebesgue-one-operator, lem:trace-ext}
Let $C_V = 1 + \|\Theta\|_1$ and, for $k \in \N$, let $V_k f = \Theta_{k+\#d} * f$, both for
$f \in L_p(\T^d)$, $1 \leq p < \infty$, and for bounded continuous $f$ on $\T^d$. For all
$k \in \N$, $1 \leq p < \infty$, $f \in L_p(\T^d)$, bounded continuous $g$ on $\T^d$ and
$a \in \mathcal{T}_{2^k}$:
\begin{enumerate}
  \item $V_k f \in \mathcal{T}_{2^{k+\#d+2}}$, $V_k g \in \mathcal{T}_{2^{k+\#d+2}}$, the class
    of $V_k g$ in $L_p(\T^d)$ is $V_k$ applied to the class of $g$, and $V_k a = a$;
  \item $\|f - V_k f\|_p \leq C_V \|f - a\|_p$ and
    $\|g - V_k g\|_\infty \leq C_V \|g - a\|_\infty$.
\end{enumerate}
For every bounded continuous $g$ on $\T^d$, $\|g - V_k g\|_\infty \to 0$ as $k \to \infty$. For
all $1 \leq p < \infty$, $\beta \geq 0$ and $f \in A^\beta_1$ of the trigonometric scheme in
$L_p(\T^d)$, $2^{k\beta} \|f - V_k f\|_p \to 0$ as $k \to \infty$.
\end{lemma}

\begin{proof}
\leanok
For $J = k + \#d$, $\mathcal{F}\Theta_J(\xi) = 0$ for $|\xi| \geq 2^{J+2}$
(\Cref{lem:lp-kernels-frequency}) and $\deg m \leq |m|$ (\Cref{lem:torus-proj}), so
$\mathcal{F}\Theta_J(m) = 0$ if $\deg m \geq 2^{J+2}$, and $\Theta_J * a = a$ for
$a \in \mathcal{T}_{2^k}$ (\Cref{thm:torus-lp-bernstein}). \Cref{thm:torus-econv-fourier} and
\Cref{lem:torus-econv-continuous} give the first item. Since
$\|\Theta_J\|_1 = \|\Theta\|_1$ (\Cref{lem:lp-kernels-basic}), $V_k$ has norm at most
$\|\Theta\|_1$ in both norms, and \Cref{lem:lebesgue-one-operator} gives the second item. By the
second item, $\|g - V_k g\|_\infty \leq C_V E_{2^k}(g)_\infty$, which tends to $0$ by
\Cref{lem:torus-uniform-density}. For the last claim, the canonical near-best approximants
$a_k \in \mathcal{T}_{2^k}$ of \Cref{lem:trace-ext} satisfy $2^{k\beta} \|f - a_k\|_p \to 0$, and
$\|f - V_k f\|_p \leq C_V \|f - a_k\|_p$.
\end{proof}

\begin{theorem}[The trace of a continuous function]
\label{thm:trace-torus-continuous}
\lean{Torus.continuous_restrictFace, Torus.toLpR_valleePoussinBCF_mem,
  Torus.periodicTraceBound_limitTrace_eq_restrictFace, Torus.trace_besov_eq_restrictFace,
  Torus.ae_trace_eq_restrictFace, Torus.exists_trace_eq_restrictFace}
\leanok
\uses{thm:trace-torus, lem:trace-ext, lem:torus-vallee-poussin, thm:torus-embeddings,
  thm:approx-embedding}
Let $1 \leq p < \infty$ and let $\operatorname{Tr}$ be the trace of \Cref{lem:trace-ext} for the
restriction to the face, as in \Cref{thm:trace-torus}.
\begin{enumerate}
  \item For every bounded continuous $f$ on $\T^{n+1}$ with
    $f \in \Besov{1/p}{p}{1}(\T^{n+1})$, $\operatorname{Tr} f = f(0, \cdot)$. In particular,
    $\operatorname{Tr} f = f(0, \cdot)$ for all $s > 0$, $1 \leq q \leq \infty$ and bounded
    continuous $f \in \Besov{s+1/p}{p}{q}(\T^{n+1})$.
  \item For all $s > n/p$, $1 \leq q \leq \infty$ and $f \in \Besov{s+1/p}{p}{q}(\T^{n+1})$
    there is a bounded continuous $g$ on $\T^{n+1}$ with $f = g$ almost everywhere and
    $\operatorname{Tr} f = g(0, \cdot)$.
\end{enumerate}
\end{theorem}

\begin{proof}
\leanok
By \Cref{lem:torus-vallee-poussin}, $V_k f$ lies in $\mathcal{T}_{2^{k+n+3}}$ and
$2^{k/p} \|f - V_k f\|_p \to 0$, so \Cref{lem:trace-ext} gives
$(V_k f)(0, \cdot) \to \operatorname{Tr} f$ in $L_p(\T^n)$. Since $V_k f \to f$ uniformly,
also $(V_k f)(0, \cdot) \to f(0, \cdot)$ uniformly, hence in $L_p(\T^n)$. For $s > 0$,
$\Besov{s+1/p}{p}{q}(\T^{n+1}) \subseteq \Besov{1/p}{p}{1}(\T^{n+1})$ by
\Cref{thm:approx-embedding}. For the second item, $s + 1/p > (n+1)/p$, and
\Cref{thm:torus-embeddings} gives the continuous representative $g$.
\end{proof}

\chapter{Classical definitions on the torus}
\label{ch:classical}

The Besov and Sobolev spaces of \Cref{ch:besov}, \Cref{ch:sobolev} and \Cref{ch:torus} are
defined through approximation. The classical definitions use moduli of smoothness and
derivatives~\cite{Triebel, DeVoreLorentz}. This chapter shows that they coincide on
$L_2(\T^d)$, and for moduli of smoothness and Littlewood--Paley blocks on every $L_p(\T^d)$,
$1 \leq p < \infty$, with the notation of \Cref{ch:torus}; the spaces by moduli of smoothness
are also real interpolation spaces between two of them (\Cref{thm:torus-besov-interpolate}).

\section{Moduli of smoothness}

\begin{definition}[Translation]
\label{def:translation}
\lean{Torus.periodicTranslate, Torus.coeFn_periodicTranslate}
\leanok
For $h \in \T^d$, the translation $\tau_h : L_2(\T^d) \to L_2(\T^d)$ is
$(\tau_h f)(x) = f(x + h)$.
\end{definition}

\begin{lemma}[Fourier coefficients of translates]
\label{lem:translation-fourier}
\lean{Torus.mFourierCoeff_periodicTranslate, Torus.mFourier_add_apply,
  Torus.enorm_periodicTranslate_sub_sq, Torus.norm_mFourier_sub_one_le, Torus.norm_prod_sub_one_le,
  Torus.norm_fourier_sub_one_le, Torus.mFourier_apply_zero, Torus.mFourier_neg_neg,
  Torus.repr_periodicTranslate}
\leanok
\uses{def:translation}
For every $h \in \T^d$, $f \in L_2(\T^d)$ and $k \in \Z^d$,
$\widehat{\tau_h f}(k) = e_k(h) \hat f(k)$ and $|e_k(h) - 1| \leq 2\pi \, \#d \, \deg k \,
\|h\|$, and
\[
  \|\tau_h f - f\|_2^2 = \sum_{k \in \Z^d} |e_k(h) - 1|^2 |\hat f(k)|^2 .
\]
\end{lemma}

\begin{proof}
\leanok
The Haar measure is invariant under translation and $e_k(x + h) = e_k(x) e_k(h)$; the second
formula is Parseval's identity. For the bound, $|\prod_j a_j - 1| \leq \sum_j |a_j - 1|$ for
complex numbers of modulus $1$, and $|e^{i\theta} - 1| \leq |\theta|$.
\end{proof}

\begin{definition}[Modulus of smoothness]
\label{def:modulus}
\lean{Torus.periodicModulus, Torus.enorm_periodicTranslate_sub_eq, Torus.periodicModulus_eq_iSup,
  Torus.periodicModulus_mono, Torus.enorm_periodicTranslate_sub_le_periodicModulus,
  Torus.norm_coe_le_abs}
\leanok
\uses{def:translation, def:rd-modulus}
For $f \in L_2(\T^d)$ and $t \in \R$,
$\omega(f, t)_2 = \sup_{\|h\| \leq t} \|\tau_h f - f\|_2$. Since
$\|\tau_h f - f\|_2 = \|\Delta_h f\|_2$, this is the modulus of order $1$ in $L_2$ of
\Cref{def:rd-modulus} for the Haar measure on $\T^d$.
\end{definition}

\begin{lemma}[Growth of the modulus]
\label{lem:modulus-linear}
\lean{Torus.enorm_periodicTranslate, Torus.periodicTranslate_zero, Torus.periodicTranslate_add,
  Torus.enorm_periodicTranslate_nsmul_sub_le, Torus.exists_coe_eq_abs_eq_norm,
  Torus.periodic_smoothnessModulus_nat_mul_le, Torus.periodic_smoothnessModulus_mul_le,
  Torus.norm_exp_I_mul_sub_one_sq, Torus.fourier_coe_eq_exp, Torus.norm_fourier_sub_one_sq}
\leanok
\uses{def:modulus, lem:rd-modulus}
For $f \in L_2(\T^d)$, $h \in \T^d$ and $n \in \N$,
$\|\tau_{n h} f - f\|_2 \leq n \|\tau_h f - f\|_2$. Let $\omega_r(f, t)_p$ be the modulus of
order $r$ of \Cref{def:rd-modulus} for the Haar measure on $\T^d$. For measurable $f$,
$1 \leq p$, $r \in \N$, $m \geq 1$ and $c, t \geq 0$,
\[
  \omega_r(f, m t)_p \leq m^r \omega_r(f, t)_p , \qquad
  \omega_r(f, c t)_p \leq (1 + c)^r \omega_r(f, t)_p ;
\]
in particular $\omega(f, n t)_2 \leq n \, \omega(f, t)_2$ for $n \geq 1$.
\end{lemma}

\begin{proof}
\leanok
$\tau_{n h} f - f = \sum_{i < n} \tau_{i h} (\tau_h f - f)$, where translations are isometries.
Every $h$ with $\|h\| \leq m t$ is $m h'$ with $\|h'\| \leq t$: divide representatives $x_j$ of
$h_j$ with $|x_j| = \|h_j\|$ by $m$. So the divisibility hypothesis of \Cref{lem:rd-modulus}
holds on $\T^d$, and \Cref{lem:rd-modulus} gives
$\omega_r(f, m t)_p \leq m^r \omega_r(f, t)_p$ and, for real $c$ with
$m = \lfloor c \rfloor + 1 \leq 1 + c$, the second bound.
\end{proof}

The classical Jackson theorem bounds the error of trigonometric approximation by the modulus
of smoothness. For $p = 2$ it follows from Parseval's identity by averaging over
translations along the coordinate axes.

\begin{theorem}[Jackson's theorem on the torus]
\label{thm:torus-jackson}
\lean{Torus.approxError_sq_le_card_mul_periodicModulus_sq, Torus.approxError_sq_le_tsum,
  Torus.integral_norm_fourier_sub_one_sq, Torus.le_integral_norm_fourier_sub_one_sq,
  Torus.approxError_le_periodicModulus, Torus.ofReal_le_lintegral_enorm_fourier_sub_one_sq,
  Torus.mFourier_apply_single}
\leanok
\uses{def:modulus, lem:translation-fourier, def:torus-sobolev}
For every $n \geq 1$ and $f \in L_2(\T^d)$, the error of approximation by $\mathcal{T}_n$ in
$L_2$ satisfies
\[
  E_n(f)_2^2 \leq \#d \cdot \omega(f, 1/n)_2^2 .
\]
\end{theorem}

\begin{proof}
\leanok
First $E_n(f)_2^2 \leq \sum_{\deg k \geq n} |\hat f(k)|^2$, by the truncation. For an integer
$a$ and $t > 0$,
$\int_0^t |e^{2\pi i a u} - 1|^2 \, du = 2t - \sin(2\pi a t)/(\pi a)$, which is at least $t$
when $|a| t \geq 1$. For $\deg k \geq n$ some coordinate has $|k_j| \geq n$, so
$1 \leq n \int_0^{1/n} |e_k(u \, e_j) - 1|^2 \, du$ with $u\,e_j$ the point $u$ in coordinate
$j$. Hence, by \Cref{lem:translation-fourier},
\[
  \sum_{\deg k \geq n} |\hat f(k)|^2 \leq \sum_{j \in d} n \int_0^{1/n} \sum_k
  |e_k(u e_j) - 1|^2 |\hat f(k)|^2 \, du
\]
and the inner sum is $\|\tau_{u e_j} f - f\|_2^2$. Since $\|u e_j\| \leq u \leq 1/n$, each
integrand is at most $\omega(f, 1/n)_2^2$.
\end{proof}

In particular $E_n(f)_2 \leq (\#d + 1) \, \omega(f, 1/n)_2$, the form used below.

\begin{lemma}[The inverse inequality on the torus]
\label{lem:torus-inverse}
\lean{Torus.mFourierCoeff_eq_zero_of_mem, Torus.enorm_periodicTranslate_sub_le_of_mem,
  Torus.periodicModulus_le_of_mem}
\leanok
\uses{def:modulus, lem:translation-fourier, def:torus-sobolev}
Let $n, j, k \in \N$, $h \in \T^d$, and let $a \in \mathcal{T}_n$ be a trigonometric
polynomial of degree $< n$. Then $\hat a(k) = 0$ for $\deg k \geq n$ and
$\|\tau_h a - a\|_2 \leq 2\pi \, \#d \, n \|h\| \, \|a\|_2$. For $a \in \mathcal{T}_{2^j}$,
\[
  \omega(a, 2^{-k})_2 \leq \max(2, 2\pi \, \#d) \, 2^{-(k - j)_+} \|a\|_2 .
\]
\end{lemma}

\begin{proof}
\leanok
By \Cref{lem:translation-fourier}, $\|\tau_h a - a\|_2^2 = \sum_k |e_k(h) - 1|^2 |\hat a(k)|^2$
with $|e_k(h) - 1| \leq 2\pi \, \#d \deg k \, \|h\|$, and $\hat a(k) = 0$ unless
$\deg k < n$; Parseval's identity gives the first bound. For $k < j$ use
$\omega(a, t)_2 \leq 2 \|a\|_2$, and for $k \geq j$ the first bound with $\|h\| \leq 2^{-k}$.
\end{proof}

\begin{definition}[Besov norm through the modulus]
\label{def:modulus-norm}
\lean{Torus.trigModulusData}
\leanok
\uses{def:modulus, def:modulus-data, thm:torus-jackson, lem:torus-inverse}
The modulus $\omega(f, k) = \omega(f, 2^{-k})_2$ is compatible with the scheme of trigonometric
polynomials in $L_2(\T^d)$ with inverse order $\theta = 1$ (\Cref{def:modulus-data}), with
$C_0 = 2$, $k_0 = 0$, $C_J = \#d + 1$ and $C_B = \max(2, 2\pi \, \#d)$. For $s \in \R$,
$1 \leq q \leq \infty$ and $f \in L_2(\T^d)$, its classical quasi-norm is
\[
  \|f\|^{\omega}_{s,q} = \|f\|_2 + \bigl\| \bigl(2^{js} \omega(f, 2^{-j})_2\bigr)_{j \in \N}
  \bigr\|_{\ell_q} .
\]
\end{definition}

\begin{theorem}[Besov spaces through moduli of smoothness]
\label{thm:moduli-coincidence}
\lean{Torus.approxNormDyadic_le_modulusNorm, Torus.modulusNorm_le_approxNormDyadic,
  Torus.memBesov_iff_modulusNorm_lt_top}
\leanok
\uses{def:modulus-norm, thm:modulus-coincidence}
Let $s \in \R$, $1 \leq q \leq \infty$ and $f \in L_2(\T^d)$. Then for $s \geq 0$
\[
  \|f\|^{\mathrm{dy}}_{\Besov{s}{2}{q}} \leq (\#d + 2) \, \|f\|^{\omega}_{s,q} ,
\]
and for $0 \leq s < 1$, with $C_B = \max(2, 2\pi \, \#d)$,
\[
  \|f\|^{\omega}_{s,q} \leq \bigl(7 + 27 \, C_B (1 - 2^{-(1-s)})^{-1}\bigr)
  \|f\|^{\mathrm{dy}}_{\Besov{s}{2}{q}} .
\]
In particular, for $0 < s < 1$, $f \in \Besov{s}{2}{q}(\T^d)$ if and only if
$\|f\|^{\omega}_{s,q} < \infty$.
\end{theorem}

\begin{proof}
\leanok
\Cref{thm:modulus-coincidence} with the constants of \Cref{def:modulus-norm}: the first
constant is $1 + C_J$, the second $1 + 3 C_0 + C_B (3 + 6 \cdot 4)(1 - 2^{-(1-s)})^{-1}$.
\end{proof}

The modulus is also controlled by the dyadic Fourier blocks
$\Delta_K f = \sum_{2^K \leq \deg k < 2^{K+1}} \hat f(k) e_k$.

\begin{lemma}[The modulus through the dyadic blocks]
\label{lem:torus-modulus-blocks}
\lean{Torus.blockMult, Torus.enorm_mFourier_sub_one_le_blockMult,
  Torus.enorm_periodicTranslate_sub_sq_le, Torus.periodicModulus_le_tsum_block, Torus.blockMult_le}
\leanok
\uses{def:modulus, lem:translation-fourier, thm:torus-coefficients}
Let $f \in L_2(\T^d)$, $t \in \R$, and let $m_K(t) = \min(2, 2\pi \, \#d \, 2^{K+1} t)$. For
$\|h\| \leq t$,
\[
  \|\tau_h f - f\|_2^2 \leq \sum_K m_K(t)^2 \, \|\Delta_K f\|_2^2 ,
  \qquad\text{so}\qquad
  \omega(f, t)_2 \leq \sum_K m_K(t) \, \|\Delta_K f\|_2 .
\]
For $j \in \N$, $m_K(2^{-j}) \leq 2$ for $K \geq j$ and $m_K(2^{-j}) \leq 4\pi \, \#d \,
2^{K - j}$ for $K \leq j$.
\end{lemma}

\begin{proof}
\leanok
On the block $K$, $|e_k(h) - 1| \leq \min(2, 2\pi \, \#d \deg k \, \|h\|)$ by
\Cref{lem:translation-fourier}; Parseval's identity and $\sum_K b_K^2 \leq (\sum_K b_K)^2$
give the bounds.
\end{proof}

For $s \geq 1$ the first-order modulus no longer describes $\Besov{s}{2}{q}$; the moduli of
higher order describe $\Besov{s}{p}{q}$ for every $s > 0$ and $1 \leq p < \infty$
(\Cref{sec:torus-lp-moduli}).

\section{Convolution with kernels on \texorpdfstring{$\R^d$}{R\^{}d}}

Let $\pi : \R^d \to \T^d$, $\pi(x) = (x_j \bmod 1)_j$, be the quotient map, and view
$m \in \Z^d$ as a point of $\R^d$.

\begin{lemma}[The quotient map]
\label{lem:torus-proj}
\lean{Torus.quotientProj, Torus.quotientProj_apply, Torus.continuous_quotientProj,
  Torus.measurable_quotientProj, Torus.norm_quotientProj_le, Torus.norm_le_card_mul,
  Torus.exists_quotientProj_eq, Torus.latticeVec, Torus.latticeVec_apply, Torus.latticeVec_neg,
  Torus.inner_latticeVec, Torus.inner_latticeVec_single, Torus.trigDeg_le_norm_latticeVec,
  Torus.norm_latticeVec_le, Torus.mFourier_neg_quotientProj,
  Torus.integral_mul_mFourier_neg_quotientProj}
\leanok
\uses{def:trig-poly}
The quotient map is an additive homomorphism with $\|\pi x\| \leq |x|$, and every $h \in \T^d$ is
$\pi x$ with $|x| \leq \#d \, \|h\|$. For $m \in \Z^d$ and $v \in \R^d$, the point $-m$ of
$\R^d$ is the negative of $m$, $\langle m, v \rangle = \sum_i m_i v_i$, and
$\langle m, e_i \rangle = m_i$ for every unit vector $e_i$. Moreover
$\deg m \leq |m| \leq \#d \deg m$ and $e_m(-\pi y) = e^{-2\pi i \langle y, m \rangle}$, so
$\int_{\R^d} \Phi(y) e_m(-\pi y) \, dy = \mathcal{F}\Phi(m)$.
\end{lemma}

\begin{proof}
\leanok
Choose representatives $x_j$ of $h_j$ with $|x_j| = \|h_j\|$; then
$|x|^2 = \sum_j x_j^2 \leq (\#d \, \|h\|)^2$. The characters multiply:
$e_m(\pi y) = \prod_j e^{2\pi i m_j y_j}$.
\end{proof}

\begin{definition}[Convolution with a kernel on $\R^d$]
\label{def:torus-econv}
\lean{Torus.periodicEconv, Torus.measurable_comp_sub_quotientProj, Torus.eLpNorm_comp_sub,
  Torus.aestronglyMeasurable_periodicEconv, Torus.periodicEconv_congr_ae}
\leanok
\uses{lem:torus-proj}
For $\Phi : \R^d \to \C$ and $f : \T^d \to \C$,
\[
  (\Phi * f)(t) = \int_{\R^d} \Phi(y) f(t - \pi y) \, dy ,
\]
the convolution of $f$ with the periodisation $\sum_{m \in \Z^d} \Phi(\cdot + m)$ of $\Phi$. It
depends on $f$ only through its class modulo null sets.
\end{definition}

\begin{lemma}[Young's inequality on the torus]
\label{lem:torus-econv-young}
\lean{Torus.eLpNorm_periodicEconv_le, Torus.memLp_periodicEconv, Torus.ae_integrable_econv,
  Torus.periodicEconv_add_ae, Torus.periodicEconv_smul, Torus.periodicEconvLp,
  Torus.periodicEconvLp_apply_ae, Torus.coeFn_sum_periodicEconvLp_ae,
  Torus.enorm_periodicEconvLp_le}
\leanok
\uses{def:torus-econv, lem:weighted-jensen}
For $1 \leq p < \infty$, $\Phi \in L_1(\R^d)$ and $f \in L_p(\T^d)$, the integral converges
absolutely for almost every $t$ and $\|\Phi * f\|_{L_p(\T^d)} \leq \|\Phi\|_{L_1(\R^d)} \|f\|_p$.
So $f \mapsto \Phi * f$ is a bounded operator on $L_p(\T^d)$. For all $f, g \in L_1(\T^d)$ and
$c \in \C$, $\Phi * (f + g) = \Phi * f + \Phi * g$ almost everywhere and
$\Phi * (c f) = c \, (\Phi * f)$ at every point. For every finite set $s$, integrable kernels
$K_j$ and $u_j \in L_p(\T^d)$, the class $\sum_{j \in s} K_j * u_j$ is represented by
$t \mapsto \sum_{j \in s} (K_j * u_j)(t)$.
\end{lemma}

\begin{proof}
\leanok
Every translate $f(\cdot - \pi y)$ has the norm of $f$, and \Cref{lem:weighted-jensen} averages
them with the weight $|\Phi|$. Tonelli's theorem gives
$\int \int |\Phi(y)| |f(t - \pi y)| \, dy \, dt = \|\Phi\|_1 \|f\|_1$, so the integrals of
$\Phi * f$ and $\Phi * g$ converge absolutely almost everywhere and add there. Constants leave the
integral, and the class of a finite sum is the sum of the representatives.
\end{proof}

\begin{theorem}[Fourier multipliers on the torus]
\label{thm:torus-econv-fourier}
\lean{Torus.periodicEconv_mFourier, Torus.periodicEconvLp_mFourierLp, Torus.trigPolyLp,
  Torus.trigPolyLp_eq_span, Torus.mFourierLp_mem_trigPolyLp, Torus.isClosed_trigPolyLp,
  Torus.mem_toLp_trigScheme_iff, Torus.periodicEconvLp_mem_trigPolyLp,
  Torus.periodicEconvLp_eq_self_of_mem}
\leanok
\uses{lem:torus-econv-young, def:trig-poly}
For $\Phi \in L_1(\R^d)$ and $m \in \Z^d$, $\Phi * e_m = \mathcal{F}\Phi(m) \, e_m$. If
$\mathcal{F}\Phi(m) = 0$ whenever $\deg m \geq n$, then $\Phi * f \in \mathcal{T}_n$ for every
$f \in L_p(\T^d)$, $1 \leq p < \infty$. If $\mathcal{F}\Phi(m) = 1$ whenever $\deg m < n$, then
$\Phi * a = a$ for every $a \in \mathcal{T}_n$.
\end{theorem}

\begin{proof}
\leanok
$e_m(t - \pi y) = e_m(t) e_m(-\pi y)$ and \Cref{lem:torus-proj}. The trigonometric polynomials of
degree $< n$ form a finite-dimensional, hence closed, subspace of $L_p(\T^d)$, the operator maps
every $e_m$ into it, and the $e_m$ span a dense subspace of $L_p(\T^d)$ (Mathlib's
\texttt{span\_\allowbreak mFourierLp\_\allowbreak closure\_\allowbreak eq\_\allowbreak top}).
\end{proof}

\begin{lemma}[Uniform density of trigonometric polynomials]
\label{lem:torus-uniform-density}
\lean{Torus.map_span_mFourier_le_iSup_trigPoly, Torus.dense_iUnion_trigPoly,
  Torus.tendsto_approxError_trigScheme}
\leanok
\uses{def:trig-poly, lem:error-tendsto}
The trigonometric polynomials are dense in the bounded continuous functions on $\T^d$ for the
sup norm. For every bounded continuous $g$ on $\T^d$, the error $E_n(g)_\infty$ of best uniform
approximation by $\mathcal{T}_n$ tends to $0$ as $n \to \infty$.
\end{lemma}

\begin{proof}
\leanok
The bounded continuous functions are the image of $C(\T^d)$ under an isometry, and the span of
the characters is dense in $C(\T^d)$ (Mathlib's
\texttt{span\_\allowbreak mFourier\_\allowbreak closure\_\allowbreak eq\_\allowbreak top}). Every
character is a trigonometric polynomial, so the image of the span lies in
$\bigcup_n \mathcal{T}_n$. \Cref{lem:error-tendsto} gives the convergence of the errors.
\end{proof}

\begin{lemma}[Convolution of continuous functions on the torus]
\label{lem:torus-econv-continuous}
\lean{Torus.integrable_mul_comp_sub_quotientProj, Torus.continuous_periodicEconv,
  Torus.norm_periodicEconv_le, Torus.periodicEconvBCF, Torus.periodicEconvBCF_apply,
  Torus.norm_periodicEconvBCF_le, Torus.toLp_periodicEconvBCF, Torus.periodicEconvBCF_mem_trigPoly,
  Torus.periodicEconvBCF_eq_self_of_mem}
\leanok
\uses{def:torus-econv, lem:torus-econv-young, thm:torus-econv-fourier, def:trig-poly}
For every $\Phi \in L_1(\R^d)$ the map $g \mapsto \Phi * g$ is a bounded linear operator on the
bounded continuous functions on $\T^d$ with $\|\Phi * g\|_\infty \leq \|\Phi\|_1 \|g\|_\infty$.
For all $1 \leq p < \infty$ and bounded continuous $g$, the class of $\Phi * g$ in $L_p(\T^d)$
is $\Phi$ convolved with the class of $g$. For every $n \in \N$: if $\mathcal{F}\Phi(m) = 0$
whenever $\deg m \geq n$, then $\Phi * g \in \mathcal{T}_n$ for every bounded continuous $g$;
if $\mathcal{F}\Phi(m) = 1$ whenever $\deg m < n$, then $\Phi * a = a$ for every
$a \in \mathcal{T}_n$.
\end{lemma}

\begin{proof}
\leanok
Dominated convergence with the bound $|\Phi(y)| \|g\|_\infty$ shows that $\Phi * g$ is
continuous, and the same bound gives the sup-norm estimate. Both convolutions are the same
integral, so the classes agree. The last two claims hold for the classes in $L_2$ by
\Cref{thm:torus-econv-fourier}, and continuous functions with the same class are equal.
\end{proof}

\begin{lemma}[Differences of the convolution]
\label{lem:torus-econv-differences}
\lean{Torus.periodicEconv_add_quotientProj, Torus.higherDiff_periodicEconv_ae}
\leanok
\uses{def:torus-econv, lem:torus-econv-young}
For $\Phi \in L_1(\R^d)$, $f \in L_1(\T^d)$, $x \in \R^d$ and $t \in \T^d$,
$(\Phi * f)(t + \pi x) = (\Phi(\cdot + x) * f)(t)$, so
$\Delta^r_{\pi x}(\Phi * f) = (\Delta^r_x \Phi) * f$ almost everywhere.
\end{lemma}

\begin{proof}
\leanok
Substitute $y \mapsto y + x$; the differences are finite sums of translates, and the integrals
converge almost everywhere (\Cref{lem:torus-econv-young}).
\end{proof}

\section{Moduli of smoothness in \texorpdfstring{$L_p$}{L\_p}}
\label{sec:torus-lp-moduli}

In this section $1 \leq p < \infty$, $r \geq 1$, and $\omega_r(f, t)_p$ is the modulus of order $r$
of \Cref{def:rd-modulus} for the Haar measure on $\T^d$.

\begin{lemma}[Jackson's approximant on the torus]
\label{lem:torus-lp-jackson-approx}
\lean{Torus.periodicEconv_dilate_apply, Torus.periodic_ae_integrable_dilate_mul,
  Torus.neg_one_pow_mul_higherDiff_quotientProj, Torus.periodicJacksonApprox,
  Torus.periodicJacksonApprox_mem_trigPolyLp, Torus.ae_sub_periodicJacksonApprox_eq}
\leanok
\uses{thm:torus-econv-fourier, def:lp-kernels, lem:fourier-dilate}
With the kernel $\Phi$ of \Cref{def:lp-kernels} and $\Phi_c(x) = c^d \Phi(c x)$, the approximant
$J_k f = \sum_{l < r} (-1)^l \binom{r}{l+1} \Phi_{2^k/(l+1)} * f$ lies in $\mathcal{T}_{2^{k+1}}$,
and for almost every $t$
\[
  f(t) - J_k f(t) = \int_{\R^d} \Phi(z) (-1)^r \Delta^r_{\pi(-2^{-k} z)} f(t) \, dz .
\]
\end{lemma}

\begin{proof}
\leanok
$\mathcal{F}\Phi_c(\xi) = \mathcal{F}\Phi(\xi / c)$ vanishes for $|\xi| \geq 2c$, and
$|m| \geq \deg m$ (\Cref{lem:torus-proj}), so \Cref{thm:torus-econv-fourier} applies. The
identity is the algebra of \Cref{thm:band-jackson} with $\int \Phi = 1$ and
$(l+1) \, \pi(-c^{-1} z) = -\pi((c/(l+1))^{-1} z)$.
\end{proof}

\begin{theorem}[Jackson's inequality in $L_p(\T^d)$]
\label{thm:torus-lp-jackson}
\lean{Torus.enorm_sub_periodicJacksonApprox_le, Torus.approxError_trigScheme_le,
  Torus.lpModulus_le_two_pow_mul_succ}
\leanok
\uses{lem:torus-lp-jackson-approx, lem:modulus-linear, lem:jackson-weighted-average}
For $f \in L_p(\T^d)$ and $k \in \N$, with the constant $C_J$ of
\Cref{lem:jackson-weighted-average} on $\R^d$,
\[
  E_{2^{k+1}}(f)_p \leq \|f - J_k f\|_p \leq C_J \, \omega_r(f, 2^{-k})_p
  \leq 2^r C_J \, \omega_r(f, 2^{-(k+1)})_p .
\]
\end{theorem}

\begin{proof}
\leanok
$\|\Delta^r_{\pi(-2^{-k} z)} f\|_p \leq \omega_r(f, 2^{-k} |z|)_p \leq
(1 + |z|)^r \omega_r(f, 2^{-k})_p$ (\Cref{lem:modulus-linear}), and
\Cref{lem:jackson-weighted-average} with $M = \omega_r(f, 2^{-k})_p$ averages with the weight
$|\Phi|$.
\end{proof}

\begin{theorem}[Bernstein's inequality in $L_p(\T^d)$]
\label{thm:torus-lp-bernstein}
\lean{Torus.periodicBernsteinConst, Torus.periodicBernsteinConst_ne_top,
  Torus.periodicEconvLp_dilatedBandKernel_eq_self, Torus.lpModulus_le_of_mem_trigPolyLp}
\leanok
\uses{thm:torus-econv-fourier, lem:torus-econv-differences, lem:lp-kernels-frequency,
  thm:band-inverse}
Let $a \in \mathcal{T}_{2^j}$ and $k \in \N$. Then
\[
  \omega_r(a, 2^{-k})_p \leq C_B \, 2^{-(k-j)_+ r} \|a\|_p ,
  \qquad C_B = 2^r + (2^{\#d} \#d)^r \|D^r \Theta\|_1 .
\]
\end{theorem}

\begin{proof}
\leanok
With $J = j + \#d$, $\mathcal{F}\Theta_J(m) = 1$ for $\deg m < 2^j$, since
$|m| \leq \#d \deg m < 2^J$ (\Cref{lem:lp-kernels-frequency}); so $\Theta_J * a = a$
(\Cref{thm:torus-econv-fourier}). For $k \geq j$ and $\|h\| \leq 2^{-k}$ write $h = \pi x$ with
$|x| \leq \#d \, \|h\|$. Then $\Delta^r_h a = (\Delta^r_x \Theta_J) * a$
(\Cref{lem:torus-econv-differences}) and
$\|\Delta^r_x \Theta_J\|_1 \leq (2^J |x|)^r \|D^r \Theta\|_1$ (\Cref{thm:band-inverse}). For
$k < j$, $\omega_r(a, t)_p \leq 2^r \|a\|_p$.
\end{proof}

\begin{theorem}[Besov spaces on the torus through moduli of smoothness]
\label{thm:torus-lp-moduli}
\lean{Torus.trigLpModulusData, Torus.approxNormDyadic_le_modNorm_lp,
  Torus.modNorm_le_approxNormDyadic_lp, Torus.memBesov_iff_modNorm_lt_top_lp}
\leanok
\uses{thm:torus-lp-jackson, thm:torus-lp-bernstein, thm:modulus-coincidence, def:modulus-data}
Let $1 \leq p < \infty$, $0 < s < r$ and $1 \leq q \leq \infty$. Then
$f \in \Besov{s}{p}{q}(\T^d)$ if and only if
\[
  \|f\|_p + \bigl\| \bigl(2^{ks} \omega_r(f, 2^{-k})_p\bigr)_{k \in \N} \bigr\|_{\ell_q} < \infty ,
\]
with equivalent quasi-norms.
\end{theorem}

\begin{proof}
\leanok
The modulus $\omega_r(f, 2^{-k})_p$ is compatible with the trigonometric polynomials in
$L_p(\T^d)$ with inverse order $\theta = r$ (\Cref{def:modulus-data}): $C_0 = 2^r$, Jackson's
inequality from $k_0 = 1$ on with the constant $2^r C_J$ (\Cref{thm:torus-lp-jackson}), and
$C_B$ of \Cref{thm:torus-lp-bernstein}. \Cref{thm:modulus-coincidence} gives the equivalence.
\end{proof}

\section{Besov spaces on the torus as interpolation spaces}
\label{sec:torus-besov-interpolation}

The periodic Besov spaces interpolate as the Besov spaces of every scheme
(\Cref{thm:besov-interpolate}). The moduli of smoothness of \Cref{thm:torus-lp-moduli} give the
same for the quasi-norms by differences.

\begin{theorem}[Besov spaces on the torus by differences interpolate]
\label{thm:torus-besov-interpolate}
\lean{Torus.periodic_memKMethod_modCouple_iff, Torus.periodic_memKMethod_leftModCouple_iff}
\leanok
\uses{thm:torus-lp-moduli, def:modulus-couples, cor:modulus-spaces-interpolate}
Let $1 \leq p < \infty$, let $r \geq 1$ be an integer, let $0 < \eta < 1$,
$q_0, q_1, q \in [1, \infty]$ and $f \in L_p(\T^d)$, and write
$N_{s, q}(f) = \|f\|_p + \|(2^{ks} \omega_r(f, 2^{-k})_p)_k\|_{\ell_q}$.
\begin{enumerate}
  \item For every $0 < s_0 < s_1 < r$, $f$ lies in the space of the continuous $K$-method with
    $(\eta, q)$ of the couple $(N_{s_0, q_0}, N_{s_1, q_1})$ if and only if
    $N_{(1-\eta)s_0 + \eta s_1, q}(f) < \infty$.
  \item For every $0 < s_1 < r$, $f$ lies in the space of the continuous $K$-method with
    $(\eta, q)$ of the couple $(\|\cdot\|_p, N_{s_1, q_1})$ if and only if
    $N_{\eta s_1, q}(f) < \infty$.
\end{enumerate}
\end{theorem}

\begin{proof}
\leanok
$N_{s, q}$ is the classical quasi-norm of the modulus $\omega_r(f, 2^{-k})_p$, which is
compatible with the trigonometric polynomials with inverse order $r$
(\Cref{thm:torus-lp-moduli}). Apply \Cref{cor:modulus-spaces-interpolate}.
\end{proof}

\section{Littlewood--Paley blocks on the torus}
\label{sec:torus-lp-blocks}

In this section $1 \leq p < \infty$, and the kernels of \Cref{def:lp-kernels} on $\R^d$ act on
$L_p(\T^d)$ by convolution (\Cref{def:torus-econv}). Every identity between these operators is
checked on the characters $e_m$, where it becomes an identity of Fourier multipliers
(\Cref{thm:torus-econv-fourier}), and density of the characters carries it to all of
$L_p(\T^d)$.

\begin{lemma}[Identities of convolution operators on the torus]
\label{lem:torus-econv-identities}
\lean{Torus.ext_mFourierLp, Torus.periodicEconvLp_comp_periodicEconvLp,
  Torus.periodicEconvLp_comp_periodicEconvLp_of_eq, Torus.periodicEconvLp_comm,
  Torus.periodicEconvLp_add_kernel, Torus.periodicEconvLp_sub_kernel,
  Torus.periodicEconv_smul_kernel, Torus.periodicEconvLp_congr_kernel,
  Torus.periodicEconvLp_eq_zero_of_mem}
\leanok
\uses{def:torus-econv, thm:torus-econv-fourier, lem:torus-econv-young, lem:fourier-conv}
For every normed space $F$, two bounded linear maps $L_p(\T^d) \to F$ that agree on every $e_m$,
$m \in \Z^d$, are equal. For all
$K, L \in L_1(\R^d)$ and every $f \in L_p(\T^d)$,
\[
  K * (L * f) = (K * L) * f = L * (K * f), \qquad (K \pm L) * f = K * f \pm L * f .
\]
For every $c \in \C$, every function $K$ on $\R^d$ and every function $g$ on $\T^d$,
$(c K) * g = K * (c g)$ at every point of $\T^d$. For every $n \in \N$ and every
$\Phi \in L_1(\R^d)$ with $\mathcal{F}\Phi(m) = 0$ whenever $\deg m < n$, $\Phi * a = 0$ for
every $a \in \mathcal{T}_n$.
\end{lemma}

\begin{proof}
\leanok
The characters span a dense subspace of $L_p(\T^d)$, so a bounded linear map is determined by
its values on them. All operators in the identities are bounded (\Cref{lem:torus-econv-young}) and
multiply $e_m$ by $\mathcal{F}K(m) \, \mathcal{F}L(m) = \mathcal{F}(K * L)(m)$, respectively by
$\mathcal{F}K(m) \pm \mathcal{F}L(m) = \mathcal{F}(K \pm L)(m)$ (\Cref{thm:torus-econv-fourier},
\Cref{lem:fourier-conv}). The identity $(c K) * g = K * (c g)$ is
$(c K(y)) \, g(t - \pi y) = K(y) \, (c g(t - \pi y))$ under the integral of
\Cref{def:torus-econv}. The last claim holds on the characters $e_m$ with
$\deg m < n$, which span $\mathcal{T}_n$.
\end{proof}

The operator $f \mapsto K * f$ is determined by the kernel $K$: equal integrable kernels $K = L$
on $\R^d$ give the same operator on $L_p(\T^d)$.

\begin{theorem}[Approximate identities on the torus]
\label{thm:torus-approx-identity}
\lean{Torus.norm_periodicEconvLp_le, Torus.tendsto_periodicEconvLp_of_fourier,
  Torus.tendsto_periodicEconvLp_dilate, Torus.tendsto_periodicEconvLp_dilatedLpKernel}
\leanok
\uses{lem:torus-econv-young, thm:torus-econv-fourier, lem:fourier-conv, lem:fourier-dilate,
  lem:dilate, lem:lp-kernels-basic}
Let $(\Phi_i)_{i \in I}$ be integrable kernels on $\R^d$, $\ell$ a filter on $I$ and
$C < \infty$ with $\|\Phi_i\|_1 \leq C$ for every $i \in I$ and $\mathcal{F}\Phi_i(m) \to 1$ along
$\ell$ for every $m \in \Z^d$. Then $\Phi_i * f \to f$ in $L_p(\T^d)$ along $\ell$ for every
$f \in L_p(\T^d)$. In particular, for every $f \in L_p(\T^d)$:
\begin{enumerate}
  \item for every $K \in L_1(\R^d)$ with $\int K = 1$ and all $c_i > 0$, $i \in I$, with
    $c_i \to \infty$ along $\ell$, $K_{c_i} * f \to f$ along $\ell$;
  \item $\Phi_J * f \to f$ as $J \to \infty$.
\end{enumerate}
\end{theorem}

\begin{proof}
\leanok
The $f$ with $\Phi_i * f \to f$ form a linear subspace. It contains every $e_m$, since
$\Phi_i * e_m = \mathcal{F}\Phi_i(m) \, e_m$ (\Cref{thm:torus-econv-fourier}). It is closed: the
operators are bounded by $C$ (\Cref{lem:torus-econv-young}), hence equicontinuous, and the set
where an equicontinuous family converges to the identity is closed. So it contains the closure
of the span of the characters, which is $L_p(\T^d)$. For (1), $\mathcal{F}K_c(m) =
\mathcal{F}K(m / c)$ (\Cref{lem:fourier-dilate}) tends to $\mathcal{F}K(0) = \int K = 1$ by the
continuity of $\mathcal{F}K$ (\Cref{lem:fourier-conv}), and $\|K_c\|_1 = \|K\|_1$
(\Cref{lem:dilate}). (2) is (1) with $K = \Phi$, $\ell$ the limit $J \to \infty$ and
$c_J = 2^J$ (\Cref{lem:lp-kernels-basic}).
\end{proof}

\begin{definition}[Littlewood--Paley blocks on the torus]
\label{def:torus-lp-blocks}
\lean{Torus.periodicLpBlock, Torus.periodicLittlewoodPaleyNorm}
\leanok
\uses{def:lp-kernels, lem:torus-econv-young, def:lq-norm}
Let $s \in \R$ and $1 \leq q \leq \infty$. The Littlewood--Paley blocks of $f \in L_p(\T^d)$ are
$\Delta_j f = \psi_j * f$, $j \in \N$, with the kernels $\psi_j$ of \Cref{def:lp-kernels} on
$\R^d$, and the Littlewood--Paley quasi-norm is
$\|f\|^{\mathrm{LP}}_{s,q} = \|(2^{js} \|\Delta_j f\|_p)_j\|_{\ell_q}$.
\end{definition}

\begin{lemma}[Bounds and reconstruction on the torus]
\label{lem:torus-lp-blocks}
\lean{Torus.enorm_periodicLpBlock_le, Torus.periodicLpBlock_succ, Torus.sum_periodicLpBlock,
  Torus.tendsto_sum_periodicLpBlock, Torus.hasSum_periodicLpBlock}
\leanok
\uses{def:torus-lp-blocks, lem:torus-econv-young, lem:torus-econv-identities,
  thm:torus-approx-identity, lem:lp-kernels-basic, lem:weighted-blocks}
Let $f \in L_p(\T^d)$. For all $j, J \in \N$, $\|\Delta_j f\|_p \leq 2 \|\Phi\|_1 \|f\|_p$ and
$\sum_{j=0}^J \Delta_j f = \Phi_J * f$, and $\sum_{j=0}^J \Delta_j f \to f$ in $L_p(\T^d)$ as
$J \to \infty$. If $s > 0$, $1 \leq q \leq \infty$ and $\|f\|^{\mathrm{LP}}_{s,q} < \infty$,
then $\sum_j \Delta_j f = f$, the series converging in $L_p(\T^d)$.
\end{lemma}

\begin{proof}
\leanok
The bound is \Cref{lem:torus-econv-young} with $\|\psi_j\|_1 \leq 2 \|\Phi\|_1$
(\Cref{lem:lp-kernels-basic}). The sum telescopes, since
$\psi_{j+1} * f = \Phi_{j+1} * f - \Phi_j * f$ (\Cref{lem:torus-econv-identities}), and
$\Phi_J * f \to f$ by \Cref{thm:torus-approx-identity}. For the last claim apply
\Cref{lem:weighted-blocks}(1).
\end{proof}

\begin{lemma}[Blocks and trigonometric polynomials]
\label{lem:torus-lp-blocks-trig}
\lean{Torus.periodicLpBlock_mem_trigPolyLp, Torus.periodicLpBlock_eq_zero_of_mem,
  Torus.enorm_periodicLpBlock_le_approxError}
\leanok
\uses{lem:torus-lp-blocks, thm:torus-econv-fourier, lem:torus-econv-identities,
  lem:lp-kernels-frequency, lem:torus-proj, lem:error-basic}
Let $f \in L_p(\T^d)$ and $i, j \in \N$. Then $\Delta_j f \in \mathcal{T}_{2^{j+1}}$. If
$j \geq i + \#d + 1$, then $\Delta_j a = 0$ for every $a \in \mathcal{T}_{2^i}$, and
\[
  \|\Delta_j f\|_p \leq 2 \|\Phi\|_1 \, E_{2^i}(f)_p .
\]
\end{lemma}

\begin{proof}
\leanok
If $\mathcal{F}\psi_j(m) \neq 0$, then $\deg m \leq |m| < 2^{j+1}$
(\Cref{lem:lp-kernels-frequency} and \Cref{lem:torus-proj}), so \Cref{thm:torus-econv-fourier} gives
the first claim. For $j = k + 1$ with $k \geq i + \#d$ and $\deg m < 2^i$,
$|m| \leq \#d \deg m < 2^{\#d} 2^i \leq 2^k$, hence $\mathcal{F}\psi_j(m) = 0$
(\Cref{lem:lp-kernels-frequency}) and $\Delta_j a = 0$ (\Cref{lem:torus-econv-identities}).
Hence $\Delta_j f = \Delta_j (f - a)$ for $a \in \mathcal{T}_{2^i}$, and \Cref{lem:torus-lp-blocks}
bounds it by $2 \|\Phi\|_1 \|f - a\|_p$.
\end{proof}

\begin{theorem}[Besov spaces on the torus by Littlewood--Paley blocks]
\label{thm:torus-lp-besov}
\lean{Torus.approxNormDyadic_trigScheme_le,
  Torus.periodicLittlewoodPaleyNorm_le_approxNormDyadic_trigScheme,
  Torus.memBesov_iff_periodicLittlewoodPaleyNorm_lt_top}
\leanok
\uses{lem:torus-lp-blocks, lem:torus-lp-blocks-trig, thm:block-decomposition, def:dyadic-norm,
  thm:dyadic-norm, def:besov, lem:shifted-blocks, lem:weighted-blocks}
Let $1 \leq q \leq \infty$, let $A$ be the scheme of trigonometric polynomials in $L_p(\T^d)$,
and let $f \in L_p(\T^d)$. For every $s > 0$,
\[
  \|f\|^{\mathrm{dy}}_{A^s_q} \leq 2 \cdot 2^s (1 - 2^{-s})^{-1} \, \|f\|^{\mathrm{LP}}_{s,q} ,
\]
and for every $s \geq 0$,
\[
  \|f\|^{\mathrm{LP}}_{s,q} \leq 2 (\#d + 1) \, 2^{(\#d + 1) s} \, \|\Phi\|_1 \,
  \|f\|^{\mathrm{dy}}_{A^s_q} .
\]
In particular, for every $s > 0$, $f \in \Besov{s}{p}{q}(\T^d)$ if and only if
$\|f\|^{\mathrm{LP}}_{s,q} < \infty$.
\end{theorem}

\begin{proof}
\leanok
Put $g_0 = 0$ and $g_{k+1} = \Delta_k f$, the blocks moved up by one level. Then
$g_k \in \mathcal{T}_{2^k}$ by \Cref{lem:torus-lp-blocks-trig}, a finite Littlewood--Paley
quasi-norm makes $f = \sum_k g_k$ a block decomposition (\Cref{lem:torus-lp-blocks}), and
$\|(2^{ks} \|g_k\|_p)_k\|_{\ell_q} \leq 2^s \|f\|^{\mathrm{LP}}_{s,q}$ by
\Cref{lem:shifted-blocks}; \Cref{thm:block-decomposition} gives the first inequality. For the
second, apply \Cref{lem:weighted-blocks}(2) with $M = \#d + 1$ and $C = 2\|\Phi\|_1$
(\Cref{lem:torus-lp-blocks}, \Cref{lem:torus-lp-blocks-trig}).
\end{proof}

For $p = 2$ the sharp dyadic blocks of the Fourier coefficients describe the same spaces
(\Cref{thm:torus-coefficients}); the smooth blocks work for every $1 \leq p < \infty$.

\section{Weak derivatives}
\label{sec:torus-weak-derivatives}

For a multi-index $\alpha \in \N^d$ the derivative $\partial^\alpha$ multiplies the Fourier
coefficients by $(2\pi i k)^\alpha$. Defining weak derivatives through this rule is the
distributional definition $\int g \, \varphi = (-1)^{|\alpha|} \int f \, \partial^\alpha
\varphi$, tested against the monomials $\varphi = \overline{e_k}$, whose span is dense in
the smooth functions. For $1 \leq p \leq \infty$ the weak derivatives are the distributional
derivatives of the periodic lifts on $\R^d$ (\Cref{thm:torus-weak-derivative-lift}).

\begin{definition}[Weak derivatives and $W^m_p(\T^d)$]
\label{def:weak-derivative}
\lean{Torus.periodicDerivSymbol, Torus.HasPeriodicWeakDeriv, Torus.MemPeriodicSobolev,
  Torus.norm_two_pi_I_mul_intCast}
\leanok
For $\alpha \in \N^d$ put $|\alpha| = \sum_j \alpha_j$ and $(2\pi i k)^\alpha = \prod_j
(2\pi i k_j)^{\alpha_j}$ for $k \in \Z^d$. Let $1 \leq p \leq \infty$. An element
$g \in L_p(\T^d)$ is the \emph{weak derivative} $\partial^\alpha f$ of $f \in L_p(\T^d)$ if
$\hat g(k) = (2\pi i k)^\alpha \hat f(k)$ for every $k \in \Z^d$. For $m \in \N$, the Sobolev
space $\Sob{m}{p}(\T^d)$ consists of the $f \in L_p(\T^d)$ that have a weak derivative
$\partial^\alpha f \in L_p(\T^d)$ for every $\alpha$ with $|\alpha| \leq m$.
\end{definition}

The remaining results of this section are for $p = 2$; \Cref{ch:sobolev-torus} treats weak
derivatives along directions in $L_p(\T^d)$ for every $1 \leq p \leq \infty$.

\begin{lemma}[Existence of weak derivatives]
\label{lem:weak-derivative}
\lean{Torus.hasPeriodicWeakDeriv_zero, Torus.exists_hasPeriodicWeakDeriv,
  Torus.HasPeriodicWeakDeriv.mem_l2}
\leanok
\uses{def:weak-derivative}
For every $1 \leq p \leq \infty$ and every $f \in L_p(\T^d)$, $\partial^0 f = f$. Let
$f \in L_2(\T^d)$ and $\alpha \in \N^d$. The weak derivative $\partial^\alpha f$ exists in
$L_2(\T^d)$ if and only if $\sum_k |(2\pi i k)^\alpha \hat f(k)|^2 < \infty$.
\end{lemma}

\begin{proof}
\leanok
Every square-summable family of coefficients is the family of some element of $L_2(\T^d)$,
since the trigonometric monomials form a Hilbert basis.
\end{proof}

\begin{lemma}[Symbols and the Sobolev weight]
\label{lem:symbol-bounds}
\lean{Torus.norm_periodicDerivSymbol, Torus.norm_periodicDerivSymbol_le,
  Torus.one_add_trigDeg_pow_le, Torus.periodicDerivSymbol_zero, Torus.periodicDerivSymbol_single,
  Torus.one_le_two_pi, Torus.abs_le_one_add_trigDeg}
\leanok
\uses{def:weak-derivative, def:trig-poly}
Let $m \in \N$ and $k \in \Z^d$. For every $\alpha \in \N^d$ with $|\alpha| \leq m$,
\[
  |(2\pi i k)^\alpha| \leq (2\pi)^m (1 + \deg k)^m ,
\]
and, writing $m e_j$ for the multi-index with $m$ in coordinate $j$ and $0$ elsewhere,
\[
  (1 + \deg k)^m \leq 2^m \Bigl(1 + \sum_{j \in d} |(2\pi i k)^{m e_j}|\Bigr) .
\]
\end{lemma}

\begin{proof}
\leanok
For the first bound, $|k_j| \leq 1 + \deg k$ for every $j$ and $2\pi (1 + \deg k) \geq 1$.
For the second, $(1 + x)^m \leq 2^m (1 + x^m)$ for $x \geq 0$; if $\deg k > 0$, then
$\deg k = |k_j|$ for some $j$, and $|k_j|^m \leq (2\pi |k_j|)^m$.
\end{proof}

\begin{theorem}[$\Sob{m}{2}(\T^d) = H^m(\T^d) = \Besov{m}{2}{2}(\T^d)$]
\label{thm:weak-sobolev}
\lean{Torus.memPeriodicSobolev_iff, Torus.memPeriodicSobolev_iff_memBesov,
  Torus.memPeriodicSobolev_iff_mem_range}
\leanok
\uses{lem:weak-derivative, lem:symbol-bounds, thm:torus-coefficients, def:torus-sobolev}
Let $m \in \N$ and $f \in L_2(\T^d)$. Then $f \in \Sob{m}{2}(\T^d)$ if and only if
$\sum_k (1 + \deg k)^{2m} |\hat f(k)|^2 < \infty$. For $m \geq 1$ this holds if and only if
$f \in \Besov{m}{2}{2}(\T^d)$, and if and only if $f \in H^m(\T^d)$.
\end{theorem}

\begin{proof}
\leanok
If $f \in \Sob{m}{2}$, the coefficient families of $f$ and of $\partial_j^m f$ are
square-summable, and by the second bound of \Cref{lem:symbol-bounds} they dominate
$((1 + \deg k)^m \hat f(k))_k$. Conversely, the first bound shows that
$((2\pi i k)^\alpha \hat f(k))_k$ is square-summable for $|\alpha| \leq m$, and
\Cref{lem:weak-derivative} gives $\partial^\alpha f$. The rest is
\Cref{thm:torus-coefficients} and \Cref{def:torus-sobolev}.
\end{proof}

Classical derivatives of functions with continuously differentiable periodic lifts are weak
derivatives (\Cref{thm:torus-classical-derivative}).

\chapter{Filtrations and martingale differences}
\label{ch:filtrations}

Throughout, $(\Omega, \mathcal{F}, \mu)$ is a finite measure space and $1 \leq p \leq \infty$.
Conditional expectations are projections of norm at most $1$ on every $L_p$, so a filtration
gives a projection scheme.

\begin{lemma}[Conditional expectation on $L_p$]
\label{lem:condexp-lp}
\lean{MeasureTheory.Lp.condExpLp, MeasureTheory.Lp.coeFn_condExpLp,
  MeasureTheory.Lp.norm_condExpLp_le, MeasureTheory.Lp.condExpLp_mem_lpMeas,
  MeasureTheory.Lp.condExpLp_of_mem, MeasureTheory.Lp.condExpLpLinear,
  MeasureTheory.Lp.coeFn_condExpLpLinear}
\leanok
Let $\mathcal{G} \subseteq \mathcal{F}$ be a sub-$\sigma$-algebra. The conditional
expectation $f \mapsto E[f \mid \mathcal{G}]$ is a bounded linear map $L_p \to L_p$ of norm at
most $1$, its values are $\mathcal{G}$-measurable, and it fixes every $\mathcal{G}$-measurable
element of $L_p$.
\end{lemma}

\begin{proof}
\leanok
The bound is the conditional Jensen inequality $\|E[f \mid \mathcal{G}]\|_p \leq \|f\|_p$,
which Mathlib provides; linearity holds almost everywhere.
\end{proof}

\begin{definition}[Filtration scheme]
\label{def:filtration-scheme}
\lean{MeasureTheory.Filtration.filtrationScheme, MeasureTheory.Filtration.mem_filtrationScheme,
  MeasureTheory.Filtration.filtrationProjectionScheme}
\leanok
\uses{lem:condexp-lp, def:projection-scheme, def:comap}
Let $\mathcal{F}_0 \subseteq \mathcal{F}_1 \subseteq \cdots \subseteq \mathcal{F}$ be a
filtration. Let $A_0 = \{0\}$ and, for $n \geq 1$, let $A_n$ be the set of
$\mathcal{F}_{\lfloor \log_2 n \rfloor}$-measurable elements of $L_p$. These sets form an
approximation scheme, the \emph{filtration scheme}, and $P_0 = 0$,
$P_n = E[\,\cdot \mid \mathcal{F}_{\lfloor \log_2 n \rfloor}]$ form a projection scheme for it
with bound $1$, with $P_{2^k} = E[\,\cdot \mid \mathcal{F}_k]$.
\end{definition}

For the dyadic filtration of $[0, 1]$ with Lebesgue measure, $\mathcal{F}_j$ is generated by
the dyadic intervals of
length $2^{-j}$, and the approximants of complexity $2^j$ are the step functions on these
$2^j$ intervals: the scheme of the Haar system.

\begin{theorem}[Martingale differences]
\label{thm:martingale}
\lean{MeasureTheory.Filtration.memApprox_iff_martingale}
\leanok
\uses{def:filtration-scheme, thm:blocks}
Let $\mathcal{F}_0 \subseteq \mathcal{F}_1 \subseteq \cdots$ be a filtration whose scheme
has $\bigcup_n A_n$ dense in $L_p$, and let $s > 0$ and $1 \leq q \leq \infty$. Then
$f \in A^s_q$ if and only if
\[
  \|E[f \mid \mathcal{F}_0]\|_p + \Bigl\| \Bigl(2^{ks} \bigl\|E[f \mid \mathcal{F}_{k+1}] -
  E[f \mid \mathcal{F}_k]\bigr\|_p\Bigr)_k \Bigr\|_{\ell_q} < \infty .
\]
\end{theorem}

\begin{proof}
\leanok
\Cref{thm:blocks} with $P_{2^k} = E[\,\cdot \mid \mathcal{F}_k]$: the dyadic blocks are the
martingale differences.
\end{proof}

For the dyadic filtration of $[0, 1]$ with Lebesgue measure the density holds for $p < \infty$,
and
$E[f \mid \mathcal{F}_{k+1}] - E[f \mid \mathcal{F}_k]$ is the part of the Haar expansion of
$f$ at level $k$ (\Cref{sec:haar}).

\section{The Haar system}
\label{sec:haar}

On $[0, 1]$ with Lebesgue measure the dyadic filtration makes \Cref{thm:martingale} explicit.
The conditional expectations are averages over dyadic intervals, which follows from a general
formula for partitions.

\begin{lemma}[Conditional expectation onto a partition]
\label{lem:condexp-partition}
\lean{MeasureTheory.comap_top_le, MeasureTheory.StronglyMeasurable.eq_of_comap_top,
  MeasureTheory.condExp_comap_ae_eq}
\leanok
Let $\mu$ be a finite measure on $\Omega$, let $\varphi : \Omega \to I$ with $I$ countable and
every fibre $\varphi^{-1}(i)$ measurable, and let $\mathcal{G}$ be the $\sigma$-algebra of
the sets $\varphi^{-1}(T)$, $T \subseteq I$. For every integrable $f$ and almost every
$x \in \Omega$,
\[
  E[f \mid \mathcal{G}](x) = \mu(C_x)^{-1} \int_{C_x} f \, d\mu ,
  \qquad C_x = \varphi^{-1}(\varphi(x)) .
\]
\end{lemma}

\begin{proof}
\leanok
A $\mathcal{G}$-measurable function is constant on every fibre. On a fibre $C$ of positive
measure, $\int_C E[f \mid \mathcal{G}] \, d\mu = \int_C f \, d\mu$ determines the constant.
The fibres of measure $0$ are countably many, so their union is a null set.
\end{proof}

\begin{definition}[Dyadic intervals and the dyadic filtration]
\label{def:dyadic-filtration}
\lean{Haar.vol, Haar.dyadicIndex, Haar.haarDyadicInterval, Haar.haarDyadicInterval_eq_Ico,
  Haar.haarDyadicInterval_eq_union, Haar.vol_haarDyadicInterval, Haar.dyadicSigma,
  Haar.dyadicFiltration, Haar.haarScheme, Haar.mem_haarDyadicInterval, Haar.measurable_dyadicIndex,
  Haar.measurableSet_haarDyadicInterval, Haar.disjoint_haarDyadicInterval, Haar.dyadicIndex_nonneg,
  Haar.dyadicIndex_lt, Haar.haarDyadicInterval_natCast, Haar.haarDyadicInterval_subset,
  Haar.vol_real_haarDyadicInterval,
  Haar.dyadicIndex_succ_div, Haar.dyadicSigma_mono}
\leanok
\uses{def:filtration-scheme, def:pp-scheme}
For $j \in \N$ and $n \in \Z$ let $I_{j,n} = [n 2^{-j}, (n + 1) 2^{-j})$, the set of $x \in \R$
with $\lfloor 2^j x \rfloor = n$. Let $\mathcal{F}_j$ be the $\sigma$-algebra generated by the
intervals $I_{j,n}$, $n \in \Z$. Then $I_{j,n} = I_{j+1,2n} \cup I_{j+1,2n+1}$, so
$\mathcal{F}_0 \subseteq \mathcal{F}_1 \subseteq \cdots$ is a filtration, and
$|I_{j,n}| = 2^{-j}$. For $n \in \N$ these intervals reappear as the dyadic intervals of the
piecewise polynomials (\Cref{def:pp-scheme}). The \emph{Haar scheme} in $L_p[0, 1]$ is its
filtration scheme.
\end{definition}

\begin{lemma}[Dyadic conditional expectations]
\label{lem:dyadic-condexp}
\lean{Haar.condExp_dyadic_ae_eq, Haar.condExpLp_dyadic_ae_eq, Haar.condExp_zero_ae_eq,
  Haar.condExpLp_zero_ae_eq, Haar.enorm_condExpLp_zero}
\leanok
\uses{lem:condexp-partition, def:dyadic-filtration}
For $f \in L_1[0, 1]$, $j \in \N$ and almost every $x \in [0, 1]$,
\[
  E[f \mid \mathcal{F}_j](x) = 2^j \int_{I_{j,n}} f , \qquad n = \lfloor 2^j x \rfloor .
\]
In particular $E[f \mid \mathcal{F}_0] = \int_0^1 f$.
\end{lemma}

\begin{proof}
\leanok
\Cref{lem:condexp-partition} with $\varphi(x) = \lfloor 2^j x \rfloor$; on $[0, 1)$ the
relevant fibres are the intervals $I_{j,n}$, $0 \leq n < 2^j$, of measure $2^{-j}$.
\end{proof}

\begin{lemma}[The normalising factor $2^{j/2}$]
\label{lem:dyadic-normalising-factor}
\lean{ofReal_two_rpow_half_mul_self, enorm_two_rpow_half, ofReal_two_rpow_half_pow_mul_self,
  enorm_two_rpow_half_pow}
\leanok
\uses{def:lq-norm}
For every $j \in \N$, $2^{j/2} \cdot 2^{j/2} = 2^j$, and the modulus of $2^{j/2}$ as a complex
number is the dyadic weight $2^{j/2}$. For every $d \in \N$, $(2^{j/2})^d (2^{j/2})^d = 2^{jd}$,
and the modulus of $(2^{j/2})^d$ is $2^{jd/2}$.
\end{lemma}

\begin{proof}
\leanok
$2^{j/2} \cdot 2^{j/2} = 2^{j/2 + j/2} = 2^j$, and $2^{j/2} > 0$. Take $d$-th powers.
\end{proof}

The Haar functions, the multiwavelets of \Cref{sec:multiwavelets} and the Daubechies wavelets
carry this factor, which makes the dilation $g \mapsto 2^{j/2} g(2^j \cdot)$ an isometry of
$L_2(\R)$.

\begin{definition}[Haar functions and coefficients]
\label{def:haar}
\lean{Haar.haarFunction, Haar.measurable_haarFunction, Haar.haarFunction_eq_zero,
  Haar.abs_haarFunction, Haar.abs_haarFunction_le, Haar.enorm_haarFunction, Haar.haarFunction_mul,
  Haar.haarCoeff, Haar.haarCoeff_eq, Haar.haarCoeff_congr_ae}
\leanok
\uses{def:dyadic-filtration}
For $j \in \N$ and $m \in \Z$ the Haar function is
$h_{j,m} = 2^{j/2} (\mathbf{1}_{I_{j+1,2m}} - \mathbf{1}_{I_{j+1,2m+1}})$. It vanishes outside
$I_{j,m}$ and has modulus $2^{j/2}$ on $I_{j,m}$, so $|h_{j,m}| \leq 2^{j/2}$. The Haar
coefficient of $f \in L_1[0, 1]$ is
\[
  \langle f, h_{j,m} \rangle = \int_0^1 h_{j,m} f
  = 2^{j/2} \Bigl( \int_{I_{j+1,2m}} f - \int_{I_{j+1,2m+1}} f \Bigr) ,
\]
and functions that agree almost everywhere on $[0, 1]$ have the same Haar coefficients.
\end{definition}

\begin{lemma}[Martingale differences are Haar sums]
\label{lem:haar-martingale}
\lean{Haar.sum_fin_eq_of_ne, Haar.setIntegral_haarDyadicInterval, Haar.condExp_succ_sub_ae_eq,
  Haar.condExpLp_succ_sub_ae_eq}
\leanok
\uses{lem:dyadic-condexp, def:haar, lem:dyadic-normalising-factor}
For $f \in L_1[0, 1]$ and $k \in \N$,
\[
  E[f \mid \mathcal{F}_{k+1}] - E[f \mid \mathcal{F}_k]
  = \sum_{m=0}^{2^k - 1} \langle f, h_{k,m} \rangle \, h_{k,m}
\]
almost everywhere.
\end{lemma}

\begin{proof}
\leanok
Let $x \in I_{k,m}$ with $0 \leq m < 2^k$, and let $a$ and $b$ be the integrals of $f$ over
$I_{k+1,2m}$ and $I_{k+1,2m+1}$. By \Cref{lem:dyadic-condexp} the left side at $x$ is
$2^{k+1} a - 2^k (a + b) = 2^k (a - b)$ on $I_{k+1,2m}$ and $-2^k (a - b)$ on $I_{k+1,2m+1}$.
On the right only the term $m$ is nonzero at $x$, and it has the same value since
$\langle f, h_{k,m} \rangle = 2^{k/2} (a - b)$.
\end{proof}

\begin{lemma}[Haar sums of one level]
\label{lem:haar-level}
\lean{Haar.eLpNorm_sum_haar}
\leanok
\uses{def:haar}
Let $0 < p < \infty$, $j \in \N$ and $c_0, \dots, c_{2^j - 1} \in \C$. Then
\[
  \Bigl\| \sum_{m=0}^{2^j - 1} c_m h_{j,m} \Bigr\|_p = 2^{j(1/2 - 1/p)} \|c\|_{\ell_p} .
\]
\end{lemma}

\begin{proof}
\leanok
At every point of $I_{j,m}$ the sum is $c_m h_{j,m}$, of modulus $|c_m| 2^{j/2}$, so the
$p$-th power of the norm is $\sum_m |c_m|^p 2^{jp/2} 2^{-j}$.
\end{proof}

\begin{lemma}[Density of step functions]
\label{lem:haar-density}
\lean{Haar.stepFun, Haar.stronglyMeasurable_stepFun, Haar.exists_norm_sub_stepFun_le,
  Haar.toLp_mem_closure, Haar.closure_iUnion_haarScheme}
\leanok
\uses{def:dyadic-filtration}
For $1 \leq p < \infty$ the Haar approximants are dense in $L_p[0, 1]$.
\end{lemma}

\begin{proof}
\leanok
A bounded continuous $g$ is uniformly continuous on $[0, 1]$, so the step functions
$x \mapsto g(\lfloor 2^j x \rfloor 2^{-j})$, which are $\mathcal{F}_j$-measurable, converge to
$g$ uniformly on $[0, 1)$, hence in $L_p$. Bounded continuous functions are dense in $L_p$ for
$p < \infty$.
\end{proof}

\begin{theorem}[Haar coefficient characterisation]
\label{thm:haar}
\lean{Haar.memApprox_iff_haarCoeff}
\leanok
\uses{thm:martingale, lem:haar-martingale, lem:haar-level, lem:haar-density}
Let $1 \leq p < \infty$, $s > 0$ and $1 \leq q \leq \infty$. Then $f \in L_p[0, 1]$ lies in the
approximation space $A^s_q$ of the Haar scheme if and only if
\[
  \Bigl| \int_0^1 f \Bigr| + \Bigl\| \Bigl(2^{k(s + 1/2 - 1/p)}
  \bigl\| (\langle f, h_{k,m} \rangle)_{0 \leq m < 2^k} \bigr\|_{\ell_p} \Bigr)_k
  \Bigr\|_{\ell_q} < \infty .
\]
\end{theorem}

\begin{proof}
\leanok
\Cref{thm:martingale}, with \Cref{lem:haar-density} for the density, the formula for
$E[f \mid \mathcal{F}_0]$ from \Cref{lem:dyadic-condexp}, and
\Cref{lem:haar-martingale} and \Cref{lem:haar-level} for the martingale differences.
\end{proof}

For $0 < s < 1/p$ this approximation space is the classical Besov space
$\Besov{s}{p}{q}[0, 1]$; for larger $s$ it is strictly smaller, since the Haar functions are
discontinuous. The comparison with moduli of smoothness is \Cref{cor:haar-besov}.

\section{The Haar basis and \texorpdfstring{$n$}{n}-term approximation}
\label{sec:haar-nterm}

In breadth-first order the Haar system is $\varphi_0 = 1$ and $\varphi_{2^k + i} = h_{k,i}$ for
$k \in \N$ and $0 \leq i < 2^k$; every $n \geq 1$ is $2^k + i$ for exactly one such pair, with
$k = \lfloor \log_2 n \rfloor$. The Lean project TauCeti~\cite{TauCeti} orders its Haar basis of
$L_p[0, 1]$, $1 \leq p < \infty$, in the same way.

\begin{proposition}[The Haar basis]
\label{prop:haar-basis}
\lean{Haar.haarIndex, Haar.haarIndex_two_pow_add, Nat.log_two_pow_add, Nat.exists_eq_two_pow_add,
  Nat.two_pow_add_eq_iff, Haar.haarSystem, Haar.haarSystem_zero, Haar.haarSystem_two_pow_add,
  Haar.measurable_haarSystem, Haar.exists_abs_haarSystem_le, Haar.memLp_haarSystem, Haar.haarVec,
  Haar.coeFn_haarVec, Haar.inner_haarVec, Haar.dyadicIndex_eq_of_le,
  Haar.haarFunction_eq_of_dyadicIndex_eq, Haar.haarFunction_ae_eq_zero, Haar.integral_haarFunction,
  Haar.integral_haarFunction_mul_self, Haar.haarFunction_mul_haarFunction_of_ne,
  Haar.integral_haarFunction_mul_of_lt, Haar.integral_haarFunction_mul,
  Haar.integral_haarSystem_mul, Haar.orthonormal_haarVec, Haar.condExpLp_mem_span_haarVec,
  Haar.iUnion_haarScheme_subset_span, Haar.span_haarVec_closure_eq_top, Haar.haarBasis,
  Haar.coe_haarBasis, Haar.haarBasis_repr_apply, Haar.haarBasis_repr_zero,
  Haar.haarBasis_repr_two_pow_add}
\leanok
\uses{def:haar, lem:dyadic-condexp, lem:haar-martingale, lem:haar-density}
The functions $\varphi_n$, $n \in \N$, form a Hilbert basis of $L_2[0, 1]$. The coefficients of
$f \in L_2[0, 1]$ in this basis are $\langle \varphi_0, f \rangle = \int_0^1 f$ and
$\langle \varphi_{2^k + i}, f \rangle = \langle f, h_{k,i} \rangle$.
\end{proposition}

\begin{proof}
\leanok
Let $0 \leq m < 2^k$ and $0 \leq n < 2^j$. The function $h_{j,n}$ has mean zero, since
$I_{j+1,2n}$ and $I_{j+1,2n+1}$ have the same length, and $h_{j,n}^2 = 2^j$ on $I_{j,n}$ and $0$
elsewhere, so $\int_0^1 h_{j,n}^2 = 1$. For $k = j$ and $m \neq n$ the supports are disjoint. For
$k < j$ the function $h_{k,m}$ depends only on the dyadic interval of level $k + 1 \leq j$
containing the point, so it is constant on $I_{j,n}$, and $\int_0^1 h_{k,m} h_{j,n}$ is this
constant times $\int_0^1 h_{j,n} = 0$. Hence the $\varphi_n$ are orthonormal. By
\Cref{lem:dyadic-condexp} and \Cref{lem:haar-martingale},
\[
  E[f \mid \mathcal{F}_j] = \Bigl(\int_0^1 f\Bigr) \varphi_0
  + \sum_{k < j} \sum_{i=0}^{2^k - 1} \langle f, h_{k,i} \rangle \, \varphi_{2^k + i}
\]
lies in the span of the $\varphi_n$. Since $E[\,\cdot \mid \mathcal{F}_j]$ fixes the
$\mathcal{F}_j$-measurable functions, every Haar approximant lies in this span, and
\Cref{lem:haar-density} with $p = 2$ shows that the span is dense. The coefficients are
$\int_0^1 \varphi_n f$, since the $\varphi_n$ are real.
\end{proof}

\begin{theorem}[$n$-term approximation from the Haar basis]
\label{thm:haar-nterm}
\lean{Haar.memApprox_nTerm_iff_haarCoeff, Haar.tsum_haarBasis_repr_rpow,
  ENNReal.tsum_eq_add_tsum_two_pow_add}
\leanok
\uses{prop:haar-basis, thm:stechkin}
Let $0 < \tau < 2$ and $\theta = 1/\tau - 1/2$. Then $f \in L_2[0, 1]$ lies in the
approximation space $A^\theta_\tau$ of $n$-term approximation from the dictionary
$\{c \, \varphi_n : n \in \N, \ c \in \C\}$ if and only if
\[
  \Bigl| \int_0^1 f \Bigr|^\tau
  + \sum_{k=0}^\infty \sum_{i=0}^{2^k - 1} |\langle f, h_{k,i} \rangle|^\tau < \infty .
\]
\end{theorem}

\begin{proof}
\leanok
\Cref{thm:stechkin} for the Haar basis of \Cref{prop:haar-basis}, with the sum over $n \in \N$
split into the term $n = 0$ and the blocks $2^k \leq n < 2^{k+1}$.
\end{proof}

DeVore~\cite[\S3.6]{DeVore1998} characterises the rate $n^{-\theta}$, that is,
$A^\theta_\infty$, by the decay $\gamma_n(f) \leq M n^{-\theta-1/2}$ of the $n$-th largest
modulus $\gamma_n(f)$ of the Haar coefficients, for every $\theta > 0$. Here, too, the range
$\tau < 1$, that is, $\theta > 1/2$, is included, by \Cref{thm:stechkin}. For $\tau \geq 1$ the
condition is the one of \Cref{thm:haar} in $L_\tau$ with $p = q = \tau$ and $s = \theta$, where
the weight $2^{k(s + 1/2 - 1/p)}$ is $1$.

\begin{theorem}[Nonlinear approximation in $L_2$ and linear approximation in $L_\tau$]
\label{thm:haar-nterm-linear}
\lean{Haar.memApprox_nTerm_iff_memApprox_haarScheme}
\leanok
\uses{thm:haar-nterm, thm:haar, lem:lq-basic}
Let $1 \leq \tau < 2$ and $\theta = 1/\tau - 1/2$. Then $f \in L_2[0, 1]$ lies in the
approximation space $A^\theta_\tau$ of $n$-term approximation from the Haar basis in
$L_2[0, 1]$ if and only if it lies in the approximation space $A^\theta_\tau$ of the Haar scheme
in $L_\tau[0, 1]$.
\end{theorem}

\begin{proof}
\leanok
By \Cref{thm:haar} with $p = q = \tau$ and $s = \theta > 0$ the second condition is
$|\int_0^1 f| + \|(\|(\langle f, h_{k,i} \rangle)_i\|_{\ell_\tau})_k\|_{\ell_\tau} < \infty$.
By \Cref{lem:lq-basic} the $\tau$-th power of the second term is
$\sum_k \sum_i |\langle f, h_{k,i} \rangle|^\tau$, so this is the condition of
\Cref{thm:haar-nterm}.
\end{proof}

\chapter{Polynomial approximation on intervals}
\label{ch:interval}

This chapter treats the approximation by piecewise polynomials on an interval without
Fourier analysis and identifies their approximation spaces with Besov spaces defined by
moduli of smoothness. Functions are complex-valued and measurable on $\R$; norms are taken
on subintervals with Lebesgue measure.

\section{Differences and moduli}

\begin{definition}[Differences of higher order]
\label{def:iter-diff}
\lean{Rd.higherDiff, Rd.higherDiff_def, Rd.measurable_higherDiff,
  Interval.interval_higherDiff_apply, Interval.higherDiff_eq_sum_smul, Interval.higherDiff_zero,
  Interval.higherDiff_succ, Interval.higherDiff_succ', Interval.interval_higherDiff_add,
  Interval.higherDiff_sub, Interval.higherDiff_const_mul, Interval.interval_higherDiff_comp_add,
  Interval.higherDiff_reflect, Interval.enorm_higherDiff_coeff,
  Interval.enorm_higherDiff_le_of_forall}
\leanok
For $f : \R \to \C$, $r \in \N$ and $h \in \R$ the difference of order $r$ is
$\Delta_h^r f(x) = \sum_{k=0}^r (-1)^{r-k} \binom{r}{k} f(x + k h)$, the $r$-fold iterate of
$\Delta_h f(x) = f(x + h) - f(x)$. The same definition applies to functions on any additive
group, such as $\R^d$, and differences of measurable functions are measurable. Since
$\sum_k \binom{r}{k} = 2^r$, $|\Delta_h^r f(x)| \leq 2^r \max_{k \leq r} |f(x + k h)|$.
\end{definition}

\begin{lemma}[Differences annihilate polynomials]
\label{lem:iter-diff-poly}
\lean{Interval.natDegree_comp_affine_le, Interval.comp_affine_mem_degreeLT,
  Interval.higherDiff_polynomial}
\leanok
\uses{def:iter-diff}
Affine substitutions $x \mapsto u + v x$ do not raise the degree of a polynomial. For every
polynomial $P$ of degree $< r$ and every $h \in \R$, $\Delta_h^r P = 0$.
\end{lemma}

\begin{proof}
\leanok
Mathlib proves this for the step $1$; the polynomial $y \mapsto P(x + h y)$ has degree $< r$ and
its differences with step $1$ at $0$ are those of $P$ with step $h$ at $x$.
\end{proof}

\begin{definition}[Modulus of smoothness on an interval]
\label{def:interval-modulus}
\lean{Interval.intervalModulus, Interval.diffStepSet_Icc_of_nonneg,
  Interval.eLpNorm_indicator_diffStepSet_Icc, Interval.intervalModulus_eq_iSup_nonneg,
  Interval.eLpNorm_higherDiff_le_intervalModulus, Interval.intervalModulus_mono,
  Interval.eLpNorm_comp_add_Icc, Interval.eLpNorm_comp_sub_Icc, Interval.eLpNorm_Icc_mono}
\leanok
\uses{def:iter-diff, def:rd-modulus-on}
For $1 \leq p \leq \infty$, $r \in \N$ and $t \geq 0$ the modulus of smoothness of $f$ on
$[a, b]$ is the modulus on the set $[a, b]$ of \Cref{def:rd-modulus-on},
\[
  \omega_r(f, t, [a, b])_p = \sup_{|h| \leq t} \|\Delta_h^r f\|_{L_p([a, b]_{rh})}
  = \sup_{0 \leq h \leq t} \|\Delta_h^r f\|_{L_p[a, b - r h]} ,
\]
with $[a, b]_{rh} = \{x : x, x + h, \dots, x + rh \in [a, b]\}$, which is $[a, b - rh]$ for
$h \geq 0$. The two suprema agree since the steps $h$ and $-h$ give the same norm
(\Cref{lem:rd-modulus-on}).
\end{definition}

\begin{lemma}[Properties of the modulus on an interval]
\label{lem:interval-modulus}
\lean{Interval.eLpNorm_higherDiff_le_of_forall, Interval.eLpNorm_comp_add_mul_le,
  Interval.interval_eLpNorm_higherDiff_le, Interval.intervalModulus_le,
  Interval.intervalModulus_add_le, Interval.intervalModulus_const_mul,
  Interval.intervalModulus_polynomial, Interval.eLpNorm_higherDiff_congr_ae,
  Interval.intervalModulus_congr_ae, Interval.eLpNorm_higherDiff_reflect,
  Interval.intervalModulus_reflect}
\leanok
\uses{def:interval-modulus, lem:iter-diff-poly}
The modulus $\omega_r(f, t, [a, b])_p$ increases in $t$, is subadditive and homogeneous in $f$,
satisfies $\omega_r(f, t, [a, b])_p \leq 2^r \|f\|_{L_p[a, b]}$, vanishes for polynomials of
degree $< r$, depends only on the restriction of $f$ to $[a, b]$ up to null sets, and is
invariant under $f \mapsto f(a + b - \cdot)$.
\end{lemma}

\begin{proof}
\leanok
Each term $f(\cdot + k h)$ is a translate of $f$ restricted to a subinterval of $[a, b]$, and
$\sum_k \binom{r}{k} = 2^r$. Reflection maps $\Delta_h^r$ to $\pm \Delta_h^r$ at the reflected
point.
\end{proof}

\section{Polynomials on intervals}

\begin{lemma}[Nikolskii inequality for polynomials on intervals]
\label{lem:interval-nikolskii}
\lean{Interval.polyFun, Interval.continuous_polyFun, Interval.measurable_polyFun,
  Interval.finrank_degreeLT, Interval.polyFun_comp_affine, Interval.eq_zero_of_ae_eq_zero,
  Interval.memLp_polyFun, Interval.polyFun_degreeLTEquiv_symm, Interval.coeffToLp,
  Interval.enorm_coeffToLp, Interval.coeffToLp_injective,
  Interval.norm_polyFun_degreeLTEquiv_symm_le, Interval.exists_enorm_linearMap_le_eLpNorm,
  Interval.exists_enorm_le_eLpNorm,
  Interval.interval_eLpNorm_comp_affine, Interval.NikolskiiBound, Interval.NikolskiiBound.enorm_le,
  Interval.enorm_polyFun_le}
\leanok
Let $r \in \N$ and $1 \leq p < \infty$. There is $K < \infty$ such that for every polynomial $P$ of
degree $< r$, every interval $[u, u + L]$ with $L > 0$ and every $x \in [u, u + 2L]$,
\[
  |P(x)| \leq K L^{-1/p} \|P\|_{L_p[u, u + L]} .
\]
Moreover, for every interval $[\alpha, \beta]$ with $\alpha < \beta$ and every linear map $\Phi$
of the polynomials that maps the polynomials of degree $< r$ into themselves, there is
$K_\Phi < \infty$ with $|\Phi P(x)| \leq K_\Phi \|P\|_{L_p[\alpha, \beta]}$ for every polynomial
$P$ of degree $< r$ and every $x \in [0, 1]$.
\end{lemma}

\begin{proof}
\leanok
The polynomials of degree $< r$ form a space of dimension $r$. The map from the coefficients to
$L_p[0, \tfrac12]$ is injective, since a nonzero polynomial has finitely many zeros, hence bounded
below in finite dimension; this bounds the coefficients and so $|\Phi P(x)|$, $x \in [0, 1]$, by
$\|P\|_{L_p[0,1/2]}$ for every linear map $\Phi$ of the polynomials of degree $< r$ into
themselves, a linear map of $\C^r$ in coordinates. The identity $\Phi P = P$ bounds $|P(x)|$.
The general interval follows by the affine change of variables
$x = u + 2L y$, which multiplies the $L_p$ norm by $(2L)^{-1/p} \leq L^{-1/p}$.
\end{proof}

The bound on the doubled interval lets one piece of a piecewise polynomial control its values on
the neighbouring interval, which the inverse inequality for splines needs.

\begin{lemma}[Polynomial approximation of smooth functions]
\label{lem:taylor-lp}
\lean{Interval.exists_derivative_eq, Interval.hasDerivAt_polyFun,
  Interval.exists_eLpNorm_sub_polyFun_le}
\leanok
Let $1 \leq p \leq \infty$, $r \in \N$, $g \in C^r(\R)$ and $u \leq v$. There is a polynomial $Q$
of degree $< r$ with
$\|g - Q\|_{L_p[u, v]} \leq (v - u)^r \|g^{(r)}\|_{L_p[u, v]}$.
\end{lemma}

\begin{proof}
\leanok
Induction on $r$. Approximate $g'$ by $Q'$ of degree $< r - 1$ and let $Q$ be an antiderivative of
$Q'$ with $Q(u) = g(u)$. By the fundamental theorem of calculus and Hölder's inequality,
$|(g - Q)(x)| \leq \int_u^v |g' - Q'| \leq (v - u)^{1 - 1/p} \|g' - Q'\|_{L_p[u, v]}$, and the
$L_p$ norm of a function bounded by $M$ is at most $(v - u)^{1/p} M$.
\end{proof}

\section{Steklov smoothing}

\begin{definition}[Steklov averages]
\label{def:steklov}
\lean{Interval.steklov, Interval.steklov_eq, Interval.steklovIter, Interval.steklovIter_succ}
\leanok
For $g : \R \to \C$ and $\eta > 0$ the Steklov average is
$M_\eta g(x) = \eta^{-1} \int_0^\eta g(x + \tau) \, d\tau$, and $M_\eta^m$ is its $m$-fold
iterate.
\end{definition}

\begin{lemma}[Averages gain derivatives]
\label{lem:steklov-deriv}
\lean{Interval.hasDerivAt_steklov, Interval.differentiable_steklov, Interval.deriv_steklov,
  Interval.steklov_comp_add, Interval.steklov_sub, Interval.steklov_fwdDiff,
  Interval.continuous_steklov, Interval.continuous_fwdDiff, Interval.contDiff_fwdDiff,
  Interval.contDiff_steklov, Interval.contDiff_steklovIter, Interval.continuous_steklovIter,
  Interval.steklovIter_fwdDiff, Interval.iteratedDeriv_steklovIter_add,
  Interval.iteratedDeriv_steklovIter}
\leanok
\uses{def:steklov, def:iter-diff}
For continuous $g$, $(M_\eta g)' = \eta^{-1} \Delta_\eta g$. Hence $M_\eta^m g \in C^m$ and
$(M_\eta^r g)^{(r)} = \eta^{-r} \Delta_\eta^r g$.
\end{lemma}

\begin{proof}
\leanok
$M_\eta g(x) = \eta^{-1} (G(x + \eta) - G(x))$ with a primitive $G$ of $g$. Averages commute with
translations and differences, so each further average adds one derivative.
\end{proof}

\begin{lemma}[Minkowski's inequality for averages]
\label{lem:minkowski-avg}
\lean{Interval.eLpNorm_avg_le, Interval.avgStep, Interval.avgStepIter, Interval.avgStepIter_succ,
  Interval.continuous_avgStep, Interval.continuous_avgStepIter, Interval.eLpNorm_avgStepIter_le,
  Interval.avgStepIter_congr_point}
\leanok
\uses{lem:weighted-jensen}
Let $X$ be a second countable topological space with an $s$-finite Borel measure $\mu$, let
$K : \R \times X \to \C$ be continuous, $\delta > 0$ and $1 \leq p < \infty$. Then
\[
  \Bigl\| \delta^{-1} \int_0^\delta K(t, \cdot) \, dt \Bigr\|_{L_p(\mu)}
  \leq \sup_{0 \leq t \leq \delta} \|K(t, \cdot)\|_{L_p(\mu)} .
\]
Iterating $m$ times, an average over $v \in [0, \delta]^m$ of $G(\sigma + v_1 + \cdots + v_m, \cdot)$
has $L_p(\mu)$ norm at most $\sup_{\sigma \leq h \leq \sigma + m \delta} \|G(h, \cdot)\|_{L_p(\mu)}$.
The average at a point $x$ depends only on the values $G(\cdot, x)$.
\end{lemma}

\begin{proof}
\leanok
The average is the weighted average with the weight
$\Phi = \delta^{-1} \mathbf{1}_{(0, \delta]}$, which has $\|\Phi\|_1 = 1$, so
\Cref{lem:weighted-jensen} bounds it by
$\sup_{0 < t \leq \delta} \|K(t, \cdot)\|_{L_p(\mu)} \leq
\sup_{0 \leq t \leq \delta} \|K(t, \cdot)\|_{L_p(\mu)}$.
\end{proof}

\begin{theorem}[Steklov smoothing]
\label{thm:steklov-smoothing}
\lean{Interval.avgStepIter_const, Interval.avgStepIter_shift, Interval.avgStepIter_sum,
  Interval.steklovSmooth, Interval.sub_steklovSmooth_eq, Interval.contDiff_steklovSmooth,
  Interval.iteratedDeriv_steklovSmooth, Interval.continuous_higherDiff_uncurry,
  Interval.eLpNorm_higherDiff_Icc_le_intervalModulus, Interval.eLpNorm_sub_steklovSmooth_le,
  Interval.iteratedDeriv_steklovSmooth_eq_sum_smul, Interval.sum_range_choose_succ_le,
  Interval.eLpNorm_smoothCoeff_smul_le, Interval.eLpNorm_sum_smoothCoeff_smul_le,
  Interval.eLpNorm_iteratedDeriv_steklovSmooth_le}
\leanok
\uses{lem:steklov-deriv, lem:minkowski-avg, def:interval-modulus}
Let $f$ be continuous, $r \in \N$, $\delta > 0$ and
$S_\delta f = \sum_{k=1}^r (-1)^{k+1} \binom{r}{k} M_{k\delta}^r f$. Then $S_\delta f \in C^r$, and for
$1 \leq p < \infty$ and $[u, w]$ with $a \leq u$ and $w + r^2 \delta \leq b$,
\[
  \|f - S_\delta f\|_{L_p[u, w]} \leq \omega_r(f, r\delta, [a, b])_p
\]
and
\[
  \|(S_\delta f)^{(r)}\|_{L_p[u, w]} \leq 2^r \delta^{-r} \omega_r(f, r\delta, [a, b])_p .
\]
At every point,
$(S_\delta f)^{(r)} = \sum_{k=1}^r (-1)^{k+1} \binom{r}{k} (k\delta)^{-r} \Delta_{k\delta}^r f$.
For every measure $\mu$ on a space $X$, every $1 \leq p \leq \infty$, every $B \in [0, \infty]$
and every family $(D_h)_{h \in \R}$ of functions $D_h : X \to \C$ with
$\|D_h\|_{L_p(\mu)} \leq B$ for $0 \leq h \leq r\delta$,
\[
  \Bigl\| \sum_{k=1}^r (-1)^{k+1} \binom{r}{k} (k\delta)^{-r} D_{k\delta} \Bigr\|_{L_p(\mu)}
  \leq 2^r \delta^{-r} B .
\]
\end{theorem}

\begin{proof}
\leanok
Averaging $f(x + k h)$ over $h = v_1 + \cdots + v_r$, $v \in [0, \delta]^r$, gives
$M_{k\delta}^r f(x)$, so $f - S_\delta f = (-1)^r$ times the average of $\Delta_h^r f$, which
\Cref{lem:minkowski-avg} bounds. The identity for the derivative is \Cref{lem:steklov-deriv} for
each $M_{k\delta}^r f$. For the general bound, $(k\delta)^{-r} \leq \delta^{-r}$,
$\sum_{k=1}^r \binom{r}{k} \leq 2^r$, and the triangle inequality. With
$D_h = \Delta_h^r f$ it gives the bound on $[u, w]$, since $[u, w] \subseteq [a, b - rh]$ and so
$\|\Delta_h^r f\|_{L_p[u, w]} \leq \omega_r(f, r\delta, [a, b])_p$ for $0 \leq h \leq r\delta$.
\end{proof}

\begin{lemma}[Smoothing on a left and a right part]
\label{lem:two-sided-smoothing}
\lean{Interval.step_div_sq_bounds, Interval.two_pow_mul_inv_div_sq_pow_mul_le,
  Interval.exists_left_smoothing, Interval.exists_eLpNorm_sub_polyFun_reflect_le,
  Interval.exists_right_smoothing}
\leanok
\uses{thm:steklov-smoothing, lem:taylor-lp, lem:interval-modulus}
Let $F$ be continuous, $r \geq 1$, $1 \leq p < \infty$, $L > 0$, $w + L \leq b$ and
$\omega = \omega_r(F, L, [0, b])_p$. On $P = [0, w]$ and on $P = [b - w, b]$ there are functions
$g$ and $D$ with $\|F - g\|_{L_p(P)} \leq \omega$, $\|D\|_{L_p(P)} \leq L^{-r} 2^r r^{2r} \omega$, and
for all $a \leq b$ a polynomial $T$ of degree $< r$ with
$\|g - T\|_{L_p[a,b]} \leq (b - a)^r \|D\|_{L_p[a,b]}$.
\end{lemma}

\begin{proof}
\leanok
On $[0, w]$ take the Steklov smoothing $g = S_\delta F$, $\delta = L/r^2$, so that
$r^2 \delta = L$, and $D = g^{(r)}$ (\Cref{thm:steklov-smoothing}, \Cref{lem:taylor-lp}). On
$[b - w, b]$ apply this to $F(b - \cdot)$, which has the same modulus, and reflect back.
\end{proof}

\section{Whitney's inequality}

\begin{lemma}[Polynomials and their middle third]
\label{lem:middle-third}
\lean{Interval.exists_eLpNorm_le_middle}
\leanok
\uses{lem:interval-nikolskii}
Let $1 \leq p < \infty$ and $r \in \N$. There is $K' < \infty$ with
$\|P\|_{L_p[0,1]} \leq K' \|P\|_{L_p[1/3, 2/3]}$ for every polynomial $P$ of degree $< r$.
\end{lemma}

\begin{proof}
\leanok
All norms on the space of polynomials of degree $< r$ are equivalent, so
$|P(x)| \leq K' \|P\|_{L_p[1/3,2/3]}$ for $x \in [0, 1]$, and $[0, 1]$ has length $1$.
\end{proof}

\begin{lemma}[Gluing two polynomial approximations]
\label{lem:poly-glue}
\lean{Interval.exists_eLpNorm_sub_polyFun_le_glue, Interval.eLpNorm_sub_polyFun_le_add}
\leanok
\uses{lem:middle-third}
Let $1 \leq p < \infty$ and $r \in \N$. There is $K < \infty$ such that for every measurable $u$
and all polynomials $T_1$, $T_2$ of degree $< r$,
\[
  \|u - T_1\|_{L_p[0,1]}
  \leq (1 + K) \bigl( \|u - T_1\|_{L_p[0, 2/3]} + \|u - T_2\|_{L_p[1/3, 1]} \bigr) .
\]
If $b - a \leq 1$ and $\|g - T\|_{L_p[a,b]} \leq (b - a)^r \|D\|_{L_p[a,b]}$, then
$\|u - T\|_{L_p[a,b]} \leq \|u - g\|_{L_p[a,b]} + \|D\|_{L_p[a,b]}$.
\end{lemma}

\begin{proof}
\leanok
On $[\tfrac13, \tfrac23]$,
$\|T_1 - T_2\| \leq \|u - T_1\|_{L_p[0,2/3]} + \|u - T_2\|_{L_p[1/3,1]}$, and
\Cref{lem:middle-third} extends this bound to $[0, 1]$ with the factor $K$. On $[0, \tfrac12]$
keep $T_1$; on $[\tfrac12, 1]$ write $u - T_1 = (u - T_2) - (T_1 - T_2)$. The second statement is
the triangle inequality.
\end{proof}

\begin{theorem}[Whitney's inequality]
\label{thm:whitney}
\lean{Interval.eLpNorm_Icc_split, Interval.eLpNorm_sub_polyFun_le_of_smoothing,
  Interval.exists_eLpNorm_sub_polyFun_le_intervalModulus}
\leanok
\uses{lem:two-sided-smoothing, lem:poly-glue, lem:interval-modulus}
Let $1 \leq p < \infty$, $r \geq 1$ and $0 < t \leq \tfrac13$. There is $W < \infty$ such that
every continuous $F$ has a polynomial $P$ of degree $< r$ with
\[
  \|F - P\|_{L_p[0, 1]} \leq W \, \omega_r(F, t, [0, 1])_p .
\]
\end{theorem}

\begin{proof}
\leanok
Let $\omega = \omega_r(F, t, [0,1])_p$ and $A = 1 + t^{-r} 2^r r^{2r}$. By
\Cref{lem:two-sided-smoothing} with $L = t$ and $w = 1 - t$, each on one interval of length at
most $1$, there are polynomials $T_1$, $T_2$ of degree $< r$ with
$\|F - T_1\|_{L_p[0, 1-t]} \leq A \omega$ and $\|F - T_2\|_{L_p[t, 1]} \leq A \omega$
(\Cref{lem:poly-glue}, second statement). Since $[0, \tfrac23] \subseteq [0, 1 - t]$ and
$[\tfrac13, 1] \subseteq [t, 1]$, \Cref{lem:poly-glue} gives
$\|F - T_1\|_{L_p[0,1]} \leq 2 (1 + K) A \omega$.
\end{proof}

Whitney's inequality bounds the best approximation by a single polynomial on the whole interval.
For the dyadic schemes below it is the Jackson inequality at level $0$, where the approximants
are the polynomials of degree $< r$.

\section{Dyadic piecewise polynomials}

\begin{definition}[Dyadic piecewise polynomials]
\label{def:pp-scheme}
\lean{Interval.unitVol, Interval.isProbabilityMeasure_unitVol, Interval.ae_mem_Ico_unitVol,
  Interval.dyadicIco, Interval.measurableSet_dyadicIco, Interval.mem_dyadicIco,
  Interval.eq_of_mem_dyadicIco,
  Interval.mem_dyadicIco_floor, Interval.floor_two_pow_succ_div, Interval.ppFun,
  Interval.ppFun_eq_of_mem, Interval.measurable_ppFun, Interval.memLp_ppFun, Interval.ppFun_add,
  Interval.ppFun_const_mul, Interval.polyFun_zero, Interval.ppFun_zero, Interval.ppFun_sum,
  Interval.dyadicIco_subset, Interval.ppFun_eq_zero, Interval.exists_norm_ppFun_le,
  Interval.integrableOn_ppFun, Interval.ppFun_refine_eq, Interval.ppFun_refine,
  Interval.ppSpace, Interval.ppSpace_mono, Interval.ppScheme, Interval.mem_ppScheme_zero,
  Interval.mem_ppScheme, Interval.mem_ppScheme_two_pow, Interval.ppScheme_isLinearScheme}
\leanok
\uses{def:comap}
Let $r \geq 1$ and $1 \leq p < \infty$. For $k, n \in \N$ the dyadic interval of level $k$ is
$I_{k,n} = [n 2^{-k}, (n+1) 2^{-k})$. A piecewise polynomial of level $k$ agrees on every
$I_{k,n}$ with $n < 2^k$ with a polynomial of degree $< r$, and vanishes outside $[0, 1)$; every
piecewise polynomial of level $k$ is bounded on $\R$ and integrable on $[0, 1]$. The space
$V_k \subseteq L_p[0,1]$ consists of the classes of piecewise polynomials of level $k$. The
scheme of dyadic piecewise polynomials is the dyadic scheme of $(V_k)_k$: $A_0 = \{0\}$ and
$A_n = V_{\lfloor \log_2 n \rfloor}$ for $n \geq 1$.
\end{definition}

\begin{lemma}[Bookkeeping for dyadic intervals]
\label{lem:pp-bookkeeping}
\lean{Interval.floor_eq_of_mem_dyadicIco,
  Interval.pairwise_disjoint_dyadicIco, Interval.ppFun_eq_polyFun,
  Interval.ite_mem_degreeLT, Interval.integrableOn_dyadicIco_of_continuous,
  Interval.lintegral_Icc_le_sum, Interval.sum_le_lintegral_Icc,
  Interval.lintegral_Icc_eq_dyadicIco, Interval.zero_add_two_pow_mul_eq_one,
  Interval.lintegral_Icc_le_sum_dyadicIco, Interval.sum_lintegral_dyadicIco_le,
  Interval.lintegral_unit_eq_sum_dyadicIco, MeasureTheory.eLpNorm_le_of_lintegral_le}
\leanok
\uses{def:pp-scheme}
For $x \in I_{k,n}$, $\lfloor 2^k x \rfloor = n$. A piecewise polynomial of level $k$ vanishes
outside $[0, 1)$, and the pieces $Q_{\lfloor n/2 \rfloor}$ of level $k + 1$ give the same function
as the pieces $Q_n$ of level $k$. For $u \in \R$, $L \geq 0$, $m \in \N$ and every
$G : \R \to [0, \infty]$,
$\int_{[u, u + mL]} G = \sum_{n < m} \int_{[u + nL, u + (n+1)L]} G$; in particular
$\int_{[0, 1]} G = \sum_{n < 2^k} \int_{I_{k,n}} G$.
\end{lemma}

\begin{proof}
\leanok
The first claim restates $n 2^{-k} \leq x < (n+1) 2^{-k}$, and the claims on piecewise
polynomials follow from their definition, since $I_{k+1,n} \subseteq I_{k,\lfloor n/2 \rfloor}$.
The identity for the integrals follows by induction on $m$, splitting off the last interval,
whose intersection with the others is a single point.
\end{proof}

\begin{lemma}[Piecewise Taylor approximation]
\label{lem:pp-taylor}
\lean{Interval.exists_taylor_pieces, Interval.exists_ppFun_approx_of_taylor,
  Interval.one_half_add_le_one}
\leanok
\uses{def:pp-scheme, lem:pp-bookkeeping}
Let $1 \leq p < \infty$, $r \geq 1$, $k \geq 1$ and $L = 2^{-k}$. Let $g_1, D_1, g_2, D_2$ be
measurable, and such that for all $a \leq b$ and $\nu = 1, 2$ there is a polynomial $T$ of degree
$< r$ with $\|g_\nu - T\|_{L_p[a,b]} \leq (b - a)^r \|D_\nu\|_{L_p[a,b]}$. For every measurable $u$
there is a piecewise polynomial $S$ of level $k$ and degree $< r$ with
\begin{multline*}
  \|u - S\|_{L_p[0,1]} \leq \bigl( \|u - g_1\|_{L_p[0,1/2]} + L^r \|D_1\|_{L_p[0,1/2]} \bigr) \\
    + \bigl( \|u - g_2\|_{L_p[1/2,1]} + L^r \|D_2\|_{L_p[1/2,1]} \bigr) .
\end{multline*}
\end{lemma}

\begin{proof}
\leanok
On each $I_{k,n} \subseteq [0, \tfrac12]$ take the polynomial $T$ of $g_1$ on the closure of
$I_{k,n}$, so that $\|g_1 - S\|_{L_p(I_{k,n})} \leq L^r \|D_1\|_{L_p(I_{k,n})}$; the $p$-th powers
add up over the $2^{k-1}$ intervals of $[0, \tfrac12]$. The same with $g_2$ on $[\tfrac12, 1]$,
where $\tfrac12$ is a knot since $k \geq 1$. The triangle inequality on each half gives the bound.
\end{proof}

The lemma asks nothing of $u$, $g_\nu$ and $D_\nu$ beyond measurability, so it applies to the
slices of functions on a cube as well (\Cref{thm:cube-jackson}).

\begin{theorem}[Jackson's inequality for piecewise polynomials]
\label{thm:pp-jackson}
\lean{Interval.intervalJacksonConst, Interval.intervalJacksonConst_ne_top,
  Interval.add_ofReal_pow_mul_le_intervalJacksonConst, Interval.exists_ppFun_approx,
  Interval.intervalLpModulus, Interval.intervalLpModulus_add_le, Interval.intervalLpModulus_le,
  Interval.intervalLpModulus_anti, Interval.interval_approxError_le_of_continuous,
  Interval.approxError_ppScheme_le, Interval.exists_approxError_ppScheme_one_le,
  Interval.exists_approxError_ppScheme_le}
\leanok
\uses{def:pp-scheme, lem:pp-taylor, lem:two-sided-smoothing, thm:whitney, lem:error-dense}
Let $1 \leq p < \infty$ and $r \geq 1$. For every $f \in L_p[0,1]$ and $k \geq 1$,
\[
  E_{2^k}(f)_p \leq 2 (1 + 2^r r^{2r}) \, \omega_r(f, 2^{-k}, [0,1])_p ,
\]
and there is $C_J < \infty$ with $E_{2^k}(f)_p \leq C_J \, \omega_r(f, 2^{-k}, [0,1])_p$ for
every $f \in L_p[0,1]$ and every $k \geq 0$.
\end{theorem}

\begin{proof}
\leanok
First let $f$ be continuous, $L = 2^{-k}$ and $\omega = \omega_r(f, L, [0,1])_p$. Since
$\tfrac12 + L \leq 1$, \Cref{lem:two-sided-smoothing} on $[0, \tfrac12]$ and on $[\tfrac12, 1]$
gives $g_1, D_1$ and $g_2, D_2$ with $\|f - g_\nu\| \leq \omega$ and
$\|D_\nu\| \leq L^{-r} 2^r r^{2r} \omega$ on the respective half, and Taylor polynomials on all
intervals. \Cref{lem:pp-taylor} with $u = f$ gives a piecewise polynomial $S$ of level $k$ with
$\|f - S\|_{L_p[0,1]} \leq 2 (1 + 2^r r^{2r}) \omega$. General $f$ follow, for every scheme in
$L_p[0,1]$, by
\Cref{lem:error-dense}: the bounded continuous functions are dense in $L_p$, and
$\omega_r(\cdot, t)_p$ is subadditive with $\omega_r(f, t)_p \leq 2^r \|f\|_p$. For $k = 0$ the approximants are the
polynomials of degree $< r$, and \Cref{thm:whitney} with $t = \tfrac14$, passed to $L_p$ in the
same way, gives $E_1(f)_p \leq W \omega_r(f, \tfrac14, [0,1])_p \leq W \omega_r(f, 1, [0,1])_p$.
\end{proof}

\begin{lemma}[Local estimates for piecewise polynomials]
\label{lem:pp-local}
\lean{Interval.localMass, Interval.tsum_localMass_ppFun_le, Interval.IsPiecewisePoly,
  Interval.isPiecewisePoly_ppFun, Interval.eLpNorm_piece_eq_localMass, Interval.enorm_le_localMass,
  Interval.floor_add_mul_le, Interval.localMass_le_sum_range,
  Interval.higherDiff_eq_of_mem_piece, Interval.higherDiff_eq_zero_of_isPiecewisePoly,
  Interval.enorm_higherDiff_le_of_isPiecewisePoly_of_le,
  Interval.enorm_higherDiff_le_of_isPiecewisePoly,
  Interval.stripBound, Interval.lintegral_stripBound_le, Interval.eLpNorm_le_of_strip,
  Interval.eLpNorm_higherDiff_le_of_pointwise, Interval.sum_sum_localMass_le}
\leanok
\uses{def:pp-scheme, lem:pp-bookkeeping, lem:interval-nikolskii, lem:iter-diff-poly}
Let $1 \leq p < \infty$, and let $a$ be a piecewise polynomial of level $j$ and degree $< r$
on $[0, \infty)$: on every $I_{j,n}$, $n \in \N$, it agrees with a polynomial of degree $< r$.
Let $m_n = \int_{I_{j,n}} |a|^p$ for $n \in \N$, $0 \leq h \leq 2^{-j}$, $x \geq 0$ and
$n = \lfloor 2^j x \rfloor$. If $k \in \N$ and $x + kh < (n+1) 2^{-j}$, then
$\Delta_h^k a(x) = \Delta_h^k P(x)$ for the piece $P$ of $a$ on $I_{j,n}$; for $k = r$ this
vanishes. In every case, for $k \leq r$,
\[
  |\Delta_h^k a(x)| \leq 2^k K 2^{j/p} \Bigl( \sum_{l=0}^r m_{n+l} \Bigr)^{1/p}
\]
with the constant $K$ of \Cref{lem:interval-nikolskii}. If $a$ vanishes outside $[0, 1)$, then
$\sum_{n < 2^j} \sum_{l = 0}^r m_{n+l} \leq (r+1) \|a\|_p^p$. Finally, let $S \subseteq [0, 1]$ be
measurable. If $|\Delta_h^r a(x)| \leq G (\sum_{l=0}^r m_{n+l})^{1/p}$ for almost every $x \in S$
with $(n+1) 2^{-j} \leq x + rh$, and $\sum_{n < 2^j} \sum_{l = 0}^r m_{n+l} \leq (r+1) B^p$, then
$\|\Delta_h^r a\|_{L_p(S)} \leq G (rh (r + 1))^{1/p} B$.
\end{lemma}

\begin{proof}
\leanok
In the first case $x, x + h, \dots, x + kh$ lie in $I_{j,n}$, where $a = P$; for $k = r$,
\Cref{lem:iter-diff-poly} gives $0$. In general, for $k \leq r$ and $i \leq k$,
$x + ih \in I_{j, n+l}$ with $0 \leq l \leq r$, where $|a| \leq K 2^{j/p} m_{n+l}^{1/p}$ by
\Cref{lem:interval-nikolskii}, and $\sum_i \binom{k}{i} = 2^k$. In the double sum every $m_n$
occurs at most $r + 1$ times. For the last statement, $|\Delta_h^r a|^p$ vanishes on $I_{j,n}$
outside the strip $[(n+1)2^{-j} - rh, (n+1)2^{-j})$ of width $rh$ and is bounded on it by
$G^p \sum_{l=0}^r m_{n+l}$.
\end{proof}

\begin{theorem}[Inverse inequality for piecewise polynomials]
\label{thm:pp-inverse}
\lean{Interval.ofReal_two_pow_rpow_mul_eq, Interval.eLpNorm_higherDiff_le_of_isPiecewisePoly,
  Interval.eLpNorm_higherDiff_ppFun_le, ENNReal.ofReal_two_pow_mul_rpow_le, Interval.ppInverseConst,
  Interval.ppInverseConst_ne_top,
  Interval.two_pow_le_ppInverseConst, Interval.le_ppInverseConst_mul,
  Interval.exists_inverse_ppScheme}
\leanok
\uses{lem:pp-local}
Let $1 \leq p < \infty$ and $r \in \N$. There is $C_B < \infty$ such that for all $j, k \in \N$
and every $a \in V_j$,
\[
  \omega_r(a, 2^{-k}, [0, 1])_p \leq C_B \, 2^{-\max(k - j, 0)/p} \|a\|_p .
\]
\end{theorem}

\begin{proof}
\leanok
For $k < j$ this is $\omega_r(a, t, [0,1])_p \leq 2^r \|a\|_p$. Let $k \geq j$ and
$0 \leq h \leq 2^{-k}$. By \Cref{lem:pp-local} with $G = 2^r K 2^{j/p}$, for every piecewise
polynomial $a$ of level $j$ on $[0, \infty)$ with
$\sum_{n < 2^j} \sum_{l=0}^r m_{n+l} \leq (r+1) B^p$,
$\|\Delta_h^r a\|_{L_p[0, 1]} \leq 2^r K (r (r+1) \, 2^j h)^{1/p} B$. Here $B = \|a\|_p$,
and $2^j h \leq 2^{-(k-j)}$.
\end{proof}

\begin{theorem}[Approximation spaces of dyadic piecewise polynomials]
\label{thm:pp-besov}
\lean{Interval.ppModulusData, Interval.memApprox_ppScheme_iff}
\leanok
\uses{thm:pp-jackson, thm:pp-inverse, def:modulus-data, thm:modulus-coincidence}
Let $1 \leq p < \infty$, $1 \leq q \leq \infty$, $r \geq 1$ and $0 \leq s < 1/p$. A function
$f \in L_p[0,1]$ lies in the approximation space $A^s_q$ of the scheme of dyadic piecewise
polynomials of degree $< r$ if and only if
\[
  \|f\|_p + \bigl\| \bigl( 2^{ks} \omega_r(f, 2^{-k}, [0,1])_p \bigr)_k \bigr\|_{\ell_q} < \infty .
\]
\end{theorem}

\begin{proof}
\leanok
By \Cref{thm:pp-jackson}, \Cref{thm:pp-inverse} and $\omega_r(f, t, [0,1])_p \leq 2^r \|f\|_p$
(\Cref{lem:interval-modulus}), $\omega_r(\cdot, 2^{-k}, [0,1])_p$ is a modulus compatible with
the scheme of inverse order $\theta = 1/p$ in the sense of \Cref{def:modulus-data}. Apply
\Cref{thm:modulus-coincidence}.
\end{proof}

For $0 < s < 1/p$ the condition says that $f$ lies in the Besov space $\Besov{s}{p}{q}[0,1]$
defined by moduli of smoothness. The bound $s < 1/p$ comes from the jumps at the knots: the
indicator of $[0, \tfrac12)$ lies in $V_1$, and its modulus is of order $t^{1/p}$. Higher
smoothness needs piecewise polynomials that are smooth across the knots, that is, splines.

\section{Piecewise polynomials on the circle}

\begin{definition}[Periodic modulus of smoothness]
\label{def:per-modulus}
\lean{Interval.perExt, Interval.perExt_periodic, Interval.perExt_of_mem,
  Interval.measurable_perExt, Interval.perExt_add, Interval.perModulus,
  Interval.eLpNorm_higherDiff_le_perModulus, Interval.perModulus_mono}
\leanok
\uses{def:interval-modulus}
For $f : \R \to \C$ the periodic extension of $f|_{[0,1)}$ is
$\tilde f(x) = f(x - \lfloor x \rfloor)$. For $1 \leq p \leq \infty$, $r \in \N$ and $t \geq 0$
the periodic modulus of smoothness is
\[
  \omega_r^{\mathrm{per}}(f, t)_p = \sup_{0 \leq h \leq t} \|\Delta_h^r \tilde f\|_{L_p[0,1]} .
\]
\end{definition}

Through $\tilde f$, functions on $[0, 1)$ are the functions on the circle $\T = \R / \Z$, and
$L_p[0,1]$ is $L_p(\T)$. Unlike $\omega_r(f, t, [0,1])_p$, the periodic modulus compares values
near $1$ with values near $0$.

\begin{lemma}[Properties of the periodic modulus]
\label{lem:per-modulus}
\lean{Interval.perExt_ae_eq, Interval.higherDiff_perExt_ae_eq, Interval.lintegral_Icc_periodic,
  Interval.perModulus_congr_ae, Interval.perModulus_add_le, Interval.eLpNorm_perExt_Icc,
  Interval.perModulus_le, Interval.intervalModulus_le_perModulus, Interval.perExt_eq_of_periodic,
  Interval.higherDiff_periodic, Interval.lintegral_Icc_nat_le, Interval.intervalModulus_perExt_le}
\leanok
\uses{def:per-modulus, lem:interval-modulus}
Let $1 \leq p < \infty$, $r \in \N$ and $t \geq 0$. The periodic modulus depends only on the class
of $f$ in $L_p[0,1]$, is subadditive in $f$, and for measurable $f$ and every $n \in \N$
\[
  \omega_r(f, t, [0,1])_p \leq \omega_r^{\mathrm{per}}(f, t)_p \leq 2^r \|f\|_{L_p[0,1]}
\]
and
\[
  \omega_r(\tilde f, t, [0, n])_p \leq n^{1/p} \, \omega_r^{\mathrm{per}}(f, t)_p .
\]
\end{lemma}

\begin{proof}
\leanok
If $f = g$ almost everywhere on $[0, 1]$, then $\{\tilde f \neq \tilde g\}$ is contained in the
union of the integer translates of a null set. The integral of a $1$-periodic function over
$[c, c + 1]$ is its integral over the circle, so it does not depend on $c$; hence each term
$\tilde f(\cdot + kh)$ of $\Delta_h^r \tilde f$ has $L_p[0,1]$ norm $\|f\|_{L_p[0,1]}$, and
$\sum_k \binom{r}{k} = 2^r$. For $x \in [0, 1 - rh)$ the points $x, \dots, x + rh$ lie in
$[0, 1)$, where $\tilde f = f$. Finally $\Delta_h^r \tilde f$ is $1$-periodic and $[0, n]$ is
covered by the $n$ intervals $[i, i + 1]$, so
$\|\Delta_h^r \tilde f\|_{L_p[0, n]}^p \leq n \|\Delta_h^r \tilde f\|_{L_p[0, 1]}^p$.
\end{proof}

\begin{theorem}[Approximation spaces on the circle]
\label{thm:pp-besov-per}
\lean{Interval.perPieces, Interval.perPieces_eq, Interval.perPieces_mem,
  Interval.perExt_ppFun_eq, Interval.isPiecewisePoly_perExt_ppFun, Interval.sum_localMass_perExt_le,
  Interval.sum_sum_localMass_perExt_le, Interval.lpPerModulus, Interval.lpPerModulus_add_le,
  Interval.lpPerModulus_le, Interval.intervalLpModulus_le_lpPerModulus,
  Interval.exists_approxError_ppScheme_le_per, Interval.exists_inverse_ppScheme_per,
  Interval.ppPerModulusData, Interval.memApprox_ppScheme_iff_per}
\leanok
\uses{thm:pp-jackson, thm:pp-inverse, lem:pp-local, lem:per-modulus, def:modulus-data,
  thm:modulus-coincidence}
Let $1 \leq p < \infty$, $1 \leq q \leq \infty$, $r \geq 1$ and $0 \leq s < 1/p$. A function
$f \in L_p[0,1]$ lies in the approximation space $A^s_q$ of the scheme of dyadic piecewise
polynomials of degree $< r$ if and only if
\[
  \|f\|_p + \bigl\| \bigl( 2^{ks} \omega_r^{\mathrm{per}}(f, 2^{-k})_p \bigr)_k \bigr\|_{\ell_q}
  < \infty .
\]
\end{theorem}

\begin{proof}
\leanok
Jackson's inequality for $\omega_r^{\mathrm{per}}$ at every level follows from
\Cref{thm:pp-jackson} and $\omega_r(f, t, [0,1])_p \leq \omega_r^{\mathrm{per}}(f, t)_p$. For the inverse inequality, the
periodic extension $\tilde a$ of $a \in V_j$ is a piecewise polynomial of level $j$ on
$[0, \infty)$: on $I_{j, m' + 2^j c}$ with $m' < 2^j$ it is the piece of $a$ on $I_{j, m'}$
shifted by $c$. Any $2^j$ consecutive pieces of $\tilde a$ make up one period, so
$\sum_{n < 2^j} \sum_{l = 0}^r m_{n+l} \leq (r + 1) \|a\|_p^p$, and the estimate in the proof of
\Cref{thm:pp-inverse} gives
$\omega_r^{\mathrm{per}}(a, 2^{-k})_p \leq C_B 2^{-\max(k-j, 0)/p} \|a\|_p$. Apply
\Cref{thm:modulus-coincidence} with $\theta = 1/p$.
\end{proof}

\begin{corollary}[Interval and circle]
\label{cor:interval-circle}
\lean{Interval.lqNorm_modulus_lt_top_iff_perModulus}
\leanok
\uses{thm:pp-besov, thm:pp-besov-per}
Let $1 \leq p < \infty$, $1 \leq q \leq \infty$, $r \geq 1$ and $0 \leq s < 1/p$. For every
$f \in L_p[0,1]$,
\begin{multline*}
  \|f\|_p + \bigl\| \bigl( 2^{ks} \omega_r(f, 2^{-k}, [0,1])_p \bigr)_k \bigr\|_{\ell_q} < \infty \\
  \iff
  \|f\|_p + \bigl\| \bigl( 2^{ks} \omega_r^{\mathrm{per}}(f, 2^{-k})_p \bigr)_k \bigr\|_{\ell_q}
  < \infty .
\end{multline*}
\end{corollary}

\begin{proof}
\leanok
By \Cref{thm:pp-besov} and \Cref{thm:pp-besov-per} both conditions say $f \in A^s_q$.
\end{proof}

So $\Besov{s}{p}{q}[0,1] = \Besov{s}{p}{q}(\T)$ for $0 < s < 1/p$: below smoothness $1/p$ a jump
between the ends of the interval costs nothing.

\section{Haar coefficients}
\label{sec:haar-coefficients}

\begin{corollary}[Besov spaces through Haar coefficients]
\label{cor:haar-besov}
\lean{Haar.polyFun_of_mem_degreeLT_one, Haar.dyadicIndex_eq_floor, Haar.dyadicIndex_intCast_div,
  Haar.dyadicIndex_natCast_div, Haar.ppFun_ae_eq_stepFun, Haar.ae_eq_ppFun_of_stronglyMeasurable,
  Haar.mem_haarScheme_of_mem_ppScheme, Haar.mem_ppScheme_of_mem_haarScheme,
  Haar.haarScheme_eq_ppScheme, Haar.memApprox_ppScheme_one_iff,
  Haar.lqNorm_modulus_lt_top_iff_haarCoeff, Haar.lqNorm_perModulus_lt_top_iff_haarCoeff}
\leanok
\uses{thm:haar, thm:pp-besov, thm:pp-besov-per}
For every $1 \leq p < \infty$ the Haar scheme in $L_p[0, 1]$ is the scheme of dyadic piecewise
polynomials of degree $< 1$. Let $1 \leq p < \infty$, $1 \leq q \leq \infty$ and $0 < s < 1/p$.
For $f \in L_p[0,1]$ the following are equivalent.
\begin{enumerate}
  \item $\|f\|_p + \bigl\| \bigl( 2^{ks} \omega_1(f, 2^{-k}, [0,1])_p \bigr)_k \bigr\|_{\ell_q}
    < \infty$.
  \item $\|f\|_p + \bigl\| \bigl( 2^{ks} \omega_1^{\mathrm{per}}(f, 2^{-k})_p \bigr)_k
    \bigr\|_{\ell_q} < \infty$.
  \item $\bigl| \int_0^1 f \bigr| + \Bigl\| \Bigl(2^{k(s + 1/2 - 1/p)}
    \bigl\| (\langle f, h_{k,m} \rangle)_{0 \leq m < 2^k} \bigr\|_{\ell_p} \Bigr)_k
    \Bigr\|_{\ell_q} < \infty$.
\end{enumerate}
\end{corollary}

\begin{proof}
\leanok
A function measurable for the dyadic $\sigma$-algebra $\mathcal{F}_j$ is constant on every
$I_{j,n}$, and a function constant on every $I_{j,n}$ is $\mathcal{F}_j$-measurable; this
gives the first claim. By it, each of the three conditions says $f \in A^s_q$, by
\Cref{thm:pp-besov}, \Cref{thm:pp-besov-per} and \Cref{thm:haar}.
\end{proof}

The case of the Haar system and $1 \leq \tau < 2$ of the theorem of DeVore, Jawerth and
Popov~\cite{DeVoreJawerthPopov} (see also \cite[\S3.6 and Remark~7.6]{DeVore1998}) reads as
follows.

\begin{corollary}[$n$-term Haar approximation and Besov spaces]
\label{cor:haar-nterm-besov}
\lean{Haar.memApprox_nTerm_iff_modulus, Haar.memApprox_nTerm_iff_perModulus}
\leanok
\uses{thm:haar-nterm-linear, cor:haar-besov}
Let $1 \leq \tau < 2$ and $\theta = 1/\tau - 1/2$. For $f \in L_2[0, 1]$ the following are
equivalent.
\begin{enumerate}
  \item $f$ lies in the approximation space $A^\theta_\tau$ of $n$-term approximation from the
    Haar basis in $L_2[0, 1]$.
  \item $\|f\|_\tau + \bigl\| \bigl( 2^{k\theta} \omega_1(f, 2^{-k}, [0,1])_\tau \bigr)_k
    \bigr\|_{\ell_\tau} < \infty$.
  \item $\|f\|_\tau + \bigl\| \bigl( 2^{k\theta} \omega_1^{\mathrm{per}}(f, 2^{-k})_\tau
    \bigr)_k \bigr\|_{\ell_\tau} < \infty$.
\end{enumerate}
\end{corollary}

\begin{proof}
\leanok
\Cref{thm:haar-nterm-linear}, and \Cref{cor:haar-besov} with $p = q = \tau$ and $s = \theta$,
which applies since $0 < \theta < 1/\tau$.
\end{proof}

The second condition says $f \in \Besov{\theta}{\tau}{\tau}[0, 1]$: best $n$-term approximation
in $L_2$ at the rate $n^{-\theta}$, measured in $\ell_\tau$ with $1/\tau = \theta + 1/2$,
describes this Besov space. For $\tau < 1$ the Besov space lies in the quasi-Banach space
$L_\tau$, outside the range of \Cref{cor:haar-besov}; that case is
\Cref{thm:haar-nterm-besov-all}.

\section{Dyadic step functions for \texorpdfstring{$0 < p \leq 1$}{0 < p <= 1}}
\label{sec:steps-below-one}

For $0 < p < 1$ the space $L_p[0, 1]$ is a quasi-Banach space: in place of the triangle
inequality, $\|f + g\|_p^p \leq \|f\|_p^p + \|g\|_p^p$ (\Cref{lem:lp-rho-triangle}). The modulus
$\omega_r(f, t, [a, b])_p$ is defined for these $p$ by the formula of
\Cref{def:interval-modulus}. Approximation by dyadic step functions in $L_p[0, 1]$ is described by
the $E$-functional of a couple (\Cref{sec:e-functional}), and the coincidence theorem for
$\rho$-normed couples (\Cref{thm:couple-modulus-coincidence}) identifies its approximation spaces
with a Besov space of the first modulus for every $0 < s < 1/p$, as in
Oswald~\cite[Theorem~6]{Oswald1980}, \cite[(7)]{Oswald2018}.

\begin{definition}[The couple of dyadic step functions]
\label{def:step-couple}
\lean{Haar.stepCouple}
\leanok
\uses{def:dyadic-filtration, def:complexity-couple}
Let $0 < p \leq \infty$ and let $V_j \subseteq L_p[0, 1]$ be the subgroup of the
$\mathcal{F}_j$-measurable elements, the classes of the step functions on the dyadic intervals
$I_{j,n}$. The \emph{step couple} in $L_p[0, 1]$ is the couple of the subgroups
$V_0 \subseteq V_1 \subseteq \cdots$ (\Cref{def:complexity-couple}) with the first norm
$\|\cdot\|_p$; its spaces $E_{s, q}$ are the approximation spaces of dyadic step functions.
\end{definition}

The couple uses only the group structure and the quasi-norm of $L_p$, so it covers
$0 < p < 1$, where $L_p$ is not a normed space.

\begin{lemma}[The step couple]
\label{lem:step-couple}
\lean{MeasureTheory.Filtration.monotone_lpMeasSubgroup, Haar.eFunctional_stepCouple_two_pow,
  Haar.mem_lpMeas_of_stepCouple_norm1_le, Haar.stepCouple_eq_complexityCouple,
  Haar.memEMethod_stepCouple_iff_memApprox}
\leanok
\uses{def:step-couple, lem:complexity-couple, def:filtration-scheme,
  thm:complexity-couple-approx}
Let $0 < p \leq \infty$. The subgroups $V_j$ increase with $j$. In the step couple an element of
complexity at most $2^j$ lies in $V_j$, and $E(2^j, f) = \inf_{b \in V_j} \|f - b\|_p$ for every
$f \in L_p[0, 1]$. For $1 \leq p \leq \infty$ the step couple is the complexity couple of the Haar
scheme, so for $0 < q \leq \infty$ and $s > 0$, and for $s = 0$ and $q = \infty$, its space
$E_{s, q}$ is the approximation space $A^s_q$ of the Haar scheme.
\end{lemma}

\begin{proof}
\leanok
For $i \leq j$, $\mathcal{F}_i \subseteq \mathcal{F}_j$, so every $\mathcal{F}_i$-measurable
element is $\mathcal{F}_j$-measurable. The next two claims are (1) and (3) of
\Cref{lem:complexity-couple}. For $p \geq 1$ the dyadic levels of the subgroups $V_j$ are the sets
of the filtration scheme (\Cref{def:filtration-scheme}), and \Cref{thm:complexity-couple-approx}
identifies the spaces.
\end{proof}

\begin{lemma}[Moduli for $0 < p \leq 1$]
\label{lem:modulus-below-one}
\lean{Interval.intervalModulus_add_rpow_le, Interval.eLpNorm_higherDiff_le_of_forall_of_le_one,
  Interval.eLpNorm_higherDiff_le_of_le_one, Interval.intervalModulus_le_of_le_one,
  Interval.intervalLpModulus_add_rpow_le, Interval.intervalLpModulus_le_of_le_one}
\leanok
\uses{def:iter-diff, def:interval-modulus, lem:interval-modulus, lem:lp-rho-triangle,
  thm:pp-jackson}
Let $0 < p \leq 1$, $r \in \N$, $a, b \in \R$ and $t \geq 0$. For $f, g : \R \to \C$
\[
  \omega_r(f + g, t, [a, b])_p^p \leq \omega_r(f, t, [a, b])_p^p + \omega_r(g, t, [a, b])_p^p .
\]
Let $\mu$ be a measure on $\R$, $h \in \R$ and $B \in [0, \infty]$. If
$\|f(\cdot + kh)\|_{L_p(\mu)} \leq B$ for $k = 0, \dots, r$, then
$\|\Delta_h^r f\|_{L_p(\mu)} \leq 2^{r/p} B$. Hence for measurable $f$ and $h \geq 0$,
$\|\Delta_h^r f\|_{L_p[a, b - rh]} \leq 2^{r/p} \|f\|_{L_p[a, b]}$, and
$\omega_r(f, t, [a, b])_p \leq 2^{r/p} \|f\|_{L_p[a, b]}$. Both bounds on $\omega_r$ hold for
$f, g \in L_p[0, 1]$, $[a, b] = [0, 1]$ and $t = 2^{-k}$.
\end{lemma}

\begin{proof}
\leanok
For each step $h$, \Cref{lem:lp-rho-triangle} gives
$\|\Delta_h^r (f + g)\|_p^p \leq \|\Delta_h^r f\|_p^p + \|\Delta_h^r g\|_p^p$; take the supremum.
For the second claim expand $\Delta_h^r f = \sum_{k \leq r} (-1)^{r-k} \binom{r}{k} f(\cdot + kh)$
(\Cref{def:iter-diff}): the $p$-th power of the quasi-norm is subadditive,
$\binom{r}{k}^p \leq \binom{r}{k}$, and $\sum_k \binom{r}{k} = 2^r$. On $[a, b - rh]$ each
translate has norm at most $\|f\|_{L_p[a, b]}$.
\end{proof}

\begin{lemma}[Jackson's inequality for steps by averaging]
\label{lem:step-jackson-below-one}
\lean{MeasureTheory.exists_setLIntegral_enorm_sub_rpow_le, Interval.volume_dyadicIco_eq_inv,
  Interval.exists_ppFun_lintegral_le, Interval.setLIntegral_dyadicIco_le,
  Interval.lintegral_indicator_add_eq_eLpNorm,
  Interval.two_pow_mul_sum_lintegral_le_intervalModulus, Haar.eFunctional_stepCouple_le}
\leanok
\uses{def:step-couple, lem:step-couple, def:interval-modulus, lem:rd-modulus-on, def:pp-scheme,
  lem:pp-bookkeeping, lem:haar-density, cor:haar-besov, thm:pp-jackson}
Let $0 < p < \infty$, let $f : \R \to \C$ be measurable and $j \in \N$. There are $c_n \in \C$
such that the step function $u = \sum_{n < 2^j} c_n \mathbf{1}_{I_{j,n}}$ satisfies
\[
  \int_0^1 |f - u|^p \leq 2^j \sum_{n < 2^j} \int_{I_{j,n}} \int_{I_{j,n}} |f(x) - f(y)|^p
  \, dx \, dy \leq 2 \, \omega_1(f, 2^{-j}, [0, 1])_p^p .
\]
Hence $E(2^j, f) \leq 2^{1/p} \omega_1(f, 2^{-j}, [0, 1])_p$ in the step couple for every
$f \in L_p[0, 1]$.
\end{lemma}

\begin{proof}
\leanok
Let $D_n(y) = \int_{I_{j,n}} |f(x) - f(y)|^p \, dx$. Some $y_n \in I_{j,n}$ has $D_n(y_n)$ at most
the mean $2^j \int_{I_{j,n}} D_n$ of $D_n$ over $I_{j,n}$ (the first moment method, Mathlib's
\texttt{exists\_\allowbreak le\_\allowbreak setLAverage}); take $c_n = f(y_n)$. For the second
inequality substitute $x = y + h$: for $x, y \in I_{j,n}$, $|h| < 2^{-j}$ and $x \in [0, 1]$, so
the sum is at most
$\int_0^1 \int_{-2^{-j}}^{2^{-j}} \mathbf{1}_{[0, 1]}(y + h) \, |f(y + h) - f(y)|^p \, dh \, dy$.
After exchanging the integrals, the integral over $y$ is $\|\Delta_h f\|_{L_p([0, 1]_h)}^p \leq
\omega_1(f, 2^{-j}, [0, 1])_p^p$ for each $h$, and the steps fill an interval of length
$2 \cdot 2^{-j}$. The step function $u$ is $\mathcal{F}_j$-measurable
(\Cref{lem:haar-density}, \Cref{cor:haar-besov}), so it lies in $V_j$.
\end{proof}

The proof averages pointwise and uses Tonelli's theorem, so it holds for every
$0 < p < \infty$ with the constant $2^{1/p}$.

\begin{lemma}[The inverse inequality for steps, $0 < p \leq 1$]
\label{lem:step-inverse-below-one}
\lean{Interval.enorm_higherDiff_one_rpow_le_stripBound, Interval.lintegral_rpow_eq_of_step,
  Interval.eLpNorm_higherDiff_one_le_of_step, Haar.intervalLpModulus_le_of_mem_lpMeas}
\leanok
\uses{def:step-couple, def:pp-scheme, lem:pp-bookkeeping, lem:pp-local,
  lem:modulus-below-one, cor:haar-besov, thm:pp-jackson, thm:pp-inverse}
Let $0 < p \leq 1$ and $j \in \N$, let $u : \R \to \C$ take the value $v_n$ on $I_{j,n}$ for
every $n \in \N$, with $v_{2^j} = 0$, and let $0 \leq h \leq 2^{-j}$. For $x \geq 0$ in
$I_{j,n}$, $\Delta_h u(x) = 0$ unless $x$ lies in the strip $[(n+1) 2^{-j} - h, (n+1) 2^{-j})$, and
there $|\Delta_h u(x)|^p \leq |v_n|^p + |v_{n+1}|^p$. Hence
\[
  \|\Delta_h u\|_{L_p[0, 1]} \leq \bigl(2 \cdot 2^j h\bigr)^{1/p} \, \|u\|_{L_p[0, 1]} .
\]
Every $a \in V_j$ satisfies, for every $k \in \N$,
\[
  \omega_1(a, 2^{-k}, [0, 1])_p \leq 2^{1/p} \, 2^{-\max(k - j, 0)/p} \, \|a\|_p .
\]
\end{lemma}

\begin{proof}
\leanok
In the strip of $I_{j,n}$, $x + h \in I_{j, n+1}$, so $\Delta_h u(x) = v_{n+1} - v_n$ and
$(\alpha + \beta)^p \leq \alpha^p + \beta^p$; outside it, $x$ and $x + h$ lie in $I_{j,n}$. The
strip integral of \Cref{lem:pp-local} with $r = 1$, which holds for every $0 < p < \infty$, gives
$\int_0^1 |\Delta_h u|^p \leq 2 h \sum_{n < 2^j} |v_n|^p$, since each $v_n$ occurs in at most two
terms, and $\|u\|_{L_p[0, 1]}^p = 2^{-j} \sum_{n < 2^j} |v_n|^p$. For $k < j$ the bound on $a$ is
$\omega_1(a, t, [0, 1])_p \leq 2^{1/p} \|a\|_p$ (\Cref{lem:modulus-below-one}). For $k \geq j$,
$a$ agrees almost everywhere on $[0, 1]$ with a step function $u$ of level $j$ that vanishes
outside $[0, 1)$ (\Cref{cor:haar-besov}), and $2^j h \leq 2^{-(k - j)}$ for $h \leq 2^{-k}$.
\end{proof}

For $p < 1$ the order $1/p$ exceeds the order $\min(r, r - 1 + 1/p)$ of DeVore and
Popov~\cite[Theorem~4.8]{DeVorePopovTAMS}, which is $1$ for $r = 1$ and $p \leq 1$: a first
difference of a step function sees only its jumps.

\begin{theorem}[Dyadic steps and the first modulus for $0 < p \leq 1$]
\label{thm:steps-besov-below-one}
\lean{Haar.stepModulusData, Haar.eNormDyadic_stepCouple_lt_top_iff_modulus,
  Haar.memEMethod_stepCouple_iff_modulus}
\leanok
\uses{def:step-couple, lem:step-couple, lem:lp-rho-triangle, lem:complexity-couple,
  lem:modulus-below-one, lem:step-jackson-below-one, lem:step-inverse-below-one,
  def:couple-modulus-data, thm:couple-modulus-coincidence}
Let $0 < p \leq 1$, $0 < q \leq \infty$ and $f \in L_p[0, 1]$. For every $0 \leq s < 1/p$,
\[
  \|f\|_p + \bigl\| \bigl( 2^{ks} E(2^k, f) \bigr)_k \bigr\|_{\ell_q} < \infty
\]
if and only if
\[
  \|f\|_p + \bigl\| \bigl( 2^{ks} \omega_1(f, 2^{-k}, [0, 1])_p \bigr)_k \bigr\|_{\ell_q}
  < \infty .
\]
For every $0 < s < 1/p$, $f$ lies in the approximation space $E_{s, q}$ of the step couple if and
only if the second condition holds.
\end{theorem}

\begin{proof}
\leanok
By \Cref{lem:modulus-below-one}, \Cref{lem:step-jackson-below-one} and
\Cref{lem:step-inverse-below-one}, $\omega_1(\cdot, 2^{-k}, [0, 1])_p$ is a modulus compatible
with the step couple of exponent $\rho = p$ and inverse order $\theta = 1/p$
(\Cref{def:couple-modulus-data}), and only $0$ has complexity below $1$. Apply
\Cref{thm:couple-modulus-coincidence}.
\end{proof}

The condition defines the Besov space $B^s_{p,q,1}[0, 1]$ of the first
modulus~\cite[(7)]{Oswald2018}, \cite[Theorem~6]{Oswald1980}. For $s < 1$ it is the Besov space
of DeVore and Popov~\cite[(2.4)]{DeVorePopovTAMS} with $r = 1$, whose coincidence
theorem~\cite[Theorem~5.1]{DeVorePopovTAMS} reaches only $s < 1$ for $r = 1$. For
$1/p - 1 < s < 1$ and $q < \infty$ it is also the Fourier-analytic Besov space
$\Besov{s}{p}{q}[0, 1]$~\cite[(2)]{Oswald2018}. For $s > 1$, and for $s = 1$ with $q < \infty$, it
differs from the space defined by the second modulus: the line $x \mapsto x$ has vanishing second
modulus and lies outside it (\Cref{lem:steps-line}).

\begin{lemma}[The line]
\label{lem:steps-line}
\lean{Interval.intervalModulus_two_id, Haar.idLp, Haar.intervalLpModulus_two_idLp,
  Interval.le_lintegral_dyadicIco_enorm_sub_rpow, Haar.le_eFunctional_stepCouple_idLp,
  Haar.not_memEMethod_stepCouple_idLp}
\leanok
\uses{lem:interval-modulus, lem:interval-nikolskii, def:pp-scheme, lem:pp-bookkeeping,
  cor:haar-besov, thm:pp-jackson, def:step-couple, lem:step-couple, lem:complexity-couple,
  lem:e-method-dyadic, def:lq-norm, lem:lq-basic, lem:dyadic-interpolation,
  lem:daub-dyadic-scales}
Let $\ell(x) = x$. Then $\omega_2(\ell, t, [0, 1])_p = 0$ for every $0 < p \leq \infty$ and
$t \geq 0$. For $0 < p < \infty$ and $j \in \N$, every step function $b$ of level $j$ has
$\|\ell - b\|_p^p \geq \tfrac12 (2^{-j}/4)^p$, so $E(2^j, \ell) \geq 2^{-1/p} \, 2^{-j}/4$ in the
step couple. Hence for $0 < p < \infty$, $0 < q \leq \infty$ and $s \geq 1$ with $s > 1$ or
$q < \infty$, $\ell \notin E_{s, q}$.
\end{lemma}

\begin{proof}
\leanok
The second modulus vanishes on polynomials of degree $1$ (\Cref{lem:interval-modulus}). For
$c \in \C$ at most half of $I_{j,n}$ lies within $2^{-j}/4$ of $\operatorname{Re} c$, so
$\int_{I_{j,n}} |x - c|^p \, dx \geq (2^{-j}/4)^p \, 2^{-j}/2$; sum over the $2^j$ intervals. By
\Cref{lem:e-method-dyadic}, $\ell \in E_{s,q}$ would bound the $\ell_q$ norm of
$2^{ks} E(2^k, \ell) \geq 2^{-1/p} \, 2^{k(s - 1)}/4$ by the finite dyadic norm. For $s > 1$ these
terms are unbounded (\Cref{lem:dyadic-interpolation}); for $s = 1$ and $q < \infty$ they are
bounded below by a positive constant, so their $\ell_q$ norm is infinite.
\end{proof}

On $\R^d$, identifying the approximation spaces of entire functions of exponential type with Besov
spaces for $0 < p < 1$ rests on Jackson's inequality in $L_p$, $p < 1$
\cite{StorozhenkoOswald1978, KolomoitsevTikhonov2020}; that identification lies outside this
project.

\section{The Haar system for \texorpdfstring{$0 < p \leq 1$}{0 < p <= 1}}
\label{sec:haar-below-one}

For $s > 1/p - 1$ the approximation spaces of the step couple in $L_p[0, 1]$, $0 < p \leq 1$, lie
in $L_1[0, 1]$, so Haar coefficients are defined on them, and they describe these spaces. For
$1/p - 1 < s < 1/p$ this describes $B^s_{p,q,1}[0, 1]$ by Haar coefficients and gives the
Schauder basis property of the Haar system there, which is due to Oswald~\cite{Oswald1981} (see
\cite[\S1 and Theorem~1]{Oswald2018}). Throughout, $0 < p \leq 1$, and $E(t, f)$ is the
$E$-functional of the step couple (\Cref{def:step-couple}).

\begin{lemma}[Fatou's lemma for telescoping sums]
\label{lem:telescoping-fatou}
\lean{MeasureTheory.TendstoInMeasure.sub_right, MeasureTheory.eLpNorm_sub_rpow_le_sum,
  MeasureTheory.eLpNorm_sub_rpow_le_tsum}
\leanok
\uses{lem:lp-rho-triangle}
Let $\mu$ be a measure, $0 < r \leq 1$, and let $u_0, u_1, \dots$ and $g$ be strongly measurable
functions with values in a normed group such that $u_n \to g$ in measure. Then for every $k \in \N$
\[
  \|g - u_k\|_{L_r(\mu)}^r \leq \sum_{i=0}^\infty \|u_{k+i+1} - u_{k+i}\|_{L_r(\mu)}^r .
\]
\end{lemma}

\begin{proof}
\leanok
By \Cref{lem:lp-rho-triangle}, $\|u_{k+n} - u_k\|_r^r \leq \sum_{i<n} \|u_{k+i+1} - u_{k+i}\|_r^r$
for every $n$. The functions $u_{k+n} - u_k$ converge in measure to $g - u_k$, and a bound of the
$L_r$ quasi-norms passes to the limit (Mathlib's
\texttt{eLpNorm\_\allowbreak le\_\allowbreak of\_\allowbreak tendstoInMeasure}, Fatou's lemma
along an almost everywhere convergent subsequence).
\end{proof}

For $r = 1$ this is the triangle inequality for a series in $L_1(\mu)$.

\begin{lemma}[Weighted tails]
\label{lem:weighted-tails}
\lean{ENNReal.le_dyadicWeight_neg_mul_lqNorm, ENNReal.tendsto_zero_of_lqNorm_dyadicWeight_ne_top,
  ENNReal.lqNorm_rpow_le_of_rpow_le_tsum, ENNReal.le_two_rpow_mul_of_rpow_le}
\leanok
\uses{def:lq-norm, lem:lq-basic, lem:hardy-every-exponent, lem:lp-rho-triangle}
Let $q \in (0, \infty]$, $s \in \R$ and let $e$, $x$, $y$ be sequences in $[0, \infty]$.
\begin{enumerate}
  \item $e_j \leq 2^{-js} \|(2^{ks} e_k)_k\|_{\ell_q}$ for every $j$; if $s > 0$ and the norm is
    finite, then $e_j \to 0$.
  \item Let $\rho > 0$, $\beta \in \R$ and $c \in [0, \infty]$. If
    $x_k^\rho \leq c \sum_{i \geq 0} 2^{i\beta} y_{k+i}^\rho$ for every $k$, then
    \[
      \|(2^{ks} x_k)_k\|_{\ell_q}^\rho \leq c \, H_{s\rho - \beta, q/\rho} \,
      \|(2^{ks} y_k)_k\|_{\ell_q}^\rho ,
    \]
    with $H$ as in \Cref{lem:hardy-every-exponent}.
  \item For $0 < p < \infty$ and $a, b \in [0, \infty]$, $a^p \leq 2 b^p$ implies
    $a \leq 2^{1/p} b$.
\end{enumerate}
\end{lemma}

\begin{proof}
\leanok
(1) $2^{js} e_j$ is one term of the sequence. (2) Multiply by $2^{ks\rho}$: with
$\alpha = s\rho - \beta$ and $b_j = 2^{j\beta} y_j^\rho$ the right side is
$c \, 2^{k\alpha} \sum_{i \geq 0} b_{k+i}$, and \Cref{lem:hardy-every-exponent} in
$\ell_{q/\rho}$ applies, since $\|(z_k^\rho)_k\|_{\ell_{q/\rho}} = \|z\|_{\ell_q}^\rho$
(\Cref{lem:lp-rho-triangle}).
\end{proof}

\begin{lemma}[Steps in $L_1$]
\label{lem:steps-into-l1}
\lean{Haar.eq_of_dyadicIndex_eq, Haar.memLp_of_stronglyMeasurable_dyadicFiltration,
  Haar.memLp_of_mem_lpMeas, Haar.mem_lpMeas_of_le, Haar.sum_setLIntegral_haarDyadicInterval,
  Haar.setLIntegral_enorm_le_of_dyadicIndex, Haar.eLpNorm_one_le_of_mem_lpMeas}
\leanok
\uses{def:dyadic-filtration, lem:condexp-partition, cor:haar-besov, def:pp-scheme,
  lem:pp-bookkeeping, def:step-couple}
Let $j \in \N$ and let $G$ be constant on every $I_{j,n}$, $n \in \Z$. Let $S$ be a measurable
union of such intervals. Then
\[
  \int_{S \cap [0, 1]} |G| \leq 2^{j(1/p - 1)} \Bigl( \int_{S \cap [0, 1]} |G|^p \Bigr)^{1/p} .
\]
Every $\mathcal{F}_j$-measurable $a \in L_p[0, 1]$ lies in $L_r[0, 1]$ for every
$0 < r \leq \infty$, and $\|a\|_1 \leq 2^{j(1/p - 1)} \|a\|_p$.
\end{lemma}

\begin{proof}
\leanok
For $x \in S \cap [0, 1)$ the interval $I$ of level $j$ containing $x$ lies in $S$, so
$|G(x)|^p 2^{-j} = \int_I |G|^p \leq J$ with $J = \int_{S \cap [0, 1]} |G|^p$. Hence
$|G| = |G|^p |G|^{1-p} \leq |G|^p (2^j J)^{(1-p)/p}$ on $S \cap [0, 1)$; integrate. An
$\mathcal{F}_j$-measurable function is constant on every $I_{j,n}$
(\Cref{lem:condexp-partition}), and on $[0, 1]$ it is a step function with $2^j$ values
(\Cref{cor:haar-besov}).
\end{proof}

With the values $c_n$ of $a$ on the intervals of level $j$, the last inequality reads
$2^{-j} \sum_n |c_n| \leq 2^{-j} (\sum_n |c_n|^p)^{1/p}$.

\begin{definition}[Level norms of the Haar expansion]
\label{def:haar-level-norm}
\lean{Haar.haarLevelNorm, Haar.haarLevelNorm_congr_ae, Haar.haarLevelNorm_rpow,
  Haar.haarLevelNorm_eq_eLpNorm}
\leanok
\uses{def:haar, lem:haar-martingale, lem:haar-level}
For $f : \R \to \C$ and $k \in \N$ let
\[
  H_k(f) = 2^{k(1/2 - 1/p)} \bigl\| (\langle f, h_{k,m} \rangle)_{0 \leq m < 2^k}
  \bigr\|_{\ell_p} .
\]
It depends only on the class of $f$ in $L_1[0, 1]$, and for integrable $f$,
$H_k(f) = \|E[f \mid \mathcal{F}_{k+1}] - E[f \mid \mathcal{F}_k]\|_p$.
\end{definition}

The equality is \Cref{lem:haar-martingale} and \Cref{lem:haar-level}.

\begin{lemma}[Properties of the level norms]
\label{lem:haar-level-norm}
\lean{Haar.integrable_haarFunction_mul, Haar.haarCoeff_add, Haar.haarCoeff_sub,
  Haar.enorm_haarCoeff_le, Haar.haarCoeff_eq_zero_of_dyadicIndex,
  Haar.haarCoeff_eq_zero_of_mem_lpMeas, Haar.haarLevelNorm_le_lintegral,
  Haar.haarLevelNorm_add_rpow_le, Haar.haarLevelNorm_le_of_mem_lpMeas,
  Haar.haarLevelNorm_eq_zero_of_mem_lpMeas}
\leanok
\uses{def:haar-level-norm, def:haar, prop:haar-basis, lem:steps-into-l1, def:lq-norm,
  lem:lq-basic}
Let $k \in \N$. For every $f$ and $m$,
$|\langle f, h_{k,m} \rangle| \leq 2^{k/2} \int_{I_{k,m}} |f|$, so
$H_k(f) \leq 2^k \int_0^1 |f|$. For integrable $f$ and $g$,
$H_k(f + g)^p \leq H_k(f)^p + H_k(g)^p$. For $j \in \N$ and an $\mathcal{F}_j$-measurable
$a \in L_p[0, 1]$, $H_k(a) = 0$ if $j \leq k$, and
\[
  H_k(a) \leq 2^{(j - k)(1/p - 1)} \|a\|_p \qquad \text{if } k \leq j .
\]
\end{lemma}

\begin{proof}
\leanok
The Haar function $h_{k,m}$ vanishes outside $I_{k,m}$ and has modulus $2^{k/2}$ on it. The
coefficients are linear on integrable functions, and $|\alpha + \beta|^p \leq |\alpha|^p +
|\beta|^p$. A function constant on $I_{k,m}$ has $\langle f, h_{k,m} \rangle = 0$, since
$h_{k,m}$ has mean zero (\Cref{prop:haar-basis}). For $k \leq j$, \Cref{lem:steps-into-l1} on
$I_{k,m}$ gives $(\int_{I_{k,m}} |a|)^p \leq 2^{j(1-p)} \int_{I_{k,m}} |a|^p$; summing over $m$,
\[
  H_k(a)^p \leq 2^{k(p/2 - 1)} \, 2^{kp/2} \, 2^{j(1-p)} \int_0^1 |a|^p
  = 2^{(j - k)(1 - p)} \|a\|_p^p .
\]
\end{proof}

\begin{lemma}[Near-best steps]
\label{lem:near-best-steps}
\lean{Haar.enorm_sub_rpow_le_of_nearBest, Haar.eLpNorm_one_sub_le_of_nearBest,
  Haar.tendstoInMeasure_of_nearBest, Haar.eLpNorm_one_sub_le_tsum_of_nearBest,
  Haar.eLpNorm_one_le_of_nearBest, Haar.integrable_of_nearBest,
  Haar.haarLevelNorm_rpow_le_tsum_of_nearBest, Haar.exists_nearBest_stepCouple,
  Haar.antitone_eFunctional_stepCouple_add}
\leanok
\uses{lem:steps-into-l1, lem:haar-level-norm, lem:telescoping-fatou, lem:weighted-tails,
  lem:lp-rho-triangle, lem:step-couple, lem:e-functional}
Let $f \in L_p[0, 1]$ and $s \geq -1$. There are $\mathcal{F}_j$-measurable $b_j \in L_p[0, 1]$
with $\|f - b_j\|_p \leq e_j = E(2^j, f) + \|f\|_p 2^{-j(s+1)}$, and $e$ decreases. Let
$b_j$ be $\mathcal{F}_j$-measurable with $\|f - b_j\|_p \leq e_j$ for a decreasing sequence $e$
with $e_j \to 0$. Then $\|b_{j+1} - b_j\|_p^p \leq 2 e_j^p$, $b_j \to f$ in measure, and for every
$k$
\[
  \|f - b_k\|_1 \leq \sum_{i=0}^\infty 2^{(k+i+1)(1/p - 1)} 2^{1/p} e_{k+i} .
\]
If moreover $\|f\|_p \leq e_0$, then
$\|f\|_1 \leq 2^{1/p} e_0 + \sum_{i \geq 0} 2^{(i+1)(1/p - 1)} 2^{1/p} e_i$. If this series
converges, then $f$ is integrable and for every $k$
\[
  H_k(f)^p \leq \sum_{i=0}^\infty \bigl(2^{(i+1)(1/p - 1)} 2^{1/p} e_{k+i}\bigr)^p .
\]
\end{lemma}

\begin{proof}
\leanok
Near-best elements exist by \Cref{lem:e-functional}; for $\|f\|_p = 0$ take $b_j = 0$. The
difference $b_{j+1} - b_j = (f - b_j) - (f - b_{j+1})$ is $\mathcal{F}_{j+1}$-measurable, so
\Cref{lem:steps-into-l1} bounds its $L_1$ norm, and \Cref{lem:telescoping-fatou} with $r = 1$
gives the bound of $\|f - b_k\|_1$; at $k = 0$ it shows $f - b_0 \in L_1$, so $f$ is integrable
when the series converges. For the level norms, telescope
$b_{k+n} = b_k + \sum_{i<n} (b_{k+i+1} - b_{k+i})$ with $H_k(b_k) = 0$ and
\Cref{lem:haar-level-norm} for each block, which is $\mathcal{F}_{k+i+1}$-measurable; then
$H_k(f)^p \leq H_k(f - b_{k+n})^p + H_k(b_{k+n})^p$, and
$H_k(f - b_{k+n}) \leq 2^k \|f - b_{k+n}\|_1 \to 0$ as $n \to \infty$.
\end{proof}

\begin{theorem}[The approximation spaces of steps embed into $L_1$]
\label{thm:besov-into-l1}
\lean{Haar.tsum_dyadicWeight_succ_mul_le_lqNorm, Haar.stepCoupleEmbedConst,
  Haar.stepCoupleEmbedConst_ne_top,
  Haar.pos_of_one_div_sub_one_lt, Haar.eLpNorm_one_le_eNormDyadic,
  Haar.integrable_of_eNormDyadic_lt_top}
\leanok
\uses{lem:near-best-steps, lem:weighted-tails, thm:couple-modulus-coincidence,
  lem:e-method-dyadic, lem:dyadic-interpolation}
Let $s > 1/p - 1$ and $q \in (0, \infty]$, and let $\delta = s - (1/p - 1)$. Every
$f \in L_p[0, 1]$ satisfies
\[
  \|f\|_1 \leq C \Bigl( \|f\|_p + \bigl\| \bigl( 2^{ks} E(2^k, f) \bigr)_k \bigr\|_{\ell_q} \Bigr)
\]
with $C = 2^{1/p} \bigl(1 + 2^{1/p - 1} (1 - 2^{-\delta})^{-1}\bigr) \, C_q (1 + \gamma_{1,q})$,
$C_q$ and $\gamma_{1,q}$ as in \Cref{thm:couple-modulus-coincidence}. In particular every element
of $E_{s,q}$ of the step couple is integrable.
\end{theorem}

\begin{proof}
\leanok
Since $p \leq 1$, $s > 0$. Take the near-best steps of \Cref{lem:near-best-steps}. By
\Cref{thm:couple-modulus-coincidence},
$L = \|(2^{ks} e_k)_k\|_{\ell_q} \leq C_q (1 + \gamma_{1,q})$ times the dyadic norm, and by
\Cref{lem:weighted-tails} $e_j \leq 2^{-js} L$, so $e_j \to 0$ and
$\sum_i 2^{(i+1)(1/p-1)} 2^{1/p} e_i \leq 2^{1/p} 2^{1/p-1} (1 - 2^{-\delta})^{-1} L$. With
$e_0 \leq L$, \Cref{lem:near-best-steps} gives the bound.
\end{proof}

\begin{lemma}[Martingale convergence in $L_1$]
\label{lem:haar-martingale-l1}
\lean{Haar.tendsto_eLpNorm_condExp_sub, Haar.ae_eq_zero_of_haarCoeff_eq_zero,
  Haar.eLpNorm_sub_condExp_rpow_le_tsum, Haar.eFunctional_stepCouple_le_eLpNorm_sub_condExp}
\leanok
\uses{def:filtration-scheme, lem:lebesgue, lem:haar-density, lem:dyadic-condexp,
  lem:haar-martingale, lem:telescoping-fatou, def:haar-level-norm, lem:step-couple,
  lem:steps-into-l1}
For integrable $f$, $E[f \mid \mathcal{F}_N] \to f$ in $L_1[0, 1]$. If moreover $\int_0^1 f = 0$
and $\langle f, h_{k,m} \rangle = 0$ for all $k$ and $0 \leq m < 2^k$, then $f = 0$ almost
everywhere on $[0, 1]$. For integrable $f$ and every $k$,
\[
  \|f - E[f \mid \mathcal{F}_k]\|_p^p \leq \sum_{j=0}^\infty H_{k+j}(f)^p ,
\]
and $E(2^k, f) \leq \|f - E[f \mid \mathcal{F}_k]\|_p$ for every $f \in L_p[0, 1]$.
\end{lemma}

\begin{proof}
\leanok
The conditional expectations form a projection scheme of bound $1$ in $L_1$
(\Cref{def:filtration-scheme}) with dense approximants (\Cref{lem:haar-density}), so they converge
(\Cref{lem:lebesgue}). Under the hypotheses on the coefficients every $E[f \mid \mathcal{F}_N]$
vanishes, by \Cref{lem:dyadic-condexp} and \Cref{lem:haar-martingale}, so $\|f\|_1 = 0$. The
convergence in $L_1$ is convergence in measure, and \Cref{lem:telescoping-fatou} with $r = p$
applies to the martingale differences (\Cref{def:haar-level-norm}). Finally
$E[f \mid \mathcal{F}_k]$ is $\mathcal{F}_k$-measurable and lies in $L_p$
(\Cref{lem:steps-into-l1}).
\end{proof}

\begin{theorem}[Haar coefficients for $0 < p \leq 1$]
\label{thm:haar-below-one}
\lean{Haar.dyadicWeight_succ_mul_rpow, Haar.lqNorm_haarLevelNorm_lt_top,
  Haar.eNormDyadic_lt_top_of_lqNorm_haarLevelNorm, Haar.dyadicWeight_mul_lqNorm_haarCoeff,
  Haar.memEMethod_stepCouple_iff_haarCoeff}
\leanok
\uses{thm:besov-into-l1, lem:near-best-steps, lem:haar-martingale-l1, lem:weighted-tails,
  lem:e-method-dyadic, def:haar-level-norm, thm:couple-modulus-coincidence}
Let $s > 1/p - 1$ and $q \in (0, \infty]$. Then $f \in L_p[0, 1]$ lies in the approximation space
$E_{s,q}$ of the step couple if and only if $f \in L_1[0, 1]$ and
\[
  \Bigl| \int_0^1 f \Bigr| + \Bigl\| \Bigl(2^{k(s + 1/2 - 1/p)}
  \bigl\| (\langle f, h_{k,m} \rangle)_{0 \leq m < 2^k} \bigr\|_{\ell_p} \Bigr)_k
  \Bigr\|_{\ell_q} < \infty .
\]
\end{theorem}

\begin{proof}
\leanok
The weighted sequence is $(2^{ks} H_k(f))_k$. If $f \in E_{s,q}$, then $f$ is integrable
(\Cref{thm:besov-into-l1}), and the near-best steps of \Cref{lem:near-best-steps} give
$H_k(f)^p \leq 2 \cdot 2^{1-p} \sum_{i \geq 0} 2^{i(1-p)} e_{k+i}^p$; \Cref{lem:weighted-tails}
with $\rho = p$, $\beta = 1 - p$ and $s p - (1 - p) > 0$, which is $s > 1/p - 1$, bounds
$\|(2^{ks} H_k(f))_k\|_{\ell_q}$ by $\|(2^{ks} e_k)_k\|_{\ell_q} < \infty$. Conversely, by
\Cref{lem:haar-martingale-l1}, $E(2^k, f)^p \leq \sum_{j \geq 0} H_{k+j}(f)^p$, and
\Cref{lem:weighted-tails} with $\beta = 0$ and $s p > 0$ bounds
$\|(2^{ks} E(2^k, f))_k\|_{\ell_q}$; finally use \Cref{lem:e-method-dyadic}.
\end{proof}

The condition $s > 1/p - 1$ cannot be dropped: for $0 < s < 1/p - 1$ and $q < \infty$, the space
$B^s_{p,q,1}[0, 1]$ has a trivial dual~\cite[\S1 and Theorem~1~(i)]{Oswald2018}; the limiting case
$s = 1/p - 1$ is the subject of~\cite{Oswald2018}.

\begin{corollary}[Besov spaces of the first modulus through Haar coefficients]
\label{cor:haar-besov-below-one}
\lean{Haar.lqNorm_modulus_lt_top_iff_haarCoeff_of_le_one}
\leanok
\uses{thm:haar-below-one, thm:steps-besov-below-one}
Let $1/p - 1 < s < 1/p$ and $q \in (0, \infty]$. For $f \in L_p[0, 1]$,
\[
  \|f\|_p + \bigl\| \bigl( 2^{ks} \omega_1(f, 2^{-k}, [0,1])_p \bigr)_k \bigr\|_{\ell_q} < \infty
\]
if and only if $f \in L_1[0, 1]$ and the Haar condition of \Cref{thm:haar-below-one} holds.
\end{corollary}

\begin{proof}
\leanok
Both conditions say $f \in E_{s,q}$, by \Cref{thm:steps-besov-below-one} and
\Cref{thm:haar-below-one}.
\end{proof}

For $p = 1$ and $q \geq 1$ this is also \Cref{cor:haar-besov} with $p = 1$.

\begin{definition}[Haar partial sums]
\label{def:haar-partial-sums}
\lean{Haar.haarSystemCoeff, Haar.haarSystemCoeff_zero, Haar.haarSystemCoeff_two_pow_add,
  Haar.haarSystemCoeff_congr_ae, Haar.haarSystemLp, Haar.coeFn_sum_smul_haarSystemLp,
  Haar.integrable_haarSystem_mul, Haar.integrable_sum_mul_haarSystem,
  Haar.haarPartialSum, Haar.coeFn_haarPartialSum, Haar.stronglyMeasurable_haarSystem,
  Haar.sum_smul_haarSystemLp_mem, Haar.haarPartialSum_mem}
\leanok
\uses{prop:haar-basis, def:step-couple}
With the Haar system $\varphi_0 = 1$, $\varphi_{2^k + i} = h_{k,i}$ of \Cref{prop:haar-basis},
the coefficients of $f \in L_1[0, 1]$ are $c_n(f) = \int_0^1 \varphi_n f$, so
$c_0(f) = \int_0^1 f$ and $c_{2^k + i}(f) = \langle f, h_{k,i} \rangle$, and the partial sums are
$S_N f = \sum_{n < N} c_n(f) \varphi_n \in L_p[0, 1]$. For $N \leq 2^l$ every combination of
$\varphi_0, \dots, \varphi_{N-1}$ is $\mathcal{F}_l$-measurable.
\end{definition}

\begin{lemma}[Remainders of Haar partial sums]
\label{lem:haar-partial-sums}
\lean{Haar.sum_haarSystem_two_pow_ae_eq, Haar.sum_fin_ite_lt, Haar.sub_sum_haarSystem_ae_eq,
  Haar.eLpNorm_sum_haar_ite_le, Haar.eLpNorm_sub_sum_haarSystem_rpow_le,
  Haar.eFunctional_stepCouple_sub_of_mem}
\leanok
\uses{def:haar-partial-sums, lem:dyadic-condexp, lem:haar-martingale, lem:haar-level,
  lem:haar-martingale-l1, def:haar-level-norm, lem:lp-rho-triangle, lem:complexity-couple}
Let $f$ be integrable, $k \in \N$ and $0 \leq i \leq 2^k$. Then
$S_{2^k} f = E[f \mid \mathcal{F}_k]$,
\[
  f - S_{2^k + i} f = \bigl(f - E[f \mid \mathcal{F}_{k+1}]\bigr) +
  \sum_{i \leq m < 2^k} \langle f, h_{k,m} \rangle h_{k,m} ,
\]
and $\|f - S_{2^k + i} f\|_p^p \leq \sum_{j \geq 0} H_{k+j}(f)^p$. For $f \in L_p[0, 1]$ and
an $\mathcal{F}_l$-measurable $a \in L_p[0, 1]$, $E(2^l, f - a) = E(2^l, f)$.
\end{lemma}

\begin{proof}
\leanok
By induction on $k$ with \Cref{lem:dyadic-condexp} and \Cref{lem:haar-martingale}. The part of
the Haar sum of level $k$ has quasi-norm at most $H_k(f)$ (\Cref{lem:haar-level}), and
$\|f - E[f \mid \mathcal{F}_{k+1}]\|_p^p \leq \sum_{j \geq 1} H_{k+j}(f)^p$
(\Cref{lem:haar-martingale-l1}). The last claim is (3) of \Cref{lem:complexity-couple}.
\end{proof}

\begin{theorem}[The Haar system as a Schauder basis for $0 < p \leq 1$]
\label{thm:haar-schauder-below-one}
\lean{Haar.eNormDyadic_sub_haarPartialSum_le, Haar.tendsto_dyadicWeight_mul_tsum_haarLevelNorm,
  Haar.eNormDyadic_sub_haarPartialSum_le_of_le_one, Haar.tendsto_eNormDyadic_sub_haarPartialSum,
  Haar.exists_enorm_haarSystemCoeff_le, Haar.haarSystemCoeff_sub_sum,
  Haar.eq_haarSystemCoeff_of_tendsto}
\leanok
\uses{lem:haar-partial-sums, thm:haar-below-one, thm:besov-into-l1, lem:weighted-tails,
  lem:hardy-every-exponent, lem:lq-quasi-triangle, lem:lq-reindex, lem:mixed-norm-tails,
  lem:e-method-dyadic, prop:haar-basis}
Let $s > 1/p - 1$ and $0 < q < \infty$, and write
$\|g\|_{E} = \|g\|_p + \|(2^{ks} E(2^k, g))_k\|_{\ell_q}$. For every $f \in E_{s,q}$,
$\|f - S_n f\|_E \to 0$. If $a_0, a_1, \dots \in \C$ and
$\|f - \sum_{n < N} a_n \varphi_n\|_E \to 0$ for some $f \in L_p[0, 1]$, then $a_n = c_n(f)$ for
every $n$.
\end{theorem}

\begin{proof}
\leanok
Let $2^k \leq n < 2^{k+1}$ and $R = f - S_n f$, and let
$W_k = 2^{ks} (\sum_{j \geq 0} H_{k+j}(f)^p)^{1/p}$. By \Cref{lem:haar-partial-sums},
$\|R\|_p \leq W_k$, $E(2^l, R) \leq \|R\|_p$ for $l \leq k$ and $E(2^l, R) = E(2^l, f)$ for
$l > k$, since $S_n f$ is $\mathcal{F}_{k+1}$-measurable. As
$\|(2^{ls})_{l \leq k}\|_{\ell_q} \leq 2^{ks} \gamma_{s,q}$ (\Cref{lem:lq-reindex}),
\[
  \|R\|_E \leq W_k + C_q \Bigl( \gamma_{s,q} W_k + \bigl\| \bigl( 2^{ls} E(2^l, f)
  \bigr)_{l > k} \bigr\|_{\ell_q} \Bigr) .
\]
By \Cref{thm:haar-below-one} and Hardy's inequality for tails in $\ell_{q/p}$ with the exponent
$sp$ (\Cref{lem:weighted-tails}), $(W_k)_k \in \ell_q$, so $W_k \to 0$ as $q < \infty$, and
the tails of the $\ell_q$ norm of $(2^{ls} E(2^l, f))_l$ tend to $0$
(\Cref{lem:mixed-norm-tails}). For the uniqueness, the
$L_1$ norm is at most $C \|\cdot\|_E$ (\Cref{thm:besov-into-l1}), $c_i$ is bounded on $L_1$, and
$c_i(f - \sum_{n<N} a_n \varphi_n) = c_i(f) - a_i$ for $N > i$ by orthonormality.
\end{proof}

For $1/p - 1 < s < 1/p$ the space $E_{s,q}$ is $B^s_{p,q,1}[0, 1]$
(\Cref{cor:haar-besov-below-one}), and the theorem is the Schauder basis property of the Haar
system in it~\cite{Oswald1981}, \cite[Theorem~1]{Oswald2018}. For $1/2 < p < 1$,
$1/p - 1 < s < 1$ and $q < \infty$ the Haar system is also a Schauder basis of the
Fourier-analytic $\Besov{s}{p}{q}[0, 1]$~\cite{TriebelHaar1978},
\cite[Theorem~2]{Oswald2018}; that space is not formalised here.

\begin{theorem}[$n$-term Haar approximation and Besov spaces for $0 < \tau < 2$]
\label{thm:haar-nterm-besov-all}
\lean{Haar.ofReal_le_two, Haar.mem_l2_of_tsum_enorm_rpow_ne_top,
  Haar.memApprox_nTerm_iff_memEMethod, Haar.exists_inclusion_eq_of_memEMethod,
  Haar.memApprox_nTerm_iff_intervalModulus_of_lt_two}
\leanok
\uses{thm:haar-nterm, thm:haar-below-one, lem:haar-martingale-l1, prop:haar-basis,
  thm:steps-besov-below-one, cor:haar-nterm-besov, lem:lq-basic, lem:inclusion}
Let $0 < \tau < 2$ and $\theta = 1/\tau - 1/2$.
\begin{enumerate}
  \item For $\tau \leq 1$, $f \in L_2[0, 1]$ lies in the approximation space $A^\theta_\tau$ of
    $n$-term approximation from the Haar basis if and only if its class in $L_\tau[0, 1]$ lies in
    the approximation space $E_{\theta, \tau}$ of the step couple in $L_\tau[0, 1]$.
  \item For $\tau \leq 1$, every element of $E_{\theta, \tau}$ is the class of a function in
    $L_2[0, 1]$.
  \item $f \in L_2[0, 1]$ lies in $A^\theta_\tau$ if and only if
    \[
      \|f\|_\tau + \bigl\| \bigl( 2^{k\theta} \omega_1(f, 2^{-k}, [0,1])_\tau \bigr)_k
      \bigr\|_{\ell_\tau} < \infty .
    \]
\end{enumerate}
\end{theorem}

\begin{proof}
\leanok
(1) By \Cref{thm:haar-nterm} and \Cref{thm:haar-below-one} with $p = q = \tau$ and $s = \theta$,
where $\theta > 1/\tau - 1$ and the weight $2^{k(s + 1/2 - 1/p)}$ is $1$, both conditions say
$\sum_k \sum_i |\langle f, h_{k,i} \rangle|^\tau < \infty$ (\Cref{lem:lq-basic}); $f \in L_2$ is
integrable. (2) The Haar coefficients of $g \in E_{\theta, \tau}$ lie in
$\ell_\tau \subseteq \ell_2$, so \Cref{prop:haar-basis} gives $F \in L_2[0, 1]$ with the same
coefficients; $F - g$ is integrable with vanishing coefficients, so $F = g$ almost everywhere
(\Cref{lem:haar-martingale-l1}). (3) For $\tau \leq 1$ combine (1) with
\Cref{thm:steps-besov-below-one}, since $0 < \theta < 1/\tau$; for $1 < \tau < 2$ it is
\Cref{cor:haar-nterm-besov}.
\end{proof}

For $\tau \leq 1$, by (1) and (2) the approximation space of $n$-term Haar approximation in
$L_2$ and the approximation space of steps in $L_\tau$ are equal as sets. For
$2/3 < \tau < 2$, that is $\theta < 1$, the condition in (3) defines the classical Besov space
$\Besov{\theta}{\tau}{\tau}[0, 1]$~\cite[\S3.6]{DeVore1998}, and (3) is the Haar case of the
theorem of DeVore, Jawerth and Popov~\cite{DeVoreJawerthPopov}; for $\tau \leq 2/3$ it is the
space $B^\theta_{\tau,\tau,1}[0, 1]$ of the first modulus (\Cref{lem:steps-line}).

\section{Local projections and multiscale details}

\begin{lemma}[Dual basis for polynomials]
\label{lem:poly-dual-basis}
\lean{Interval.integrableOn_polyFun_mul, Interval.polyFun_add, Interval.polyFun_smul,
  Interval.polyPairing, Interval.polyPairing_apply, Interval.polyPairing_comm,
  Interval.polyFun_map_conj, Interval.polyPairing_separatingLeft,
  Interval.polyPairing_nondegenerate, Interval.monoBasis, Interval.dualPoly,
  Interval.sum_polyPairing_smul_dualPoly, Interval.polyFun_eq_sum_dualPoly}
\leanok
Let $r \in \N$ and $b_i(t) = t^i$, $i < r$. There are polynomials $d_0, \dots, d_{r-1}$ of degree
$< r$ such that every polynomial $P$ of degree $< r$ satisfies
\[
  P = \sum_{i < r} \Bigl( \int_0^1 P \, b_i \Bigr) d_i .
\]
\end{lemma}

\begin{proof}
\leanok
The pairing $\langle P, Q \rangle = \int_0^1 P Q$ on polynomials of degree $< r$ is
nondegenerate: $\langle P, \bar P \rangle = \int_0^1 |P|^2$, where $\bar P$ has the conjugate
coefficients, vanishes only for $P = 0$. So $(b_i)$ has a dual basis $(d_i)$ with
$\langle d_i, b_k \rangle = \delta_{ik}$, and the coordinates of $P$ in it are
$\langle P, b_i \rangle$.
\end{proof}

\begin{definition}[Local projections]
\label{def:loc-proj}
\lean{Interval.exists_enorm_polyFun_le_of_fintype, Interval.locCoeff, Interval.locProj,
  Interval.locProj_mem, Interval.polyFun_locProj, Interval.integrableOn_mul_monoBasis,
  Interval.locCoeff_add, Interval.locCoeff_const_mul, Interval.locCoeff_congr,
  Interval.locProj_add, Interval.locProj_const_mul, Interval.locProj_congr}
\leanok
\uses{lem:poly-dual-basis, def:pp-scheme}
For $g \in L_1(I_{j,m})$ the local projection of $g$ on $I_{j,m}$ is the polynomial
\[
  \Pi_{j,m} g(x) = \sum_{i < r} c_i \, d_i(2^j x - m), \qquad
  c_i = 2^j \int_{I_{j,m}} g(y) \, b_i(2^j y - m) \, dy .
\]
It is linear in $g$ and depends only on the class of $g$ in $L_1(I_{j,m})$.
\end{definition}

\begin{lemma}[Properties of the local projections]
\label{lem:loc-proj}
\lean{Interval.integral_dyadicIco_eq, Interval.polyFun_locProj_of_eqOn,
  Interval.ofReal_two_pow_mul_inv_rpow, Interval.volume_dyadicIco, Interval.enorm_locCoeff_le,
  Interval.enorm_polyFun_locProj_le, Interval.eLpNorm_locProj_le,
  Interval.integral_locProj_mul_monoBasis, Interval.integral_sub_locProj_mul}
\leanok
\uses{def:loc-proj}
Let $1 \leq p < \infty$ and let $M$ bound all $|b_i|$ and $|d_i|$ on $[0, 1]$. If $g$ agrees on
$I_{j,m}$ with a polynomial $Q$ of degree $< r$, then $\Pi_{j,m} g = Q$. For every
$g \in L_p(I_{j,m})$,
\[
  \|\Pi_{j,m} g\|_{L_p(I_{j,m})} \leq r M^2 \|g\|_{L_p(I_{j,m})} ,
\]
and $\int_{I_{j,m}} (g - \Pi_{j,m} g) Q = 0$ for every polynomial $Q$ of degree $< r$: the local
projection is the $L_2(I_{j,m})$-orthogonal projection.
\end{lemma}

\begin{proof}
\leanok
With $y = (m + t) 2^{-j}$ the coefficients of $Q$ are $c_i = \int_0^1 \tilde Q \, b_i$ for
$\tilde Q(t) = Q((m + t) 2^{-j})$, so $\Pi_{j,m} Q(x) = \tilde Q(2^j x - m) = Q(x)$ by
\Cref{lem:poly-dual-basis}. Hölder's inequality on an interval of length $2^{-j}$ gives
$|c_i| \leq M 2^{j/p} \|g\|_{L_p(I_{j,m})}$, hence
$|\Pi_{j,m} g| \leq r M^2 2^{j/p} \|g\|_{L_p(I_{j,m})}$ on $I_{j,m}$, and the $L_p$ norm of a
function bounded by $B$ on $I_{j,m}$ is at most $2^{-j/p} B$. For the orthogonality, $\Pi_{j,m} g$
and $g$ have the same moments against $b_l(2^j x - m)$, since
$\langle d_i, b_l \rangle = \delta_{il}$, and $Q$ is a combination of the $b_l(2^j x - m)$.
\end{proof}

\begin{theorem}[Projections onto piecewise polynomials]
\label{thm:pp-proj}
\lean{Interval.ppFun_congr, Interval.ae_dyadicIco_of_unitVol, Interval.integrableOn_dyadicIco,
  Interval.basisBound, Interval.basisBound_ne_top, Interval.enorm_monoBasis_le,
  Interval.enorm_dualPoly_le, Interval.ppProjConst, Interval.ppProjConst_ne_top,
  Interval.eLpNorm_ppFun_locProj_le, Interval.ppProjLinear, Interval.ppProj,
  Interval.coeFn_ppProj, Interval.norm_ppProj_le, Interval.ppProj_mem, Interval.ppProj_of_mem,
  Interval.ppProjectionScheme, Interval.ppProjectionScheme_two_pow, Interval.ppFun_locProj_ppFun}
\leanok
\uses{lem:loc-proj, lem:pp-bookkeeping, def:projection-scheme}
Let $1 \leq p < \infty$ and $r \in \N$. For $j \in \N$ let $P_j f = \Pi_{j,m} f$ on every
$I_{j,m}$, $m < 2^j$. Then $P_j$ is a bounded projection of $L_p[0,1]$ onto $V_j$, and
$\|P_j\| \leq r M^2$ for all $j$. As functions, $P_j a = a$ for every piecewise polynomial $a$ of
level $j$ and degree $< r$.
\end{theorem}

\begin{proof}
\leanok
$P_j f \in V_j$ by construction, and $P_j a = a$ for $a \in V_j$ by \Cref{lem:loc-proj}. Summing
the $p$-th powers of the local bounds of \Cref{lem:loc-proj} over the intervals of level $j$
gives $\|P_j f\|_p \leq r M^2 \|f\|_p$.
\end{proof}

\begin{theorem}[Multiscale characterisation]
\label{thm:pp-multiscale}
\lean{Interval.ppScheme_mono_degree, Interval.closure_iUnion_ppScheme,
  Interval.memApprox_ppScheme_iff_block, Interval.lqNorm_modulus_lt_top_iff_block,
  Interval.lqNorm_perModulus_lt_top_iff_block}
\leanok
\uses{thm:pp-proj, lem:haar-density, cor:haar-besov, thm:pp-besov, thm:pp-besov-per}
Let $1 \leq p < \infty$, $r \geq 1$, $s > 0$ and $1 \leq q \leq \infty$. A function
$f \in L_p[0,1]$ lies in the approximation space $A^s_q$ of the scheme of dyadic piecewise
polynomials of degree $< r$ if and only if
\[
  \|P_0 f\|_p + \bigl\| \bigl( 2^{ks} \|P_{k+1} f - P_k f\|_p \bigr)_k \bigr\|_{\ell_q} < \infty .
\]
For $0 < s < 1/p$ this condition is equivalent to $f \in \Besov{s}{p}{q}[0,1]$ and to
$f \in \Besov{s}{p}{q}(\T)$, defined by $\omega_r$ and $\omega_r^{\mathrm{per}}$.
\end{theorem}

\begin{proof}
\leanok
The piecewise polynomials are dense in $L_p[0,1]$, since they contain the step functions of
the Haar system, which are dense (\Cref{lem:haar-density}); the Haar scheme is the scheme of
piecewise polynomials of degree $< 1$ (\Cref{cor:haar-besov}). The
projections of \Cref{thm:pp-proj} form a projection scheme, and the abstract block
characterisation applies. The last statement combines this with \Cref{thm:pp-besov} and
\Cref{thm:pp-besov-per}.
\end{proof}

The details $P_{k+1} f - P_k f$ are what multiwavelets of degree $< r$ expand at level $k$, and
the local projections $P_j$ are those of discontinuous Galerkin methods.

\section{Multiwavelets}
\label{sec:multiwavelets}

\begin{definition}[Detail space and multiwavelets]
\label{def:multiwavelets}
\lean{Alpert.pieceOf, Alpert.pieceOf_mem, Alpert.pieceOf_add, Alpert.pieceOf_smul,
  Alpert.pieceOf_degreeLTEquiv, Alpert.twoPieces, Alpert.twoPieces_mem, Alpert.twoPieces_add,
  Alpert.twoPieces_smul, Alpert.levelOne, Alpert.coeFn_levelOne, Alpert.eq_zero_of_ae_dyadicIco,
  Alpert.levelOne_injective, Alpert.polyL2, Alpert.range_polyL2_le, Alpert.finrank_range_polyL2,
  Alpert.finrank_range_levelOne, Alpert.detail, Alpert.finiteDimensional_detail,
  Alpert.finrank_detail, Alpert.detailBasis, Alpert.detailBasis_mem_range,
  Alpert.multiwaveletPieces, Alpert.levelOne_multiwaveletPieces, Alpert.multiwavelet,
  Alpert.coeFn_detailBasis, Alpert.inner_eq_integral}
\leanok
\uses{def:pp-scheme, lem:poly-dual-basis}
Let $r \geq 1$. The \emph{detail space} $W \subseteq L_2[0,1]$ consists of the piecewise
polynomials of degree $< r$ on $[0, \tfrac12)$ and $[\tfrac12, 1)$ that are $L_2$-orthogonal to the
polynomials of degree $< r$. It has dimension $2r - r = r$. The \emph{multiwavelets}
$\psi_1, \dots, \psi_r$ are an orthonormal basis of $W$, and
$\psi_{k,m,i}(x) = 2^{k/2} \psi_i(2^k x - m)$ for $k \in \N$, $m < 2^k$.
\end{definition}

B.~Alpert chose such a basis explicitly through Legendre polynomials; any orthonormal basis of
$W$ serves. For $r = 1$, $\psi_1$ is the Haar function.

\begin{lemma}[Properties of the multiwavelets]
\label{lem:multiwavelets}
\lean{Alpert.integral_conj_multiwavelet_mul_multiwavelet,
  Alpert.integral_conj_polyFun_mul_multiwavelet, Alpert.integral_polyFun_mul_multiwavelet,
  Alpert.ppFun_ae_eq_sum_multiwavelet, Interval.conj_ppFun, Alpert.integrable_conj_multiwavelet_mul,
  Alpert.integral_conj_multiwavelet_mul_sum, Alpert.memLp_sum_multiwavelet, Alpert.sumMultiwavelet,
  Alpert.coeFn_sumMultiwavelet, Alpert.sumMultiwavelet_injective, Alpert.exists_stable}
\leanok
\uses{def:multiwavelets}
The multiwavelets are orthonormal in $L_2[0,1]$, have $r$ vanishing moments,
$\int_0^1 \psi_i(x) x^l \, dx = 0$ for $l < r$, and every $u \in W$ is
$u = \sum_i \langle u, \psi_i \rangle \psi_i$. For $1 \leq p \leq \infty$ there are $K, M$ with
\[
  \|\beta\|_{\ell_p} \leq K \Bigl\| \sum_i \beta_i \psi_i \Bigr\|_{L_p[0,1]} \leq K M \|\beta\|_{\ell_p}
\]
for all $\beta \in \C^r$.
\end{lemma}

\begin{proof}
\leanok
The first three statements hold in $W$. The map $\beta \mapsto \sum_i \beta_i \psi_i$ into
$L_p[0,1]$ is injective, since $\beta_j = \int_0^1 \overline{\psi_j} \sum_i \beta_i \psi_i$, and
it is defined on a space of dimension $r$, so its inverse is bounded; the upper bound is the
triangle inequality.
\end{proof}

\begin{lemma}[Details in the multiwavelets]
\label{lem:multiwavelet-details}
\lean{Interval.integrableOn_mul_polyFun, Interval.dyadicIco_eq_union,
  Interval.disjoint_dyadicIco_halves, Alpert.integral_sub_ppFun_succ_mul, Alpert.alpertDetailFun,
  Alpert.integral_alpertDetailFun_mul, Alpert.exists_conj_multiwavelet_eq_polyFun,
  Alpert.integral_ppFun_mul_conj_multiwavelet,
  Alpert.integral_alpertDetailFun_mul_conj_multiwavelet, Alpert.alpertCoeff, Alpert.detailPieces,
  Alpert.ppFun_detailPieces, Alpert.levelOne_detailPieces_mem,
  Alpert.integral_conj_multiwavelet_mul_detailPieces, Alpert.ppFun_detailPieces_ae_eq}
\leanok
\uses{lem:multiwavelets, lem:loc-proj, thm:pp-proj, lem:dyadic-normalising-factor}
Let $g \in L_1[0,1]$, $k \in \N$ and $m < 2^k$. On $I_{k,m}$,
\[
  P_{k+1} g - P_k g = \sum_{i \leq r} \langle g, \psi_{k,m,i} \rangle \, \psi_{k,m,i}
  \qquad \text{almost everywhere,}
\]
with $\langle g, \psi_{k,m,i} \rangle = 2^{k/2} \int_{I_{k,m}} g(x) \,
\overline{\psi_i(2^k x - m)} \, dx$.
\end{lemma}

\begin{proof}
\leanok
On $I_{k,m}$ the detail $\Delta_k g = P_{k+1} g - P_k g$ is a polynomial of degree $< r$ on each
half, and it is orthogonal to the polynomials of degree $< r$: on each half
$P_{k+1} g - g$ is, and on $I_{k,m}$ so is $g - P_k g$ (\Cref{lem:loc-proj}). Moved to $[0, 1]$
it therefore lies in $W$ and expands in the $\psi_i$. Its coefficients are those of $g$: the
moved multiwavelet is a polynomial on each half, so $\int_{I_{k,m}} (P_{k+1} g - g)
\overline{\psi_i(2^k \cdot - m)} = 0$, and $\int_{I_{k,m}} P_k g \, \overline{\psi_i(2^k \cdot - m)} =
0$ by the vanishing moments.
\end{proof}

\begin{lemma}[Norms of the details]
\label{lem:multiwavelet-norms}
\lean{ENNReal.ofReal_two_pow_eq, Alpert.eLpNorm_alpertDetailFun_dyadicIco,
  Alpert.exists_eLpNorm_alpertDetailFun_equiv, Alpert.enorm_ppProj_sub}
\leanok
\uses{lem:multiwavelet-details, lem:multiwavelets, lem:dyadic-normalising-factor,
  lem:pp-bookkeeping}
Let $1 \leq p < \infty$. There are $K, M$ with
\[
  2^{k(1/2 - 1/p)} \bigl\| (\langle g, \psi_{k,m,i} \rangle)_{m,i} \bigr\|_{\ell_p}
  \leq K \|P_{k+1} g - P_k g\|_p
\]
and
\[
  \|P_{k+1} g - P_k g\|_p
  \leq M \, 2^{k(1/2 - 1/p)} \bigl\| (\langle g, \psi_{k,m,i} \rangle)_{m,i} \bigr\|_{\ell_p}
\]
for every $g \in L_p[0,1]$ and every $k \in \N$.
\end{lemma}

\begin{proof}
\leanok
By \Cref{lem:multiwavelet-details} and a change of variables,
$2^{k/p} \|P_{k+1} g - P_k g\|_{L_p(I_{k,m})} = \|\sum_i 2^{k/2} \langle g, \psi_{k,m,i} \rangle
\psi_i\|_{L_p[0,1]}$. Apply \Cref{lem:multiwavelets} on each interval and sum the $p$-th powers
over $m < 2^k$.
\end{proof}

\begin{theorem}[Multiwavelet characterisation]
\label{thm:multiwavelet-besov}
\lean{Alpert.memApprox_ppScheme_iff_alpertCoeff, Alpert.lqNorm_modulus_lt_top_iff_alpertCoeff,
  Alpert.lqNorm_perModulus_lt_top_iff_alpertCoeff}
\leanok
\uses{thm:pp-multiscale, lem:multiwavelet-norms, thm:pp-besov, thm:pp-besov-per}
Let $1 \leq p < \infty$, $r \geq 1$, $s > 0$ and $1 \leq q \leq \infty$. A function
$f \in L_p[0,1]$ lies in the approximation space $A^s_q$ of the scheme of dyadic piecewise
polynomials of degree $< r$ if and only if
\[
  \|P_0 f\|_p + \Bigl\| \Bigl( 2^{k(s + 1/2 - 1/p)} \bigl\| (\langle f, \psi_{k,m,i} \rangle)_{m < 2^k,
  i \leq r} \bigr\|_{\ell_p} \Bigr)_k \Bigr\|_{\ell_q} < \infty .
\]
For $0 < s < 1/p$ this condition is equivalent to $f \in \Besov{s}{p}{q}[0,1]$ and to
$f \in \Besov{s}{p}{q}(\T)$.
\end{theorem}

\begin{proof}
\leanok
By \Cref{lem:multiwavelet-norms} the sequences $(2^{ks} \|P_{k+1} f - P_k f\|_p)_k$ and
$(2^{k(s + 1/2 - 1/p)} \|(\langle f, \psi_{k,m,i} \rangle)_{m,i}\|_{\ell_p})_k$ are comparable, so
their $\ell_q$ norms are finite together; apply \Cref{thm:pp-multiscale}.
\end{proof}

\section{Cardinal B-splines}

\begin{definition}[Cardinal B-splines]
\label{def:bspline}
\lean{Interval.bsplineCoeff, Interval.bsplinePiece, Interval.bspline, Interval.bspline_eq_piece,
  Interval.bspline_eq_zero_of_lt, Interval.bspline_eq_zero_of_ge}
\leanok
Let $r \geq 1$. For $0 \leq j \leq r$ let
\[
  B_j(x) = \frac{1}{(r-1)!} \sum_{l = 0}^{j} (-1)^l \binom{r}{l} (x - l)^{r-1} .
\]
The cardinal B-spline of order $r$ is $N_r(x) = B_j(x)$ for $x \in [j, j+1)$, $0 \leq j < r$, and
$N_r(x) = 0$ for $x \notin [0, r)$.
\end{definition}

\begin{lemma}[The pieces of the B-spline]
\label{lem:bspline-pieces}
\lean{Interval.bsplinePiece_succ, Interval.bsplinePiece_zero, Interval.pow_mem_degreeLT,
  Interval.bsplinePiece_mem, Interval.bsplinePiece_top}
\leanok
\uses{def:bspline, lem:iter-diff-poly}
Every $B_j$ is a polynomial of degree $< r$, $B_{j+1} - B_j$ is a multiple of $(x - j - 1)^{r-1}$, and
$B_r = 0$.
\end{lemma}

\begin{proof}
\leanok
The first two statements are immediate. $B_r(z)$ is $\pm 1/(r-1)!$ times the $r$-th difference
with step $1$ at $0$ of the polynomial $t \mapsto (z - t)^{r-1}$ of degree $< r$, which vanishes
(\Cref{lem:iter-diff-poly}, here directly Mathlib's statement for the step $1$).
\end{proof}

So consecutive pieces of $N_r$, including the zero pieces outside $[0, r)$, agree to order $r - 2$
at the knots: $N_r$ is a spline of order $r$.

\begin{lemma}[B-splines on a unit interval]
\label{lem:bspline-basis}
\lean{Interval.bsplineLocal, Interval.bsplineLocal_eq, Interval.bsplineLocal_mem,
  Interval.linearIndependent_X_add_C_pow, Interval.shiftedPow, Interval.top_le_span_shiftedPow,
  Interval.bsplineVec, Interval.bsplineVec_eq, Interval.top_le_span_bsplineVec,
  Interval.bsplineBasis, Interval.bsplineBasis_apply}
\leanok
\uses{lem:bspline-pieces}
The $r$ B-splines $N_r(\cdot + j)$, $0 \leq j < r$, restricted to $[0, 1)$ are the polynomials
$\beta_j(x) = B_j(x + j)$, and these form a basis of the polynomials of degree $< r$.
\end{lemma}

\begin{proof}
\leanok
The shifted powers $(x + m)^{r-1}$, $0 \leq m < r$, are linearly independent: the coefficients
of a vanishing combination $\sum_m g_m (x + m)^{r-1}$ give $\sum_m g_m m^e = 0$ for $e < r$, a
Vandermonde system. They form a basis, since there are $r$ of them. In this basis
$\beta_j = \frac{1}{(r-1)!} (x + j)^{r-1} + \sum_{m < j} (\ldots) (x + m)^{r-1}$ is triangular,
so the $\beta_j$ span as well.
\end{proof}

\begin{theorem}[Marsden's identity]
\label{thm:marsden}
\lean{Interval.sum_range_eq_add_last, Interval.sum_range_eq_first_add, Interval.bcoord,
  Interval.sum_bcoord, Interval.bcoord_eq_of_sum, Interval.bsplineLocal_comp_add_one,
  Interval.bsplineLocal_zero, Interval.bsplineLocal_top, Interval.bcoord_comp_add_one,
  Interval.comp_X_add_C_mem, Interval.bsplineCoef, Interval.bcoord_comp_eq_bsplineCoef,
  Interval.bspline_sub_eq, Interval.polyFun_eq_sum_bsplineCoef, Interval.bcoord_add,
  Interval.bcoord_C_mul, Interval.bsplineCoef_add, Interval.bsplineCoef_C_mul,
  Interval.bsplineCoef_sum}
\leanok
\uses{lem:bspline-basis}
Let $r \geq 1$. For every polynomial $P$ of degree $< r$ there are coefficients
$c_n(P) \in \C$, $n \in \Z$, such that for all $i \in \Z$ and $x \in [i, i+1)$
\[
  P(x) = \sum_{j < r} c_{i - j}(P) \, N_r(x - (i - j)) ,
\]
that is, $P = \sum_{n \in \Z} c_n(P) N_r(\cdot - n)$ on $\R$. The coefficients $c_n(P)$ depend
linearly on $P$.
\end{theorem}

\begin{proof}
\leanok
On $[i, i+1)$ expand $P(\cdot + i)$ in the basis of \Cref{lem:bspline-basis}; let $c_i(P)$ be the
coordinate of $\beta_0$ there. Since $\beta_j(x + 1) = \beta_{j+1}(x) - c\, x^{r-1}$ with
$\beta_r = 0$ and $x^{r-1} = (r-1)! \beta_0$, the coordinate $j + 1$ of $P(\cdot + i + 1)$ is the
coordinate $j$ of $P(\cdot + i)$. By induction the coordinate $j$ on $[i, i+1)$ is
$c_{i-j}(P)$, and $\beta_j(x - i) = N_r(x - (i - j))$ there. Coordinates in a basis are linear.
\end{proof}

\begin{lemma}[Coefficients of shifted powers]
\label{lem:bspline-coef-powers}
\lean{Interval.bsplinePiece_eq_neg_sum, Interval.bsplineLocal_eq_neg_sum,
  Interval.X_sub_C_pow_mem_span_bsplineLocal, Interval.bcoord_bsplineLocal,
  Interval.bcoord_zero_of_mem_span, Interval.bsplineCoef_X_sub_C_pow}
\leanok
\uses{thm:marsden, lem:bspline-pieces}
Let $r \geq 1$ and $m, n \in \Z$ with $n < m < n + r$. Then $c_n\bigl((x - m)^{r-1}\bigr) = 0$.
\end{lemma}

\begin{proof}
\leanok
By definition $c_n(P)$ is the coordinate of $\beta_0$ in the expansion of $P(\cdot + n)$, and
$(x + n - m)^{r-1} = (x - a)^{r-1}$ with $a = m - n \in \{1, \dots, r - 1\}$. Since $B_r = 0$
(\Cref{lem:bspline-pieces}), $B_j = -\sum_{j < l \leq r} \frac{(-1)^l}{(r-1)!} \binom{r}{l}
(x - l)^{r-1}$, so
\[
  \beta_{r-a}(x) = -\frac{(-1)^r}{(r-1)!} (x - a)^{r-1}
    - \sum_{1 \leq b < a} \frac{(-1)^{b + r - a}}{(r-1)!} \binom{r}{b + r - a} (x - b)^{r-1} .
\]
By induction on $a$, $(x - a)^{r-1}$ is a combination of $\beta_1, \dots, \beta_{r-1}$, so its
coordinate at $\beta_0$ vanishes.
\end{proof}

\section{Dyadic splines}

\begin{lemma}[Spline joins]
\label{lem:spline-joins}
\lean{Interval.exists_eq_C_mul_pow_iff}
\leanok
Let $r \geq 1$, $\xi \in \C$ and $P$ a polynomial of degree $< r$. Then
$P = c (x - \xi)^{r-1}$ for some $c \in \C$ if and only if $P^{(i)}(\xi) = 0$ for all $i < r - 1$.
\end{lemma}

\begin{proof}
\leanok
If the derivatives vanish, $\xi$ is a root of multiplicity $\geq r - 1$, so $(x - \xi)^{r-1}$
divides $P$, and the quotient has degree $0$.
\end{proof}

\begin{definition}[Dyadic splines]
\label{def:spline-space}
\lean{Interval.SplineJoins, Interval.SplineJoinsBelow, Interval.SplineJoins.below,
  Interval.splineJoins_zero, Interval.SplineJoins.add,
  Interval.SplineJoins.const_mul, Interval.SplineJoins.sum, Interval.SplineJoins.refine,
  Interval.splineSpace, Interval.splineSpace_le_ppSpace, Interval.splineSpace_zero,
  Interval.splineSpace_mono, Interval.splineScheme, Interval.mem_splineScheme_zero,
  Interval.mem_splineScheme, Interval.mem_splineScheme_two_pow,
  Interval.splineScheme_isLinearScheme}
\leanok
\uses{def:pp-scheme, lem:spline-joins}
Let $r \geq 1$. The space $S_k \subseteq V_k$ of splines of order $r$ and level $k$ consists of
the piecewise polynomials of degree $< r$ on the intervals $I_{k,n}$ whose pieces $Q_n$ satisfy
$Q_n - Q_{n-1} = c_n (x - n 2^{-k})^{r-1}$ at every interior knot, $0 < n < 2^k$. By
\Cref{lem:spline-joins} these are the piecewise polynomials that are $(r-2)$ times continuously
differentiable. At level $0$ there is no interior knot, so $S_0 = V_0$. The spaces $S_k$ increase
in $k$, and the spline scheme is their dyadic scheme, $A_n = S_{\lfloor \log_2 n \rfloor}$ for
$n \geq 1$.
\end{definition}

\begin{lemma}[Dyadic B-splines are splines]
\label{lem:bspline-dyadic}
\lean{Interval.bsplinePieces, Interval.bsplinePieces_mem, Interval.X_sub_C_pow_comp_affine,
  Interval.bsplinePieces_join, Interval.splineJoins_bsplinePieces, Interval.polyFun_bsplinePieces}
\leanok
\uses{def:spline-space, lem:bspline-pieces}
For $n \in \Z$ the B-spline $N_r(2^k x - n)$ is on $I_{k,m}$, $m \in \N$, the polynomial
$B_{m-n}(2^k x - n)$ if $0 \leq m - n \leq r$ and $0$ otherwise, and these pieces have spline
joins at every knot $m 2^{-k}$, $m \geq 1$.
\end{lemma}

\begin{proof}
\leanok
At the knot $m 2^{-k}$ the pieces differ by $(B_d - B_{d-1})(2^k x - n)$, $d = m - n$, a multiple
of $(2^k x - m)^{r-1}$ by \Cref{lem:bspline-pieces}; at the ends of the support one of the two
pieces is $0$ and the other is $B_0$ or $B_{r-1} = B_{r-1} - B_r$.
\end{proof}

\begin{lemma}[B-spline coefficients of the pieces of a spline]
\label{lem:spline-coef-pieces}
\lean{Interval.X_sub_C_pow_comp_scale, Interval.bsplineCoef_comp_eq_of_splineJoins}
\leanok
\uses{def:spline-space, lem:bspline-coef-powers}
Let $r \geq 1$, let $s \in S_k$ have the pieces $Q_m$, and let $\tilde Q_m(y) = Q_m(y 2^{-k})$. For
$n \in \Z$ and $j \leq i < 2^k$ with $n \leq j$ and $i < n + r$,
$c_n(\tilde Q_i) = c_n(\tilde Q_j)$.
\end{lemma}

\begin{proof}
\leanok
$\tilde Q_{m} - \tilde Q_{m-1}$ is a multiple of $(y - m)^{r-1}$, and for $j < m \leq i$ we have
$n < m < n + r$, so \Cref{lem:bspline-coef-powers} and the linearity of $c_n$ give
$c_n(\tilde Q_m) = c_n(\tilde Q_{m-1})$.
\end{proof}

\section{The inverse inequality for splines}

\begin{lemma}[Differences of truncated powers]
\label{lem:spline-inverse-tools}
\lean{Interval.truncPow, Interval.enorm_higherDiff_truncPow_le_of_lt,
  Interval.enorm_higherDiff_truncPow_le}
\leanok
\uses{def:iter-diff}
Let $r \geq 1$, $m \in \N$, $x < \xi$ and $h \geq 0$. Then
$|\Delta_h^m (\cdot - \xi)_+^{r-1}(x)| \leq 2^m (mh)^{r-1}$, in particular
$|\Delta_h^r (\cdot - \xi)_+^{r-1}(x)| \leq 2^r (rh)^{r-1}$.
\end{lemma}

\begin{proof}
\leanok
$0 \leq (x + ih - \xi)_+ \leq mh$ for $i \leq m$, and $\sum_i \binom{m}{i} = 2^m$.
\end{proof}

\begin{lemma}[Splines across the knots]
\label{lem:spline-across}
\lean{Interval.eq_add_sum_Ioc_sub, Interval.jumpCoeff, Interval.sub_eq_jumpCoeff,
  Interval.polyFun_eq_add_sum_jumps, Interval.higherDiff_eq_add_sum_jumps,
  Interval.higherDiff_eq_sum_jumps, Interval.enorm_jumpCoeff_le,
  Interval.enorm_jumpCoeff_mul_higherDiff_truncPow_le_any_order}
\leanok
\uses{def:spline-space, lem:interval-nikolskii, lem:pp-local, lem:iter-diff-poly,
  lem:spline-inverse-tools}
Let $N \in \N$ and let $u$ agree on $I_{j,m}$, $m < N$, with polynomials $P_m$ of degree $< r$ such
that $P_m - P_{m-1} = c_m (x - \xi_m)^{r-1}$, $\xi_m = m 2^{-j}$, for $0 < m < N$; let
$m_n = \int_{I_{j,n}} |u|^p$. For $0 \leq h \leq 2^{-j}$ and $0 \leq x$ with $x + rh < N 2^{-j}$,
$n_0 = \lfloor 2^j x \rfloor$,
\[
  \Delta_h^r u(x) = \sum_{n_0 < m \leq n_0 + r} c_m \, \Delta_h^r (\cdot - \xi_m)_+^{r-1}(x),
\]
and $|c_m| 2^{-j(r-1)} \leq K 2^{j/p} (m_m^{1/p} + m_{m-1}^{1/p})$ for $0 < m < N$. For a spline in
$S_j$, $N = 2^j$. For differences of order $k \leq r$ with $x + kh < N 2^{-j}$ the same holds
with the piece added, $\Delta_h^k u(x) = \Delta_h^k P_{n_0}(x) +
\sum_{n_0 < m \leq n_0 + r} c_m \, \Delta_h^k (\cdot - \xi_m)_+^{r-1}(x)$, and every jump term is
at most $2 \cdot 2^k k^{r-1} K (2^j h)^{r-1} 2^{j/p} (\sum_{l \leq r} m_{n_0+l})^{1/p}$.
\end{lemma}

\begin{proof}
\leanok
On $[x, x + rh]$, $u = Q_{n_0} + \sum_m c_m (\cdot - \xi_m)_+^{r-1}$ by telescoping, and
$\Delta_h^r Q_{n_0} = 0$. With $L = 2^{-j}$,
$c_m L^{r-1} = Q_m(\xi_m + L) - Q_{m-1}(\xi_m + L)$; the first value is bounded by Nikolskii's
inequality (\Cref{lem:interval-nikolskii}) on $I_{j,m}$, the second by the same inequality on
the doubled interval of $I_{j,m-1}$. For order $k \leq r$ the same telescoping on $[x, x + kh]$
keeps the piece $Q_{n_0}$, and each jump term is $|c_m|$ times
$|\Delta_h^k (\cdot - \xi_m)_+^{r-1}(x)| \leq 2^k (kh)^{r-1}$
(\Cref{lem:spline-inverse-tools}). Since $2^{j(r-1)} (kh)^{r-1} = k^{r-1} (2^j h)^{r-1}$ and
$m_m$, $m_{m-1}$ are among the $m_{n_0+l}$, $l \leq r$, the bound on $c_m$ gives the stated
bound.
\end{proof}

\begin{theorem}[Inverse inequality for splines]
\label{thm:spline-inverse}
\lean{Interval.splineStripConst, Interval.enorm_higherDiff_spline_le,
  Interval.eLpNorm_higherDiff_spline_le, Interval.eLpNorm_higherDiff_ppFun_spline_le,
  Interval.splineInverseConst,
  Interval.splineInverseConst_ne_top, Interval.splineStripConst_mul_eq,
  Interval.splineStripConst_mul_le, Interval.two_pow_le_two_pow_mul_one_add,
  Interval.le_two_pow_mul_one_add, Interval.exists_inverse_splineScheme}
\leanok
\uses{lem:spline-across, lem:spline-inverse-tools, lem:pp-local}
Let $1 \leq p < \infty$ and $r \geq 1$. There is $C_B < \infty$ such that for all $j, k \in \N$ and
every spline $a \in S_j$ of order $r$,
\[
  \omega_r(a, 2^{-k}, [0, 1])_p \leq C_B \, 2^{-\max(k - j, 0)(r - 1 + 1/p)} \|a\|_p .
\]
\end{theorem}

\begin{proof}
\leanok
For $k \geq j$ and $0 \leq h \leq 2^{-k}$, \Cref{lem:spline-across} and
\Cref{lem:spline-inverse-tools} give
$|\Delta_h^r a(x)| \leq C (2^j h)^{r-1} 2^{j/p} (\sum_{l \leq r} m_{n_0 + l})^{1/p}$ on the
strips. The last statement of \Cref{lem:pp-local} gives
$\|\Delta_h^r a\|_{L_p[0, 1 - rh]} \leq C' (2^j h)^{r - 1 + 1/p} \|a\|_p$, and
$2^j h \leq 2^{-(k-j)}$.
\end{proof}

\section{A spline quasi-interpolant}

\begin{definition}[Quasi-interpolant]
\label{def:quasi-interp}
\lean{Interval.eq_of_polyFun_eq, Interval.locProjScaled, Interval.qiIdx, Interval.qiCoeff,
  Interval.qiPieces, Interval.qiPieces_apply, Interval.qiPieces_mem, Interval.qiPieces_join,
  Interval.splineJoins_qiPieces, Interval.polyFun_qiPieces}
\leanok
\uses{thm:marsden, def:loc-proj, lem:bspline-dyadic}
Let $r \geq 1$ and $k, N \in \N$. For $-r < n < N$ let $J(n) = I_{k, \max(n, 0)}$, an interval of
the support of $N_r(2^k x - n)$ in $[0, \infty)$, and let $\lambda_n(g) = c_n(\tilde P)$ be the
coefficient of \Cref{thm:marsden} of the rescaled local projection
$\tilde P(y) = \Pi_{J(n)} g(y 2^{-k})$. The quasi-interpolant is
\[
  Q_k g(x) = \sum_{-r < n < N} \lambda_n(g) \, N_r(2^k x - n) ,
\]
with spline joins at every knot by \Cref{lem:bspline-dyadic}. On $[0, 1)$, $N = 2^k$, it is a
spline of order $r$ and level $k$.
\end{definition}

\begin{lemma}[Properties of the quasi-interpolant]
\label{lem:quasi-interp}
\lean{Interval.bspline_eq_zero_of_not_mem, Interval.polyFun_qiPieces_of_coeff,
  Interval.polyFun_qiPieces_of_eq,
  Interval.exists_enorm_bspline_le, Interval.exists_enorm_bcoord_le, Interval.qiCoeff_eq_bcoord,
  Interval.enorm_qiCoeff_le, Interval.bsplinePieces_eq_zero_of_not_mem,
  Interval.qiPieces_eq_sum_Ioc, Interval.polyFun_qiPieces_eq_sum_Ioc, Interval.eLpNorm_qiPieces_le}
\leanok
\uses{def:quasi-interp, lem:loc-proj, lem:bspline-basis}
Let $1 \leq p < \infty$, $r \geq 1$, $k, N \in \N$ and $i < N$.
\begin{enumerate}
  \item If $\lambda_n(g) = c_n(T(\cdot 2^{-k}))$ for $i - r < n \leq i$ and a polynomial $T$ of
    degree $< r$, then $Q_k g = T$ on $I_{k,i}$. This holds in particular if $g$ agrees on the
    intervals $J(n)$, $i - r < n \leq i$, with $T$.
  \item There is $C$, independent of $k$ and $i$, with
    $\|Q_k g\|_{L_p(I_{k,i})} \leq C \sum_{i - r < n \leq i} \|g\|_{L_p(J(n))}$.
\end{enumerate}
\end{lemma}

\begin{proof}
\leanok
(1) On $I_{k,i}$ only the B-splines $i - r < n \leq i$ are nonzero, and Marsden's identity sums
them to $T$. If $g = T$ on $J(n)$, the local projections reproduce $T$ (\Cref{lem:loc-proj}), so
$\lambda_n(g) = c_n(T(\cdot 2^{-k}))$. (2) The coordinates in the B-spline basis of \Cref{lem:bspline-basis}
are linear functionals on a space of dimension $r$, hence bounded by $C_b \|\cdot\|_{L_p[0,1]}$;
by the coordinate consistency of \Cref{thm:marsden}, $\lambda_n(g)$ is such a coordinate of
$\tilde P(\cdot + \max(n, 0))$, whose norm on $[0, 1]$ is $2^{k/p} \|\Pi_{J(n)} g\|_{L_p(J(n))}
\leq 2^{k/p} r M^2 \|g\|_{L_p(J(n))}$. With $|N_r| \leq M_N$ and $|I_{k,i}| = 2^{-k}$ this gives (2).
\end{proof}

\begin{theorem}[The quasi-interpolant reproduces splines]
\label{thm:qi-splines}
\lean{Interval.polyFun_qiPieces_ppFun, Interval.ppFun_qiPieces_ppFun,
  Interval.ppFun_eq_sum_bsplinePieces}
\leanok
\uses{lem:quasi-interp, lem:spline-coef-pieces, lem:loc-proj}
Let $r \geq 1$ and $k \in \N$. Then $Q_k s = s$ for every $s \in S_k$, with $N = 2^k$. In
particular, every spline $s \in S_k$ is a combination
$s = \sum_{-r < n < 2^k} \lambda_n(s) N_r(2^k x - n)$ of the B-splines that meet $[0, 1)$.
\end{theorem}

\begin{proof}
\leanok
Let $s$ have the pieces $Q_m$ and let $i < 2^k$ and $i - r < n \leq i$. The local projection on
$J(n) = I_{k, \max(n, 0)}$ reproduces the piece $Q_{\max(n,0)}$ (\Cref{lem:loc-proj}), and
$n \leq \max(n, 0) \leq i < n + r$, so \Cref{lem:spline-coef-pieces} gives
$\lambda_n(s) = c_n(\tilde Q_{\max(n,0)}) = c_n(\tilde Q_i)$. By part (1) of
\Cref{lem:quasi-interp}, $Q_k s = Q_i$ on $I_{k,i}$.
\end{proof}

\section{Jackson's inequality for splines and the coincidence theorem}
\label{sec:spline-jackson-coincidence}

\begin{lemma}[The local error of the quasi-interpolant]
\label{lem:qi-local-error}
\lean{Interval.sum_lintegral_window_le, Interval.locProjScaled_mem,
  Interval.qiCoeff_add_const_mul, Interval.polyFun_qiPieces_sub, Interval.qiWindow,
  Interval.dyadicIco_subset_qiWindow, Interval.dyadicIco_subset_qiWindow_self,
  Interval.natCast_add_one_sub_le, Interval.qiWindow_subset, Interval.eLpNorm_sub_qiPieces_le,
  Interval.sum_eLpNorm_sub_qiPieces_le, Interval.exists_window_poly, Interval.exists_qiWindow_poly,
  Interval.qiBasisPoly, Interval.qiBasisPoly_mem, Interval.locProjScaled_eq_sum,
  Interval.qiCoeff_eq_sum}
\leanok
\uses{lem:quasi-interp, thm:marsden}
The coefficients $\lambda_n(g)$ are linear combinations of the coefficients of the local
projection of $g$ on $J(n)$. Let $W_i = [(i + 1 - r) 2^{-k}, (i + 1) 2^{-k}]$ (truncated at
$0$). There is $C$ with
$\|g - Q_k g\|_{L_p(I_{k,i})} \leq C \|g - T\|_{L_p(W_i)}$ for every $i < N$ and every polynomial
$T$ of degree $< r$, and the windows $W_i$, $i < N$, cover every point at most $r$ times. Hence,
if for every $i$ in a set $s \subseteq \{0, \dots, N - 1\}$ some $T_i$ has
$\|g - T_i\|_{L_p(W_i)}^p \leq \int_{W_i} H$, then
$\sum_{i \in s} \|g - Q_k g\|_{L_p(I_{k,i})}^p \leq C^p r \int H$.
\end{lemma}

\begin{proof}
\leanok
$Q_k$ is linear and reproduces $T$ on $I_{k,i}$, so $g - Q_k g = (g - T) - Q_k(g - T)$ there,
and part (2) of \Cref{lem:quasi-interp} applies with $J(n) \subseteq W_i$.
\end{proof}

\begin{lemma}[Spline approximation from smoothings]
\label{lem:qi-taylor}
\lean{Interval.exists_sum_eLpNorm_sub_qiPieces_le_of_taylor, Interval.lintegral_sub_ppFun_le_sum,
  Interval.natCast_div_two_pow_le, Interval.exists_eLpNorm_sub_qi_le_of_taylor,
  Interval.add_ofReal_div_pow_mul_le, Interval.exists_eLpNorm_qi_le}
\leanok
\uses{lem:qi-local-error}
Let $1 \leq p < \infty$ and $r \geq 1$. There is $A < \infty$ with the following property. Let
$k, N \in \N$, $s \subseteq \{0, \dots, N - 1\}$ a set of cells whose windows lie in $[a, b]$,
$u$ measurable and integrable on the dyadic intervals $I_{k,m}$, $m < N$, and $g, D$ measurable
such that for all $v \leq w$ some polynomial $T$ of degree $< r$ has
$\|g - T\|_{L_p[v,w]} \leq (w - v)^r \|D\|_{L_p[v,w]}$. Then, with $c = (r 2^{-k})^r$,
\[
  \sum_{i \in s} \|u - Q_k u\|_{L_p(I_{k,i})}^p \leq
  \bigl(A (\|u - g\|_{L_p[a,b]} + c \|D\|_{L_p[a,b]})\bigr)^p .
\]
Consequently, for $k \geq r + 2$, every such $u$ with $N = 2^k$, and measurable
$g_1, D_1, g_2, D_2$ such that for all $a \leq b$ and $\nu = 1, 2$ some polynomial $T$ of degree
$< r$ has $\|g_\nu - T\|_{L_p[a,b]} \leq (b - a)^r \|D_\nu\|_{L_p[a,b]}$,
\begin{multline*}
  \|u - Q_k u\|_{L_p[0,1]} \leq A \bigl( \bigl( \|u - g_1\|_{L_p[0,3/4]}
    + c \|D_1\|_{L_p[0,3/4]} \bigr) \\
  + \bigl( \|u - g_2\|_{L_p[1/4,1]} + c \|D_2\|_{L_p[1/4,1]} \bigr) \bigr) ,
  \qquad c = (r 2^{-k})^r .
\end{multline*}
There is $C < \infty$ with $\|Q_k u\|_{L_p[0,1]} \leq C \|u\|_{L_p[0,1]}$ for every $k$ and every
such $u$.
\end{lemma}

\begin{proof}
\leanok
On every window $W_i$, $i \in s$, a Taylor polynomial $T_i$ of $g$ satisfies
$\|u - T_i\|_{L_p(W_i)}^p \leq 2^{p-1} \int_{W_i} (|u - g|^p + c^p |D|^p)$, since $W_i$ has length
at most $r 2^{-k}$. \Cref{lem:qi-local-error} with $H = 2^{p-1} (|u - g|^p + c^p |D|^p)$ on
$[a, b]$ bounds the sum by $C^p r 2^{p-1} (\|u - g\|^p + (c \|D\|)^p)$, and
$a^p + b^p \leq (a + b)^p$. Every window has length $\leq r 2^{-k} \leq \tfrac14$, so it lies in
$[0, \tfrac34]$ or in $[\tfrac14, 1]$; the first statement for the two sets of cells gives the
second. For the third take $g = D = 0$ on $[0, 1]$.
\end{proof}

\begin{theorem}[Jackson's inequality for splines]
\label{thm:spline-jackson}
\lean{Interval.exists_eLpNorm_sub_qi_le, Interval.approxError_splineScheme_one,
  Interval.exists_approxError_splineScheme_le}
\leanok
\uses{lem:two-sided-smoothing, lem:qi-taylor, def:spline-space, thm:whitney,
  thm:pp-jackson}
Let $1 \leq p < \infty$ and $r \geq 1$. There is $C_J$ with
$E_{2^k}(f)_p \leq C_J \, \omega_r(f, 2^{-k}, [0, 1])_p$ for every $f \in L_p[0,1]$ and $k \geq 0$,
where $E_{2^k}$ is the error of best approximation by splines of order $r$ and level $k$.
\end{theorem}

\begin{proof}
\leanok
For continuous $f$ and $k \geq r + 2$ take the smoothings of \Cref{lem:two-sided-smoothing} with
$L = 2^{-k}$ and $w = \tfrac34$ on $[0, \tfrac34]$ and on $[\tfrac14, 1]$. \Cref{lem:qi-taylor}
with $u = f$ gives $\|f - Q_k f\|_p \leq C' \omega_r(f, 2^{-k}, [0,1])_p$, since
$(r 2^{-k})^r 2^{kr} = r^r$. General
$f$ follow by density of the bounded continuous functions, as in \Cref{thm:pp-jackson}. For
$k \leq r + 1$, $E_{2^k}(f)_p \leq E_1(f)_p$, the splines of level $0$ are the polynomials of
degree $< r$, and \Cref{thm:whitney} with $t = 2^{-(r+1)}$, passed to $L_p$ as in
\Cref{thm:pp-jackson}, gives $E_1(f)_p \leq W \omega_r(f, 2^{-(r+1)}, [0,1])_p \leq W \omega_r(f, 2^{-k}, [0,1])_p$.
\end{proof}

\begin{theorem}[Approximation spaces of dyadic splines]
\label{thm:spline-besov}
\lean{Interval.splineModulusData, Interval.memApprox_splineScheme_iff}
\leanok
\uses{thm:spline-jackson, thm:spline-inverse, def:modulus-data, thm:modulus-coincidence}
Let $1 \leq p < \infty$, $1 \leq q \leq \infty$, $r \geq 1$ and $0 \leq s < r - 1 + 1/p$. A function
$f \in L_p[0,1]$ lies in the approximation space $A^s_q$ of the scheme of dyadic splines of
order $r$ if and only if
\[
  \|f\|_p + \bigl\| \bigl( 2^{ks} \omega_r(f, 2^{-k}, [0,1])_p \bigr)_k \bigr\|_{\ell_q} < \infty .
\]
\end{theorem}

\begin{proof}
\leanok
\Cref{thm:spline-jackson} and \Cref{thm:spline-inverse} make
$\omega_r(\cdot, 2^{-k}, [0,1])_p$ a modulus compatible with the spline scheme of inverse order
$\theta = r - 1 + 1/p$; apply \Cref{thm:modulus-coincidence}.
\end{proof}

For $0 < s < r - 1 + 1/p$ the condition says $f \in \Besov{s}{p}{q}[0, 1]$. For $r = 1$ the
splines are the step functions and this is \Cref{thm:pp-besov}; for $r \geq 2$ it reaches
smoothness above $1$, for example $s < \tfrac32$ for continuous piecewise linear splines in $L_2$.

\section{Splines on the circle}

\begin{definition}[Periodic splines]
\label{def:per-spline-space}
\lean{Interval.PerJoin, Interval.perJoin_zero, Interval.PerJoin.add, Interval.PerJoin.const_mul,
  Interval.PerJoin.refine, Interval.perSplineSpace, Interval.perSplineSpace_le_splineSpace,
  Interval.perSplineSpace_mono, Interval.perSplineScheme, Interval.mem_perSplineScheme_zero,
  Interval.mem_perSplineScheme, Interval.mem_perSplineScheme_two_pow,
  Interval.perSplineScheme_isLinearScheme}
\leanok
\uses{def:spline-space, def:per-modulus}
Let $r \geq 1$ and $k \in \N$. A spline $u \in S_k$ with pieces $Q_0, \dots, Q_{2^k - 1}$ is
\emph{periodic} if in addition $Q_0(x - 1) - Q_{2^k - 1}(x) = c (x - 1)^{r-1}$ for some
$c \in \C$. The periodic splines form subspaces $S_k(\T) \subseteq S_k$ of $L_p(\T) = L_p[0,1]$
increasing in $k$, and the scheme of periodic dyadic splines has the approximants $\{0\}$ at
complexity $0$ and $S_{\lfloor \log_2 n \rfloor}(\T)$ at complexity $n \geq 1$.
\end{definition}

The additional join is the spline join at the point $0 = 1$ of the circle: the periodic extension
$\tilde u$ has the piece $Q_0(x - 1)$ on $[1, 1 + 2^{-k})$.

\begin{lemma}[Periodic extensions of periodic splines]
\label{lem:per-spline-ext}
\lean{Interval.perPieces_join, Interval.splineJoinsBelow_perPieces}
\leanok
\uses{def:per-spline-space, thm:pp-besov-per}
Let $u \in S_k(\T)$ have pieces $Q_0, \dots, Q_{2^k - 1}$. On $I_{k, m' + 2^k c}$, $m' < 2^k$,
$c \in \N$, the periodic extension $\tilde u$ is the polynomial $P_{m' + 2^k c}(x) = Q_{m'}(x - c)$,
and $P_m - P_{m-1} = c_m (x - m 2^{-k})^{r-1}$ for every $m \geq 1$.
\end{lemma}

\begin{proof}
\leanok
For $m = e + 2^k d$ with $0 < e < 2^k$ this is the join of $Q_e$ and $Q_{e-1}$ shifted by $d$;
for $m = 2^k d$, $d \geq 1$, it is the join across $0 = 1$ shifted by $d - 1$.
\end{proof}

\begin{theorem}[The inverse inequality for periodic splines]
\label{thm:per-spline-inverse}
\lean{Interval.exists_inverse_perSplineScheme}
\leanok
\uses{lem:per-spline-ext, lem:spline-across, thm:spline-inverse, thm:pp-besov-per}
Let $1 \leq p < \infty$ and $r \geq 1$. There is $C_B$ with
\[
  \omega_r^{\mathrm{per}}(a, 2^{-k})_p \leq C_B \, 2^{-\max(k - j, 0)(r - 1 + 1/p)} \|a\|_p
\]
for all $j, k \in \N$ and every $a \in S_j(\T)$.
\end{theorem}

\begin{proof}
\leanok
By \Cref{lem:per-spline-ext}, $\tilde a$ satisfies the hypotheses of \Cref{lem:spline-across}
with $N = 2^j + r + 1$, and $x + rh < N 2^{-j}$ for $x \in [0, 1]$ and $h \leq 2^{-j}$. Any $2^j$
consecutive pieces of $\tilde a$ make up one period, so the proof of \Cref{thm:spline-inverse}
runs on all of $[0, 1]$.
\end{proof}

\begin{lemma}[Quasi-interpolants of periodic functions]
\label{lem:per-qi}
\lean{Interval.locCoeff_add_two_pow, Interval.locProj_add_two_pow,
  Interval.locProjScaled_add_two_pow, Interval.bsplineCoef_comp_sub, Interval.qiCoeff_add_two_pow,
  Interval.bsplinePieces_add_two_pow_comp, Interval.qiPieces_add_two_pow_comp,
  Interval.perQiPieces, Interval.perQiPieces_mem, Interval.polyFun_perQiPieces,
  Interval.splineJoins_perQiPieces, Interval.perJoin_perQiPieces}
\leanok
\uses{def:quasi-interp, lem:quasi-interp, def:per-spline-space}
Let $r \geq 1$, $k \in \N$ and let $g$ be $1$-periodic. Then $\lambda_{n + 2^k}(g) = \lambda_n(g)$
for $n \geq 0$, and the pieces of the quasi-interpolant on $[0, 2)$, $N = 2^{k+1}$, satisfy
$Q_{i + 2^k}(x + 1) = Q_i(x)$ for $r - 1 \leq i < 2^k$. The \emph{periodic quasi-interpolant}
$Q_k^{\mathrm{per}} g(x) = Q_k g(x + 1)$, $x \in [0, 1)$, is a spline of order $r$ and level $k$,
and a periodic one when $r \leq 2^k$.
\end{lemma}

\begin{proof}
\leanok
The local projections of $g$ on $I_{k, m + 2^k} = I_{k,m} + 1$ and on $I_{k,m}$ are translates,
and B-spline coefficients and B-spline pieces move with the translation. On $I_{k, i + 2^k}$ only
the coefficients $\lambda_n$, $i + 2^k - r < n \leq i + 2^k$, enter, all with $n - 2^k \geq 0$. The
joins of $Q_k^{\mathrm{per}} g$ inside $[0, 1)$ are those of $Q_k g$ inside $[1, 2)$. Across
$0 = 1$, its first piece moved back by $1$ is the piece of $Q_k g$ on $I_{k, 2^k}$, and its last
piece is the piece on $I_{k, 2^{k+1} - 1}$ moved by $1$, which is the piece on $I_{k, 2^k - 1}$;
so the join is that of $Q_k g$ at the knot $1$.
\end{proof}

\begin{theorem}[Jackson's inequality for periodic splines]
\label{thm:per-spline-jackson}
\lean{Interval.eLpNorm_comp_add_one_dyadicIco, Interval.qiWindow_subset_Icc_zero_two,
  Interval.eLpNorm_sub_perQiPieces_eq, Interval.exists_eLpNorm_sub_perQi_le,
  Interval.intervalCutoff, Interval.continuous_intervalCutoff, Interval.intervalCutoff_zero,
  Interval.intervalCutoff_one, Interval.intervalCutoff_nonneg, Interval.intervalCutoff_le_one,
  Interval.intervalCutoff_eq_one, Interval.dense_periodic,
  Interval.exists_approxError_perSplineScheme_le}
\leanok
\uses{lem:per-qi, lem:qi-taylor, lem:two-sided-smoothing, lem:per-modulus, lem:error-dense}
Let $1 \leq p < \infty$ and $r \geq 1$. There is $C_J$ with
$E_{2^k}(f)_p \leq C_J \, \omega_r^{\mathrm{per}}(f, 2^{-k})_p$ for every $f \in L_p(\T)$ and
every $k \geq r$, where $E_{2^k}$ is the error of best approximation by periodic splines of order
$r$ and level $k$.
\end{theorem}

\begin{proof}
\leanok
First let $f$ be continuous and $1$-periodic. On $[0, 1)$ the error of $Q_k^{\mathrm{per}} f$ is
the error of $Q_k f$ on $[1, 2)$, whose windows $W_i$, $2^k \leq i < 2^{k+1}$, lie in $[0, 2]$.
\Cref{lem:two-sided-smoothing} on $[0, 2]$ with $L = 2^{-k}$ and $b = 3$ gives polynomials on all
windows, and \Cref{lem:qi-taylor} and \Cref{lem:per-modulus} give
$\|f - Q_k^{\mathrm{per}} f\|_p \leq C \omega_r(f, 2^{-k}, [0, 3])_p
\leq 3^{1/p} C \omega_r^{\mathrm{per}}(f, 2^{-k})_p$. The continuous $1$-periodic functions are
dense in $L_p[0,1]$: for a bounded continuous $g$ and a continuous cut-off $\varphi_\delta$ with
values in $[0, 1]$, equal to $1$ on $[\delta, 1 - \delta]$ and to $0$ at $0$ and $1$, the function
$\varphi_\delta g$ extends to a continuous $1$-periodic function, and
$\|g - \varphi_\delta g\|_p \leq \|g\|_\infty (2\delta)^{1/p}$. \Cref{lem:error-dense} gives the
bound for all $f$, since $Q_k^{\mathrm{per}} f \in S_k(\T)$ for $k \geq r$ by \Cref{lem:per-qi}.
\end{proof}

\begin{theorem}[Approximation spaces of periodic splines]
\label{thm:per-spline-besov}
\lean{Interval.perSplineModulusData, Interval.memApprox_perSplineScheme_iff}
\leanok
\uses{thm:per-spline-jackson, thm:per-spline-inverse, def:modulus-data, thm:modulus-coincidence}
Let $1 \leq p < \infty$, $1 \leq q \leq \infty$, $r \geq 1$ and $0 \leq s < r - 1 + 1/p$. A
function $f \in L_p(\T)$ lies in the approximation space $A^s_q$ of the scheme of periodic dyadic
splines of order $r$ if and only if
\[
  \|f\|_p + \bigl\| \bigl( 2^{ks} \omega_r^{\mathrm{per}}(f, 2^{-k})_p \bigr)_k \bigr\|_{\ell_q}
  < \infty .
\]
\end{theorem}

\begin{proof}
\leanok
\Cref{thm:per-spline-jackson}, from level $k_0 = r$ on, and \Cref{thm:per-spline-inverse} make
$\omega_r^{\mathrm{per}}(\cdot, 2^{-k})_p$ a modulus compatible with the periodic spline scheme of
inverse order $\theta = r - 1 + 1/p$; apply \Cref{thm:modulus-coincidence}.
\end{proof}

For $0 < s < r - 1 + 1/p$ the condition says $f \in \Besov{s}{p}{q}(\T)$.

\section{Differences of every order}

The isotropic modulus on the cube (\Cref{sec:iso-modulus}) needs, in each variable, differences of
every order $m \leq r$ of the approximants, also of order below the order of the pieces, where
they do not vanish inside a piece.

\begin{lemma}[Differences of polynomials by derivatives]
\label{lem:poly-diff-deriv}
\lean{Interval.iterate_derivative_mem_degreeLT, Interval.hasDerivAt_higherDiff_polyFun,
  Interval.norm_higherDiff_polyFun_le}
\leanok
\uses{def:iter-diff}
Let $P$ be a polynomial, $m \in \N$, $\sigma \geq 0$, $y \in \R$ and $M \geq 0$ with
$|P^{(m)}(z)| \leq M$ for all $z \in [y, y + m \sigma]$. Then
$|\Delta_\sigma^m P(y)| \leq \sigma^m M$.
\end{lemma}

\begin{proof}
\leanok
Induction on $m$. The derivative of $\Delta_\sigma^m P$ is $\Delta_\sigma^m P'$, and
$\Delta_\sigma^{m+1} P(y) = \Delta_\sigma^m P(y + \sigma) - \Delta_\sigma^m P(y)$, so the mean
value theorem and the induction hypothesis for $P'$ at the points of $[y, y + \sigma]$ give the
bound $\sigma \cdot \sigma^m M$.
\end{proof}

\begin{lemma}[Difference bound for polynomials]
\label{lem:diff-bound}
\lean{Interval.DiffBound, Interval.exists_diffBound}
\leanok
\uses{lem:poly-diff-deriv, lem:interval-nikolskii}
Let $1 \leq p < \infty$ and $r, m \in \N$. There is $K$ such that for every polynomial $P$ of
degree $< r$, all $u \in \R$, $L > 0$, $t \geq 0$ and $x$ with $u \leq x$ and
$x + mt \leq u + 2L$,
\[
  |\Delta_t^m P(x)| \leq K \Bigl(\frac{t}{L}\Bigr)^m L^{-1/p} \, \|P\|_{L_p[u, u + L]} .
\]
\end{lemma}

\begin{proof}
\leanok
Every linear map of the polynomials of degree $< r$ into themselves, in particular
$P \mapsto P^{(m)}$, satisfies $|\Phi P(x)| \leq K \|P\|_{L_p[0, 1/2]}$ for $x \in [0, 1]$
(\Cref{lem:interval-nikolskii}). With
\Cref{lem:poly-diff-deriv}, $|\Delta_\sigma^m Q(y)| \leq \sigma^m K \|Q\|_{L_p[0,1/2]}$ for
$y, y + m\sigma \in [0, 1]$; the substitution $x = u + 2Ly$ rescales.
\end{proof}

\begin{lemma}[Steps of either sign]
\label{lem:step-regions}
\lean{Interval.interval_higherDiff_neg, Interval.stepRegion, Interval.mem_stepRegion,
  Interval.stepRegion_of_nonneg, Interval.stepRegion_of_neg, Interval.stepRegion_eq_diffStepSet,
  Interval.eLpNorm_higherDiff_stepRegion_eq}
\leanok
\uses{def:iter-diff, lem:rd-modulus-on}
For $m \in \N$ and $t \in \R$ let $J(m, t) = \{x \in [0, 1] : x + mt \in [0, 1]\}$, the set
$[0, 1]_{mt}$ of \Cref{def:rd-modulus-on}. Then $\Delta_{-t}^m f(x) = (-1)^m \Delta_t^m f(x - mt)$,
and for every $f$
\[
  \|\Delta_t^m f\|_{L_p(J(m,t))} = \|\Delta_{|t|}^m f\|_{L_p[0, 1 - m|t|]} .
\]
\end{lemma}

\begin{proof}
\leanok
Reindex $k \mapsto m - k$ in the sum of the difference. For $t < 0$, $J(m, t) = [-mt, 1]$, and the
translation by $mt$ maps it onto $[0, 1 + mt]$.
\end{proof}

\begin{theorem}[Differences of every order of piecewise polynomials]
\label{thm:pp-lower-diff}
\lean{Interval.eLpNorm_le_of_lintegral_le_mul, Interval.mul_rpow_one_div_rpow,
  Interval.lintegral_localBound_le, Interval.sum_localMass_ppFun_le,
  Interval.enorm_higherDiff_le_of_mem_piece,
  Interval.enorm_higherDiff_rpow_le_strip_add_piece, Interval.rpow_strip_mul_eq,
  Interval.rpow_piece_mul_le, Interval.exists_eLpNorm_higherDiff_ppFun_le}
\leanok
\uses{lem:diff-bound, thm:pp-inverse, lem:pp-local}
Let $1 \leq p < \infty$ and $1 \leq m \leq r$. There is $C$ such that for every $j \in \N$, every
piecewise polynomial $u$ of level $j$ and degree $< r$ on $[0, 1]$ and $0 \leq h \leq 2^{-j}$,
\[
  \|\Delta_h^m u\|_{L_p[0, 1 - mh]} \leq C (2^j h)^{1/p} \|u\|_p .
\]
\end{theorem}

\begin{proof}
\leanok
If $x + mh$ lies right of the right end of the dyadic interval of $x$, then $|\Delta_h^m u(x)|$ is
bounded by the local masses of the $r + 1$ neighbouring pieces, as in the proof of
\Cref{thm:pp-inverse}; these points form strips of width $\leq rh$ left of the knots. Otherwise
$x, \dots, x + mh$ lie in one piece $I_{j,n}$, and \Cref{lem:diff-bound} gives
$|\Delta_h^m u(x)| \leq K (2^j h)^m 2^{j/p} \|u\|_{L_p(I_{j,n})}$. Integrating,
$\int |\Delta_h^m u|^p \leq C (2^j h + (2^j h)^{mp}) \|u\|_p^p$, and $(2^j h)^{mp} \leq 2^j h$
since $2^j h \leq 1$ and $mp \geq 1$.
\end{proof}

\begin{theorem}[Differences of lower order of splines]
\label{thm:spline-lower-diff}
\lean{Interval.enorm_jumpCoeff_mul_higherDiff_truncPow_le,
  Interval.enorm_higherDiff_spline_le_of_lt,
  Interval.eLpNorm_higherDiff_le_of_lt_mul, Interval.exists_eLpNorm_higherDiff_spline_le_of_lt}
\leanok
\uses{lem:diff-bound, lem:spline-across, lem:spline-inverse-tools}
Let $1 \leq p < \infty$ and $1 \leq m < r$. There is $C$ such that for every $j \in \N$, every
spline $u$ of order $r$ and level $j$ on $[0, 1]$ and $0 \leq h \leq 2^{-j}$,
\[
  \|\Delta_h^m u\|_{L_p[0, 1 - mh]} \leq C (2^j h)^m \|u\|_p .
\]
\end{theorem}

\begin{proof}
\leanok
Let first $mh \leq 2^{-j}$ and $n_0 = \lfloor 2^j x \rfloor$. On $[x, x + mh]$ the spline is the
piece $Q_{n_0}$ plus jumps $c_l (\cdot - \xi_l)_+^{r-1}$ at knots $\xi_l > x$
(\Cref{lem:spline-across}). The piece contributes
$|\Delta_h^m Q_{n_0}(x)| \leq K_m (2^j h)^m 2^{j/p} \|u\|_{L_p(I_{j,n_0})}$ by
\Cref{lem:diff-bound}, since $x, x + mh$ lie in the doubled interval of $I_{j, n_0}$. A jump
contributes $|c_l| \, 2^m (mh)^{r-1}$ (as in \Cref{lem:spline-inverse-tools}), with
$|c_l| 2^{-j(r-1)} \leq K 2^{j/p} (\dots)^{1/p}$ by \Cref{lem:spline-across}, and
$(2^j h)^{r-1} \leq (2^j h)^m$. So
$|\Delta_h^m u(x)| \leq C (2^j h)^m 2^{j/p} (\sum_{l \leq r} \int_{I_{j,n_0+l}} |u|^p)^{1/p}$ at
every point, and integrating gives the claim. For $mh > 2^{-j}$,
$\|\Delta_h^m u\|_p \leq 2^m \|u\|_p \leq 2^m (m 2^j h)^m \|u\|_p$.
\end{proof}

The jump of $u^{(r-1)}$ at a knot limits the gain for $m = r$ to $(2^j h)^{r-1+1/p}$
(\Cref{thm:spline-inverse}); below order $r$ the smoothness $C^{r-2}$ of splines gives the full
order $m$.

\chapter{Approximation on the cube}
\label{ch:cube}

This chapter carries the approximation by piecewise polynomials from the interval to the cube
$Q = [0, 1]^d$. Everything passes through slices: a bound in one variable for a fixed step passes
to the cube by Tonelli's theorem in one coordinate. The tensor approximants are built from
operators of finite rank that act on the slices in one coordinate at a time, which avoids
Fubini arguments in several variables.

\section{The cube and slices}

\begin{definition}[Cube and slices]
\label{def:cube}
\lean{Cube.cubeVol, Cube.slice, Cube.measurable_slice,
  Cube.measurable_slice_uncurry, Cube.slice_update, Cube.regionVol,
  Cube.unitVol_restrict_Icc_of_subset}
\leanok
The cube $Q = [0, 1]^d$ carries the product $\lambda_d$ of Lebesgue measure on $[0, 1]$, a
probability measure on $\R^d$. For $f : \R^d \to \C$, $x \in \R^d$ and a coordinate $i$, the
slice of $f$ through $x$ in coordinate $i$ is $f_{x,i}(t) = f(x_1, \dots, x_{i-1}, t, x_{i+1},
\dots, x_d)$; it does not depend on $x_i$. For $[a, b] \subseteq [0, 1]$ let
$Q_i[a, b] = \{x \in Q : x_i \in [a, b]\}$.
\end{definition}

\begin{lemma}[Integration through slices]
\label{lem:cube-slices}
\lean{Cube.lintegral_cubeVol_eq_slice, Cube.lintegral_region_eq_slice,
  Cube.eLpNorm_rpow_region_eq_slice, Cube.eLpNorm_rpow_cubeVol_eq_slice,
  Cube.eLpNorm_rpow_regionVol_eq_slice, Cube.measurable_eLpNorm_slice, Cube.ae_slice_ae_eq,
  Cube.lintegral_inter_eq_slice, Cube.eLpNorm_inter_eq_slice, Cube.eLpNorm_cubeVol_eq_slice,
  Cube.eLpNorm_regionVol_eq_slice, Cube.eLpNorm_inter_le_of_slice_le,
  Cube.eLpNorm_cubeVol_le_of_slice_le}
\leanok
\uses{def:cube}
Let $G : \R^d \to [0, \infty]$ and $f : \R^d \to \C$ be measurable, $i$ a coordinate,
$0 < p < \infty$ and $[a, b] \subseteq [0, 1]$. Then
\begin{gather*}
  \int_Q G \, d\lambda_d = \int_Q \int_0^1 G_{x,i}(t) \, dt \, d\lambda_d(x) , \\
  \|f\|_{L_p(Q_i[a,b])}^p = \int_Q \|f_{x,i}\|_{L_p[a,b]}^p \, d\lambda_d(x) ,
\end{gather*}
and $x \mapsto \|f_{x,i}\|_{L_p(\mu)}$ is measurable for every $s$-finite measure $\mu$ on $\R$.
If $f = g$ almost everywhere on $Q$, then for almost every $x \in Q$ the slices $f_{x,i}$ and
$g_{x,i}$ agree almost everywhere on $[0, 1]$. For every measurable $R \subseteq \R^d$ such that
$x \in R \Leftrightarrow x' \in R$ whenever $x$ and $x'$ differ only in the coordinate $i$, and
all measurable $S, S' \subseteq \R$, with integrals over subsets of $\R^d$ taken with respect to
$\lambda_d$,
\begin{gather*}
  \int_{R \cap \{x_i \in S\}} G \, d\lambda_d
  = \int_R \int_{S \cap [0, 1]} G_{x,i}(t) \, dt \, d\lambda_d(x) , \\
  \|f\|_{L_p(R \cap \{x_i \in S\})} = \bigl\| x \mapsto \|f_{x,i}\|_{L_p(S \cap [0, 1])}
  \bigr\|_{L_p(R)} ,
\end{gather*}
in particular, for $R = \R^d$, $\|f\|_{L_p(Q)} = \| x \mapsto \|f_{x,i}\|_{L_p[0,1]} \|_{L_p(Q)}$
and $\|f\|_{L_p(Q_i[a,b])} = \| x \mapsto \|f_{x,i}\|_{L_p[a,b]} \|_{L_p(Q)}$. For all measurable
$F, H : \R^d \to \C$ and $C \in [0, \infty]$ with
$\|F_{x,i}\|_{L_p(S \cap [0,1])} \leq C \|H_{x,i}\|_{L_p(S' \cap [0,1])}$ for every $x$,
$\|F\|_{L_p(R \cap \{x_i \in S\})} \leq C \|H\|_{L_p(R \cap \{x_i \in S'\})}$, in particular
$\|F\|_{L_p(Q)} \leq C \|H\|_{L_p(Q)}$ for $S = S' = \R$ and $R = \R^d$.
\end{lemma}

\begin{proof}
\leanok
The first identity is Tonelli's theorem for the coordinate $i$, in the form of marginals; the
inner integral does not depend on $x_i$, and $\lambda_d$ is a probability measure. The second
applies it to $|f|^p$ times the indicator of $Q_i[a, b]$. For the statement on $f = g$ apply the
first to the indicator of $\{f \neq g\}$. The integral over $R \cap \{x_i \in S\}$ is the first
identity for $G$ times the indicator of $R \cap \{x_i \in S\}$, since the indicator of $R$ does
not depend on $x_i$; for $|f|^p$ its $p$-th root is the identity of norms, and the bound follows
from it by the monotonicity and the homogeneity of the $L_p$ norm.
\end{proof}

\section{Moduli in the coordinates}

\begin{definition}[Moduli of smoothness on the cube]
\label{def:cube-modulus}
\lean{Cube.cubeDiff, Cube.cubeDiff_eq, Cube.measurable_cubeDiff, Cube.slice_cubeDiff,
  Cube.cubeDiff_add, Cube.diffRegion, Cube.coordModulus,
  Cube.cubeModulus, Cube.lpCubeModulus}
\leanok
\uses{def:cube, def:iter-diff}
For $r \in \N$, $h \geq 0$ and a coordinate $i$, the difference in coordinate $i$ is
$\Delta^r_{h e_i} f(x) = \Delta_h^r f_{x,i}(x_i) = \sum_{k=0}^r (-1)^{r-k} \binom{r}{k}
f(x + k h e_i)$. For $1 \leq p < \infty$ and $t \geq 0$ the modulus in coordinate $i$ and the
modulus on the cube are
\[
  \omega_r^{(i)}(f, t)_p = \sup_{0 \leq h \leq t} \|\Delta^r_{h e_i} f\|_{L_p(Q_i[0, 1 - rh])} ,
  \qquad
  \omega_r(f, t)_p = \sum_{i=1}^d \omega_r^{(i)}(f, t)_p .
\]
\end{definition}

These are the moduli in the coordinates of Nikol'ski\u{\i}'s definition of Besov spaces
\cite{Nikolskii}; they are equivalent to the isotropic modulus
$\sup_{|h| \leq t} \|\Delta_h^r f\|_p$, which is not formalised here. The supremum over the step
stays outside the integral over the other coordinates.

\begin{lemma}[Properties of the moduli on the cube]
\label{lem:cube-modulus}
\lean{Cube.eLpNorm_rpow_cubeDiff_eq_slice, Cube.eLpNorm_cubeDiff_le_coordModulus,
  Cube.coordModulus_mono, Cube.coordModulus_add_le, Cube.eLpNorm_cubeDiff_le_of_slice_le,
  Cube.eLpNorm_cubeDiff_le, Cube.coordModulus_le, Cube.coordModulus_eq_of_slice,
  Cube.coordModulus_congr_ae, Cube.cubeModulus_mono, Cube.cubeModulus_add_le, Cube.cubeModulus_le,
  Cube.cubeModulus_congr_ae, Cube.lpCubeModulus_add_le, Cube.lpCubeModulus_le,
  Cube.lpCubeModulus_anti, Cube.lpCubeModulus_le_of_coord}
\leanok
\uses{def:cube-modulus, lem:cube-slices, lem:interval-modulus}
Let $f$ be measurable, $1 \leq p < \infty$ and $h \geq 0$. Then
\[
  \|\Delta^r_{h e_i} f\|_{L_p(Q_i[0, 1 - rh])}^p
  = \int_Q \|\Delta_h^r f_{x,i}\|_{L_p[0, 1 - rh]}^p \, d\lambda_d(x) .
\]
For every $B \in [0, \infty]$ with
$\|\Delta_h^r f_{x,i}\|_{L_p[0, 1 - rh]} \leq B \|f_{x,i}\|_{L_p[0,1]}$ for every $x$,
\[
  \|\Delta^r_{h e_i} f\|_{L_p(Q_i[0, 1 - rh])} \leq B \|f\|_{L_p(Q)} .
\]
If $g$ is measurable and, for
every $0 \leq h \leq t$ and almost every $x$, the differences $\Delta_h^r f_{x,i}$ and
$\Delta_h^r g_{x,i}$ have the same norm in $L_p[0, 1 - rh]$, then
$\omega_r^{(i)}(f, t)_p = \omega_r^{(i)}(g, t)_p$.
The moduli $\omega_r^{(i)}(f, t)_p$ and $\omega_r(f, t)_p$ increase in $t$, are subadditive in $f$,
satisfy $\omega_r^{(i)}(f, t)_p \leq 2^r \|f\|_{L_p(Q)}$ and
$\omega_r(f, t)_p \leq d \, 2^r \|f\|_{L_p(Q)}$, and depend only on the class of $f$ in $L_p(Q)$.
For all $j, k \in \N$, $\theta \in \R$ and $B \geq 2^r$: if $a \in L_p(Q)$ is represented by a
measurable $g$ and, in case $j \leq k$,
$\|\Delta^r_{h e_i} g\|_{L_p(Q_i[0, 1 - rh])} \leq B \, 2^{-(k-j)\theta} \|g\|_{L_p(Q)}$ for every
coordinate $i$ and every $0 \leq h \leq 2^{-k}$, then
$\omega_r(a, 2^{-k})_p \leq d B \, 2^{-\max(k-j, 0)\theta} \|a\|_p$.
\end{lemma}

\begin{proof}
\leanok
The identity is \Cref{lem:cube-slices} for the difference, whose slices in coordinate $i$ are the
differences of the slices. The bounds hold slice by slice by \Cref{lem:interval-modulus}, and
\Cref{lem:cube-slices} passes them to $Q$. In the last statement take the supremum over $h$ and
sum over the $d$ coordinates if $j \leq k$, and use $\omega_r(a, t)_p \leq d \, 2^r \|a\|_p$ if
$k < j$.
\end{proof}

\section{Smoothing in one coordinate}

\begin{lemma}[Steklov smoothing in one coordinate]
\label{lem:cube-smoothing}
\lean{Cube.continuous_slice, Cube.continuous_cubeDiff_uncurry, Cube.continuous_cubeDiff,
  Cube.cubeSteklovSmooth, Cube.cubeSteklovSmoothDeriv, Cube.slice_cubeSteklovSmooth,
  Cube.slice_cubeSteklovSmoothDeriv, Cube.sub_cubeSteklovSmooth_eq, Cube.cubeSteklovSmoothDeriv_eq,
  Cube.continuous_cubeSteklovSmooth, Cube.cubeSteklovSmoothDeriv_eq_sum_smul,
  Cube.continuous_cubeSteklovSmoothDeriv, Cube.eLpNorm_sub_cubeSteklovSmooth_le_of_cubeDiff_le,
  Cube.eLpNorm_cubeDiff_regionVol_le, Cube.eLpNorm_sub_cubeSteklovSmooth_le,
  Cube.eLpNorm_cubeSteklovSmoothDeriv_le}
\leanok
\uses{thm:steklov-smoothing, lem:minkowski-avg, def:cube-modulus}
Let $f$ be continuous on $\R^d$, $r \in \N$, $\delta > 0$, and let
$S^{(i)}_\delta f(x) = S_\delta f_{x,i}(x_i)$ with the smoothing $S_\delta$ of
\Cref{thm:steklov-smoothing}. Let $1 \leq p < \infty$. For every s-finite measure $\mu$ on
$\R^d$ and every $B \in [0, \infty]$ with $\|\Delta^r_{h e_i} f\|_{L_p(\mu)} \leq B$ for all
$0 \leq h \leq r\delta$, $\|f - S^{(i)}_\delta f\|_{L_p(\mu)} \leq B$. For $w + r^2 \delta \leq 1$,
\begin{gather*}
  \|f - S^{(i)}_\delta f\|_{L_p(Q_i[0, w])} \leq \omega_r^{(i)}(f, r\delta)_p , \\
  \|\partial_i^r S^{(i)}_\delta f\|_{L_p(Q_i[0, w])} \leq 2^r \delta^{-r} \omega_r^{(i)}(f, r\delta)_p .
\end{gather*}
\end{lemma}

\begin{proof}
\leanok
Slice by slice, \Cref{thm:steklov-smoothing} gives $f - S^{(i)}_\delta f = (-1)^r$ times the
average over $v \in [0, \delta]^r$ of $\Delta^r_{(v_1 + \cdots + v_r) e_i} f$, and
\[
  \partial_i^r S^{(i)}_\delta f
  = \sum_{k=1}^r (-1)^{k+1} \binom{r}{k} (k\delta)^{-r} \Delta^r_{k\delta e_i} f ,
\]
with the averages over the steps only. \Cref{lem:minkowski-avg} with
the point space $\R^d$ and the measure $\mu$ bounds the average by
$\sup_h \|\Delta^r_{h e_i} f\|_{L_p(\mu)} \leq B$. For the cube measure on $Q_i[0, w]$,
$Q_i[0, w] \subseteq Q_i[0, 1 - rh]$ for $h \leq r \delta$. The derivative bound of
\Cref{thm:steklov-smoothing}, for the cube measure on
$Q_i[0, w]$ and $D_h = \Delta^r_{h e_i} f$, bounds the derivative.
\end{proof}

\begin{lemma}[Smoothing near the two ends of a coordinate]
\label{lem:cube-two-sided-smoothing}
\lean{Cube.exists_left_cubeSmoothing, Cube.reflectCoord, Cube.reflectCoord_reflectCoord,
  Cube.continuous_reflectCoord, Cube.slice_reflectCoord, Cube.coordModulus_reflectCoord,
  Cube.eLpNorm_reflectCoord_regionVol, Cube.exists_right_cubeSmoothing}
\leanok
\uses{lem:cube-smoothing, lem:two-sided-smoothing, lem:taylor-lp, lem:interval-modulus}
Let $f$ be continuous on $\R^d$, $r \geq 1$, $1 \leq p < \infty$, $L > 0$, $w + L \leq 1$ and
$\omega = \omega_r^{(i)}(f, L)_p$. On $R = Q_i[0, w]$ and on $R = Q_i[1 - w, 1]$ there are measurable
$G$ and $D$ with $\|f - G\|_{L_p(R)} \leq \omega$, $\|D\|_{L_p(R)} \leq L^{-r} 2^r r^{2r} \omega$,
and for every $x$ and all $a \leq b$ a polynomial $T$ of degree $< r$ with
$\|G_{x,i} - T\|_{L_p[a,b]} \leq (b - a)^r \|D_{x,i}\|_{L_p[a,b]}$.
\end{lemma}

\begin{proof}
\leanok
On $Q_i[0, w]$ take $G = S^{(i)}_\delta f$ with $\delta = L/r^2$, $D = \partial_i^r G$, and
\Cref{lem:taylor-lp} on every slice. On $Q_i[1 - w, 1]$ apply this to $f$ composed with the
reflection $x_i \mapsto 1 - x_i$, which preserves $\omega_r^{(i)}$ slice by slice
(\Cref{lem:interval-modulus}), and reflect back, as on the interval
(\Cref{lem:two-sided-smoothing}).
\end{proof}

\section{Operators of finite rank on slices}

\begin{definition}[Kernel operators and their action in one coordinate]
\label{def:kernel-op}
\lean{Cube.BddMeas, Cube.BddMeas.add, Cube.BddMeas.const_mul, Cube.bddMeas_zero, Cube.bddMeas_sum,
  Cube.BddMeas.integrableOn_Icc, Cube.BddMeas.memLp, Cube.BddMeas.slice, Cube.KernelOp,
  Cube.KernelOp.kernel, Cube.KernelOp.rangeFn, Cube.KernelOp.bddMeas_kernel,
  Cube.KernelOp.bddMeas_rangeFn, Cube.KernelOp.kernelCoef, Cube.KernelOp.apply,
  Cube.KernelOp.integrableOn_mul_kernel, Cube.KernelOp.kernelCoef_add, Cube.KernelOp.kernelCoef_sub,
  Cube.KernelOp.kernelCoef_const_mul, Cube.KernelOp.kernelCoef_sum, Cube.KernelOp.apply_sub,
  Cube.KernelOp.norm_kernelCoef_le, Cube.KernelOp.slicewise, Cube.KernelOp.slicewise_eq,
  Cube.KernelOp.slice_slicewise, Cube.KernelOp.measurable_kernelCoef_slice,
  Cube.KernelOp.measurable_slicewise, Cube.KernelOp.bddMeas_slicewise,
  Cube.KernelOp.slicewise_sub}
\leanok
\uses{def:cube}
A kernel operator on $[0, 1]$ is $T g = \sum_{j \in J} \bigl( \int_0^1 g \, \kappa_j \bigr)
\varphi_j$ with a finite index set $J$ and bounded measurable $\kappa_j$, $\varphi_j$. Its action in
coordinate $i$ on functions on $\R^d$ is $T^{(i)} g(x) = T g_{x,i}(x_i)$. For bounded measurable
$g$, $T^{(i)} g$ is bounded and measurable, its slices in coordinate $i$ are the $T g_{x,i}$, and
$T^{(i)}$ is linear on bounded measurable functions.
\end{definition}

\begin{lemma}[Forms of kernel operators]
\label{lem:kernel-forms}
\lean{Cube.KernelOp.IsForm, Cube.KernelOp.isForm_slicewise, Cube.KernelOp.IsForm.slice_eq,
  Cube.KernelOp.IsForm.slicewise, Cube.KernelOp.IsForm.add, Cube.KernelOp.IsForm.const_mul,
  Cube.KernelOp.isForm_zero, Cube.KernelOp.IsForm.congr, Cube.KernelOp.IsForm.refine,
  Cube.KernelOp.foldSlicewise, Cube.KernelOp.foldSlicewise_nil, Cube.KernelOp.foldSlicewise_cons,
  Cube.KernelOp.bddMeas_foldSlicewise, Cube.KernelOp.isForm_foldSlicewise}
\leanok
\uses{def:kernel-op}
A function $g$ on $\R^d$ has the form of $T$ in coordinate $i$ if
$g(x) = \sum_j c_j(x) \varphi_j(x_i)$ with bounded measurable $c_j$ that do not depend on $x_i$.
These functions form a linear space, and their slices in coordinate $i$ are combinations of the
$\varphi_j$. For bounded measurable $g$:
\begin{enumerate}
  \item $T^{(i)} g$ has the form of $T$ in coordinate $i$;
  \item if $g$ has the form of $T$ in coordinate $i$ and $j \neq i$, then $S^{(j)} g$ has the form
    of $T$ in coordinate $i$ for every kernel operator $S$;
  \item if every $\varphi_j$ is a combination of the functions $\varphi'_{j'}$ of a kernel operator
    $T'$, then every function of the form of $T$ in coordinate $i$ has the form of $T'$ there.
\end{enumerate}
In particular, $T^{(i_1)} \cdots T^{(i_m)} g$ has the form of $T$ in each of the coordinates
$i_1, \dots, i_m$.
\end{lemma}

\begin{proof}
\leanok
(1) Take $c_j(x) = \int_0^1 g_{x,i}(s) \kappa_j(s) \, ds$. (2) The factors $\varphi_a(x_i)$ are
constant along coordinate $j$, so by linearity of the integral over the finite sum
$S^{(j)} g(x) = \sum_a S^{(j)} c_a(x) \varphi_a(x_i)$, and $S^{(j)} c_a$ does not depend on $x_i$.
(3) Regroup the double sum. The last statement follows by induction on $m$.
\end{proof}

\begin{lemma}[Tensor products of a kernel operator]
\label{lem:kernel-tensor}
\lean{Cube.bddMeas_prod_coord, Cube.KernelOp.Reproduces, Cube.KernelOp.Reproduces.apply_sum,
  Cube.KernelOp.IsForm.slicewise_self, Cube.KernelOp.IsForm.kernelCoef_slice,
  Cube.KernelOp.partialTensor, Cube.KernelOp.partialTensor_subset_span,
  Cube.KernelOp.mem_span_prod_of_isForm, Cube.KernelOp.isForm_prod,
  Cube.KernelOp.isForm_of_mem_span_prod}
\leanok
\uses{lem:kernel-forms}
Let $T$ be a kernel operator with $T \varphi_j = \varphi_j$ for every $j$. A function $g$ on $\R^d$
is bounded, measurable and has the form of $T$ in every coordinate if and only if it is a linear
combination of the products $x \mapsto \prod_{i=1}^d \varphi_{j_i}(x_i)$.
\end{lemma}

\begin{proof}
\leanok
A product has the form of $T$ in coordinate $i$: the coefficient of $\varphi_{j_i}(x_i)$ is the
product of the other factors, and the other coefficients vanish. Conversely, $T$ is the identity
on the combinations of the $\varphi_j$, so $T^{(i)} c = c$ for every $c$ of the form of $T$ in
coordinate $i$, that is,
$c(x) = \sum_m c_m(x) \varphi_m(x_i)$ with $c_m(x) = \int_0^1 c_{x,i}(s) \kappa_m(s) \, ds$. For a
set $S$ of coordinates call $c(x) \prod_{i \in S} \varphi_{j_i}(x_i)$ a partial tensor product if
$c$ is bounded, measurable, does not depend on the coordinates in $S$ and has the form of $T$ in
every coordinate outside $S$. For $i_0 \notin S$ the expansion of $c$ in coordinate $i_0$ writes a
partial tensor product for $S$ as a combination of partial tensor products for $S \cup \{i_0\}$:
the coefficients $c_m$ do not depend on $x_{i_0}$ or on the coordinates in $S$, and they keep the
form of $T$ in the coordinates $i \neq i_0$, since integration in coordinate $i_0$ acts on the
coefficients of that form. Starting from $g$ itself with $S = \emptyset$, induction over the
coordinates ends with $S = \{1, \dots, d\}$, where $c$ is constant.
\end{proof}

\begin{lemma}[Telescoping over the coordinates of the cube]
\label{lem:kernel-telescoping}
\lean{Cube.KernelOp.eLpNorm_slicewise_le, Cube.KernelOp.eLpNorm_sub_foldSlicewise_le}
\leanok
\uses{def:kernel-op, lem:cube-slices}
Let $1 \leq p < \infty$ and $C \in [0, \infty]$ with
$\|T u\|_{L_p[0,1]} \leq C \|u\|_{L_p[0,1]}$ for every measurable $u$. Then
$\|T^{(i)} g\|_{L_p(Q)} \leq C \|g\|_{L_p(Q)}$ for measurable $g$, and for bounded measurable
$g$ and coordinates $i_1, \dots, i_m$,
\[
  \bigl\| g - T^{(i_1)} \cdots T^{(i_m)} g \bigr\|_{L_p(Q)}
  \leq (1 + C)^m \sum_{l=1}^m \|g - T^{(i_l)} g\|_{L_p(Q)} .
\]
\end{lemma}

\begin{proof}
\leanok
The first bound holds slice by slice, and \Cref{lem:cube-slices} passes it to $Q$. For the second,
$g - T^{(i)} h = (g - T^{(i)} g) + T^{(i)} (g - h)$; induction on $m$.
\end{proof}

\section{Piecewise polynomials on the cube}

\begin{lemma}[The local projection as a kernel operator]
\label{lem:pp-kernel}
\lean{Cube.norm_indicator_polyFun_le, Cube.ppKernel, Cube.ppKernel_kernelCoef, Cube.ppKernel_apply,
  Cube.eLpNorm_ppKernel_apply_le, Cube.ppKernel_apply_ppFun, Cube.bddMeas_ppFun,
  Cube.ppKernel_rangeFn_eq_ppFun, Cube.ppKernel_rangeFn_eq_sum, Cube.eLpNorm_sub_ppKernel_apply_le}
\leanok
\uses{def:kernel-op, def:loc-proj, thm:pp-proj, lem:pp-bookkeeping}
Let $r \geq 1$, $k \in \N$, and let $b_l$ and $d_l$ be the monomials and their dual basis of
\Cref{lem:poly-dual-basis}. The projection $P_k$ of \Cref{thm:pp-proj} is the kernel operator with
index set $\{(n, l) : n < 2^k, l < r\}$ and
\[
  \kappa_{n,l}(s) = \mathbf{1}_{I_{k,n}}(s) \, 2^k b_l(2^k s - n) ,
  \qquad
  \varphi_{n,l}(t) = \mathbf{1}_{I_{k,n}}(t) \, d_l(2^k t - n) .
\]
Every $\varphi_{n,l}$ of level $k$ is a combination of the functions $\varphi$ of level $k + 1$.
For $1 \leq p < \infty$, every $u$ integrable on $[0, 1]$ and every piecewise polynomial $S$ of
level $k$ and degree $< r$,
\[
  \|u - P_k u\|_{L_p[0,1]} \leq (1 + C_r) \|u - S\|_{L_p[0,1]} ,
\]
with the bound $C_r = r M^2$ of \Cref{thm:pp-proj}.
\end{lemma}

\begin{proof}
\leanok
The coefficient of $(n, l)$ is the local coefficient of \Cref{def:loc-proj}. The function
$\varphi_{n,l}$ is a piecewise polynomial of level $k$, hence of level $k + 1$
(\Cref{lem:pp-bookkeeping}), and $P_{k+1}$ reproduces it (\Cref{thm:pp-proj}). Since $P_k S = S$,
$u - P_k u = (u - S) - P_k (u - S)$.
\end{proof}

\begin{definition}[Piecewise polynomials on the cube]
\label{def:cube-pp}
\lean{Cube.IsCubePP, Cube.IsCubePP.succ, Cube.IsCubePP.add, Cube.IsCubePP.const_mul,
  Cube.isCubePP_zero, Cube.IsCubePP.exists_slice_eq_ppFun, Cube.cubePPSpace, Cube.cubePPSpace_mono,
  Cube.cubePPScheme, Cube.mem_cubePPScheme_two_pow, Cube.cubePPScheme_isLinearScheme, Cube.cubeProj,
  Cube.isCubePP_cubeProj}
\leanok
\uses{lem:pp-kernel, lem:kernel-forms}
Let $r \geq 1$ and $k \in \N$. A function on $\R^d$ is a piecewise polynomial of level $k$ and
degree $< r$ in each variable if it is bounded, measurable, and has the form of $P_k$ in every
coordinate. Their classes form subspaces $V^d_k \subseteq L_p(Q)$ increasing in $k$, and the
scheme of dyadic piecewise polynomials on the cube has the approximants $\{0\}$ at complexity $0$
and $V^d_{\lfloor \log_2 n \rfloor}$ at complexity $n \geq 1$. The tensor projection
$P^d_k g = P_k^{(1)} \cdots P_k^{(d)} g$ of a bounded measurable $g$ is a piecewise polynomial of
level $k$, and the slices of a piecewise polynomial of level $k$ on the cube are piecewise
polynomials of level $k$ and degree $< r$ on $[0, 1]$.
\end{definition}

The definition is coordinatewise: on every dyadic cube of level $k$ such a function is a
polynomial of degree $< r$ in each variable separately. The spaces increase since the forms of
$P_k$ are forms of $P_{k+1}$ (\Cref{lem:pp-kernel}, \Cref{lem:kernel-forms}).

\begin{lemma}[Piecewise polynomials on the cube are tensor products]
\label{lem:cube-pp-tensor}
\lean{Cube.ppKernel_reproduces, Cube.isCubePP_iff_mem_span, Cube.isCubePP_iff_mem_span_prod_ppFun}
\leanok
\uses{def:cube-pp, lem:kernel-tensor, thm:pp-proj}
Let $r \geq 1$ and $k \in \N$. For a function $g$ on $\R^d$ the following are equivalent:
\begin{enumerate}
  \item $g$ is a piecewise polynomial of level $k$ and degree $< r$ in each variable;
  \item $g$ is a linear combination of products $\prod_{i=1}^d \varphi_{n_i, l_i}(x_i)$ of the
    functions of $P_k$ (\Cref{lem:pp-kernel});
  \item $g$ is a linear combination of products $\prod_{i=1}^d u_i(x_i)$ of piecewise polynomials
    $u_i$ of level $k$ and degree $< r$ on $[0, 1]$.
\end{enumerate}
\end{lemma}

\begin{proof}
\leanok
$P_k$ reproduces the piecewise polynomials of level $k$ (\Cref{thm:pp-proj}), in particular its
functions $\varphi_{n,l}$, so \Cref{lem:kernel-tensor} gives the equivalence of (1) and (2). The
$\varphi_{n,l}$ are piecewise polynomials of level $k$, which gives (2) $\Rightarrow$ (3). For the
converse, $u_i = P_k u_i$ is a combination of the $\varphi_{n,l}$, and the product of these
combinations expands into a combination of products of the $\varphi_{n,l}$.
\end{proof}

\begin{lemma}[From slices to the cube]
\label{lem:slices-to-cube}
\lean{Cube.cube_eLpNorm_le_of_slice_le, Cube.slice_sub_slicewise,
  Cube.eLpNorm_sub_slicewise_le_of_smoothing,
  Cube.KernelOp.eLpNorm_sub_foldSlicewise_le_of_coord, Cube.eLpNorm_sub_cubeProj_le_of_coord}
\leanok
\uses{lem:cube-slices, lem:cube-two-sided-smoothing, lem:kernel-telescoping}
Let $1 \leq p < \infty$, $c, M \in [0, \infty]$, and let $F, G_1, D_1, G_2, D_2$ be measurable
on $\R^d$
and $[a_\nu, b_\nu] \subseteq [0, 1]$. If for every $x$
\[
  \|F_{x,i}\|_{L_p[0,1]} \leq c \sum_{\nu = 1}^2 \bigl( \|(G_\nu)_{x,i}\|_{L_p[a_\nu, b_\nu]}
    + M \|(D_\nu)_{x,i}\|_{L_p[a_\nu, b_\nu]} \bigr) ,
\]
then the same bound holds for $\|F\|_{L_p(Q)}$ with the norms on $Q_i[a_\nu, b_\nu]$.

Let $T$ be a kernel operator, $f$ continuous on $\R^d$, $r \geq 1$, $L > 0$, $w + L \leq 1$ and
$c, M, J \in [0, \infty]$. Assume that every continuous $u$ on $\R$ satisfies
\begin{multline*}
  \|u - T u\|_{L_p[0,1]} \leq c \bigl( \|u - g_1\|_{L_p[0, w]} + M \|D_1\|_{L_p[0, w]} \\
    + \|u - g_2\|_{L_p[1 - w, 1]} + M \|D_2\|_{L_p[1 - w, 1]} \bigr)
\end{multline*}
for all measurable $g_\nu, D_\nu$ such that for all $a \leq b$ some polynomial $P$ of degree
$< r$ has $\|g_\nu - P\|_{L_p[a,b]} \leq (b - a)^r \|D_\nu\|_{L_p[a,b]}$, and that
$\alpha + M \beta \leq J \omega$ whenever $\alpha \leq \omega$ and
$\beta \leq L^{-r} 2^r r^{2r} \omega$ in $[0, \infty]$. Then
$\|f - T^{(i)} f\|_{L_p(Q)} \leq 2 c J \, \omega_r^{(i)}(f, L)_p$ for every coordinate $i$.

Let $T$ be a kernel operator with $\|T u\|_{L_p[0,1]} \leq C \|u\|_{L_p[0,1]}$ for measurable
$u$. If $\|f - T^{(i)} f\|_{L_p(Q)} \leq B \, \omega_r^{(i)}(f, t)_p$ for every coordinate $i$
and a bounded measurable $f$, then
$\|f - T^{(1)} \cdots T^{(d)} f\|_{L_p(Q)} \leq (1 + C)^d B \, \omega_r(f, t)_p$; for $T = P_k$
this is $\|f - P^d_k f\|_{L_p(Q)} \leq (1 + C_r)^d B \, \omega_r(f, t)_p$.
\end{lemma}

\begin{proof}
\leanok
The norms of the slices are measurable functions of $x$, and the norms on $Q$ and on
$Q_i[a_\nu, b_\nu]$ are the $L_p(Q)$ norms of these functions (\Cref{lem:cube-slices}); the
bound follows from the monotonicity, the homogeneity and the triangle inequality of the
$L_p(Q)$ norm. For the second statement, \Cref{lem:cube-two-sided-smoothing} on $Q_i[0, w]$ and
on $Q_i[1 - w, 1]$ gives $G_1, D_1$ and $G_2, D_2$; the assumption on the slices
$u = f_{x,i}$, $g_\nu = (G_\nu)_{x,i}$, $D_\nu = (D_\nu)_{x,i}$ and the first statement give
$\|f - T^{(i)} f\|_{L_p(Q)} \leq c \sum_{\nu} (\|f - G_\nu\| + M \|D_\nu\|) \leq 2 c J \omega$.
The third statement is \Cref{lem:kernel-telescoping}.
\end{proof}

\begin{theorem}[Jackson's inequality on the cube]
\label{thm:cube-jackson}
\lean{Cube.eLpNorm_sub_slicewise_ppKernel_le, Cube.cubeJacksonConst, Cube.cubeJacksonConst_ne_top,
  Cube.eLpNorm_sub_cubeProj_le, Cube.approxError_le_of_continuous, Cube.approxError_cubePPScheme_le,
  Cube.exists_eLpNorm_sub_ppKernel_zero_apply_le,
  Cube.exists_eLpNorm_sub_slicewise_ppKernel_zero_le, Cube.exists_approxError_cubePPScheme_one_le,
  Cube.exists_approxError_cubePPScheme_le}
\leanok
\uses{def:cube-pp, lem:pp-kernel, lem:pp-taylor, lem:poly-glue, lem:cube-two-sided-smoothing,
  lem:slices-to-cube, lem:cube-modulus, lem:error-dense}
Let $1 \leq p < \infty$ and $r \geq 1$. For every $f \in L_p(Q)$ and $k \geq 1$,
\[
  E_{2^k}(f)_p \leq (1 + C_r)^{d+1} \, 2 (1 + 2^r r^{2r}) \, \omega_r(f, 2^{-k})_p ,
\]
where $E$ is the error of best approximation by the scheme of \Cref{def:cube-pp}, and there is
$C_J < \infty$ with $E_{2^k}(f)_p \leq C_J \, \omega_r(f, 2^{-k})_p$ for every $f \in L_p(Q)$ and
every $k \geq 0$.
\end{theorem}

\begin{proof}
\leanok
First let $f$ be bounded and continuous, $k \geq 1$ and $L = 2^{-k}$. In coordinate $i$,
\Cref{lem:cube-two-sided-smoothing} on $Q_i[0, \tfrac12]$ and on $Q_i[\tfrac12, 1]$ gives
$G_1, D_1$ and $G_2, D_2$. For every $x$, \Cref{lem:pp-taylor} for the slices of $f$, $G_\nu$ and
$D_\nu$ through $x$ gives a piecewise polynomial $S_x$ of level $k$; it need not depend
measurably on $x$, since only the bound enters. With \Cref{lem:pp-kernel},
$\|f_{x,i} - P_k f_{x,i}\|_{L_p[0,1]}$ is at most $1 + C_r$ times the sum of
$\|(f - G_\nu)_{x,i}\|$ and $L^r \|(D_\nu)_{x,i}\|$ on the two halves, and
\Cref{lem:slices-to-cube} gives
\begin{multline*}
  \|f - P_k^{(i)} f\|_{L_p(Q)}
  \leq (1 + C_r) \sum_{\nu = 1}^2 \bigl( \|f - G_\nu\| + L^r \|D_\nu\| \bigr) \\
  \leq (1 + C_r) \, 2 (1 + 2^r r^{2r}) \, \omega_r^{(i)}(f, L)_p ,
\end{multline*}
the norms on the respective halves of $Q$. The second part of \Cref{lem:slices-to-cube} sums
these bounds over the coordinates, and $P^d_k f \in V^d_k$. For $k = 0$ take the smoothings on
$Q_i[0, \tfrac34]$ and $Q_i[\tfrac14, 1]$ with $L = \tfrac14$; on every slice their Taylor
polynomials on $[0, \tfrac34]$ and $[\tfrac14, 1]$ glue by \Cref{lem:poly-glue} to one
polynomial, and the same steps give $\|f - P_0^d f\|_{L_p(Q)} \leq W \omega_r(f, \tfrac14)_p
\leq W \omega_r(f, 1)_p$. General $f$ follow by \Cref{lem:error-dense}: the bounded continuous
functions are dense in $L_p(Q)$, and $\omega_r(\cdot, t)_p$ is subadditive with
$\omega_r(f, t)_p \leq d \, 2^r \|f\|_p$ (\Cref{lem:cube-modulus}).
\end{proof}

\begin{theorem}[The inverse inequality on the cube]
\label{thm:cube-inverse}
\lean{Cube.eLpNorm_cubeDiff_le_of_isCubePP, Cube.exists_inverse_cubePPScheme}
\leanok
\uses{def:cube-pp, thm:pp-inverse, lem:cube-modulus}
Let $1 \leq p < \infty$ and $r \geq 1$. There is $C_B$ with
\[
  \omega_r(a, 2^{-k})_p \leq C_B \, 2^{-\max(k - j, 0)/p} \|a\|_p
\]
for all $j, k \in \N$ and every $a \in V^d_j$.
\end{theorem}

\begin{proof}
\leanok
Let $0 \leq h \leq 2^{-j}$. Every slice $a_{x,i}$ is a piecewise polynomial of level $j$ on
$[0, 1]$ (\Cref{def:cube-pp}), so the local estimate in the proof of \Cref{thm:pp-inverse} gives
$\|\Delta_h^r a_{x,i}\|_{L_p[0, 1 - rh]} \leq 2^r K (r (r + 1) 2^j h)^{1/p} \|a_{x,i}\|_{L_p[0,1]}$,
with the Nikolskii constant $K$ of \Cref{lem:interval-nikolskii}. Integrating over $x$
(\Cref{lem:cube-modulus}) gives the same bound for $\Delta^r_{h e_i} a$ on the cube; take the
supremum over $h \leq 2^{-k}$ and sum over $i$. For $k < j$ use
$\omega_r(a, t)_p \leq d \, 2^r \|a\|_p$.
\end{proof}

\begin{theorem}[Approximation spaces of piecewise polynomials on the cube]
\label{thm:cube-besov}
\lean{Cube.cubePPModulusData, Cube.memApprox_cubePPScheme_iff}
\leanok
\uses{thm:cube-jackson, thm:cube-inverse, def:modulus-data, thm:modulus-coincidence}
Let $1 \leq p < \infty$, $1 \leq q \leq \infty$, $r \geq 1$ and $0 \leq s < 1/p$. A function
$f \in L_p(Q)$ lies in the approximation space $A^s_q$ of the scheme of dyadic piecewise
polynomials of degree $< r$ in each variable on the cube if and only if
\[
  \|f\|_p + \bigl\| \bigl( 2^{ks} \omega_r(f, 2^{-k})_p \bigr)_k \bigr\|_{\ell_q} < \infty .
\]
\end{theorem}

\begin{proof}
\leanok
\Cref{thm:cube-jackson}, at every level, and \Cref{thm:cube-inverse} make
$\omega_r(\cdot, 2^{-k})_p$ a modulus compatible with the scheme of inverse order $\theta = 1/p$;
apply \Cref{thm:modulus-coincidence}.
\end{proof}

For $0 < s < 1/p$ the condition says $f \in \Besov{s}{p}{q}(Q)$, with the Besov space defined by
the moduli in the coordinates.

\section{Tensor splines}

\begin{lemma}[The quasi-interpolant as a kernel operator]
\label{lem:qi-kernel}
\lean{Cube.toNat_lt_of_mem_qiIdx, Cube.qiKernel, Cube.qiKernel_kernelCoef, Cube.qiKernel_apply,
  Cube.exists_eLpNorm_qiKernel_apply_le}
\leanok
\uses{def:quasi-interp, lem:qi-local-error, lem:pp-kernel, lem:qi-taylor}
Let $r \geq 1$ and $k \in \N$. The quasi-interpolant $Q_k$ of \Cref{def:quasi-interp} with
$N = 2^k$ is a kernel operator with index set $\{n \in \Z : -r < n < 2^k\}$: the kernel
$\kappa_n$ is a combination of the kernels of the local projection on $J(n)$, and $\varphi_n$ is the
B-spline $N_r(2^k x - n)$ on $[0, 1)$. For $1 \leq p < \infty$ there is $C < \infty$ with
$\|Q_k u\|_{L_p[0,1]} \leq C \|u\|_{L_p[0,1]}$ for every $k$ and every measurable $u$.
\end{lemma}

\begin{proof}
\leanok
The coefficients $\lambda_n(g)$ are combinations of the local coefficients on $J(n)$
(\Cref{lem:qi-local-error}), which are integrals against the kernels of \Cref{lem:pp-kernel}. The
bound is \Cref{lem:qi-taylor} when $u \in L_p[0,1]$, and trivial otherwise.
\end{proof}

\begin{definition}[Splines on the cube]
\label{def:cube-spline}
\lean{Cube.IsCubeSpline, Cube.IsCubeSpline.succ, Cube.IsCubeSpline.add,
  Cube.IsCubeSpline.const_mul, Cube.isCubeSpline_zero_fun, Cube.IsCubePP.isCubeSpline_zero,
  Cube.isCubeSpline_of_isForm, Cube.cubeQI, Cube.isCubeSpline_cubeQI, Cube.cubeSplineSpace,
  Cube.cubeSplineSpace_mono,
  Cube.cubeSplineScheme, Cube.mem_cubeSplineScheme_two_pow, Cube.cubeSplineScheme_isLinearScheme,
  Cube.approxError_cubeSplineScheme_one_le}
\leanok
\uses{def:spline-space, lem:qi-kernel, lem:kernel-forms, def:cube-pp}
Let $r \geq 1$ and $k \in \N$. A function on $\R^d$ is a spline of order $r$ and level $k$ in
each variable if it is bounded and measurable and each of its slices, in every coordinate and
through every point, is a dyadic spline of order $r$ and level $k$ on $[0, 1]$. Their classes form
subspaces of $L_p(Q)$ increasing in $k$, and the scheme of tensor splines has the approximants
$\{0\}$ at complexity $0$ and the splines of level $\lfloor \log_2 n \rfloor$ at complexity
$n \geq 1$. At level $0$ they contain the polynomials of degree $< r$ in each variable. The tensor
quasi-interpolant $Q_k^{(1)} \cdots Q_k^{(d)} g$ of a bounded measurable $g$ is a spline of level
$k$ on the cube.
\end{definition}

The tensor quasi-interpolant has the form of $Q_k$ in every coordinate
(\Cref{lem:kernel-forms}), so its slices are combinations of B-splines, hence splines.

\begin{lemma}[Splines on the cube are tensor products]
\label{lem:cube-spline-tensor}
\lean{Cube.qiKernel_reproduces, Cube.IsCubeSpline.slicewise_qiKernel, Cube.IsCubeSpline.isForm,
  Cube.isCubeSpline_iff_mem_span, Cube.isCubeSpline_iff_mem_span_prod_ppFun}
\leanok
\uses{def:cube-spline, thm:qi-splines, lem:kernel-tensor, lem:bspline-dyadic}
Let $r \geq 1$ and $k \in \N$. For a function $g$ on $\R^d$ the following are equivalent:
\begin{enumerate}
  \item $g$ is a spline of order $r$ and level $k$ in each variable;
  \item $g$ is a linear combination of products $\prod_{i=1}^d N_r(2^k x_i - n_i)|_{[0,1)}$ of the
    B-splines with $-r < n_i < 2^k$;
  \item $g$ is a linear combination of products $\prod_{i=1}^d u_i(x_i)$ of dyadic splines
    $u_i \in S_k$ of order $r$ on $[0, 1]$.
\end{enumerate}
\end{lemma}

\begin{proof}
\leanok
$Q_k$ reproduces the splines (\Cref{thm:qi-splines}), in particular its functions
$N_r(2^k x - n)$. So $Q_k^{(i)} g = g$ for a spline $g$ on the cube, and $g$ has the form of $Q_k$
in every coordinate (\Cref{lem:kernel-forms}). \Cref{lem:kernel-tensor} gives the equivalence of
(1) and (2). The B-splines are splines (\Cref{lem:bspline-dyadic}), which gives (2)
$\Rightarrow$ (3). For the converse, $u_i = Q_k u_i$ is a combination of B-splines, and the
product of these combinations expands into a combination of products of B-splines.
\end{proof}

\begin{theorem}[Jackson's inequality for tensor splines]
\label{thm:cube-spline-jackson}
\lean{Cube.exists_eLpNorm_sub_slicewise_qiKernel_le, Cube.exists_approxError_cubeSplineScheme_le}
\leanok
\uses{def:cube-spline, lem:qi-taylor, lem:qi-kernel, lem:cube-two-sided-smoothing,
  lem:slices-to-cube, thm:cube-jackson, lem:error-dense}
Let $1 \leq p < \infty$ and $r \geq 1$. There is $C_J < \infty$ with
$E_{2^k}(f)_p \leq C_J \, \omega_r(f, 2^{-k})_p$ for every $f \in L_p(Q)$ and every $k \geq 0$,
where $E$ is the error of best approximation by the scheme of tensor splines.
\end{theorem}

\begin{proof}
\leanok
Let $f$ be bounded and continuous and $k \geq r + 2$. In coordinate $i$,
\Cref{lem:cube-two-sided-smoothing} with $L = 2^{-k}$ on $Q_i[0, \tfrac34]$ and $Q_i[\tfrac14, 1]$
and \Cref{lem:qi-taylor} on every slice give, through \Cref{lem:slices-to-cube},
$\|f - Q_k^{(i)} f\|_{L_p(Q)} \leq B \, \omega_r^{(i)}(f, 2^{-k})_p$, and the second part of
\Cref{lem:slices-to-cube} with the bound of \Cref{lem:qi-kernel} sums over the coordinates. General
$f$ follow by \Cref{lem:error-dense}. For $k \leq r + 1$, $E_{2^k}(f)_p$ is at most the error of
the polynomials of degree $< r$ in each variable, and \Cref{thm:cube-jackson} at level $0$, with
Whitney's estimate at the scale $2^{-(r+1)}$, bounds it by $W \omega_r(f, 2^{-(r+1)})_p \leq
W \omega_r(f, 2^{-k})_p$.
\end{proof}

\begin{theorem}[The inverse inequality for tensor splines]
\label{thm:cube-spline-inverse}
\lean{Cube.splineStripConst_ne_top, Cube.eLpNorm_cubeDiff_le_of_isCubeSpline,
  Cube.exists_inverse_cubeSplineScheme}
\leanok
\uses{def:cube-spline, thm:spline-inverse, lem:cube-modulus}
Let $1 \leq p < \infty$ and $r \geq 1$. There is $C_B$ with
\[
  \omega_r(a, 2^{-k})_p \leq C_B \, 2^{-\max(k - j, 0)(r - 1 + 1/p)} \|a\|_p
\]
for all $j, k \in \N$ and every spline $a$ of order $r$ and level $j$ on the cube.
\end{theorem}

\begin{proof}
\leanok
Every slice is a spline of level $j$, so the local estimate in the proof of
\Cref{thm:spline-inverse} holds on every slice for each step $0 \leq h \leq 2^{-j}$; integrate
over the point (\Cref{lem:cube-modulus}), take the supremum over $h \leq 2^{-k}$ and sum over the
coordinates. For $k < j$ use $\omega_r(a, t)_p \leq d \, 2^r \|a\|_p$.
\end{proof}

\begin{theorem}[Approximation spaces of tensor splines]
\label{thm:cube-spline-besov}
\lean{Cube.cubeSplineModulusData, Cube.memApprox_cubeSplineScheme_iff}
\leanok
\uses{thm:cube-spline-jackson, thm:cube-spline-inverse, def:modulus-data,
  thm:modulus-coincidence}
Let $1 \leq p < \infty$, $1 \leq q \leq \infty$, $r \geq 1$ and $0 \leq s < r - 1 + 1/p$. A
function $f \in L_p(Q)$ lies in the approximation space $A^s_q$ of the scheme of tensor splines of
order $r$ if and only if
\[
  \|f\|_p + \bigl\| \bigl( 2^{ks} \omega_r(f, 2^{-k})_p \bigr)_k \bigr\|_{\ell_q} < \infty .
\]
\end{theorem}

\begin{proof}
\leanok
\Cref{thm:cube-spline-jackson} and \Cref{thm:cube-spline-inverse} make
$\omega_r(\cdot, 2^{-k})_p$ a modulus compatible with the scheme of tensor splines of inverse
order $\theta = r - 1 + 1/p$; apply \Cref{thm:modulus-coincidence}.
\end{proof}

For $0 < s < r - 1 + 1/p$ the condition says $f \in \Besov{s}{p}{q}(Q)$; for $r \geq 2$ this
reaches smoothness beyond $1$.

\section{The isotropic modulus}
\label{sec:iso-modulus}

\begin{definition}[The translation algebra]
\label{def:translation-algebra}
\lean{TranslationAlgebra.transl, TranslationAlgebra.transl_apply, TranslationAlgebra.translHom,
  TranslationAlgebra.diffAct, TranslationAlgebra.diffAct_single,
  TranslationAlgebra.diffAct_single_one, TranslationAlgebra.higherDiff_eq_diffAct,
  TranslationAlgebra.single_sum_one}
\leanok
\uses{def:iter-diff}
Let $G$ be an additive commutative monoid. The algebra $\C[G]$ of finite formal combinations of
the translations $\delta_v$, $v \in G$, acts on functions on $G$ by $\delta_v f(x) = f(x + v)$.
This action is an algebra homomorphism, $\delta_{v + w} = \delta_v \delta_w$, and
$\Delta_h^r f = (\delta_h - 1)^r f$.
\end{definition}

The algebra $\C[G]$ is commutative, so identities between differences are ring identities.

\begin{definition}[Isotropic differences]
\label{def:iso-diff}
\lean{Cube.cube_higherDiff_apply, Cube.higherDiff_add, Cube.higherDiff_single}
\leanok
\uses{def:cube-modulus, def:translation-algebra}
For $h \in \R^d$ and $r \in \N$ the isotropic difference is
\[
  \Delta_h^r f(x) = \sum_{k=0}^r (-1)^{r-k} \binom{r}{k} f(x + kh) ,
\]
the difference of \Cref{def:iter-diff} on $\R^d$, and $\Delta_h^r f = (\delta_h - 1)^r f$ in the
translation algebra $\C[\R^d]$. For $h = t e_i$ the isotropic difference is the difference
$\Delta^r_{t e_i}$ in coordinate $i$ of \Cref{def:cube-modulus}.
\end{definition}

\begin{lemma}[Expansion into mixed coordinate differences]
\label{lem:iso-expansion}
\lean{Cube.cubeMixedDiff, Cube.cubeMixedDiff_nil, Cube.cubeMixedDiff_cons,
  Cube.measurable_cubeMixedDiff, TranslationAlgebra.coordElem, Cube.diffAct_coordElem_pow,
  Cube.diffAct_list_prod, TranslationAlgebra.letterCount, TranslationAlgebra.sum_letterCount,
  Cube.letterCount_le, TranslationAlgebra.lowerVec, TranslationAlgebra.shiftVec,
  TranslationAlgebra.shiftVec_apply, TranslationAlgebra.card_lt_add_letterCount_le,
  TranslationAlgebra.prod_lt_single, TranslationAlgebra.single_sub_one_eq_sum,
  TranslationAlgebra.single_sub_one_pow_eq_sum, Cube.higherDiff_eq_sum, Cube.slice_cubeDiff_of_ne,
  Cube.slice_cubeMixedDiff_mem}
\leanok
\uses{def:iso-diff}
For a word $\alpha \colon \{1, \dots, r\} \to \{1, \dots, d\}$ let $m_i(\alpha) = \#\alpha^{-1}(i)$
and $(s_\alpha)_i = h_i \, \#\{l : i < \alpha_l\}$, and let
$M_m f = \Delta^{m_1}_{h_1 e_1} \cdots \Delta^{m_d}_{h_d e_d} f$ be the mixed coordinate
difference. Then $\sum_i m_i(\alpha) = r$, $\#\{l : i < \alpha_l\} + m_i(\alpha) \leq r$, and
\[
  \Delta_h^r f(x) = \sum_{\alpha} M_{m(\alpha)} f(x + s_\alpha) .
\]
If every slice of $f$ in coordinate $i$ lies in a subspace $V$ of functions of one variable, so
does every slice in coordinate $i$ of a mixed difference in the other coordinates.
\end{lemma}

\begin{proof}
\leanok
With $\delta_i = \delta_{h_i e_i}$, $\delta_h = \prod_i \delta_i = \prod_i (1 + (\delta_i - 1))$, so
$\delta_h - 1 = \sum_i (\delta_i - 1) \prod_{j < i} \delta_j$. Expanding the $r$-th power of this
sum gives one term $\prod_l (\delta_{\alpha_l} - 1) \prod_{j < \alpha_l} \delta_j
= \delta_{s_\alpha} \prod_i (\delta_i - 1)^{m_i(\alpha)}$ per word $\alpha$. A difference in
coordinate $c \neq i$ is a combination of values at points shifted in coordinate $c$, so it
combines slices in coordinate $i$.
\end{proof}

\begin{definition}[The isotropic modulus]
\label{def:iso-modulus}
\lean{Cube.cubeSet, Cube.measurableSet_cubeSet, Cube.cubeVol_eq_restrict,
  Cube.cubeVol_restrict_congr, Cube.isoRegion, Cube.measurableSet_isoRegion,
  Cube.isoRegion_inter_cubeSet, Cube.eLpNorm_higherDiff_isoRegion_eq, Cube.isoModulus,
  Cube.isoModulus_eq_iSup, Cube.eLpNorm_higherDiff_le_isoModulus, Cube.lpIsoModulus}
\leanok
\uses{def:iso-diff, def:cube, def:rd-modulus-on}
Let $Q_{rh} = \{x \in Q : x + rh \in Q\}$, the set of the points $x$ with
$x, x + h, \dots, x + rh \in Q$. The isotropic modulus of smoothness on the cube is the modulus on
the set $Q$ of \Cref{def:rd-modulus-on},
\[
  \omega^{\mathrm{iso}}_r(f, t)_p = \sup_{\|h\|_\infty \leq t} \|\Delta_h^r f\|_{L_p(Q_{rh})} ,
\]
the supremum over steps of every direction and sign.
\end{definition}

\begin{lemma}[Properties of the isotropic modulus]
\label{lem:iso-modulus}
\lean{Cube.eLpNorm_comp_add_le, Cube.ae_comp_add_of_ae, Cube.add_mul_mem_Icc,
  Cube.add_smul_mem_cubeSet, Cube.isoModulus_mono, Cube.isoModulus_add_le,
  Cube.cube_eLpNorm_higherDiff_le, Cube.isoModulus_le, Cube.isoModulus_congr_ae,
  Cube.diffRegion_eq_isoRegion, Cube.coordModulus_le_isoModulus, Cube.cubeModulus_le_isoModulus,
  Cube.lpIsoModulus_add_le, Cube.lpIsoModulus_le, Cube.lpIsoModulus_anti,
  Cube.lpCubeModulus_le_lpIsoModulus}
\leanok
\uses{def:iso-modulus, lem:cube-modulus, lem:rd-modulus-on}
Let $1 \leq p < \infty$. The isotropic modulus increases in $t$, is subadditive in $f$, satisfies
$\omega^{\mathrm{iso}}_r(f, t)_p \leq 2^r \|f\|_p$ and depends only on the class of $f$ in
$L_p(Q)$. The moduli in the coordinates satisfy
$\omega_r^{(i)}(f, t)_p \leq \omega^{\mathrm{iso}}_r(f, t)_p$, so
$\omega_r(f, t)_p \leq d \, \omega^{\mathrm{iso}}_r(f, t)_p$.
\end{lemma}

\begin{proof}
\leanok
For $x \in Q_{rh}$ and $0 \leq n \leq r$ the point $x + nh$ lies in $Q$, since each coordinate
$x_i + n h_i$ lies between $x_i$ and $x_i + r h_i$. Translations that stay in $Q$ preserve
$L_p$ norms and null sets. For the coordinate moduli take $h = s e_i$ with $0 \leq s \leq t$; on
$Q$ the region $Q_{r s e_i}$ is $\{x_i \leq 1 - rs\}$.
\end{proof}

\begin{theorem}[Isotropic differences from the coordinates]
\label{thm:iso-diff-bound}
\lean{Cube.eLpNorm_cubeDiff_inter_le, Cube.listRegion,
  Cube.listRegion_nil, Cube.listRegion_cons, Cube.measurableSet_listRegion,
  Cube.update_mem_listRegion_iff, Cube.eLpNorm_cubeMixedDiff_le,
  Cube.eLpNorm_shift_cubeMixedDiff_le, Cube.eLpNorm_higherDiff_le_sum, Cube.prod_rpow_eq_rpow_sum,
  Cube.eLpNorm_higherDiff_le_of_exponents, Cube.lpIsoModulus_le_of_exponents}
\leanok
\uses{lem:iso-expansion, lem:iso-modulus, lem:cube-slices, lem:step-regions}
Let $1 \leq p < \infty$, $r, j \in \N$, and let $V$ be a subspace of functions of one variable with
\[
  \|\Delta_t^m u\|_{L_p(J(m,t))} \leq C (2^j |t|)^{\theta_m} \|u\|_p
  \qquad (u \in V, \ m \leq r, \ |t| \leq 2^{-j})
\]
for exponents $\theta_m \geq 0$, with $J(m, t)$ as in \Cref{lem:step-regions}. Let
$\Theta \leq \sum_i \theta_{m_i}$ for all $m \in \N^d$ with $\sum_i m_i = r$. Then for every
measurable $f$ whose slices all lie in $V$ and every $h$ with $\|h\|_\infty \leq 2^{-j}$,
\[
  \|\Delta_h^r f\|_{L_p(Q_{rh})} \leq d^r C^d (2^j \|h\|_\infty)^{\Theta} \|f\|_p .
\]
If $\Theta \geq 0$, then for every $k \in \N$ every $a \in L_p(Q)$ represented by a measurable
function whose slices all lie in $V$ satisfies
\[
  \omega^{\mathrm{iso}}_r(a, 2^{-k})_p \leq (d^r C^d + 2^r) \, 2^{-\max(k - j, 0)\Theta} \|a\|_p .
\]
\end{theorem}

\begin{proof}
\leanok
By \Cref{lem:iso-expansion} it suffices to bound each of the $d^r$ terms. For $x \in Q_{rh}$
the point $x + s_\alpha$ lies in $Q$ and in the box $\{y : y_i + m_i h_i \in [0, 1] \ \forall i\}$,
since $\#\{l : i < \alpha_l\} + m_i \leq r$; by translation invariance the term is bounded by the
norm of $M_m f$ on the box. Apply the coordinate differences one at a time: the box constraint in
the other coordinates does not depend on the coordinate of the step, so Tonelli's theorem in that
coordinate (\Cref{lem:cube-slices}) and the bound on the slices, which stay in $V$, give the factor
$C (2^j |h_i|)^{\theta_{m_i}}$. The product is at most
$C^d (2^j \|h\|_\infty)^{\sum_i \theta_{m_i}} \leq C^d (2^j \|h\|_\infty)^\Theta$, since
$2^j \|h\|_\infty \leq 1$. For the modulus with $k \geq j$ take the supremum over
$\|h\|_\infty \leq 2^{-k}$, where $2^j \|h\|_\infty \leq 2^{-(k-j)}$; for $k < j$ use
$\omega^{\mathrm{iso}}_r(a, t)_p \leq 2^r \|a\|_p$ (\Cref{lem:iso-modulus}).
\end{proof}

\begin{lemma}[The gains of piecewise polynomials and splines]
\label{lem:iso-exponents}
\lean{Cube.ppSubmodule, Cube.splineSubmodule, Cube.IsCubePP.slice_mem_ppSubmodule,
  Cube.IsCubeSpline.slice_mem_splineSubmodule, Cube.ppExponent, Cube.splineExponent,
  Cube.ppExponent_nonneg, Cube.splineExponent_nonneg, Cube.le_sum_ppExponent,
  Cube.le_sum_splineExponent, Cube.exists_eLpNorm_higherDiff_stepRegion_pp_le,
  Cube.exists_eLpNorm_higherDiff_stepRegion_spline_le}
\leanok
\uses{thm:pp-lower-diff, thm:spline-lower-diff, thm:spline-inverse, lem:step-regions,
  def:cube-pp, def:cube-spline}
Let $1 \leq p < \infty$ and $r \geq 1$.
\begin{enumerate}
  \item The piecewise polynomials of level $j$ and degree $< r$ on $[0, 1)$ satisfy the hypothesis
    of \Cref{thm:iso-diff-bound} with $\theta_0 = 0$ and $\theta_m = 1/p$ for $m \geq 1$, and
    $\Theta = 1/p$ is admissible.
  \item The splines of order $r$ and level $j$ satisfy it with $\theta_m = m$ for $m < r$ and
    $\theta_r = r - 1 + 1/p$, and $\Theta = r - 1 + 1/p$ is admissible.
\end{enumerate}
The slices of piecewise polynomials (splines) on the cube are piecewise polynomials (splines).
\end{lemma}

\begin{proof}
\leanok
The bounds are \Cref{thm:pp-lower-diff}, \Cref{thm:spline-lower-diff} and \Cref{thm:spline-inverse}
for steps $t \geq 0$, and \Cref{lem:step-regions} for $t < 0$. If $\sum_i m_i = r \geq 1$, some
$m_i \geq 1$, which gives (1). For (2), either some $m_i = r$, or all $m_i < r$ and
$\sum_i \theta_{m_i} = \sum_i m_i = r \geq r - 1 + 1/p$.
\end{proof}

\begin{theorem}[The inverse inequalities with the isotropic modulus]
\label{thm:iso-inverse}
\lean{Cube.exists_inverse_cubePPScheme_iso, Cube.exists_inverse_cubeSplineScheme_iso}
\leanok
\uses{thm:iso-diff-bound, lem:iso-exponents, lem:iso-modulus}
Let $1 \leq p < \infty$ and $r \geq 1$. There is $C_B$ with
\[
  \omega^{\mathrm{iso}}_r(a, 2^{-k})_p \leq C_B \, 2^{-\max(k - j, 0)\theta} \|a\|_p
\]
for all $j, k \in \N$ and every piecewise polynomial $a$ of level $j$ and degree $< r$ on the cube,
with $\theta = 1/p$, and likewise for every spline $a$ of order $r$ and level $j$ on the cube, with
$\theta = r - 1 + 1/p$.
\end{theorem}

\begin{proof}
\leanok
The slices of $a$ lie in the spaces of \Cref{lem:iso-exponents}, so the bound for the modulus in
\Cref{thm:iso-diff-bound} applies with $\Theta = \theta$ and $C_B = d^r C^d + 2^r$.
\end{proof}

\begin{theorem}[Approximation spaces on the cube and the isotropic modulus]
\label{thm:iso-besov}
\lean{Cube.cubePPIsoModulusData, Cube.memApprox_cubePPScheme_iff_iso,
  Cube.cubeSplineIsoModulusData, Cube.memApprox_cubeSplineScheme_iff_iso,
  Cube.lqNorm_cubeModulus_lt_top_iff_iso}
\leanok
\uses{thm:iso-inverse, lem:iso-modulus, thm:cube-jackson, thm:cube-spline-jackson,
  thm:cube-besov, thm:cube-spline-besov, def:modulus-data, thm:modulus-coincidence}
Let $1 \leq p < \infty$, $1 \leq q \leq \infty$ and $r \geq 1$.
\begin{enumerate}
  \item For $0 \leq s < 1/p$, $f \in L_p(Q)$ lies in the approximation space $A^s_q$ of the
    piecewise polynomials of degree $< r$ in each variable if and only if
    $\|f\|_p + \| ( 2^{ks} \omega^{\mathrm{iso}}_r(f, 2^{-k})_p )_k \|_{\ell_q} < \infty$.
  \item For $0 \leq s < r - 1 + 1/p$, $f$ lies in the approximation space $A^s_q$ of the tensor
    splines of order $r$ if and only if the same condition holds.
  \item For $0 \leq s < r - 1 + 1/p$,
    $\|f\|_p + \| ( 2^{ks} \omega_r(f, 2^{-k})_p )_k \|_{\ell_q} < \infty$ with the coordinate
    moduli if and only if the same holds with the isotropic modulus.
\end{enumerate}
\end{theorem}

\begin{proof}
\leanok
Jackson's inequalities \Cref{thm:cube-jackson} and \Cref{thm:cube-spline-jackson} hold with the
isotropic modulus by \Cref{lem:iso-modulus}, at the cost of a factor $d$. With
\Cref{thm:iso-inverse} the isotropic modulus is compatible with both schemes; apply
\Cref{thm:modulus-coincidence}. Part (3) combines (2) with \Cref{thm:cube-spline-besov}.
\end{proof}

For $0 < s < r - 1 + 1/p$ the Besov spaces $\Besov{s}{p}{q}(Q)$ defined by the moduli in the
coordinates and by the isotropic modulus therefore coincide.

\section{The cube in Euclidean space}
\label{sec:cube-euclidean}

The theory of domains in \Cref{ch:lipschitz-domains} works in the Euclidean space $\R^d$, the
theory of the cube in $\R^d$ with the maximum norm. This section connects the two on the cube.

\begin{definition}[The cube in the Euclidean space]
\label{def:euclidean-cube}
\lean{Cube.euclideanCube, Cube.preimage_toLp_euclideanCube, Cube.preimage_ofLp_cubeSet,
  Cube.measurableSet_euclideanCube, Cube.measurePreserving_toLp_cubeVol,
  Cube.measurePreserving_ofLp_euclideanCube, Cube.norm_euclideanEquiv_symm_le,
  Cube.restrictEuclidean, Cube.extendEuclidean}
\leanok
\uses{def:cube, lem:modulus-on-linear-equiv, lem:zero-extension-lp}
Let $Q_E = [0, 1]^d$ in the Euclidean space $\R^d$. The identity of $\R^d$ carries $Q_E$ onto
the cube $Q$ of \Cref{def:cube} and Lebesgue measure on $Q_E$ to the cube measure, and
$\|x\|_\infty \leq \|x\|_2 \leq \sqrt{d} \, \|x\|_\infty \leq 2^d \|x\|_\infty$ for
$x \in \R^d$. For $1 \leq p \leq \infty$ the restriction $R \colon L_p(\R^d) \to L_p(Q)$,
$R F = F|_Q$, and the extension by zero $E \colon L_p(Q) \to L_p(\R^d)$ are bounded linear
operators.
\end{definition}

\begin{lemma}[Restriction and extension by zero]
\label{lem:euclidean-cube-transport}
\lean{Cube.restrictEuclidean_ae, Cube.extendEuclidean_ae, Cube.restrictEuclidean_extendEuclidean,
  Cube.indicatorLp_extendEuclidean_restrictEuclidean, Cube.restrictEuclidean_eq_of_indicatorLp_eq,
  Cube.enorm_restrictEuclidean, Cube.enorm_extendEuclidean}
\leanok
\uses{def:euclidean-cube, lem:indicator-lp}
Let $1 \leq p \leq \infty$, $g \in L_p(Q)$ and $F, G \in L_p(\R^d)$. Then $R(E g) = g$,
$\mathbf 1_{Q_E} E(R F) = \mathbf 1_{Q_E} F$, $\|R F\|_p = \|\mathbf 1_{Q_E} F\|_p$ and
$\|E g\|_p = \|g\|_p$, and $\mathbf 1_{Q_E} F = \mathbf 1_{Q_E} G$ implies $R F = R G$.
\end{lemma}

\begin{proof}
\leanok
\uses{lem:indicator-lp-eq-iff}
$R F = F$ and $E g = g$ almost everywhere on the cube, and $E g = 0$ outside of it.
\end{proof}

\begin{lemma}[The Euclidean and the isotropic modulus on the cube]
\label{lem:euclidean-cube-moduli}
\lean{Cube.lpModulusOn_euclideanCube_le_lpIsoModulus, Cube.lpIsoModulus_add_le_lpModulusOn}
\leanok
\uses{def:euclidean-cube, def:iso-modulus, def:rd-modulus-on}
Let $1 \leq p \leq \infty$, $r, k \in \N$ and $F \in L_p(\R^d)$, and let
$\omega_r(F, t, Q_E)_p$ be the modulus on $Q_E$ for the Euclidean norm
(\Cref{def:rd-modulus-on}). Then
\begin{align*}
  \omega_r(F, 2^{-k}, Q_E)_p &\leq \omega^{\mathrm{iso}}_r(R F, 2^{-k})_p , \\
  \omega^{\mathrm{iso}}_r(R F, 2^{-(k+d)})_p &\leq \omega_r(F, 2^{-k}, Q_E)_p .
\end{align*}
\end{lemma}

\begin{proof}
\leanok
\uses{lem:modulus-on-linear-equiv, lem:iso-modulus, lem:euclidean-cube-transport}
Both are moduli on the same set, with the norms $\|h\|_2$ and $\|h\|_\infty$ of the steps. Apply
\Cref{lem:modulus-on-linear-equiv} to the identity of $\R^d$ in both directions, with the norms
at most $1$ and at most $2^d$, and use that the moduli increase in $t$ and depend only on the
class of the function on the cube.
\end{proof}

\begin{proposition}[Intrinsic and isotropic quasi-norms on the cube]
\label{prop:euclidean-cube-norms}
\lean{Cube.modNormOn_euclideanCube_le_isoNorm, Cube.isoNorm_restrictEuclidean_le_modNormOn}
\leanok
\uses{def:euclidean-cube, def:iso-modulus, lem:modnorm-on}
Let $1 \leq p \leq \infty$, $r \in \N$, $s \geq 0$, $1 \leq q \leq \infty$ and
$F \in L_p(\R^d)$, and let
$N_{Q_E}(F) = \|F\|_{L_p(Q_E)} + \|(2^{ks} \omega_r(F, 2^{-k}, Q_E)_p)_k\|_{\ell_q}$. Then
\[
  N_{Q_E}(F) \leq \|R F\|_p
    + \bigl\|\bigl(2^{ks} \omega^{\mathrm{iso}}_r(R F, 2^{-k})_p\bigr)_k\bigr\|_{\ell_q}
  \leq \bigl(1 + 2^{ds} (d \, 2^r + 1)\bigr) N_{Q_E}(F) .
\]
\end{proposition}

\begin{proof}
\leanok
\uses{lem:euclidean-cube-moduli, lem:lq-basic, lem:euclidean-cube-transport, lem:iso-modulus}
By \Cref{lem:euclidean-cube-transport}, $\|R F\|_p = \|F\|_{L_p(Q_E)}$. The first bound is the
first comparison of \Cref{lem:euclidean-cube-moduli}. For the second, the isotropic modulus at
level $k + d$ is at most the Euclidean one at level $k$, the first $d$ levels are at most
$2^r \|R F\|_p$ (\Cref{lem:iso-modulus}), and a shift by $d$ levels costs the factor $2^{ds}$
(\Cref{lem:lq-basic}).
\end{proof}

\begin{theorem}[Tensor splines and the Euclidean modulus on the cube]
\label{thm:cube-splines-euclidean}
\lean{Cube.memApprox_cubeSplineScheme_restrictEuclidean_iff,
  Cube.exists_approxNormDyadic_restrictEuclidean_le}
\leanok
\uses{thm:iso-besov, prop:euclidean-cube-norms, thm:modulus-coincidence}
Let $1 \leq p < \infty$, $1 \leq q \leq \infty$, $r \geq 1$ and $F \in L_p(\R^d)$.
\begin{enumerate}
  \item For $0 \leq s < r - 1 + 1/p$, $R F$ lies in the approximation space $A^s_q$ of the
    tensor splines of order $r$ if and only if $N_{Q_E}(F) < \infty$.
  \item For $s \geq 0$ there is $C < \infty$, independent of $F$, with
    $\|R F\|^{\mathrm{dy}}_{A^s_q} \leq C \, N_{Q_E}(F)$.
\end{enumerate}
\end{theorem}

\begin{proof}
\leanok
By \Cref{prop:euclidean-cube-norms} the quasi-norm $N_{Q_E}(F)$ and the isotropic quasi-norm of
$R F$ are equivalent. Part (1) is then part (2) of \Cref{thm:iso-besov}; part (2) is Jackson's
direction of \Cref{thm:modulus-coincidence} for the isotropic modulus, which holds for every
$s \geq 0$.
\end{proof}

\chapter{Besov spaces on \texorpdfstring{$\R^d$}{R\^{}d}}
\label{ch:rd}

This chapter works on a finite-dimensional real inner product space $V$ of dimension $n$, such
as $\R^d$, with Lebesgue measure, and with $L_p = L_p(V)$ for $1 \leq p < \infty$; convolution,
approximate identities and the Littlewood--Paley kernels are those of \Cref{ch:fourier-rd}. The
Besov space $\Besov{s}{p}{q}$ is the approximation space $A^s_q$ of the band-limited scheme in
$L_p$. The chapter describes it by Littlewood--Paley blocks for $s > 0$ (\Cref{thm:lp-band}); by
differences for $0 < s < r$, through Jackson's inequality and the inverse inequality
(\Cref{thm:band-jackson}, \Cref{thm:band-inverse} and \Cref{thm:rd-besov}); for $q = p$ and
$0 < s < 1$ by the Gagliardo form of the Sobolev--Slobodeckij space $W^{s,p}(V)$
(\Cref{thm:sobolev-slobodeckij-besov}); by local means (\Cref{thm:local-means}); by atomic
decompositions with regular and with smooth atoms (\Cref{thm:atomic-characterisation} and
\Cref{thm:smooth-atomic-decomposition}); and as a real interpolation space between two Besov
spaces by differences (\Cref{thm:rd-besov-interpolate}).
Moduli of smoothness on subsets and pointwise multipliers prepare the domains of
\Cref{ch:lipschitz-domains}, and smooth atoms give the invariance of the Besov spaces under
smooth changes of variables (\Cref{thm:diffeo-invariance}), on which the manifolds of
\Cref{ch:manifolds} rest.

\section{Band-limited approximation}

\begin{definition}[Band-limited classes]
\label{def:band-space}
\lean{Rd.bandSpace, Rd.mem_bandSpace, Rd.bandScheme, Rd.mem_bandScheme_two_pow}
\leanok
\uses{def:lp-kernels, lem:conv-op, def:comap}
Let $1 \leq p < \infty$. The band-limited class of level $k \in \N$ is
$E_k = \{ g \in L_p : \Theta_k * g = g \}$. The band-limited scheme has $A_0 = \{0\}$ and
$A_n = E_{\lfloor \log_2 n \rfloor}$ for $n \geq 1$.
\end{definition}

On the Fourier side $\Theta_k * g = g$ says $\widehat g = \widehat{\Theta_k} \widehat g$: the
spectrum of $g$ lies in the ball of radius $2^{k+2}$, and every $g$ with spectrum in the ball of
radius $2^{k+1}$ belongs to $E_k$. These are the functions of exponential type of
\cite{Nikolskii}, described by a kernel, so that the definition works in every $L_p$ without
distributions.

\begin{lemma}[The band-limited classes]
\label{lem:band-space}
\lean{Rd.smul_mem_bandSpace, Rd.isClosed_bandSpace, Rd.bandSpace_mono}
\leanok
\uses{def:band-space, thm:kernel-identities}
Let $1 \leq p < \infty$. Every $E_k$ is a closed complex linear subspace of $L_p$, and
$E_k \subseteq E_{k+1}$ for every $k \in \N$. The band-limited scheme is an approximation scheme
with $A_{2^k} = E_k$.
\end{lemma}

\begin{proof}
\leanok
$E_k$ is the kernel of $T_{\Theta_k} - I$. For $g \in E_k$,
$\Theta_{k+1} * g = \Theta_{k+1} * (\Theta_k * g) = (\Theta_{k+1} * \Theta_k) * g
= \Theta_k * g = g$ by \Cref{lem:conv-op} and \Cref{thm:kernel-identities}. The last claim is
\Cref{def:comap}.
\end{proof}

\section{The Littlewood--Paley decomposition}

\begin{definition}[Littlewood--Paley blocks]
\label{def:lp-blocks}
\lean{Rd.lpBlock, Rd.littlewoodPaleyNorm}
\leanok
\uses{def:lp-kernels, lem:conv-op, def:lq-norm}
Let $1 \leq p < \infty$, $s \in \R$ and $1 \leq q \leq \infty$. The Littlewood--Paley blocks of
$f \in L_p$ are $\Delta_j f = \psi_j * f$, $j \in \N$, and the Littlewood--Paley quasi-norm is
$\|f\|^{\mathrm{LP}}_{s,q} = \|(2^{js} \|\Delta_j f\|_p)_j\|_{\ell_q}$.
\end{definition}

\begin{lemma}[Bounds and reconstruction]
\label{lem:lp-blocks}
\lean{Rd.enorm_lpBlock_le, Rd.lpBlock_succ, Rd.sum_lpBlock, Rd.tendsto_sum_lpBlock,
  Rd.hasSum_lpBlock}
\leanok
\uses{def:lp-blocks, lem:lp-kernels-basic, thm:approx-identity, lem:weighted-blocks}
Let $1 \leq p < \infty$ and $f \in L_p$. For all $j, J \in \N$,
$\|\Delta_j f\|_p \leq 2 \|\Phi\|_1 \|f\|_p$ and $\sum_{j=0}^J \Delta_j f = \Phi_J * f$, and
$\sum_{j=0}^J \Delta_j f \to f$ in $L_p$ as $J \to \infty$. If $s > 0$, $1 \leq q \leq \infty$
and $\|f\|^{\mathrm{LP}}_{s,q} < \infty$, then $\sum_j \Delta_j f = f$, the series converging
in $L_p$.
\end{lemma}

\begin{proof}
\leanok
The bound is \Cref{lem:conv-op} with $\|\psi_j\|_1 \leq 2 \|\Phi\|_1$, the sum telescopes by
\Cref{lem:lp-kernels-basic}, and $\Phi_J * f \to f$ by \Cref{thm:approx-identity}. For the last
claim apply \Cref{lem:weighted-blocks}(1).
\end{proof}

\begin{lemma}[Blocks and band-limited functions]
\label{lem:lp-blocks-band}
\lean{Rd.lpBlock_mem_bandSpace, Rd.lpBlock_eq_zero_of_mem, Rd.enorm_lpBlock_le_approxError}
\leanok
\uses{lem:lp-blocks, lem:band-space, thm:kernel-identities}
Let $1 \leq p < \infty$, $f \in L_p$ and $i, j \in \N$. Then $\Delta_j f \in E_j$. If
$j \geq i + 3$, then $\Delta_j g = 0$ for every $g \in E_i$, and
$\|\Delta_j f\|_p \leq 2 \|\Phi\|_1 \inf_{g \in E_i} \|f - g\|_p$.
\end{lemma}

\begin{proof}
\leanok
$\Theta_j * (\psi_j * f) = (\Theta_j * \psi_j) * f = \psi_j * f$ and
$\psi_j * g = \psi_j * (\Theta_i * g) = (\psi_j * \Theta_i) * g = 0$ by
\Cref{thm:kernel-identities}. Hence $\Delta_j f = \Delta_j (f - g)$ for $g \in E_i$, and
\Cref{lem:lp-blocks} bounds it by $2 \|\Phi\|_1 \|f - g\|_p$.
\end{proof}

\begin{theorem}[Band-limited approximation and Littlewood--Paley blocks]
\label{thm:lp-band}
\lean{Rd.approxNormDyadic_bandScheme_le, Rd.littlewoodPaleyNorm_le_approxNormDyadic,
  Rd.memApprox_bandScheme_iff}
\leanok
\uses{lem:lp-blocks, lem:lp-blocks-band, thm:block-decomposition, def:dyadic-norm,
  lem:weighted-blocks}
Let $1 \leq p < \infty$, $s > 0$ and $1 \leq q \leq \infty$, and let $A$ be the band-limited
scheme. For every $f \in L_p$,
\[
  \|f\|^{\mathrm{dy}}_{A^s_q} \leq 2 (1 - 2^{-s})^{-1} \, \|f\|^{\mathrm{LP}}_{s,q}
  \quad \text{and} \quad
  \|f\|^{\mathrm{LP}}_{s,q} \leq 8 \cdot 8^s \, \|\Phi\|_1 \, \|f\|^{\mathrm{dy}}_{A^s_q},
\]
and $f \in A^s_q$ if and only if $\|f\|^{\mathrm{LP}}_{s,q} < \infty$.
\end{theorem}

\begin{proof}
\leanok
By \Cref{lem:lp-blocks-band}, $\Delta_j f \in E_j = A_{2^j}$, and by \Cref{lem:lp-blocks} a
finite Littlewood--Paley quasi-norm makes $f = \sum_j \Delta_j f$ a block decomposition, so
\Cref{thm:block-decomposition} gives the first inequality. For the second, apply
\Cref{lem:weighted-blocks}(2) with $M = 3$ and $C = 2\|\Phi\|_1$ (\Cref{lem:lp-blocks},
\Cref{lem:lp-blocks-band}), and $3 \cdot 2 \leq 8$.
\end{proof}

With the Besov space $\Besov{s}{p}{q}(V)$ defined through Littlewood--Paley blocks,
\Cref{thm:lp-band} identifies it with the approximation space of the band-limited scheme.

\section{Differences and moduli of smoothness}

\begin{definition}[Differences and the modulus of smoothness]
\label{def:rd-modulus}
\lean{Rd.fwdDiff_eq_sub, Rd.higherDiff_apply, Rd.smoothnessModulus, Rd.lpModulus}
\leanok
For $f : V \to \C$, $h \in V$ and $r \in \N$ let $\Delta_h f = f(\cdot + h) - f$ and
$\Delta_h^r = (\Delta_h)^r$, so that
$\Delta_h^r f(x) = \sum_{l=0}^r (-1)^{r-l} \binom{r}{l} f(x + l h)$. For $0 < p \leq \infty$ and
$t \in \R$, the modulus of smoothness is
$\omega_r(f, t)_p = \sup_{\|h\| \leq t} \|\Delta_h^r f\|_p$.
\end{definition}

\begin{lemma}[Properties of differences]
\label{lem:rd-differences}
\lean{Rd.higherDiff_comp_add, Rd.fwdDiff_nsmul, Rd.aestronglyMeasurable_comp_add,
  Rd.eLpNorm_comp_add_right, Rd.ae_comp_add_right, Rd.aestronglyMeasurable_higherDiff,
  Rd.higherDiff_congr_ae, Rd.eLpNorm_higherDiff_le_two_pow_mul, Rd.integrable_higherDiff,
  Rd.conv_comp_add, Rd.higherDiff_conv_ae, Rd.higherDiff_dilate, Rd.higherDiff_neg,
  Rd.measurable_enorm_higherDiff_uncurry}
\leanok
\uses{def:rd-modulus, lem:conv-assoc, lem:dilate}
Let $1 \leq p < \infty$, $f \in L_p$, $K \in L_1$, $\Phi : V \to \C$, $h, c \in V$, $m, r \in \N$
and $\lambda > 0$. Then:
\begin{enumerate}
  \item $\Delta_h^r (f(\cdot + c)) = (\Delta_h^r f)(\cdot + c)$ and
    $\Delta_{mh} f = \sum_{i<m} (\Delta_h f)(\cdot + i h)$;
  \item $\|f(\cdot + c)\|_p = \|f\|_p$ and $\|\Delta_h^r f\|_p \leq 2^r \|f\|_p$;
  \item $\Delta_h^r K$ is integrable, and $\Delta_h^r (K * f) = (\Delta_h^r K) * f$ almost
    everywhere;
  \item $\Delta_h^r (\Phi_\lambda) = (\Delta_{\lambda h}^r \Phi)_\lambda$;
  \item $\Delta_{-h}^r f(x) = (-1)^r \Delta_h^r f(x - r h)$ for every $x \in V$;
  \item for every measurable $u : V \to \C$, $(x, v) \mapsto |\Delta_v^r u(x)|$ is measurable on
    $V \times V$.
\end{enumerate}
\end{lemma}

\begin{proof}
\leanok
Parts (1) and (4) follow by induction on $r$, and (2) holds since translations preserve Lebesgue
measure. For (3), $(K * f)(x + c) = (K(\cdot + c) * f)(x)$ by the substitution $y \mapsto y + c$,
and \Cref{lem:conv-assoc} gives the linearity in the kernel. Part (5) reflects the sum of
\Cref{def:rd-modulus} at $l \mapsto r - l$, with $(-1)^r (-1)^{r-l} = (-1)^l$, and in (6) that
sum is a finite sum of measurable functions of $(x, v)$.
\end{proof}

\begin{lemma}[Properties of the modulus on $\R^d$]
\label{lem:rd-modulus}
\lean{Rd.eLpNorm_higherDiff_nsmul_higherDiff_le, Rd.eLpNorm_higherDiff_le_smoothnessModulus,
  Rd.smoothnessModulus_mono, Rd.smoothnessModulus_add_le, Rd.smoothnessModulus_le,
  Rd.smoothnessModulus_congr_ae, Rd.smoothnessModulus_nat_mul_le_of_forall_exists_nsmul,
  Rd.smoothnessModulus_nat_mul_le, Rd.smoothnessModulus_mul_le_of_nat_mul_le,
  Rd.smoothnessModulus_mul_le, Rd.lpModulus_add_le, Rd.lpModulus_le, Rd.lpModulus_zero,
  Rd.smoothnessModulus_const_smul, Rd.lpModulus_const_smul}
\leanok
\uses{def:rd-modulus, lem:rd-differences}
Let $1 \leq p < \infty$, $r \in \N$ and $f, g \in L_p$. The modulus $\omega_r(f, t)_p$ increases
in $t$, depends only on the class of $f$ modulo null sets, and satisfies
$\omega_r(f + g, t)_p \leq \omega_r(f, t)_p + \omega_r(g, t)_p$,
$\omega_r(a f, t)_p = |a| \, \omega_r(f, t)_p$ for every $a \in \C$ and
$\omega_r(f, t)_p \leq 2^r \|f\|_p$. For every integer $m \geq 1$ and every $t \geq 0$,
$\omega_r(f, m t)_p \leq m^r \, \omega_r(f, t)_p$, and for all $c, t \geq 0$,
\[
  \omega_r(f, c t)_p \leq (1 + c)^r \, \omega_r(f, t)_p .
\]
More generally, let $\mu$ be an additive Haar measure on a normed group $V$, let $f$ be almost
everywhere strongly measurable and let $t \geq 0$. For every integer $m \geq 1$: if every $h$
with $\|h\| \leq m t$ is $m h'$ for some $h'$ with $\|h'\| \leq t$, then
$\omega_r(f, m t)_p \leq m^r \omega_r(f, t)_p$. If this bound holds for every integer
$m \geq 1$, then $\omega_r(f, c t)_p \leq (1 + c)^r \omega_r(f, t)_p$ for every $c \geq 0$.
\end{lemma}

\begin{proof}
\leanok
For an integer $m \geq 1$, part (1) of \Cref{lem:rd-differences} gives
$\Delta_{mh}^{a+1} \Delta_h^b f = \sum_{i<m} (\Delta_{mh}^a \Delta_h^{b+1} f)(\cdot + i h)$, and
induction on $a$ yields $\|\Delta_{mh}^r f\|_p \leq m^r \|\Delta_h^r f\|_p$, so
$\omega_r(f, m t)_p \leq m^r \omega_r(f, t)_p$. For real $c$ take
$m = \lfloor c \rfloor + 1 \leq 1 + c$.
\end{proof}

\begin{lemma}[Moduli of convolutions]
\label{lem:rd-modulus-conv}
\lean{Rd.lpModulus_convLp_le}
\leanok
\uses{lem:rd-differences, thm:young}
Let $1 \leq p < \infty$, $r, k \in \N$, $B \geq 0$ and $K \in L_1$ with
$\|\Delta_v^r K\|_1 \leq B$ for all $\|v\| \leq 2^{-k}$. Then
$\omega_r(K * h, 2^{-k})_p \leq B \, \|h\|_p$ for every $h \in L_p$.
\end{lemma}

\begin{proof}
\leanok
$\Delta_v^r (K * h) = (\Delta_v^r K) * h$ almost everywhere by \Cref{lem:rd-differences}, and
Young's inequality gives $\|\Delta_v^r (K * h)\|_p \leq \|\Delta_v^r K\|_1 \|h\|_p$.
\end{proof}

\begin{lemma}[The modulus at small scales]
\label{lem:modulus-to-zero}
\lean{Rd.tendsto_smoothnessModulus, Rd.tendsto_lpModulus}
\leanok
\uses{def:rd-modulus}
Let $1 \leq p < \infty$, $r \geq 1$, let $\mu$ be an additive Haar measure on $V$ and
$f \in L_p(\mu)$. Then $\omega_r(f, t)_p \to 0$ as $t \to 0$; in particular
$\omega_r(f, 2^{-k})_p \to 0$ as $k \to \infty$.
\end{lemma}

\begin{proof}
\leanok
\uses{lem:rd-differences, lem:translation-continuity}
For $\|h\| \leq t$, part (2) of \Cref{lem:rd-differences} applied to $\Delta_h f$ gives
$\|\Delta_h^r f\|_p \leq 2^{r-1} \|\Delta_h f\|_p = 2^{r-1} \|f(\cdot + h) - f\|_p$. Translation is
continuous in $L_p$ (\Cref{lem:translation-continuity}), so for every $\varepsilon > 0$ there is
$\delta > 0$ with $\|f(\cdot + h) - f\|_p \leq 2^{1-r} \varepsilon$ for $\|h\| < \delta$, and then
$\omega_r(f, t)_p \leq \varepsilon$ for $|t| < \delta$.
\end{proof}

\begin{lemma}[Moduli under linear changes of variables]
\label{lem:modulus-linear-equiv}
\lean{Rd.higherDiff_comp_addHom, Rd.smoothnessModulus_comp_le,
  Rd.lpModulus_compMeasurePreserving_le, MeasureTheory.Lp.enorm_compMeasurePreserving,
  Rd.modNorm_compMeasurePreserving_le, Rd.modNorm_compMeasurePreserving_lt_top_iff}
\leanok
\uses{def:rd-modulus}
Let $V$ and $W$ be finite-dimensional real normed spaces with additive Haar measures $\mu$ and
$\nu$, and let $e : V \to W$ be linear and continuous with $\nu(B) = \mu(e^{-1}(B))$ for every
Borel set $B$. For every $f : W \to \C$ and $r \in \N$,
$\Delta_h^r (f \circ e) = (\Delta_{e h}^r f) \circ e$ for every $h \in V$, and for measurable $f$,
every $0 < p \leq \infty$ and every $t$, $\omega_r(f \circ e, t)_p \leq \omega_r(f, \|e\| t)_p$.
Let $1 \leq p \leq \infty$, $s \in \R$, $1 \leq q \leq \infty$ and $F \in L_p(\nu)$. Then
$\|F \circ e\|_{L_p(\mu)} = \|F\|_{L_p(\nu)}$ and
\begin{align*}
  &\|F \circ e\|_p
    + \bigl\|\bigl(2^{ks} \omega_r(F \circ e, 2^{-k})_p\bigr)_{k \in \N}\bigr\|_{\ell_q} \\
  &\qquad \leq (1 + \|e\|)^r \Bigl(\|F\|_p
    + \bigl\|\bigl(2^{ks} \omega_r(F, 2^{-k})_p\bigr)_{k \in \N}\bigr\|_{\ell_q}\Bigr) .
\end{align*}
If $e$ is bijective, the two quasi-norms are finite together.
\end{lemma}

\begin{proof}
\leanok
\uses{lem:rd-modulus}
The identity for differences follows by induction on $r$ from $e(x + h) = e x + e h$. Since $e$
carries $\mu$ to $\nu$, $\|G \circ e\|_{L_p(\mu)} = \|G\|_{L_p(\nu)}$ for every measurable $G$, so
$\|\Delta_h^r (f \circ e)\|_p = \|\Delta_{e h}^r f\|_p \leq \omega_r(f, \|e\| t)_p$ for
$\|h\| \leq t$, as $\|e h\| \leq \|e\| \|h\|$. By \Cref{lem:rd-modulus},
$\omega_r(F, \|e\| 2^{-k})_p \leq (1 + \|e\|)^r \omega_r(F, 2^{-k})_p$, and
$(1 + \|e\|)^r \geq 1$. For bijective $e$ the inverse $e^{-1}$ carries $\nu$ to $\mu$ and
$(F \circ e) \circ e^{-1} = F$, so the bound for $e^{-1}$ gives the converse.
\end{proof}

\section{Moduli of smoothness on subsets}

\begin{definition}[The modulus of smoothness on a subset]
\label{def:rd-modulus-on}
\lean{Rd.diffStepSet, Rd.diffStepSet_subset, Rd.diffStepSet_univ, Rd.modulusOn, Rd.lpModulusOn}
\leanok
\uses{def:rd-modulus}
For $U \subseteq V$, $h \in V$ and $r \in \N$ let
$U_{rh} = \{x \in V : x, x + h, \dots, x + rh \in U\}$, so that $U_{rh} \subseteq U$ and
$V_{rh} = V$. For $f : V \to \C$, $0 < p \leq \infty$ and $t \in \R$, the modulus of smoothness
on $U$ is
\[
  \omega_r(f, t, U)_p = \sup_{\|h\| \leq t} \|\Delta_h^r f\|_{L_p(U_{rh})} .
\]
\end{definition}

This is the modulus of \cite[(2.1)]{DeVoreSharpley}, written with the restricted differences
$\Delta_h^r(f, x, U) = \Delta_h^r f(x)$ for $x \in U_{rh}$ and $\Delta_h^r(f, x, U) = 0$ otherwise.

\begin{lemma}[Properties of the modulus on a subset]
\label{lem:rd-modulus-on}
\lean{Rd.indicator_diffStepSet_higherDiff_indicator, Rd.mem_diffStepSet_neg,
  Rd.indicator_diffStepSet_neg, Rd.eLpNorm_indicator_diffStepSet_neg, Rd.measurableSet_diffStepSet,
  Rd.diffStepSet_mono, Rd.eLpNorm_indicator_higherDiff_le_modulusOn, Rd.modulusOn_mono,
  Rd.modulusOn_univ, Rd.modulusOn_indicator, Rd.modulusOn_add_le, Rd.modulusOn_le_smoothnessModulus,
  Rd.modulusOn_le, Rd.modulusOn_congr_ae, Rd.modulusOn_mono_set, Rd.lpModulusOn_succ_le,
  Rd.lpModulusOn_le_lpModulus, Rd.lpModulusOn_mono_set, Rd.lpModulusOn_congr,
  Rd.lpModulusOn_add_le, Rd.lpModulusOn_le, Rd.modulusOn_const_smul, Rd.lpModulusOn_const_smul}
\leanok
\uses{def:rd-modulus-on, lem:rd-modulus}
Let $U \subseteq V$ be measurable, $1 \leq p < \infty$, $r \in \N$ and $f, g \in L_p$. Then
$\omega_r(f, t, U)_p$ increases in $t$, depends only on the class of $f|_U$ modulo null sets,
and satisfies $\omega_r(f + g, t, U)_p \leq \omega_r(f, t, U)_p + \omega_r(g, t, U)_p$,
\[
  \omega_r(f, t, U)_p \leq \omega_r(f, t)_p
  \qquad\text{and}\qquad
  \omega_r(f, t, U)_p \leq 2^r \|f\|_{L_p(U)} .
\]
Moreover $\omega_r(f, t, V)_p = \omega_r(f, t)_p$,
$\omega_r(f, t, U)_p \leq \omega_r(f, t, O)_p$ for measurable $U \subseteq O$, and the steps $h$
and $-h$ give the same norm: $x \in U_{r(-h)}$ if and only if $x - rh \in U_{rh}$,
$\Delta_{-h}^r f(x) = (-1)^r \Delta_h^r f(x - rh)$, and
$\|\Delta_{-h}^r f\|_{L_p(U_{r(-h)})} = \|\Delta_h^r f\|_{L_p(U_{rh})}$ for every function $f$.
\end{lemma}

\begin{proof}
\leanok
For $x \in U_{rh}$ all points $x + lh$, $l \leq r$, lie in $U$, so $\Delta_h^r (1_U f) = \Delta_h^r f$
on $U_{rh}$. Hence $\omega_r(f, t, U)_p = \omega_r(1_U f, t, U)_p \leq \omega_r(1_U f, t)_p \leq
2^r \|1_U f\|_p$ by \Cref{lem:rd-modulus}. The reflection $l \mapsto r - l$ maps the points
$x - lh$ onto $(x - rh) + lh$, and the translation by $rh$ preserves the measure; the other
statements follow from the definition.
\end{proof}

\begin{lemma}[The modulus on $\Omega$ through a function on $V$]
\label{lem:modulus-on-extension}
\lean{Rd.lpModulusOn_le_add_lpModulus}
\leanok
\uses{lem:rd-modulus-on, lem:indicator-lp}
For every measurable $\Omega \subseteq V$, every $p \geq 1$, every order $r$, all
$f, G \in L_p(V)$ and every $k \in \N$,
\[
  \omega_r(f, 2^{-k}, \Omega)_p
    \leq 2^r \|\mathbf 1_\Omega (f - G)\|_p + \omega_r(G, 2^{-k})_p .
\]
\end{lemma}

\begin{proof}
\leanok
Split $f = (f - G) + G$. The modulus on $\Omega$ of $f - G$ depends only on
$\mathbf 1_\Omega (f - G)$ and is at most $2^r \|\mathbf 1_\Omega (f - G)\|_p$, and that of $G$ is
at most the modulus of $G$ on $V$ (\Cref{lem:rd-modulus-on}).
\end{proof}

\begin{lemma}[The intrinsic quasi-norm]
\label{lem:modnorm-on}
\lean{Rd.modNormOn, Rd.modNormOn_zero, Rd.modNormOn_add_le, Rd.modNormOn_const_smul,
  Rd.modNormOn_finsetSum_le}
\leanok
\uses{def:rd-modulus-on, lem:rd-modulus-on}
On a measurable set $\Omega$, $N_\Omega(G) = \|G\|_{L_p(\Omega)} + \|(2^{ks} \omega_r(G, 2^{-k},
\Omega)_p)_k\|_{\ell_q}$ vanishes at $0$, is homogeneous and, for $q \geq 1$, subadditive, also
for finite sums.
\end{lemma}

\begin{proof}
\leanok
The modulus on $\Omega$ is homogeneous and subadditive (\Cref{lem:rd-modulus-on}), and so are the
$L_p$ and $\ell_q$ norms.
\end{proof}

\begin{lemma}[Null sets do not matter]
\label{lem:null-set-domains}
\lean{Rd.diffStepSet_ae_eq, Rd.modulusOn_congr_set, Rd.lpModulusOn_congr_set,
  Rd.indicatorLp_congr_set, Rd.modNormOn_congr_set}
\leanok
\uses{def:rd-modulus-on, lem:modnorm-on}
If $U$ and $U'$ agree up to a null set, so do their step sets $U_{rh}$ and $U'_{rh}$, and
$\omega_r(f, t, U)_p = \omega_r(f, t, U')_p$, $1_U f = 1_{U'} f$ in $L_p$ and
$N_U(f) = N_{U'}(f)$.
\end{lemma}

\begin{proof}
\leanok
The step sets are intersections of translates, and translations preserve the Lebesgue measure.
\end{proof}

\begin{lemma}[Moduli on subsets under linear changes of variables]
\label{lem:modulus-on-linear-equiv}
\lean{Rd.diffStepSet_preimage, Rd.modulusOn_comp_le, PiLp.norm_continuousLinearEquiv_le,
  PiLp.norm_continuousLinearEquiv_symm_le}
\leanok
\uses{def:rd-modulus-on, lem:modulus-linear-equiv}
Let $V$, $W$, $\mu$, $\nu$ and $e$ be as in \Cref{lem:modulus-linear-equiv}, and let
$U \subseteq W$ be measurable. Then $(e^{-1}U)_{rh} = e^{-1}(U_{r\,eh})$ for every $h \in V$ and
$r \in \N$, and for every measurable $f : W \to \C$, every $0 < p \leq \infty$ and every $t$,
\[
  \omega_r(f \circ e, t, e^{-1}U)_{L_p(\mu)} \leq \omega_r(f, \|e\| t, U)_{L_p(\nu)} .
\]
For $1 \leq p \leq \infty$ and normed spaces $\beta_i$, $i$ in a finite set $I$, the identity map
from $\ell_p^I(\beta)$ to the product $\prod_i \beta_i$ with the maximum norm has norm at most
$1$, and its inverse has norm at most $(\# I)^{1/p}$.
\end{lemma}

\begin{proof}
\leanok
\uses{lem:rd-modulus-on}
Since $e(x + lh) = e x + l \, e h$, the points $x, x + h, \dots, x + rh$ lie in $e^{-1}U$ if and
only if $e x, e x + e h, \dots, e x + r \, e h$ lie in $U$. Hence the restricted difference of
$f \circ e$ with step $h$ is the restricted difference of $f$ with step $e h$, composed with $e$,
and $e$ preserves the measures; moreover $\|e h\| \leq \|e\| t$ for $\|h\| \leq t$. The norm
bounds are $\max_i \|x_i\| \leq \|x\|_p \leq (\# I)^{1/p} \max_i \|x_i\|$.
\end{proof}

\begin{lemma}[Moduli on subsets under affine maps]
\label{lem:modulus-on-affine}
\lean{Rd.measurePreserving_affine, Rd.eLpNorm_comp_affine, Rd.ae_eq_comp_affine,
  Rd.memLp_comp_affine, Rd.diffStepSet_preimage_affine, Rd.higherDiff_comp_affine,
  Rd.modulusOn_comp_affine_le, Rd.compAffineLp, Rd.compAffineLp_ae, Rd.enorm_compAffineLp,
  Rd.compAffineLp_inv_compAffineLp, Rd.compAffineLp_compAffineLp_inv,
  Rd.lpModulusOn_compAffineLp_le, Rd.enorm_indicatorLp_compAffineLp,
  Rd.exists_modNormOn_compAffineLp_le}
\leanok
\uses{def:rd-modulus-on, lem:modnorm-on}
Let $V$ be a finite-dimensional real normed space of dimension $n$ with an additive Haar
measure $\mu$, let $a \in V$, $L > 0$ and $A x = a + L x$. Then $\mu(A^{-1} B) = L^{-n} \mu(B)$
for every Borel set $B$. For $1 \leq p \leq \infty$ and every measurable $G$,
$\|G \circ A\|_p = L^{-n/p} \|G\|_p$, with $L^{-n/p} = 1$ for $p = \infty$, and
$G \mapsto G \circ A$ is a bounded operator on $L_p(\mu)$, inverted on both sides by the
composition with $A^{-1} y = -L^{-1} a + L^{-1} y$. For every
$U \subseteq V$, $r \in \N$ and $h \in V$, $(A^{-1}U)_{rh} = A^{-1}(U_{r\,Lh})$ and
$\Delta_h^r (f \circ A) = (\Delta_{Lh}^r f) \circ A$; for measurable $U$ and $f$ and every $t$,
\[
  \omega_r(f \circ A, t, A^{-1}U)_p \leq L^{-n/p} \, \omega_r(f, L t, U)_p .
\]
For measurable $U$, $r \in \N$, $s \geq 0$ and $1 \leq q \leq \infty$ there is $C < \infty$ with
$N_{A^{-1}U}(G \circ A) \leq C \, N_U(G)$ for every $G \in L_p(\mu)$, with the intrinsic
quasi-norms of \Cref{lem:modnorm-on}.
\end{lemma}

\begin{proof}
\leanok
\uses{lem:modulus-on-linear-equiv, lem:rd-modulus-on, lem:lq-basic}
The dilation $x \mapsto L x$ carries $\mu$ to $L^{-n} \mu$, and translations preserve
$L^{-n} \mu$; hence $\int |G \circ A|^p \, d\mu = L^{-n} \int |G|^p \, d\mu$ for $p < \infty$,
and for $p = \infty$ the measures $\mu$ and $L^{-n} \mu$ have the same null sets. As in
\Cref{lem:modulus-on-linear-equiv}, since $A(x + lh) = A x + l \, L h$, the points
$x, \dots, x + rh$ lie in $A^{-1}U$ if and only if $A x, \dots, A x + r \, L h$ lie in $U$, and
$\Delta_h^r (f \circ A) = (\Delta_{Lh}^r f) \circ A$; with $\|L h\| \leq L t$ for
$\|h\| \leq t$ this gives the bound of the moduli. For the
quasi-norms choose $m \in \N$ with $L < 2^m$. Then $L \, 2^{-(k+m)} \leq 2^{-k}$, so
$\omega_r(G \circ A, 2^{-(k+m)}, A^{-1}U)_p \leq L^{-n/p} \omega_r(G, 2^{-k}, U)_p$, while the
first $m$ levels are at most $2^r L^{-n/p} \|G\|_{L_p(U)}$ by \Cref{lem:rd-modulus-on}. The shift of
$m$ levels costs $2^{ms}$ (\Cref{lem:lq-basic}), and $C = L^{-n/p} (1 + 2^{ms} (m 2^r + 1))$.
The compositions with $A$ and with $A^{-1}$ undo each other, since $A(A^{-1} y) = y$ and
$A^{-1}(A x) = x$.
\end{proof}

\section{Moment kernels}

\begin{lemma}[The coefficients of the moment kernels]
\label{lem:moment-coefficients}
\lean{Rd.jacksonCoeff, Rd.neg_one_pow_mul_neg_one_pow_sub, Rd.neg_one_pow_mul_sum_higherDiff_coeff,
  Rd.sum_jacksonCoeff}
\leanok
Let $L \in \N$ and $c_l = (-1)^{l+1} \binom{L}{l}$ for $1 \leq l \leq L$. For all
$F_0, \dots, F_L \in \C$,
\[
  (-1)^L \sum_{l=0}^L (-1)^{L-l} \binom{L}{l} F_l = F_0 - \sum_{l=1}^L c_l F_l ,
\]
and $\sum_{l=1}^L c_l = 1$ for $L \geq 1$.
\end{lemma}

\begin{proof}
\leanok
$(-1)^L (-1)^{L-l} = (-1)^l$ for $l \leq L$, and $\sum_{l=0}^L (-1)^l \binom{L}{l} = (1 - 1)^L = 0$
for $L \geq 1$.
\end{proof}

\begin{lemma}[Moments of the coefficients of the moment kernels]
\label{lem:jackson-moments}
\lean{Rd.sum_jacksonCoeff_mul_pow}
\leanok
\uses{lem:moment-coefficients}
For all $r, j \in \N$ with $1 \leq j < r$,
$\sum_{l < r} (-1)^l \binom{r}{l+1} (l + 1)^j = 0$.
\end{lemma}

\begin{proof}
\leanok
By \Cref{lem:moment-coefficients} the sum is $0^j - (-1)^r \Delta^r_1 (x^j)(0)$, and differences
of order $r$ annihilate polynomials of degree $j < r$ (Mathlib's
\texttt{fwdDiff\_\allowbreak iter\_\allowbreak pow\_\allowbreak eq\_\allowbreak zero\_\allowbreak
of\_\allowbreak lt}).
\end{proof}

\begin{definition}[Moment kernels]
\label{def:moment-kernel}
\lean{Rd.momentKernel, Rd.dilate_dilate, Rd.dilate_momentKernel, Rd.integrable_momentKernel,
  Rd.eLpNorm_momentKernel_ne_top, Rd.eLpNorm_momentKernel_lt_top,
  Rd.integral_momentKernel, Rd.lintegral_enorm_ne_zero_of_integral_eq_one}
\leanok
\uses{lem:moment-coefficients, lem:dilate}
Let $g \in L_1$ with $\int_V g = 1$ and $L \in \N$. The moment kernel of order $L$ is
\[
  m = \sum_{l=1}^L (-1)^{l+1} \binom{L}{l} \, g_{1/l} .
\]
Its dilations are $m_c = \sum_{l=1}^L (-1)^{l+1} \binom{L}{l} \, g_{c/l}$ for $c > 0$,
$\|m\|_1 < \infty$, and $\int_V m = 1$ for $L \geq 1$.
\end{definition}

\begin{lemma}[First moments of the moment kernels of order two]
\label{lem:moment-kernel-first-moment}
\lean{Rd.integral_inner_mul_momentKernel_two}
\leanok
\uses{def:moment-kernel, lem:first-moments-dilate}
For every $w \in V$ and every $g \colon V \to \C$ such that $x \mapsto \langle w, x \rangle g(x)$
is integrable, the moment kernel $m = 2 g - g_{1/2}$ of order two satisfies
$\int_V \langle w, x \rangle \, m(x) \, dx = 0$.
\end{lemma}

\begin{proof}
\leanok
By \Cref{lem:first-moments-dilate} the first moment is $2 M - 2 M = 0$ with
$M = \int_V \langle w, x \rangle g(x) \, dx$.
\end{proof}

\begin{lemma}[Averaged differences]
\label{lem:averaged-differences}
\lean{Rd.conv_finsetSum_kernel_ae, Rd.ae_sub_sum_conv_dilate_eq,
  Rd.ae_sub_conv_dilate_momentKernel_eq}
\leanok
\uses{def:moment-kernel, lem:conv-assoc, def:rd-modulus}
Let $1 \leq p < \infty$, $g \in L_1$ with $\int_V g = 1$, $L \in \N$, $c > 0$ and $f \in L_p$.
For almost every $x \in V$,
\[
  f(x) - (m_c * f)(x) = (-1)^L \int_V g(z) \, \Delta^L_{-c^{-1} z} f(x) \, dz .
\]
\end{lemma}

\begin{proof}
\leanok
Write $(g_{c/l} * f)(x) = \int_V g(z) f(x - l c^{-1} z) \, dz$ and $f(x) = \int_V g(z) f(x) \, dz$,
and apply \Cref{lem:moment-coefficients} with $F_l = f(x - l c^{-1} z)$.
\end{proof}

\begin{theorem}[Moment kernels approximate with the modulus]
\label{thm:moment-approx}
\lean{Rd.eLpNorm_indicator_sub_convLp_momentKernel_le, Rd.enorm_sub_convLp_momentKernel_le}
\leanok
\uses{lem:averaged-differences, lem:weighted-jensen, lem:rd-modulus-on}
Let $1 \leq p < \infty$, $L \in \N$, $c > 0$, and let $g \in L_1$ with $\int_V g = 1$ vanish
outside the ball of radius $R$. Let $U, T \subseteq V$ with $U$ measurable and $U \subseteq T_{Lh}$
for every $h = -c^{-1} z$ with $g(z) \neq 0$. Then for every $f \in L_p$,
\[
  \|f - m_c * f\|_{L_p(U)} \leq \|g\|_1 \, \omega_L(f, R/c, T)_p ,
\]
and in particular $\|f - m_c * f\|_p \leq \|g\|_1 \, \omega_L(f, R/c)_p$.
\end{theorem}

\begin{proof}
\leanok
By \Cref{lem:averaged-differences},
$1_U (f - m_c * f) = (-1)^L \int_V g(z) \, 1_U \Delta^L_{-c^{-1} z} f \, dz$ almost everywhere,
and $|1_U \Delta^L_h f| \leq |1_{T_{Lh}} \Delta^L_h f|$ for the steps $h = -c^{-1} z$ with
$g(z) \neq 0$, which have length at most $R/c$. The weighted Jensen inequality
(\Cref{lem:weighted-jensen}) gives
$\|f - m_c * f\|_{L_p(U)}^p \leq \|g\|_1^{p-1} \int_V |g(z)| \, \omega_L(f, R/c, T)_p^p \, dz$.
The case $U = T = V$ is the second statement, by \Cref{lem:rd-modulus-on}.
\end{proof}

\section{Jackson's inequality}

\begin{lemma}[The Jackson approximant]
\label{lem:jackson-approx}
\lean{Rd.conv_dilatedBandKernel_dilate_lpKernel, Rd.jacksonScale, Rd.jacksonScale_pos,
  Rd.jacksonScale_le, Rd.jacksonApprox, Rd.jacksonApprox_mem_bandSpace,
  Rd.ae_sub_jacksonApprox_eq}
\leanok
\uses{lem:averaged-differences, lem:lp-kernels-frequency, lem:conv-of-fourier, lem:band-space,
  lem:dilate}
Let $1 \leq p < \infty$, $r, k \in \N$ and $f \in L_p$, and let
\[
  J_k f = \sum_{l=1}^r (-1)^{l+1} \binom{r}{l} \, \Phi_{2^k/l} * f .
\]
Then $J_k f \in E_k$, and for almost every $x \in V$,
\[
  f(x) - J_k f(x) = (-1)^r \int_V \Phi(z) \, \Delta^r_{-2^{-k} z} f(x) \, dz .
\]
\end{lemma}

\begin{proof}
\leanok
The transform $\widehat{\Phi_c} = \chi_{1,2}(c^{-1} \cdot)$ vanishes outside the ball of radius $2c$,
and $\widehat{\Theta_k} = 1$ on the ball of radius $2^{k+1}$, so $\Theta_k * \Phi_c = \Phi_c$ for
$c \leq 2^k$ by \Cref{lem:conv-of-fourier}. The identity is \Cref{lem:averaged-differences} for
$g = \Phi$, $L = r$ and $c = 2^k$, since $\int_V \Phi = \widehat{\Phi}(0) = 1$.
\end{proof}

Jackson's inequality averages differences with the weight $|\Phi|$. The weighted average is a
lemma of its own, since it serves the torus as well (\Cref{thm:torus-lp-jackson}).

\begin{lemma}[Weighted averages with the Jackson weight]
\label{lem:jackson-weighted-average}
\lean{one_add_pow_le_two_pow_mul, Rd.bandJacksonConst, Rd.bandJacksonConst_ne_top,
  Rd.eLpNorm_integral_lpKernel_mul_le}
\leanok
\uses{lem:weighted-jensen, lem:lp-kernels-basic}
Let $1 \leq p < \infty$ and $r \in \N$, and let $C \geq 0$ with
$C^p = \|\Phi\|_1^{p-1} \int_V |\Phi(z)| (1 + |z|)^{rp} \, dz$. Then $C < \infty$. For every
s-finite measure space $(X, \mu)$, every $M < \infty$ and every measurable
$G \colon X \times V \to \C$ with $\|G(\cdot, z)\|_{L_p(\mu)} \leq (1 + |z|)^r M$ for every
$z \in V$,
\[
  \Bigl\|\int_V \Phi(z) \, G(\cdot, z) \, dz\Bigr\|_{L_p(\mu)} \leq C M .
\]
\end{lemma}

\begin{proof}
\leanok
The constant is finite, since $\Phi$ is a Schwartz function,
$(1 + t)^N \leq 2^N (1 + t^N)$ for $t \geq 0$, and $|\Phi(z)| (1 + |z|)^N$ is integrable for every
$N$. By \Cref{lem:weighted-jensen},
\begin{align*}
  \Bigl\|\int_V \Phi(z) \, G(\cdot, z) \, dz\Bigr\|_{L_p(\mu)}^p
    &\leq \|\Phi\|_1^{p-1} \int_V |\Phi(z)| \, \|G(\cdot, z)\|_{L_p(\mu)}^p \, dz \\
    &\leq \|\Phi\|_1^{p-1} \int_V |\Phi(z)| (1 + |z|)^{rp} \, dz \; M^p = C^p M^p .
\end{align*}
\end{proof}

\begin{theorem}[Jackson's inequality for band-limited functions]
\label{thm:band-jackson}
\lean{Rd.enorm_sub_jacksonApprox_le, Rd.approxError_bandScheme_le}
\leanok
\uses{lem:jackson-approx, lem:rd-modulus, lem:jackson-weighted-average}
Let $1 \leq p < \infty$ and $r \in \N$. There is $C < \infty$ with
\[
  \inf_{g \in E_k} \|f - g\|_p \leq \|f - J_k f\|_p \leq C \, \omega_r(f, 2^{-k})_p
\]
for all $f \in L_p$ and $k \in \N$; one can take
$C^p = \|\Phi\|_1^{p-1} \int_V |\Phi(z)| (1 + |z|)^{rp} \, dz$.
\end{theorem}

\begin{proof}
\leanok
By \Cref{lem:jackson-approx},
$f - J_k f = (-1)^r \int_V \Phi(z) \, \Delta^r_{-2^{-k} z} f \, dz$ almost everywhere, and
$\|\Delta^r_{-2^{-k} z} f\|_p \leq \omega_r(f, 2^{-k} |z|)_p \leq (1 + |z|)^r
\omega_r(f, 2^{-k})_p$ by \Cref{lem:rd-modulus}. \Cref{lem:jackson-weighted-average} with
$M = \omega_r(f, 2^{-k})_p$ gives the bound.
\end{proof}

\section{The inverse inequality}

\begin{lemma}[Differences of Schwartz functions in $L_p$]
\label{lem:schwartz-differences-lp}
\lean{Rd.fwdDiff_schwartz_eq_integral, Rd.enorm_fwdDiff_schwartz_le_lintegral,
  Rd.lintegral_lintegral_Ioc_add_smul, Rd.rpow_lintegral_Ioc_le,
  Rd.enorm_fwdDiff_schwartz_rpow_le, Rd.lintegral_enorm_fwdDiff_rpow_le,
  Rd.lintegral_enorm_fwdDiff_le, Rd.eLpNorm_fwdDiff_schwartz_le}
\leanok
\uses{def:rd-modulus, lem:weighted-jensen}
For every Schwartz function $G$ on $V$ and all $v, x \in V$,
\[
  G(x + v) - G(x) = \int_0^1 \partial_v G(x + t v) \, dt
\]
and
\[
  |G(x + v) - G(x)| \leq \int_0^1 |\partial_v G(x + t v)| \, dt .
\]
For every measurable $\varphi \colon V \to [0, \infty]$ and every $v \in V$,
$\int_V \int_0^1 \varphi(x + t v) \, dt \, dx = \int_V \varphi$. For every real $p \geq 1$ and every
measurable $\varphi \colon (0, 1] \to [0, \infty]$,
$\bigl(\int_0^1 \varphi\bigr)^p \leq \int_0^1 \varphi^p$. For every real $p \geq 1$, every
Schwartz function $G$ on $V$ and all $v, x \in V$,
\[
  |G(x + v) - G(x)|^p \leq \int_0^1 |\partial_v G(x + t v)|^p \, dt
  \quad \text{and} \quad
  \|G(\cdot + v) - G\|_p \leq \|\partial_v G\|_p .
\]
\end{lemma}

\begin{proof}
\leanok
The identity is the fundamental theorem of calculus along the segment from $x$ to $x + v$; the
pointwise bound follows. The integral identity is Tonelli's theorem with the translation
invariance of Lebesgue measure. The bound $(\int_0^1 \varphi)^p \leq \int_0^1 \varphi^p$ is the
weighted Jensen inequality of \Cref{lem:weighted-jensen} with the weight $1$; with the pointwise
bound it gives the bound of
$|G(x + v) - G(x)|^p$. Integrating it over $x$ and averaging the translates of
$|\partial_v G|^p$ along the segment gives $\|G(\cdot + v) - G\|_p \leq \|\partial_v G\|_p$.
\end{proof}

\begin{lemma}[Higher differences of Schwartz functions]
\label{lem:schwartz-differences}
\lean{Rd.fwdDiffSchwartz, Rd.coe_fwdDiffSchwartz, Rd.lineDerivOp_fwdDiffSchwartz,
  Rd.iteratedLineDerivOp_fwdDiffSchwartz, Rd.lintegral_enorm_higherDiff_le,
  Rd.enorm_iteratedLineDerivOp_le, Rd.lintegral_enorm_higherDiff_schwartz_le,
  Rd.lintegral_enorm_iteratedFDeriv_ne_top}
\leanok
\uses{lem:schwartz-differences-lp}
For every Schwartz function $G$ on $V$, $v \in V$ and $r \in \N$,
\[
  \|\Delta_v^r G\|_1 \leq \|v\|^r \int_V \|D^r G(x)\| \, dx < \infty .
\]
\end{lemma}

\begin{proof}
\leanok
The difference $\Delta_v$ and the line derivative $\partial_v$ commute and preserve Schwartz
functions. With $\|\Delta_v G\|_1 \leq \|\partial_v G\|_1$ (\Cref{lem:schwartz-differences-lp}),
induction gives $\|\Delta_v^r G\|_1 \leq \|\partial_v^r G\|_1$, and
$|\partial_v^r G(x)| \leq \|v\|^r \|D^r G(x)\|$.
\end{proof}

\begin{lemma}[The difference operator on $L_2$]
\label{lem:difference-operator-l2}
\lean{Rd.fwdDiffL2, Rd.fwdDiffL2_ae_eq, Rd.enorm_fwdDiffL2, Rd.enorm_fwdDiffL2_toLp_le}
\leanok
\uses{lem:schwartz-differences-lp, lem:rd-differences}
For every $v \in V$, $\Delta_v f = f(\cdot + v) - f$ is a bounded linear operator on $L_2(V)$, and
$\|\Delta_v G\|_2 \leq \|\partial_v G\|_2$ for every Schwartz function $G$ on $V$.
\end{lemma}

\begin{proof}
\leanok
$\Delta_v$ is the translation by $v$, an isometry, minus the identity. The bound is
\Cref{lem:schwartz-differences-lp} with $p = 2$.
\end{proof}

\begin{theorem}[The inverse inequality for band-limited functions]
\label{thm:band-inverse}
\lean{Rd.bandInverseConst, Rd.bandInverseConst_ne_top, Rd.eLpNorm_higherDiff_dilatedBandKernel_le,
  Rd.lpModulus_le_of_mem_bandSpace}
\leanok
\uses{lem:schwartz-differences, lem:rd-differences, lem:rd-modulus, thm:young, def:band-space}
Let $1 \leq p < \infty$, $r \in \N$ and $C = 2^r + \int_V \|D^r \Theta(x)\| \, dx$. For all
$j, k \in \N$ and $a \in E_j$,
\[
  \omega_r(a, 2^{-k})_p \leq C \, 2^{-\max(k - j, 0) r} \, \|a\|_p .
\]
\end{theorem}

\begin{proof}
\leanok
For $k < j$ use $\omega_r(a, t)_p \leq 2^r \|a\|_p$. For $k \geq j$ and $\|h\| \leq 2^{-k}$,
$a = \Theta_j * a$ gives $\Delta_h^r a = (\Delta_h^r \Theta_j) * a$ almost everywhere, and
$\Delta_h^r \Theta_j = (\Delta_{2^j h}^r \Theta)_{2^j}$ (\Cref{lem:rd-differences}). Young's
inequality and \Cref{lem:schwartz-differences} give
$\|\Delta_h^r a\|_p \leq (2^j \|h\|)^r \|D^r \Theta\|_1 \|a\|_p$, and $2^j \|h\| \leq 2^{-(k-j)}$.
\end{proof}

\section{Besov spaces by differences}
\label{sec:rd-besov-differences}

\begin{theorem}[Approximation, Littlewood--Paley blocks and differences]
\label{thm:rd-besov}
\lean{Rd.bandModulusData, Rd.modNorm_bandModulusData, Rd.memApprox_bandScheme_iff_modulus,
  Rd.littlewoodPaleyNorm_lt_top_iff_modulus, Rd.modNorm_le_modNorm_of_lt}
\leanok
\uses{thm:band-jackson, thm:band-inverse, lem:rd-modulus, def:modulus-data,
  thm:modulus-coincidence, thm:lp-band}
Let $1 \leq p < \infty$, $r \geq 1$, $1 \leq q \leq \infty$ and $f \in L_p$.
\begin{enumerate}
  \item For $0 \leq s < r$, $f$ lies in the approximation space $A^s_q$ of the band-limited
    scheme if and only if $\|f\|_p + \|(2^{ks} \omega_r(f, 2^{-k})_p)_k\|_{\ell_q} < \infty$.
  \item For $0 < s < r$, $\|f\|^{\mathrm{LP}}_{s,q} < \infty$ if and only if
    $\|f\|_p + \|(2^{ks} \omega_r(f, 2^{-k})_p)_k\|_{\ell_q} < \infty$.
  \item For $0 \leq s < r'$ and every order $r$ there is $C < \infty$ with
    $\|f\|_p + \|(2^{ks} \omega_{r'}(f, 2^{-k})_p)_k\|_{\ell_q} \leq
    C (\|f\|_p + \|(2^{ks} \omega_r(f, 2^{-k})_p)_k\|_{\ell_q})$.
\end{enumerate}
\end{theorem}

\begin{proof}
\leanok
\Cref{thm:band-jackson} and \Cref{thm:band-inverse} make the modulus compatible with the
band-limited scheme, with inverse order $r$; apply \Cref{thm:modulus-coincidence}. Part (2)
combines (1) with \Cref{thm:lp-band}. Part (3) combines the inverse direction of
\Cref{thm:modulus-coincidence} for the order $r'$ with Jackson's direction for the order $r$.
\end{proof}

For $0 < s < r$ the Besov space $\Besov{s}{p}{q}(V)$ thus has three equivalent descriptions: by
approximation with band-limited functions, by Littlewood--Paley blocks, and by differences.

\section{Sobolev--Slobodeckij norms}
\label{sec:sobolev-slobodeckij}

For $0 < s < 1$ and $1 \leq p < \infty$ the Sobolev--Slobodeckij space $W^{s,p}(V)$ consists of
the $f \in L_p$ with finite Gagliardo seminorm
$(\int\!\int |f(y) - f(x)|^p / |y - x|^{n + sp} \, dy \, dx)^{1/p}$, which Gagliardo introduced
to describe the boundary values of $\Sob{1}{p}$, with $s = 1 - 1/p$ on the
boundary~\cite{Gagliardo1957}. This section compares the $p$-th power of the seminorm with the
first-difference Besov norm of \Cref{thm:rd-besov}, with explicit constants in both directions,
and so identifies $W^{s,p}(V)$ with $\Besov{s}{p}{p}(V)$. Both directions sort the steps $h$
into dyadic annuli: from the form to the modulus, $\|\Delta_h f\|_p^p$ is an average of
$\|\Delta_z f\|_p^p$ over a ball, and Hardy's inequality sums the annuli; from the modulus to
the form, each annulus inside the unit ball contributes one term of the $\ell_p$ norm, and the
steps outside it give a radial integral.

\paragraph{The form.}

\begin{definition}[The Gagliardo form]
\label{def:slobodeckij-form}
\lean{Rd.slobodeckijForm}
\leanok
For every $0 < p < \infty$, every $s \in \R$ and every $f \colon V \to \C$, the Gagliardo
form is
\[
  \mathcal{E}_{s,p}(f) = \int_V \int_V \frac{|f(y) - f(x)|^p}{|y - x|^{n + sp}} \, dy \, dx
  \in [0, \infty] .
\]
\end{definition}

The Lean definition takes the exponent $p$ as a real number, which is $0$ for $p = \infty$; the
results below use the form for $1 \leq p < \infty$.

\begin{lemma}[The Gagliardo form by the step]
\label{lem:slobodeckij-form-step}
\lean{Rd.slobodeckijForm_congr_ae, Rd.slobodeckijForm_of_subsingleton, Rd.fwdDiff_congr_ae,
  Rd.eLpNorm_fwdDiff_rpow_eq, Rd.eLpNorm_fwdDiff_zero_rpow, Rd.measurable_eLpNorm_fwdDiff_rpow,
  Rd.slobodeckijForm_eq_lintegral_of_measurable, Rd.slobodeckijForm_eq_lintegral_eLpNorm}
\leanok
\uses{def:slobodeckij-form, def:rd-modulus, lem:rd-differences}
Let $0 < p < \infty$ and $s \in \R$. For all $f, g \colon V \to \C$ with $f = g$ almost
everywhere, $\Delta_h f = \Delta_h g$ almost everywhere for every $h \in V$ and
$\mathcal{E}_{s,p}(f) = \mathcal{E}_{s,p}(g)$; if $V = \{0\}$, then
$\mathcal{E}_{s,p}(f) = 0$ for every $f$. For every almost everywhere strongly measurable
$f \colon V \to \C$, the function $h \mapsto \|\Delta_h f\|_p^p$ is measurable,
$\|\Delta_0 f\|_p = 0$, and
\[
  \mathcal{E}_{s,p}(f) = \int_V \frac{\|\Delta_h f\|_p^p}{|h|^{n + sp}} \, dh .
\]
\end{lemma}

\begin{proof}
\leanok
Replace $f$ by a measurable representative. The substitution $y = x + h$ in the inner integral
and Tonelli's theorem for the measurable function
$(x, h) \mapsto |f(x + h) - f(x)|^p / |h|^{n + sp}$ exchange the order of integration, and
$\int |f(x + h) - f(x)|^p \, dx = \|\Delta_h f\|_p^p$. The same Tonelli argument gives the
measurability of $h \mapsto \|\Delta_h f\|_p^p$.
\end{proof}

\begin{lemma}[The Gagliardo measure]
\label{lem:gagliardo-measure}
\lean{Rd.gagliardoMeasure, Rd.lintegral_gagliardoMeasure, Rd.ae_gagliardoMeasure,
  Rd.eLpNorm_sub_sq_gagliardoMeasure}
\leanok
\uses{def:slobodeckij-form}
Let $s \in \R$ and let $\mu_s$ be the measure $|y - x|^{-(n + 2s)} \, dx \, dy$ on $V \times V$.
For every measurable $F \colon V \times V \to [0, \infty]$,
$\int F \, d\mu_s = \int_V \int_V F(x, y) \, |y - x|^{-(n + 2s)} \, dy \, dx$. A property that
holds almost everywhere on $V$ holds for both coordinates of $\mu_s$-almost every pair. For every
measurable $f \colon V \to \C$,
\[
  \bigl\|(x, y) \mapsto f(y) - f(x)\bigr\|_{L_2(\mu_s)}^2 = \mathcal{E}_{s,2}(f) .
\]
\end{lemma}

\begin{proof}
\leanok
$\mu_s$ has a measurable density with respect to the product of the Lebesgue measures, and
Tonelli's theorem gives the double integral. Null sets of the product measure are null for
$\mu_s$, and the preimages of null sets of $V$ under the coordinate maps are null for the product
measure. For $F(x, y) = |f(y) - f(x)|^2$ the double integral is the Gagliardo form, since
$a / b^c = a b^{-c}$.
\end{proof}

So $f \mapsto (f, (x, y) \mapsto f(y) - f(x))$ maps $W^{s,2}(V)$ with the norm
$(\|f\|_2^2 + \mathcal{E}_{s,2}(f))^{1/2}$ isometrically into $L_2(V) \times L_2(\mu_s)$ with the
$L_2$ norm on the product, which gives $W^{s,2}(V)$ an inner product.

\paragraph{Dyadic annuli and radial integrals.}

\begin{lemma}[Dyadic annuli of a ball]
\label{lem:dyadic-annuli}
\lean{Rd.dyadicAnnulus, Rd.measurableSet_dyadicAnnulus, Rd.pairwise_disjoint_dyadicAnnulus,
  Rd.closedBall_subset_union_iUnion_dyadicAnnulus,
  Rd.setLIntegral_closedBall_le_tsum_dyadicAnnulus,
  Rd.setLIntegral_closedBall_le_tsum_dyadicAnnulus_of_zero,
  Rd.tsum_setLIntegral_dyadicAnnulus_le, ENNReal.ofReal_div_two_pow_rpow,
  Rd.volume_dyadicAnnulus_le}
\leanok
\uses{lem:dyadic-ball-volume}
For every $\rho > 0$ and $j \in \N$ let
$A_{\rho, j} = \{ z \in V : \rho 2^{-(j+1)} < |z| \leq \rho 2^{-j} \}$. These sets are
measurable and pairwise disjoint, and $|A_{\rho, j}| \leq \rho^n 2^{-jn} |B(0, 1)|$. For every
$k \in \N$ the closed ball $\bar B(0, \rho 2^{-k})$ is contained in
$\{0\} \cup \bigcup_{i \in \N} A_{\rho, k + i}$, and for every $G \colon V \to [0, \infty]$
\[
  \int_{\bar B(0, \rho 2^{-k})} G \leq G(0) \, |\{0\}| + \sum_{i \in \N} \int_{A_{\rho, k+i}} G ,
  \qquad
  \sum_{j \in \N} \int_{A_{\rho, j}} G \leq \int_V G .
\]
\end{lemma}

\begin{proof}
\leanok
For $z \neq 0$ in the ball choose $i$ with $2^i \leq \rho 2^{-k} / |z| < 2^{i+1}$; then
$z \in A_{\rho, k+i}$. For $i < j$, the points of $A_{\rho, j}$ satisfy
$|z| \leq \rho 2^{-j} \leq \rho 2^{-(i+1)}$, so they lie outside $A_{\rho, i}$. The integral
bounds follow from the countable subadditivity of the integral over sets and from its
additivity over disjoint measurable sets, and $A_{\rho, j}$ lies in the closed ball of radius
$\rho 2^{-j}$, whose volume is $\rho^n 2^{-jn} |B(0, 1)|$.
\end{proof}

\begin{lemma}[Radial integrals of powers]
\label{lem:radial-power-integrals}
\lean{ENNReal.ofReal_two_pow_mul_rpow, ENNReal.rpow_le_rpow_add_rpow_of_le_of_le,
  Rd.complBallRpowConst,
  Rd.complBallRpowConst_ne_top, Rd.setLIntegral_compl_ball_inter_dyadicShell_le,
  Rd.lintegral_compl_ball_norm_rpow_neg_le, Rd.ballRpowConst, Rd.ballRpowConst_ne_top,
  Rd.setLIntegral_dyadicAnnulus_norm_rpow_neg_le, Rd.lintegral_ball_norm_rpow_neg_le}
\leanok
\uses{lem:dyadic-annuli, lem:maximal-function-kernels}
For every $R > 0$ and every real $a$:
\begin{enumerate}
  \item if $a > n$, then
    $\int_{|z| \geq R} |z|^{-a} \, dz \leq 2^n |B(0, 1)| \bigl(1 - 2^{-(a - n)}\bigr)^{-1}
    R^{n - a}$;
  \item if $a < n$, then
    $\int_{|z| < R} |z|^{-a} \, dz \leq (1 + 2^a) |B(0, 1)| \bigl(1 - 2^{-(n - a)}\bigr)^{-1}
    R^{n - a}$.
\end{enumerate}
\end{lemma}

\begin{proof}
\leanok
(1) On the part of the shell $2^i R \leq |z| < 2^{i+1} R$ of \Cref{lem:maximal-function-kernels}
outside $B(0, R)$, $|z|^{-a} \leq 2^{-ia} R^{-a}$, and the shell lies in a ball of volume
$2^{(i+1)n} R^n |B(0, 1)|$; sum the geometric series in $2^{-(a-n)}$. (2) On the annulus
$A_{R, i}$ of \Cref{lem:dyadic-annuli},
$|z|^{-a} \leq (R 2^{-(i+1)})^{-a} + (R 2^{-i})^{-a} = (1 + 2^a) R^{-a} 2^{ia}$ for either sign
of $a$, and $|A_{R, i}| \leq R^n 2^{-in} |B(0, 1)|$; sum the geometric series in
$2^{-(n-a)}$. For $n \geq 1$ the point $0$ is a null set; for $n = 0$ the integrand vanishes,
since then $-a > 0$.
\end{proof}

Part (2) holds for exponents $a$ of either sign.

\paragraph{Averages over balls.}

\begin{lemma}[A difference as an average of differences]
\label{lem:difference-average}
\lean{Rd.eLpNorm_fwdDiff_le_add, Rd.eLpNorm_fwdDiff_rpow_le, Rd.setLIntegral_ball_comp_sub_le,
  Rd.eLpNorm_fwdDiff_rpow_le_average}
\leanok
\uses{lem:rd-differences, lem:slobodeckij-form-step}
Let $1 \leq p < \infty$ and let $f \colon V \to \C$ be almost everywhere strongly measurable.
For all $h, z \in V$,
$\|\Delta_h f\|_p \leq \|\Delta_{h - z} f\|_p + \|\Delta_z f\|_p$. For every
$G \colon V \to [0, \infty]$, all $t \in \R$ and $h \in V$ with $|h| \leq t$,
$\int_{|z| < t} G(h - z) \, dz \leq \int_{|z| < 2t} G(z) \, dz$. For all $t > 0$ and $h \in V$
with $|h| \leq t$,
\[
  \|\Delta_h f\|_p^p \leq 2^p \, |B(0, t)|^{-1} \int_{|z| < 2t} \|\Delta_z f\|_p^p \, dz .
\]
\end{lemma}

\begin{proof}
\leanok
$\Delta_h f = (\Delta_{h - z} f)(\cdot + z) + \Delta_z f$, and translations keep $L_p$ norms.
Hence $\|\Delta_h f\|_p^p \leq 2^{p-1} (\|\Delta_{h-z} f\|_p^p + \|\Delta_z f\|_p^p)$ by the
convexity of $x \mapsto x^p$. Integrate over $|z| < t$. The reflection $z \mapsto h - z$
preserves Lebesgue measure and maps $B(0, t)$ into $B(0, 2t)$, since
$|h - z| \leq |h| + |z| < 2t$, so both integrals are at most
$\int_{|z| < 2t} \|\Delta_z f\|_p^p \, dz$, and dividing by $|B(0, t)|$ gives the claim.
\end{proof}

\paragraph{The two bounds.}

\begin{theorem}[The modulus by the Gagliardo form]
\label{thm:modulus-by-slobodeckij-form}
\lean{Rd.lpModulus_rpow_le_average, Rd.inv_volume_ball_one_div_two_pow,
  Rd.dyadicWeight_mul_lpModulus_rpow_le, Rd.dyadicWeight_mul_setLIntegral_dyadicAnnulus_le,
  Rd.tsum_dyadicWeight_mul_setLIntegral_closedBall_le, Rd.slobodeckijConstUpper,
  Rd.slobodeckijConstUpper_ne_top, Rd.lqNorm_modulus_rpow_le_slobodeckijForm}
\leanok
\uses{lem:difference-average, lem:dyadic-annuli, lem:slobodeckij-form-step, lem:hardy,
  lem:rd-modulus, def:lq-norm}
Let $1 \leq p < \infty$, $s > 0$ and $b = n + sp$, and let
$C_1 = 2^p |B(0, 1)|^{-1} (1 - 2^{-b})^{-1} 2^b$. Then $C_1 < \infty$, and for every
$f \in L_p$
\[
  \bigl\| \bigl(2^{ks} \omega_1(f, 2^{-k})_p\bigr)_k \bigr\|_{\ell_p}^p
  \leq C_1 \, \mathcal{E}_{s,p}(f) .
\]
\end{theorem}

\begin{proof}
\leanok
By \Cref{lem:difference-average} with $t = 2^{-k}$ for every step of the supremum, and
$|B(0, 2^{-k})|^{-1} = 2^{kn} |B(0, 1)|^{-1}$,
\[
  \bigl(2^{ks} \omega_1(f, 2^{-k})_p\bigr)^p \leq 2^p |B(0, 1)|^{-1} \, 2^{kb}
  \int_{|z| \leq 2^{1-k}} \|\Delta_z f\|_p^p \, dz .
\]
Write $G(z) = \|\Delta_z f\|_p^p$, so that $G(0) = 0$, and $I_j = \int_{A_{2, j}} G$. By
\Cref{lem:dyadic-annuli} the ball of radius $2^{1-k}$ contributes at most $\sum_{i} I_{k+i}$,
and Hardy's inequality (\Cref{lem:hardy}, $q = 1$) gives
$\sum_k 2^{kb} \sum_i I_{k+i} \leq (1 - 2^{-b})^{-1} \sum_j 2^{jb} I_j$. On $A_{2, j}$,
$|z| \leq 2^{1-j}$, so $2^{jb} \leq 2^b |z|^{-b}$ and
$2^{jb} I_j \leq 2^b \int_{A_{2, j}} G(z) |z|^{-b} \, dz$. The annuli are disjoint, so the sum
is at most $2^b \int_V G(z) |z|^{-b} \, dz = 2^b \mathcal{E}_{s,p}(f)$ by
\Cref{lem:slobodeckij-form-step}.
\end{proof}

\begin{theorem}[The Gagliardo form by the modulus]
\label{thm:slobodeckij-form-by-modulus}
\lean{Rd.setLIntegral_dyadicAnnulus_le_lpModulus, Rd.setLIntegral_compl_ball_le_enorm,
  Rd.slobodeckijConstLower, Rd.slobodeckijConstLower_ne_top, Rd.slobodeckijForm_le}
\leanok
\uses{lem:dyadic-annuli, lem:radial-power-integrals, lem:slobodeckij-form-step,
  lem:rd-modulus, def:lq-norm}
Let $1 \leq p < \infty$, $s > 0$ and $b = n + sp$, and let
$C_2 = 2^b |B(0, 1)| + 2^p \cdot 2^n |B(0, 1)| (1 - 2^{-sp})^{-1}$. Then $C_2 < \infty$, and for
every $f \in L_p$
\[
  \mathcal{E}_{s,p}(f) \leq C_2 \Bigl( \|f\|_p^p
  + \bigl\| \bigl(2^{ks} \omega_1(f, 2^{-k})_p\bigr)_k \bigr\|_{\ell_p}^p \Bigr) .
\]
\end{theorem}

\begin{proof}
\leanok
Split the integral of \Cref{lem:slobodeckij-form-step} into $|z| \leq 1$ and $|z| \geq 1$. Since
the integrand vanishes at $z = 0$, \Cref{lem:dyadic-annuli} with $\rho = 1$ bounds the first
part by the sum over the annuli $2^{-(i+1)} < |z| \leq 2^{-i}$. There
$\|\Delta_z f\|_p \leq \omega_1(f, 2^{-i})_p$ and $|z|^{-b} \leq 2^{(i+1) b}$, and the annulus
has volume at most $2^{-in} |B(0, 1)|$, so it contributes at most
$2^b |B(0, 1)| \, 2^{isp} \omega_1(f, 2^{-i})_p^p$. For $|z| \geq 1$,
$\|\Delta_z f\|_p \leq 2 \|f\|_p$, and part (1) of \Cref{lem:radial-power-integrals} with
$a = b > n$ and $R = 1$ bounds the second part by
$2^p \|f\|_p^p \cdot 2^n |B(0, 1)| (1 - 2^{-sp})^{-1}$.
\end{proof}

\begin{theorem}[Sobolev--Slobodeckij spaces are Besov spaces]
\label{thm:sobolev-slobodeckij-besov}
\lean{Rd.slobodeckijForm_lt_top_iff_memApprox}
\leanok
\uses{thm:modulus-by-slobodeckij-form, thm:slobodeckij-form-by-modulus, thm:rd-besov}
For every $1 \leq p < \infty$, every $0 < s < 1$ and every $f \in L_p$,
$\mathcal{E}_{s,p}(f) < \infty$ if and only if $f$ lies in the approximation space $A^s_p$ of
the band-limited scheme.
\end{theorem}

\begin{proof}
\leanok
By \Cref{thm:rd-besov} with $r = 1$ and $q = p$, $f \in A^s_p$ if and only if
$\|f\|_p + \|(2^{ks} \omega_1(f, 2^{-k})_p)_k\|_{\ell_p} < \infty$. The two directions are
\Cref{thm:modulus-by-slobodeckij-form} and \Cref{thm:slobodeckij-form-by-modulus}, whose
constants are finite.
\end{proof}

Thus $W^{s,p}(V) = \Besov{s}{p}{p}(V)$ for $0 < s < 1$ and $1 \leq p < \infty$, with
$\|f\|_p^p + \mathcal{E}_{s,p}(f)$ and
$\|f\|_p^p + \|(2^{ks} \omega_1(f, 2^{-k})_p)_k\|_{\ell_p}^p$ equivalent: each is at most
$\max(1, C_1)$, respectively $1 + C_2$, times the other.

\section{Local means}
\label{sec:rd-local-means}

\begin{definition}[Local means]
\label{def:local-means}
\lean{Rd.LocalMeansBump, Rd.LocalMeansBump.hasCompactSupport, Rd.LocalMeansBump.integrable,
  Rd.LocalMeansBump.contDiff_momentKernel, Rd.LocalMeansBump.momentKernel_eq_zero,
  Rd.LocalMeansBump.hasCompactSupport_momentKernel, Rd.LocalMeansBump.momentSchwartz,
  Rd.LocalMeansBump.derivNorm, Rd.LocalMeansBump.scaledMoment, Rd.LocalMeansBump.meansKernel,
  Rd.LocalMeansBump.dualKernel, Rd.LocalMeansBump.integrable_scaledMoment,
  Rd.LocalMeansBump.integrable_meansKernel, Rd.LocalMeansBump.integrable_dualKernel,
  Rd.LocalMeansBump.localMean, Rd.LocalMeansBump.localMeansNorm}
\leanok
\uses{def:moment-kernel, lem:conv-op}
Let $g : V \to \C$ be smooth with $\int_V g = 1$ and $g = 0$ outside the ball of radius $R$, let
$L \geq 1$ and let $m$ be the moment kernel of order $L$. The local means kernels and the dual
kernels are
\[
  k_0 = \eta_0 = m, \qquad k_j = m_{2^j} - m_{2^{j-1}}, \qquad \eta_j = m_{2^j} + m_{2^{j-1}}
  \qquad (j \geq 1) .
\]
For $1 \leq p < \infty$, $s \in \R$, $1 \leq q \leq \infty$ and $f \in L_p$, the local means
quasi-norm is $\|(2^{js} \|k_j * f\|_p)_j\|_{\ell_q}$.
\end{definition}

\begin{lemma}[The reproducing formula of local means]
\label{lem:reproducing-formula}
\lean{Rd.LocalMeansBump.sum_dual_localMean, Rd.LocalMeansBump.tendsto_sum_dual_localMean,
  Rd.tendsto_apply_apply_of_tendsto}
\leanok
\uses{def:local-means, lem:conv-op, lem:conv-assoc, thm:approx-identity}
With the notation of \Cref{def:local-means}, for every $J \in \N$ and $f \in L_p$,
\[
  \sum_{j=0}^J \eta_j * k_j * f = m_{2^J} * m_{2^J} * f ,
\]
and $\sum_{j=0}^J \eta_j * k_j * f \to f$ in $L_p$ as $J \to \infty$.
\end{lemma}

\begin{proof}
\leanok
The convolution operators commute (\Cref{lem:conv-op}), so
$T_{\eta_j} T_{k_j} = T_{m_{2^j}}^2 - T_{m_{2^{j-1}}}^2$ for $j \geq 1$, from
$a^2 - b^2 = (a + b)(a - b)$, and the sum telescopes. Since $m_{2^J} * f \to f$
(\Cref{thm:approx-identity}) and $\|m_{2^J} * h\|_p \leq \|m\|_1 \|h\|_p$, also
$m_{2^J} * m_{2^J} * f \to f$: for uniformly bounded additive operators $T_J$ with $T_J f \to f$,
$\|T_J T_J f - f\| \leq \|T_J (T_J f - f)\| + \|T_J f - f\| \leq (C + 1) \|T_J f - f\|$.
\end{proof}

\begin{lemma}[Local means are bounded by the modulus]
\label{lem:local-means-upper}
\lean{Rd.LocalMeansBump.enorm_localMean_zero_le, Rd.LocalMeansBump.enorm_localMean_succ_le,
  Rd.LocalMeansBump.upperConst, Rd.LocalMeansBump.upperConst_ne_top,
  Rd.LocalMeansBump.localMeansNorm_le}
\leanok
\uses{def:local-means, thm:moment-approx, lem:rd-modulus}
With the notation of \Cref{def:local-means}, for $f \in L_p$ and $j \in \N$,
$\|k_0 * f\|_p \leq \|m\|_1 \|f\|_p$ and
\[
  \|k_{j+1} * f\|_p \leq 2 \|g\|_1 (1 + R)^L \, \omega_L(f, 2^{-j})_p .
\]
Hence for all $s \in \R$ and $1 \leq q \leq \infty$ there is $C < \infty$ with
\[
  \|(2^{js} \|k_j * f\|_p)_j\|_{\ell_q}
  \leq C \bigl( \|f\|_p + \|(2^{js} \omega_L(f, 2^{-j})_p)_j\|_{\ell_q} \bigr)
\]
for every $f \in L_p$.
\end{lemma}

\begin{proof}
\leanok
$k_{j+1} * f = (f - m_{2^j} * f) - (f - m_{2^{j+1}} * f)$, and \Cref{thm:moment-approx} bounds
both terms by $\|g\|_1 \omega_L(f, R 2^{-j})_p \leq \|g\|_1 (1 + R)^L \omega_L(f, 2^{-j})_p$
(\Cref{lem:rd-modulus}).
\end{proof}

\begin{lemma}[Moduli of the reproducing terms]
\label{lem:dual-kernel-moduli}
\lean{Rd.LocalMeansBump.eLpNorm_scaledMoment_one, Rd.LocalMeansBump.eLpNorm_dualKernel_one_le,
  Rd.LocalMeansBump.eLpNorm_meansKernel_one_le,
  Rd.LocalMeansBump.eLpNorm_higherDiff_scaledMoment_le,
  Rd.LocalMeansBump.eLpNorm_higherDiff_scaledMoment_le', Rd.LocalMeansBump.dualConst,
  Rd.LocalMeansBump.dualConst_ne_top, Rd.LocalMeansBump.eLpNorm_higherDiff_scaledMoment_le_dyadic,
  Rd.LocalMeansBump.eLpNorm_higherDiff_dualKernel_le, Rd.LocalMeansBump.lpModulus_dual_le}
\leanok
\uses{def:local-means, lem:schwartz-differences, lem:rd-differences, lem:rd-modulus-conv}
With the notation of \Cref{def:local-means}, $\|\eta_j\|_1 \leq 2 \|m\|_1$ and
$\|k_j\|_1 \leq 2 \|m\|_1$, and there is
$C < \infty$ with
\[
  \omega_L(\eta_j * h, 2^{-i})_p \leq C \, \min\bigl(1, 2^{-(i - j) L}\bigr) \, \|h\|_p
\]
for all $i, j \in \N$ and $h \in L_p$.
\end{lemma}

\begin{proof}
\leanok
$\eta_j$ and $k_j$ are sums or differences of at most two dilations $m_{2^k}$, each with
$\|m_{2^k}\|_1 = \|m\|_1$. The moment kernel $m$ is smooth with compact support, so
$\|\Delta_v^L m_{2^j}\|_1 \leq (2^j \|v\|)^L \|D^L m\|_1$ by \Cref{lem:rd-differences} and
\Cref{lem:schwartz-differences}, and $\|\Delta_v^L m_{2^j}\|_1 \leq 2^L \|m\|_1$. For
$\|v\| \leq 2^{-i}$ this gives $\|\Delta_v^L \eta_j\|_1 \leq C \min(1, 2^{-(i-j)L})$, and
\Cref{lem:rd-modulus-conv} concludes.
\end{proof}

\begin{lemma}[An almost diagonal estimate]
\label{lem:almost-diagonal}
\lean{ENNReal.lqNorm_almostDiagonal_le, ENNReal.almostDiagonalConst,
  ENNReal.almostDiagonalConst_ne_top, ENNReal.lqNorm_le_almostDiagonalConst_mul_of_le}
\leanok
\uses{lem:hardy}
For all $s, L \in \R$ let $A_{s,L} = (1 - 2^{-(L-s)})^{-1} + (1 - 2^{-s})^{-1}$ in $[0, \infty]$;
then $A_{s,L} < \infty$ for $0 < s < L$. Let $1 \leq q \leq \infty$, $s, L \in \R$, and let
$c_{ij} = 2^{-(i-j)L}$ for $j \leq i$ and $c_{ij} = 1$ for $j > i$. For every sequence $a$ in
$[0, \infty]$,
\[
  \Bigl\| \Bigl( 2^{is} \sum_j c_{ij} \, a_j \Bigr)_i \Bigr\|_{\ell_q}
  \leq A_{s,L} \, \|(2^{js} a_j)_j\|_{\ell_q} .
\]
For every $C \in [0, \infty]$ and all sequences $a, m$ in $[0, \infty]$ with
$a_j \leq C \sum_i c_{ji} m_i$ for every $j$,
$\|(2^{js} a_j)_j\|_{\ell_q} \leq C A_{s,L} \, \|(2^{is} m_i)_i\|_{\ell_q}$.
\end{lemma}

\begin{proof}
\leanok
For $0 < s < L$ both $2^{-(L-s)}$ and $2^{-s}$ are less than $1$, so $A_{s,L} < \infty$. For
$s \leq 0$ or $s \geq L$ the constant is $\infty$, and the right side vanishes only for
$a = 0$, where both sides are $0$. For $0 < s < L$, split the sum at $j = i$: the part
$j \leq i$ is a partial sum with the weights $2^{-i(L-s)}$, the part $j > i$ a tail sum with the
weights $2^{is}$, and \Cref{lem:hardy} bounds both. For the last claim, multiply the hypothesis
by $2^{js}$, take the $\ell_q$ norm and apply the displayed bound to $(m_i)_i$.
\end{proof}

\begin{lemma}[The synthesis estimate]
\label{lem:local-means-synthesis}
\lean{ENNReal.le_tsum_of_tendsto_sum, ENNReal.enorm_le_of_tendsto_sum,
  Rd.LocalMeansBump.bumpSynthConst, Rd.LocalMeansBump.bumpSynthConst_ne_top,
  Rd.LocalMeansBump.modNorm_le_of_tendsto, Rd.LocalMeansBump.modNorm_le_localMeansNorm}
\leanok
\uses{lem:dual-kernel-moduli, lem:almost-diagonal, lem:reproducing-formula, lem:rd-modulus}
Let $E$ be a normed group and let the partial sums of $\sum_j u_j$ converge to $F$ in $E$. For
every $C < \infty$ and every subadditive $Q \colon E \to [0, \infty]$ with $Q \leq C \|\cdot\|$,
$Q(F) \leq \sum_j Q(u_j)$. For every real $s$, every $1 \leq q \leq \infty$, every
$K \in [0, \infty]$ and all $b_j \in [0, \infty]$ with $\|u_j\| \leq K b_j$ for every $j$,
$\|F\| \leq K (1 - 2^{-s})^{-1} \|(2^{js} b_j)_j\|_{\ell_q}$.

With the notation of \Cref{def:local-means}, let $0 < s < L$ and $1 \leq q \leq \infty$. There is
$C < \infty$ such that for all $h_j \in L_p$ and $F \in L_p$ with
$\sum_{j=0}^J \eta_j * h_j \to F$ in $L_p$,
\[
  \|F\|_p + \|(2^{is} \omega_L(F, 2^{-i})_p)_i\|_{\ell_q}
  \leq C \, \|(2^{js} \|h_j\|_p)_j\|_{\ell_q} .
\]
In particular $\|f\|_p + \|(2^{is} \omega_L(f, 2^{-i})_p)_i\|_{\ell_q}
\leq C \|(2^{js} \|k_j * f\|_p)_j\|_{\ell_q}$ for every $f \in L_p$.
\end{lemma}

\begin{proof}
\leanok
The functional $Q$ passes to the limit of the partial sums, which it bounds by
$\sum_{j \leq J} Q(u_j) + C \|F - \sum_{j \leq J} u_j\|$. For the second claim take $Q = \|\cdot\|$
and $u_j \leq 2^{-js} \|(2^{js} b_j)_j\|_{\ell_q}$ termwise, with
$\sum_j 2^{-js} = (1 - 2^{-s})^{-1}$.

The norm and the modulus $\omega_L(\cdot, 2^{-i})_p$ are subadditive and bounded by multiples of
$\|\cdot\|_p$, so they pass to the series:
$\omega_L(F, 2^{-i})_p \leq \sum_j \omega_L(\eta_j * h_j, 2^{-i})_p \leq C \sum_j c_{ij}
\|h_j\|_p$ by \Cref{lem:dual-kernel-moduli}, and \Cref{lem:almost-diagonal} bounds the
$\ell_q$ norm. Likewise $\|F\|_p \leq 2 \|m\|_1 \sum_j \|h_j\|_p \leq 2 \|m\|_1 (1 - 2^{-s})^{-1}
\|(2^{js} \|h_j\|_p)_j\|_{\ell_q}$. The case $h_j = k_j * f$ is \Cref{lem:reproducing-formula}.
\end{proof}

\begin{theorem}[Local means and differences]
\label{thm:local-means}
\lean{Rd.LocalMeansBump.besov_iff_localMeans,
  Rd.LocalMeansBump.memApprox_bandScheme_iff_localMeans}
\leanok
\uses{lem:local-means-upper, lem:local-means-synthesis, thm:rd-besov}
With the notation of \Cref{def:local-means}, let $0 < s < L$, $1 \leq q \leq \infty$ and
$f \in L_p$. Then
\[
  \|f\|_p + \|(2^{ks} \omega_L(f, 2^{-k})_p)_k\|_{\ell_q} < \infty
  \iff \|(2^{js} \|k_j * f\|_p)_j\|_{\ell_q} < \infty ,
\]
and both hold if and only if $f$ lies in the approximation space $A^s_q$ of the band-limited
scheme.
\end{theorem}

\begin{proof}
\leanok
\Cref{lem:local-means-upper} and \Cref{lem:local-means-synthesis}; the last statement is
\Cref{thm:rd-besov}.
\end{proof}

Every estimate runs through the averaged differences of \Cref{lem:averaged-differences}, applied
to $f$ or to the smooth kernels. With \Cref{thm:rd-besov}, local means are a fourth description of
$\Besov{s}{p}{q}(V)$ for $0 < s < L$.

\begin{definition}[Bump mollifiers]
\label{def:bump-mollifier}
\lean{Rd.bumpMollifier, Rd.contDiff_bumpMollifier, Rd.tsupport_bumpMollifier_subset,
  Rd.hasCompactSupport_bumpMollifier, Rd.integral_bumpMollifier}
\leanok
For every smooth bump function $f$ around $c \in V$ with outer radius $r$ (Mathlib's
\texttt{ContDiffBump}), the bump mollifier of $f$ is $f / \int f$ with complex values. It is
smooth with support in $\bar B(c, r)$, compactly supported, and has integral $1$.
\end{definition}

\begin{lemma}[Bumps from smooth bump functions]
\label{lem:bump-of-smooth-bump}
\lean{Rd.LocalMeansBump.ofContDiffBump, Rd.centeredBumpBase, Rd.LocalMeansBump.centered,
  Rd.LocalMeansBump.centered_radius}
\leanok
\uses{def:bump-mollifier, def:local-means}
For every smooth bump function around $z_0 \in V$ with outer radius $r$, its bump mollifier is a
bump for local means (\Cref{def:local-means}) of every radius $R \geq \|z_0\| + r$. For every
$R > 0$ the bump mollifier of the smooth bump function around $0$ with outer radius $R$ is a bump
for local means of radius $R$.
\end{lemma}

\begin{proof}
\leanok
The bump mollifier vanishes outside $\bar B(z_0, r) \subseteq \bar B(0, \|z_0\| + r)$, since
$\|z\| \leq \|z - z_0\| + \|z_0\|$ (\Cref{def:bump-mollifier}).
\end{proof}

\section{Pointwise multipliers}

\begin{lemma}[Differences along a segment]
\label{lem:differences-segment}
\lean{Rd.norm_iteratedFDeriv_fwdDiff_le_of_segment, Rd.norm_higherDiff_le_of_segment}
\leanok
\uses{def:rd-modulus}
Let $E$ be a real normed space, $l \in \N$, $G : V \to E$ of class $C^l$, $h, x \in V$ and
$C \in \R$ with $\|D^l G(x + t h)\| \leq C$ for every $t \in [0, l]$. Then
$\|\Delta_h^l G(x)\| \leq C \|h\|^l$.
\end{lemma}

\begin{proof}
\leanok
Induction on $l$: the derivatives of order $l$ of $\Delta_h G$ at $x + t h$, $0 \leq t \leq l$,
are the differences of those of $G$ at $x + (t + 1) h$ and $x + t h$, and the mean value
inequality on the segment between these points bounds them by $C \|h\|$.
\end{proof}

The bound needs the derivatives only on the segment from $x$ to $x + l h$, which matters for
functions that decay.

\begin{lemma}[Differences of smooth functions]
\label{lem:smooth-differences}
\lean{Rd.norm_higherDiff_le_of_iteratedFDeriv, Rd.norm_le_of_iteratedFDeriv_le}
\leanok
\uses{lem:differences-segment}
Let $E$ be a normed space, $l \in \N$, $C \geq 0$ and $G : V \to E$ smooth with
$\|D^l G(x)\| \leq C$ for all $x$. Then $\|\Delta_h^l G(x)\| \leq C \|h\|^l$ for all $h, x \in V$.
\end{lemma}

\begin{proof}
\leanok
This is \Cref{lem:differences-segment} with the bound on all of $V$.
\end{proof}

\begin{lemma}[Leibniz's rule for differences]
\label{lem:leibniz-differences}
\lean{Rd.sum_choose_mul_succ_add, Rd.higherDiff_mul}
\leanok
\uses{def:rd-modulus}
For $u, v : V \to \C$, $h, x \in V$ and $r \in \N$,
\[
  \Delta_h^r (u v)(x) = \sum_{l=0}^r \binom{r}{l} \, \Delta_h^l u(x) \, \Delta_h^{r-l} v(x + l h) .
\]
\end{lemma}

\begin{proof}
\leanok
Induction on $r$, from $\Delta_h(u v)(x) = \Delta_h u(x) \, v(x + h) + u(x) \, \Delta_h v(x)$ and
Pascal's rule.
\end{proof}

\begin{lemma}[The modulus of a product]
\label{lem:product-modulus}
\lean{Rd.eLpNorm_higherDiff_mul_le, Rd.smoothnessModulus_mul_smooth_le,
  ENNReal.ofReal_one_div_two_pow_pow, Rd.memLp_mul_of_bound, Rd.mulLp, Rd.mulLp_apply_ae,
  Rd.mulLp_sub, Rd.mulLp_add, Rd.mulLp_smul, Rd.enorm_mulLp_le, Rd.mulLpL, Rd.mulLpL_apply}
\leanok
\uses{lem:smooth-differences, lem:leibniz-differences, lem:rd-differences}
Let $1 \leq p < \infty$, $r \in \N$, $C \geq 0$ and $\psi : V \to \C$ smooth with
$\|D^l \psi\| \leq C$ for $l \leq r$. For $F \in L_p$, $\psi F \in L_p$ with
$\|\psi F\|_p \leq C \|F\|_p$, the map $F \mapsto \psi F$ is complex linear, and for $t \geq 0$,
\[
  \omega_r(\psi F, t)_p \leq C \sum_{l=0}^r \binom{r}{l} \, t^l \, \omega_{r-l}(F, t)_p .
\]
\end{lemma}

\begin{proof}
\leanok
By \Cref{lem:leibniz-differences} and \Cref{lem:smooth-differences},
\[
  |\Delta_h^r(\psi F)(x)| \leq C \sum_l \binom{r}{l} \|h\|^l \, |\Delta_h^{r-l} F(x + l h)| ,
\]
and translations preserve $L_p$ norms.
\end{proof}

\begin{lemma}[Moduli of lower order]
\label{lem:lower-order-moduli}
\lean{Rd.lqNorm_lpModulus_sub_le, Rd.multiplierConst, Rd.multiplierConst_ne_top,
  Rd.lqNorm_lpModulus_le_multiplierConst}
\leanok
\uses{thm:rd-besov, def:modulus-data, thm:modulus-coincidence}
Let $1 \leq p < \infty$, $s < r$, $1 \leq q \leq \infty$ and $0 \leq l \leq r$. There is
$K_l < \infty$ with
\[
  \bigl\| \bigl( 2^{k(s-l)} \omega_{r-l}(F, 2^{-k})_p \bigr)_k \bigr\|_{\ell_q}
  \leq K_l \bigl( \|F\|_p + \|(2^{ks} \omega_r(F, 2^{-k})_p)_k\|_{\ell_q} \bigr)
\]
for every $F \in L_p$.
\end{lemma}

\begin{proof}
\leanok
For $1 \leq l \leq s$, both sides are equivalent to the norms of the band-limited approximation
spaces $A^{s-l}_q$ and $A^s_q$ (\Cref{thm:rd-besov} with the orders $r - l$ and $r$), and
$A^s_q \subseteq A^{s-l}_q$. For $l > s$, $\omega_{r-l}(F, t)_p \leq 2^{r-l} \|F\|_p$ and
$\sum_k 2^{k(s-l)q} < \infty$.
\end{proof}

\begin{theorem}[Pointwise multipliers]
\label{thm:multipliers}
\lean{Rd.multiplierTotal, Rd.modNorm_mulLp_le, Rd.multiplierTotal_ne_top,
  Rd.modNorm_mulLp_lt_top}
\leanok
\uses{lem:product-modulus, lem:lower-order-moduli}
Let $1 \leq p < \infty$, $s < r$, $1 \leq q \leq \infty$ and $\psi : V \to \C$ smooth with bounded
derivatives up to order $r$. There is $K < \infty$ with
\[
  \|\psi F\|_p + \|(2^{ks} \omega_r(\psi F, 2^{-k})_p)_k\|_{\ell_q}
  \leq K \bigl( \|F\|_p + \|(2^{ks} \omega_r(F, 2^{-k})_p)_k\|_{\ell_q} \bigr)
\]
for every $F \in L_p$.
\end{theorem}

\begin{proof}
\leanok
By \Cref{lem:product-modulus}, $2^{ks} \omega_r(\psi F, 2^{-k})_p \leq C \sum_l \binom{r}{l} \,
2^{k(s-l)} \omega_{r-l}(F, 2^{-k})_p$; Minkowski's inequality and \Cref{lem:lower-order-moduli}
bound the $\ell_q$ norm.
\end{proof}

\begin{lemma}[Smooth cutoffs are Besov multipliers]
\label{lem:besov-multiplier}
\lean{Rd.exists_mul_modNorm_le}
\leanok
\uses{thm:multipliers, lem:resolution-of-unity}
For every smooth $\chi \colon V \to \C$ with compact support, every $1 \leq p < \infty$, all
$r \in \N$ and $s \in \R$ with $s < r$, and every $q \in [1, \infty]$, there is $C < \infty$ such
that for every $F \in L_p(V)$,
\[
  \|\chi F\|_p + \|(2^{ks} \omega_r(\chi F, 2^{-k})_p)_k\|_{\ell_q}
  \leq C \bigl( \|F\|_p + \|(2^{ks} \omega_r(F, 2^{-k})_p)_k\|_{\ell_q} \bigr) .
\]
\end{lemma}

\begin{proof}
\leanok
The derivatives of $\chi$ up to order $r$ are bounded (\Cref{lem:resolution-of-unity}), so
\Cref{thm:multipliers} applies.
\end{proof}

\section{Atoms}

Atoms are functions of small support with size bounds at a dyadic scale. As in Schneider and
Vyb\'{\i}ral~\cite{SchneiderVybiral}, they need no moment conditions and no smoothness beyond the
order of the differences that measure them.

\begin{definition}[Dyadic cells]
\label{def:dyadic-cells}
\lean{Rd.cellIndex, Rd.dyadicCell, Rd.cellCenter, Rd.zcellIndex_natCast, Rd.zcell_natCast,
  Rd.zcellCenter_natCast, Rd.mem_dyadicCell_cellIndex, Rd.repr_cellCenter, Rd.measurable_cellIndex,
  Rd.measurableSet_dyadicCell}
\leanok
\uses{def:dyadic-cubes}
For $j \in \N$ and $m \in \Z^n$ the dyadic cell of level $j$ and index $m$ is the dyadic cube
$P_{j,m} = Q_{j,m}$ of \Cref{def:dyadic-cubes}, that is,
$P_{j,m} = \{x \in V : \lfloor 2^j \langle b_i, x \rangle \rfloor = m_i \text{ for all } i\}$,
with centre $c_{j,m} = \sum_i (m_i + 1/2) 2^{-j} b_i$. Every $x$ lies in the cell of index
$(\lfloor 2^j \langle b_i, x \rangle \rfloor)_i$, and the cells are measurable.
\end{definition}

\begin{lemma}[Cells and centres]
\label{lem:dyadic-cells}
\lean{Rd.abs_repr_sub_cellCenter_le, Rd.norm_sub_cellCenter_le, Rd.volume_dyadicCell,
  Rd.card_le_of_norm_cellCenter_sub_le}
\leanok
\uses{def:dyadic-cells, lem:dyadic-cubes}
For every $j \in \N$, every $m \in \Z^n$ and every $x \in P_{j,m}$, every coordinate of
$x - c_{j,m}$ is at most $2^{-j}/2$ in modulus, $\|x - c_{j,m}\| \leq n 2^{-j}$ and
$|P_{j,m}| = 2^{-jn}$. For every $j \in \N$ and every $\rho \geq 0$ at most $(2\rho + 2)^n$
centres of level $j$ lie in a closed ball of radius $\rho 2^{-j}$.
\end{lemma}

\begin{proof}
\leanok
The bounds and the volume are the case $k = j$ of \Cref{lem:dyadic-cubes}. The coordinates of
the centres in such a ball lie in intervals of length $2 \rho 2^{-j}$, and such
an interval contains at most $2\rho + 2$ points of $(\Z + 1/2) 2^{-j}$.
\end{proof}

\begin{lemma}[Bounded overlap]
\label{lem:bounded-overlap}
\lean{Rd.exists_finset_of_card_le, Rd.enorm_tsum_rpow_le_of_overlap, Rd.eLpNorm_rpow_le_of_overlap}
\leanok
Let $1 \leq p < \infty$ and let $g_m$, $m$ in a countable set, be measurable functions such that at
every point at most $N$ of them are nonzero. If $F = \sum_m g_m$ almost everywhere, then
$\|F\|_p^p \leq N^{p-1} \sum_m \|g_m\|_p^p$.
\end{lemma}

\begin{proof}
\leanok
At every point the sum has at most $N$ nonzero terms, and
$|\sum_{m \in T} a_m|^p \leq |T|^{p-1} \sum_{m \in T} |a_m|^p$ by H\"older's inequality. Integrate
and exchange the sum with the integral.
\end{proof}

\begin{definition}[Regular atoms and the coefficient quasi-norm]
\label{def:regular-atoms}
\lean{Rd.IsRegularAtomOn, Rd.IsRegularAtom, Rd.levelNorm, Rd.levelNorm_eq_tsum, Rd.levelNorm_mono,
  Rd.coeffNorm}
\leanok
\uses{def:rd-modulus-on}
Let $1 \leq p < \infty$, $\sigma \geq 0$, $r \in \N$, $d, C \geq 0$ and $\Omega \subseteq V$. A
regular atom on $\Omega$ of level $j$, exponent $\sigma$ and order $r$ with centre $c$ is a
measurable $a : V \to \C$ with $a = 0$ outside $\overline{B}(c, d 2^{-j})$,
\begin{gather*}
  \|a\|_{L_p(\Omega)} \leq C 2^{-jn/p}, \
  \|\Delta_h^r a\|_{L_p(\Omega_{rh})} \leq C 2^{-jn/p} \min(1, (2^j \|h\|)^\sigma) \quad (h \in V).
\end{gather*}
A regular atom is a regular atom on $V$.
For coefficients $\lambda = (\lambda_{j,m})$ let
$\|\lambda_{j,\cdot}\|_{\ell_p} = (\sum_m |\lambda_{j,m}|^p)^{1/p}$ and
$\|\lambda\|_{b^s_{p,q}} = \|(2^{j(s - n/p)} \|\lambda_{j,\cdot}\|_{\ell_p})_j\|_{\ell_q}$.
\end{definition}

The Lip atoms of \cite[Definition 2.1]{SchneiderVybiral}, with $|a| \leq 1$ and
$|a(x) - a(y)| \leq 2^j |x - y|$, are regular atoms with $\sigma = r = 1$, and their smooth atoms of
\cite[Definition 1.3]{SchneiderVybiral} are regular atoms with $\sigma = r = K$
(\Cref{lem:pointwise-atoms}).

\begin{lemma}[Level sums]
\label{lem:level-sums}
\lean{Rd.higherDiff_tsum_apply, Rd.le_of_rpow_le_overlap, Rd.one_sub_one_div_toReal_nonneg,
  Rd.summable_levelSum, Rd.eLpNorm_levelSum_le, Rd.card_le_of_higherDiff_ne_zero,
  Rd.smoothnessModulus_levelSum_le, Rd.eLpNorm_indicator_levelSum_le, Rd.modulusOn_levelSum_le}
\leanok
\uses{lem:bounded-overlap, def:regular-atoms, lem:rd-modulus, lem:rd-modulus-on}
Let $1 \leq p < \infty$, and let $a_m$ be regular atoms of level $j$ with centres $c_m$ and
coefficients $\lambda_m$, such that at most $N$ of the centres $c_m$ with $\lambda_m \neq 0$ lie in a
ball of radius $(d + r) 2^{-j}$. Then $f = \sum_m \lambda_m a_m$ converges at every point, and
\[
  \|f\|_p \leq N^{1-1/p} C 2^{-jn/p} \|\lambda\|_{\ell_p}, \qquad
  \omega_r(f, 2^{-k})_p \leq 2^r N^{1-1/p} C 2^{-jn/p} \|\lambda\|_{\ell_p} c_{kj},
\]
with $c_{kj} = 2^{-(k-j)\sigma}$ for $j \leq k$ and $c_{kj} = 1$ for $j > k$. For regular atoms
on a measurable set $\Omega$ the same bounds hold for $\|f\|_{L_p(\Omega)}$ and
$\omega_r(f, 2^{-k}, \Omega)_p$.
\end{lemma}

\begin{proof}
\leanok
At every point at most $N$ of the terms are nonzero, and \Cref{lem:bounded-overlap} applies. For
$j \leq k$ and $\|h\| \leq 2^{-k}$ the differences $\Delta_h^r a_m$ vanish outside
$B(c_m, (d + r) 2^{-j})$, and \Cref{lem:bounded-overlap} applies to
$\Delta_h^r f = \sum_m \lambda_m \Delta_h^r a_m$ with $(2^j \|h\|)^\sigma \leq 2^{-(k-j)\sigma}$; for
$j > k$, $\omega_r(f, t)_p \leq 2^r \|f\|_p$.
\end{proof}

\begin{theorem}[Atomic synthesis on $\R^d$]
\label{thm:rd-atomic-synthesis}
\lean{Rd.atomSynthConst, Rd.atomSynthConst_ne_top, Rd.modNorm_le_coeffNorm,
  Rd.eLpNorm_indicator_le_tsum, Rd.lpModulusOn_le_tsum, Rd.modNormOn_le_coeffNorm}
\leanok
\uses{lem:level-sums, lem:almost-diagonal, lem:rd-modulus-on, lem:local-means-synthesis}
Let $1 \leq p < \infty$, $1 \leq q \leq \infty$ and $0 < s < \sigma$. Let $a_{j,m}$ be regular atoms
of level $j$ such that the centres of the atoms with $\lambda_{j,m} \neq 0$ satisfy the overlap
condition of \Cref{lem:level-sums} at every level, and let $F = \sum_j f_j$ in $L_p$ with
$f_j = \sum_m \lambda_{j,m} a_{j,m}$ almost everywhere. Then
\[
  \|F\|_p + \|(2^{ks} \omega_r(F, 2^{-k})_p)_k\|_{\ell_q} \leq K \|\lambda\|_{b^s_{p,q}},
\]
with $K = N^{1-1/p} C \bigl((1 - 2^{-s})^{-1} + 2^r ((1 - 2^{-(\sigma - s)})^{-1} +
(1 - 2^{-s})^{-1})\bigr)$. For regular atoms on a measurable set $\Omega$,
\[
  \|F\|_{L_p(\Omega)} + \|(2^{ks} \omega_r(F, 2^{-k}, \Omega)_p)_k\|_{\ell_q} \leq
  K \|\lambda\|_{b^s_{p,q}} .
\]
\end{theorem}

\begin{proof}
\leanok
This is \cite[Theorem 2.6 and Corollary 2.8, direction $\Leftarrow$]{SchneiderVybiral} for
$p \geq 1$. The norm and the moduli are subadditive and bounded on $L_p$, so they are bounded by the
sums over the levels (\Cref{lem:local-means-synthesis}). With \Cref{lem:level-sums} and $b_j =
2^{-jn/p} \|\lambda_{j,\cdot}\|_{\ell_p}$, the almost diagonal lemma (\Cref{lem:almost-diagonal})
with exponent $\sigma$ sums the levels, since $2^{js} b_j = 2^{j(s - n/p)}
\|\lambda_{j,\cdot}\|_{\ell_p}$.
\end{proof}

\begin{lemma}[Synthesis of one level]
\label{lem:single-level}
\lean{Rd.isRegularAtomOn_zero, Rd.modNorm_level_le}
\leanok
\uses{thm:rd-atomic-synthesis, lem:lq-basic}
For regular atoms of level $j$ with the overlap of \Cref{lem:level-sums}, $F = \sum_m \lambda_m a_m$
in $L_p$ and $0 < s < \sigma$,
$\|F\|_p + \|(2^{ks} \omega_r(F, 2^{-k})_p)_k\|_{\ell_q} \leq K 2^{j(s - n/p)} \|\lambda\|_{\ell_p}$,
with the constant $K$ of \Cref{thm:rd-atomic-synthesis}.
\end{lemma}

\begin{proof}
\leanok
\Cref{thm:rd-atomic-synthesis} with the coefficients and atoms of the other levels equal to $0$; the
zero function is a regular atom, and a sequence with one nonzero term $x$ has $\ell_q$-norm $x$
(\Cref{lem:lq-basic}).
\end{proof}

\begin{lemma}[Atoms from pointwise bounds]
\label{lem:pointwise-atoms}
\lean{Rd.eLpNorm_le_of_bound_of_support, ENNReal.rpow_mul_dyadicWeight_neg,
  Rd.isRegularAtomOn_of_pointwise, Rd.isRegularAtom_of_pointwise,
  Rd.enorm_higherDiff_le_of_iteratedFDeriv}
\leanok
\uses{def:regular-atoms, lem:smooth-differences, lem:dyadic-ball-volume}
Let $a$ be measurable with $a = 0$ outside $\overline{B}(c, d 2^{-j})$, $|a| \leq A$ and
$|\Delta_h^r a| \leq A \min(1, (2^j \|h\|)^\sigma)$. Then $a$ is a regular atom with the constant
$A ((r + 1) d^n |B_1|)^{1/p}$; on a measurable set $\Omega$ the bounds are needed only on $\Omega$
and on $\Omega_{rh}$. A smooth $a$ with $\|D^l a\| \leq A 2^{jl}$ for $l \leq r$ has
$|\Delta_h^r a| \leq 2^r A \min(1, (2^j \|h\|)^r)$.
\end{lemma}

\begin{proof}
\leanok
\uses{lem:conv-pointwise}
A function with $|g| \leq B$ that vanishes outside $S$ has $\|g\|_p \leq B |S|^{1/p}$. The
difference $\Delta_h^r a$ vanishes outside $r + 1$ translates of the ball, and
$|\overline{B}(c, d 2^{-j})| = d^n |B_1| 2^{-jn}$ (\Cref{lem:dyadic-ball-volume}). For smooth $a$,
\Cref{lem:smooth-differences} gives $|\Delta_h^r a| \leq A (2^j \|h\|)^r$, and
$|\Delta_h^r a| \leq 2^r \sup |a|$ (\Cref{lem:conv-pointwise}).
\end{proof}

\section{Atomic decompositions}
\label{sec:rd-atomic-decompositions}

\begin{lemma}[The cells partition the space]
\label{lem:cells-partition}
\lean{Rd.dyadicCell_disjoint, Rd.iUnion_dyadicCell, Rd.lintegral_eq_tsum_dyadicCell,
  Rd.eLpNorm_rpow_eq_tsum_dyadicCell, Rd.card_le_of_norm_sub_cellCenter_le}
\leanok
\uses{def:dyadic-cells, lem:dyadic-cells}
The cells of one level are disjoint and cover $V$. Hence $\int F = \sum_m \int_{P_{j,m}} F$ for
$F \geq 0$ and $\|g\|_p^p = \sum_m \|1_{P_{j,m}} g\|_p^p$. At most
$\lceil (2\rho + 2)^n \rceil$ centres of level $j$ lie in a ball of radius $\rho 2^{-j}$.
\end{lemma}

\begin{proof}
\leanok
Countable additivity of the integral, and \Cref{lem:dyadic-cells}.
\end{proof}

\begin{lemma}[Pointwise bounds of convolutions]
\label{lem:conv-pointwise}
\lean{Rd.lintegral_closedBall_le, Rd.lintegral_kernel_mul_le, Rd.enorm_conv_le_of_bound,
  Rd.integrable_translate_mul_of_locallyIntegrable, Rd.conv_add_apply, Rd.higherDiff_conv_apply,
  Rd.enorm_higherDiff_le_sum, Rd.enorm_higherDiff_le_of_bound_steps,
  Rd.enorm_higherDiff_le_of_bound}
\leanok
\uses{def:conv}
Let $1 \leq p < \infty$, $g \in L_p$, and $|K| \leq A$ with $K = 0$ outside $\overline{B}(0, \rho)$.
Then $\int_{\overline{B}(x, \rho)} |g| \leq |\overline{B}(x, \rho)|^{1 - 1/p} \|g\|_p$ and
\[
  |(K * g)(x)| \leq A |\overline{B}(0, \rho)|^{1-1/p} \|g\|_p \quad (x \in V).
\]
For a continuous $K$ with compact support the integrals $\int K(y + v) g(x - y) \, dy$ converge
absolutely, they equal $(K * g)(x + v)$, and $\Delta_h^r (K * g) = (\Delta_h^r K) * g$ at every point.
For every $f : V \to \C$ and all $h, x \in V$,
$|\Delta_h^r f(x)| \leq \sum_{l \leq r} \binom{r}{l} |f(x + l h)|$; if $|f(x + l h)| \leq B$
for $l = 0, \dots, r$, then $|\Delta_h^r f(x)| \leq 2^r B$, and if $|f| \leq B$, then
$|\Delta_h^r f| \leq 2^r B$.
\end{lemma}

\begin{proof}
\leanok
H\"older's inequality on the ball, and the translation invariance of the Lebesgue measure; the
differences are finite sums of translates.
\end{proof}

\begin{definition}[Cell atoms]
\label{def:cell-atoms}
\lean{Rd.cellCoef, Rd.cellAtom, Rd.enorm_cellCoef, Rd.cellCoef_mul_cellAtom,
  Rd.cellCoef_indicator, Rd.cellAtom_indicator, Rd.cellCoef_mul_cellAtom_of_indicator}
\leanok
\uses{def:dyadic-cells}
For $1 \leq p < \infty$, a kernel $K$, a function $g : V \to \C$ and $j \in \N$, the
coefficients of the cells of level $j$ are $\lambda_{j,m} = 2^{jn/p} \|1_{P_{j,m}} g\|_p$ when this
is finite and $\lambda_{j,m} = 0$ otherwise, and the cell atoms are
$a_{j,m} = \lambda_{j,m}^{-1} K * (1_{P_{j,m}} g)$, with $a_{j,m} = 0$ if $\lambda_{j,m} = 0$.
The coefficient and the atom of a cell depend only on the values of $g$ on that cell, and
$\lambda_{j,m} a_{j,m} = K * (1_{P_{j,m}} g)$ for every $m \in \Z^n$ with
$1_{P_{j,m}} g \in L_p$.
\end{definition}

\begin{lemma}[Cell decomposition of a convolution]
\label{lem:cell-decomposition}
\lean{Rd.volume_dyadicCell_ne_top, Rd.memLp_one_indicator_dyadicCell, Rd.tsum_indicator_dyadicCell,
  Rd.conv_tsum_of_finset, Rd.conv_eq_tsum_cellAtom_of_indicator, Rd.conv_eq_tsum_cellAtom,
  Rd.continuous_cellAtom}
\leanok
\uses{def:cell-atoms, lem:cells-partition, lem:bounded-overlap, lem:conv-pointwise,
  lem:dyadic-cells}
Let $1 \leq p < \infty$ and $j \in \N$.
\begin{enumerate}
  \item For every $m \in \Z^n$ the cell $P_{j,m}$ has finite volume, and
    $1_{P_{j,m}} g \in L_1$ for every $g \in L_p$.
  \item For every $g : V \to \C$ and every $y \in V$, $g(y) = \sum_m (1_{P_{j,m}} g)(y)$.
  \item For every continuous $K$ with compact support, all locally integrable
    $u_i : V \to \C$, $i \in I$, every $x \in V$ and every finite $T \subseteq I$ with
    $u_i(x - y) = 0$ whenever $i \notin T$ and $K(y) \neq 0$,
    $\bigl(K * \sum_i u_i\bigr)(x) = \sum_i (K * u_i)(x)$.
  \item For every $\rho \geq 0$, every continuous $K$ with $K = 0$ outside
    $\overline{B}(0, \rho)$ and every $g : V \to \C$ with $1_{P_{j,m}} g \in L_p$ for every
    $m \in \Z^n$, $(K * g)(x) = \sum_m \lambda_{j,m} a_{j,m}(x)$ at every $x \in V$, with finitely
    many nonzero terms. If moreover $g \in L_p$, the cell atoms are continuous.
\end{enumerate}
\end{lemma}

\begin{proof}
\leanok
(1) A cell lies in a closed ball by \Cref{lem:dyadic-cells}. (2) Only the cell containing $y$
contributes (\Cref{lem:cells-partition}). (3) Only the finitely many terms of $T$ contribute, and
the integral of a finite sum is the sum of the integrals (\Cref{lem:conv-pointwise}). (4) Only the
finitely many cells that meet $\overline{B}(x, \rho)$ contribute (\Cref{lem:cells-partition},
\Cref{lem:bounded-overlap}); apply (2) and (3) to $u_m = 1_{P_{j,m}} g$.
\end{proof}

Item (4) applies in particular to every $g \in L_p$.

\begin{lemma}[Cell atoms are atoms]
\label{lem:cell-atoms-regular}
\lean{Rd.cellAtom_eq_zero, Rd.volume_closedBall_div_two_pow_rpow,
  Rd.enorm_inv_cellCoef_mul_conv_le, Rd.enorm_cellAtom_le, Rd.enorm_higherDiff_cellAtom_le,
  Rd.cellAtomConst, Rd.cellAtomConst_ne_top, Rd.isRegularAtom_cellAtom}
\leanok
\uses{lem:cell-decomposition, lem:conv-pointwise, lem:pointwise-atoms}
Let $K$ be continuous with $K = 0$ outside $\overline{B}(0, \rho 2^{-j})$ and
$|\Delta_h^l K| \leq A 2^{jn} \min(1, (2^j \|h\|)^l)$ for $l \in \{0, r\}$. Then the cell atoms of
level $j$ vanish outside $\overline{B}(c_{j,m}, (\rho + n) 2^{-j})$,
\begin{gather*}
  |a_{j,m}| \leq A (\rho^n |B_1|)^{1 - 1/p}, \\
  |\Delta_h^r a_{j,m}| \leq 2^r A ((\rho + r)^n |B_1|)^{1 - 1/p} \min(1, (2^j \|h\|)^r),
\end{gather*}
so they are regular atoms of order and exponent $r$ with radius factor $\rho + n$.
\end{lemma}

\begin{proof}
\leanok
By \Cref{lem:conv-pointwise} with $\lambda_{j,m} = 2^{jn/p} \|1_{P_{j,m}} g\|_p$: the factors
$2^{jn}$, $|\overline{B}(0, \rho 2^{-j})|^{1-1/p} = (\rho^n |B_1|)^{1-1/p} 2^{-jn(1-1/p)}$ and
$2^{-jn/p}$ cancel. For $\|h\| \leq 2^{-j}$ the kernel $\Delta_h^r K$ of $\Delta_h^r a_{j,m}$
vanishes outside $\overline{B}(0, (\rho + r) 2^{-j})$; for larger $h$,
$|\Delta_h^r a_{j,m}| \leq 2^r \sup |a_{j,m}|$. \Cref{lem:pointwise-atoms} concludes.
\end{proof}

\begin{lemma}[The coefficients of the cells]
\label{lem:cell-coefficients}
\lean{Rd.levelNorm_cellCoef, Rd.coeffNorm_cellCoef}
\leanok
\uses{lem:cells-partition, def:cell-atoms}
$\|\lambda_{j,\cdot}\|_{\ell_p} = 2^{jn/p} \|g\|_p$. Hence the cell coefficients of $g_j \in L_p$
have $\|\lambda\|_{b^s_{p,q}} = \|(2^{js} \|g_j\|_p)_j\|_{\ell_q}$.
\end{lemma}

\begin{proof}
\leanok
\Cref{lem:cells-partition}.
\end{proof}

\begin{lemma}[Bounds of the dual kernels]
\label{lem:dual-kernel-bounds}
\lean{ENNReal.ofReal_two_pow_pow, Rd.LocalMeansBump.scaledMoment_eq_zero_of_lt,
  Rd.LocalMeansBump.dualKernel_eq_zero_of_lt, Rd.LocalMeansBump.hasCompactSupport_dualKernel,
  Rd.LocalMeansBump.continuous_dualKernel, Rd.LocalMeansBump.exists_higherDiff_dualKernel_bound}
\leanok
\uses{def:local-means, lem:smooth-differences, lem:resolution-of-unity, lem:rd-differences}
For a bump $g$ of radius $R$ and $L, r \in \N$, the kernels $m_{2^j}$ and $\eta_j$ are continuous
and vanish outside $\overline{B}(0, L R 2^{-j})$ and $\overline{B}(0, 2 L R 2^{-j})$, and there is
$A < \infty$ with $|\Delta_h^l \eta_j| \leq A 2^{jn} \min(1, (2^j \|h\|)^l)$ for all $l \leq r$.
\end{lemma}

\begin{proof}
\leanok
The moment kernel is smooth with compact support, so its derivatives are bounded
(\Cref{lem:resolution-of-unity}). \Cref{lem:smooth-differences} bounds its differences, and
$\Delta_h^l (\Phi_c) = (\Delta_{ch}^l \Phi)_c$ for the dilations (\Cref{lem:rd-differences}).
\end{proof}

\begin{lemma}[The cell atoms of the dual kernels]
\label{lem:cell-atoms-dual}
\lean{Rd.LocalMeansBump.exists_cellAtom_pointwise_bound}
\leanok
\uses{lem:cell-atoms-regular, lem:dual-kernel-bounds}
For every order $r$ there is $A < \infty$ such that the cell atoms of $\eta_j * g$, $g \in L_p$,
satisfy $|\Delta_h^l a_{j,m}| \leq A \min(1, (2^j \|h\|)^l)$ for all $l \leq r$, $j$ and $m$.
\end{lemma}

\begin{proof}
\leanok
\Cref{lem:cell-atoms-regular} with the bounds of \Cref{lem:dual-kernel-bounds}, for each $l \leq r$.
\end{proof}

\begin{theorem}[Atomic decomposition]
\label{thm:atomic-decomposition}
\lean{Rd.LocalMeansBump.atomic_decomposition}
\leanok
\uses{lem:reproducing-formula, lem:cell-decomposition, lem:cell-atoms-regular,
  lem:cell-coefficients, lem:dual-kernel-bounds}
Let $g$ be a bump of radius $R$, $1 \leq p < \infty$ and $L, r \in \N$. There is $C < \infty$ such
that for every $F \in L_p$ the terms of the reproducing formula $F = \sum_j \eta_j * (k_j * F)$
(\Cref{lem:reproducing-formula}) are $\eta_j * (k_j * F) = \sum_m \lambda_{j,m} a_{j,m}$ almost
everywhere, with the cell atoms of $k_j * F$; these are regular atoms of order and exponent $r$
with radius factor $2LR + n$ and constant $C$, and
$\|\lambda\|_{b^s_{p,q}} = \|(2^{js} \|k_j * F\|_p)_j\|_{\ell_q}$ for all $s$ and $q$.
\end{theorem}

\begin{proof}
\leanok
\Cref{lem:cell-decomposition} for the kernels $\eta_j$, whose bounds
(\Cref{lem:dual-kernel-bounds}) make the cell atoms regular (\Cref{lem:cell-atoms-regular}), and
\Cref{lem:cell-coefficients} for the coefficients.
\end{proof}

\begin{theorem}[Atomic characterisation of Besov spaces]
\label{thm:atomic-characterisation}
\lean{Rd.LocalMeansBump.besov_iff_atomic}
\leanok
\uses{thm:atomic-decomposition, thm:rd-atomic-synthesis, thm:local-means, lem:cells-partition}
Let $1 \leq p < \infty$, $L \geq 1$, $0 < s < L$, $1 \leq q \leq \infty$ and $F \in L_p$. Then
$\|F\|_p + \|(2^{ks} \omega_L(F, 2^{-k})_p)_k\|_{\ell_q} < \infty$ if and only if
$F = \sum_j f_j$ in $L_p$ with $f_j = \sum_m \lambda_{j,m} a_{j,m}$ almost everywhere, for regular
atoms of order and exponent $L$ whose centres satisfy the overlap condition of
\Cref{lem:level-sums}, and $\|\lambda\|_{b^s_{p,q}} < \infty$.
\end{theorem}

\begin{proof}
\leanok
This is \cite[Theorem 2.6]{SchneiderVybiral} for $p \geq 1$. If the norm is finite, so is the local
means quasi-norm (\Cref{thm:local-means}), and \Cref{thm:atomic-decomposition} with $r = L$ gives
the decomposition, with the overlap of \Cref{lem:cells-partition}. The converse is
\Cref{thm:rd-atomic-synthesis} with $\sigma = r = L$.
\end{proof}

\section{Smooth atoms}

On a convex set $\Omega$ with nonempty interior, derivatives within $\Omega$ bound the differences
whose steps stay in $\Omega$. Smooth atoms on $\Omega$ are therefore regular atoms on $\Omega$, and
the cell atoms of the atomic decomposition are smooth atoms of every order. Smooth atoms keep their
form under smooth changes of variables, which is how Besov spaces move between charts.

\begin{lemma}[Differences through derivatives within a convex set]
\label{lem:convex-differences}
\lean{Rd.fwdDiff_iter_comp_ofReal, Rd.fwdDiff_iter_taylorWithinEval_eq_zero,
  Rd.norm_fwdDiff_iter_le_of_iteratedDerivWithin, Rd.norm_iteratedFDeriv_line_le,
  Rd.add_smul_mem_of_mem_diffStepSet, Rd.norm_higherDiff_le_of_iteratedFDerivWithin}
\leanok
\uses{def:rd-modulus-on, lem:rd-differences}
Let $\Omega \subseteq V$ be convex with nonempty interior, $r \geq 1$, and let $a$ be $r$ times
continuously differentiable within $\Omega$ with $\|D^i a(y)\| \leq C$ for $i \leq r$ and
$y \in \Omega$, the derivatives taken within $\Omega$. Then for $x \in \Omega_{rh}$
\[
  |\Delta_h^r a(x)| \leq 2^r r^{r+1} C \|h\|^r .
\]
\end{lemma}

\begin{proof}
\leanok
By convexity the segment $[x, x + rh]$ lies in $\Omega$. The function $g(t) = a(x + th)$ on
$[0, r]$ has $|g^{(r)}| \leq r!\, C \|h\|^r$ by the chain rule within $\Omega$ (Mathlib's
\texttt{norm\_iteratedFDerivWithin\_comp\_le}). The $r$-th difference of step $1$ annihilates the
Taylor polynomial of $g$ of degree $r - 1$ at $0$, and Taylor's theorem (Mathlib's
\texttt{taylor\_mean\_remainder\_bound}) bounds the remainder on $[0, r]$ by
$\sup |g^{(r)}|\, r^r / (r-1)!$. The binomial coefficients add up to $2^r$, and
$2^r r^r r! / (r - 1)! = 2^r r^{r+1}$.
\end{proof}

\begin{definition}[Smooth atoms]
\label{def:smooth-atoms}
\lean{Rd.IsSmoothAtom, Rd.isSmoothAtom_zero, Rd.IsSmoothAtom.of_univ}
\leanok
\uses{def:regular-atoms}
Let $\Omega \subseteq V$, $K \in \N$ and $d, A \geq 0$. A smooth atom on $\Omega$ of level $j$
and order $K$ with centre $c$ is a measurable $a : V \to \C$ with $a = 0$ outside
$\overline{B}(c, d 2^{-j})$ that is $K$ times continuously differentiable within $\Omega$ with
\[
  \|D^l a(y)\| \leq A 2^{jl} \qquad (l \leq K,\ y \in \Omega).
\]
The function $0$ is a smooth atom, and a smooth atom on $V$ is a smooth atom on every set of
unique differentiability, where the derivatives within the set are the derivatives.
\end{definition}

These are the $K$-atoms of \cite[Definition 1.3]{SchneiderVybiral} on a set, without moment
conditions.

\begin{lemma}[Smooth atoms are regular atoms]
\label{lem:smooth-atoms-regular}
\lean{Rd.IsSmoothAtom.enorm_higherDiff_le, Rd.isRegularAtomOn_of_isSmoothAtom}
\leanok
\uses{def:smooth-atoms, lem:convex-differences, lem:pointwise-atoms}
Let $\Omega$ be measurable and convex with nonempty interior, and let $a$ be a smooth atom on
$\Omega$ of order $K$ and level $j$. For $1 \leq r \leq K$ and $x \in \Omega_{rh}$,
$|\Delta_h^r a(x)| \leq 2^r r^{r+1} A \min(1, (2^j \|h\|)^r)$, and $a$ is a regular atom on
$\Omega$ of order and exponent $r$ with the constant
$2^r r^{r+1} A ((r + 1) d^n |B_1|)^{1/p}$.
\end{lemma}

\begin{proof}
\leanok
\uses{lem:conv-pointwise}
\Cref{lem:convex-differences} with $C = A 2^{jr}$ gives the bound $(2^j \|h\|)^r$, and
$|\Delta_h^r a| \leq 2^r A$ on $\Omega_{rh}$ since all its points lie in $\Omega$
(\Cref{lem:conv-pointwise}).
\Cref{lem:pointwise-atoms} on $\Omega$ concludes.
\end{proof}

\begin{lemma}[Derivatives of convolutions]
\label{lem:conv-derivatives}
\lean{Rd.conv_eq_convolution_mul_real, Rd.contDiff_conv, Rd.fderiv_conv_apply,
  Rd.iteratedFDeriv_conv_apply, Rd.norm_iteratedFDeriv_const_mul_conv_le}
\leanok
\uses{def:conv}
Let $K$ be smooth with compact support and $g$ locally integrable. Then $K * g$ is smooth and
\[
  D^l (K * g)(x)(v_1, \dots, v_l) = \bigl((D^l K)(\cdot)(v_1, \dots, v_l) * g\bigr)(x) .
\]
Hence $\|D^l (c\, K * g)(x)\| \leq B$ as soon as
$|c\, ((D^l K)(\cdot)(v_1, \dots, v_l) * g)(x)| \leq B \prod_i \|v_i\|$ for all $v_i$.
\end{lemma}

\begin{proof}
\leanok
Mathlib's \texttt{HasCompactSupport.hasFDerivAt\_convolution\_left} differentiates under the
integral; induction on $l$ applies it to the directional derivatives of $K$.
\end{proof}

\begin{lemma}[Derivatives of the dual kernels]
\label{lem:dual-kernel-derivatives}
\lean{Rd.contDiff_dilate, Rd.ofReal_two_pow_mul, Rd.norm_iteratedFDeriv_dilate_le,
  Rd.LocalMeansBump.contDiff_scaledMoment, Rd.LocalMeansBump.contDiff_dualKernel,
  Rd.LocalMeansBump.dualDerivConst, Rd.LocalMeansBump.dualDerivConst_nonneg,
  Rd.LocalMeansBump.norm_iteratedFDeriv_momentKernel_le,
  Rd.LocalMeansBump.norm_iteratedFDeriv_dualKernel_le,
  Rd.LocalMeansBump.exists_iteratedFDeriv_dualKernel_bound}
\leanok
\uses{def:local-means, lem:dual-kernel-bounds}
Dilations of smooth functions are smooth, with
$\|D^l (\Phi_c)(x)\| \leq |c|^{n+l} \sup_y \|D^l \Phi(y)\|$. For a bump $g$ and $L, K \in \N$ the
kernels $m_{2^j}$ and $\eta_j$ are smooth, and with
$A_{L,K} = 2 \sum_{l=0}^K \sup_y \|D^l m(y)\| < \infty$,
$\|D^l \eta_j(y)\| \leq A_{L,K} 2^{j(n+l)}$ for all $l \leq K$, $j \in \N$ and $y \in V$.
\end{lemma}

\begin{proof}
\leanok
$\Phi_c(x) = c^n \Phi(cx)$, and the chain rule for the linear map $x \mapsto cx$. The moment kernel
is smooth with compact support, so its derivatives are bounded by its Schwartz seminorms, and
$\eta_j$ is a sum of two dilations.
\end{proof}

\begin{lemma}[Cell atoms are smooth atoms]
\label{lem:cell-atoms-smooth}
\lean{Rd.norm_iteratedFDeriv_cellAtom_le, Rd.LocalMeansBump.isSmoothAtom_cellAtom,
  Rd.LocalMeansBump.exists_isSmoothAtom_cellAtom}
\leanok
\uses{lem:conv-derivatives, lem:dual-kernel-derivatives, lem:cell-atoms-regular}
Let $K$ be smooth with $K = 0$ outside $\overline{B}(0, \rho 2^{-j})$ and
$\|D^l K\| \leq A 2^{j(n+l)}$. Then the cell atoms of $K * g$, $g \in L_p$, have
$\|D^l a_{j,m}\| \leq A 2^{jl} (\rho^n |B_1|)^{1-1/p}$. For a bump of radius $R$, every order
$K$, every level $j$ and every $g \in L_p$, every cell atom of $\eta_j * g$ is a smooth atom on
$V$ of order $K$ with radius factor $2LR + n$ and constant
$A_{L,K} ((2LR)^n |B_1|)^{1-1/p}$, with $A_{L,K}$ of \Cref{lem:dual-kernel-derivatives}.
\end{lemma}

\begin{proof}
\leanok
\Cref{lem:conv-derivatives} moves the derivatives onto the kernel, and the bound of
\Cref{lem:conv-pointwise} for the kernel $(D^l K)(\cdot)(v_1, \dots, v_l)$ cancels the factor
$2^{jn}$ against the coefficient as in \Cref{lem:cell-atoms-regular}.
\end{proof}

\begin{theorem}[Atomic decomposition into smooth atoms]
\label{thm:smooth-atomic-decomposition}
\lean{Rd.LocalMeansBump.atomic_decomposition_smooth, Rd.stdBumpBase, Rd.LocalMeansBump.std,
  Rd.LocalMeansBump.eLpNorm_std_one}
\leanok
\uses{thm:atomic-decomposition, lem:cell-atoms-smooth, lem:bump-of-smooth-bump}
Let $g$ be a bump of radius $R$, $1 \leq p < \infty$ and $L, K \in \N$. There is $A$ such that for
every $F \in L_p$ the terms of the reproducing formula are
$\eta_j * (k_j * F) = \sum_m \lambda_{j,m} a_{j,m}$ almost everywhere, with cell atoms that are
smooth atoms on $V$ of order $K$, radius factor $2LR + n$ and constant $A$, and
$\|\lambda\|_{b^s_{p,q}} = \|(2^{js} \|k_j * F\|_p)_j\|_{\ell_q}$. The normalised smooth bump of
radius $1$ around $0$ is a bump of $L_1$ norm $1$.
\end{theorem}

\begin{proof}
\leanok
\Cref{thm:atomic-decomposition} for the decomposition and the coefficients, and
\Cref{lem:cell-atoms-smooth} for the atoms. The normalised bump is nonnegative with integral
$1$, so its $L_1$ norm is $1$.
\end{proof}

\begin{definition}[Smooth changes of variables with a cut-off]
\label{def:cutoff-diffeo}
\lean{Rd.CutoffDiffeo, Rd.CutoffDiffeo.apply, Rd.CutoffDiffeo.imageCompact,
  Rd.CutoffDiffeo.isCompact_imageCompact, Rd.CutoffDiffeo.imageCompact_subset,
  Rd.CutoffDiffeo.apply_eq_zero_of_notMem, Rd.CutoffDiffeo.apply_eq_zero_of_forall_far,
  Rd.CutoffDiffeo.uniqueDiffOn_source, Rd.CutoffDiffeo.uniqueDiffOn_target,
  Rd.CutoffDiffeo.continuousOn_apply}
\leanok
Let $\Omega \subseteq V$. A smooth change of variables with a cut-off on $\Omega$ consists of open
sets $U, U'$, a bijection $\Phi : \Omega \cap U \to \Omega \cap U'$ that is smooth within
$\Omega \cap U$ with inverse $\Psi$ smooth within $\Omega \cap U'$, a compact
$K_0 \subseteq \Omega \cap U'$, and a function $\chi$ that is smooth within $\Omega \cap U'$ and
vanishes on $\Omega \setminus K_0$. It acts on functions by
\[
  T f = 1_{\Omega \cap U}\, (\chi f) \circ \Phi .
\]
$T f$ vanishes outside the compact set $\Psi(K_0) \subseteq \Omega \cap U$, and $T a = 0$ if $a$
vanishes on a ball that contains no point of $K_0$, and $T f$ is continuous on $\Omega$ for
continuous $f$. If $\Omega$ is convex with nonempty interior, $\Omega \cap U$ and $\Omega \cap U'$
are sets of unique differentiability.
\end{definition}

The pieces of a function in two overlapping charts of a manifold are related in this way; for
$\Omega = V$ and $\Phi = \mathrm{id}$, $T$ is a pointwise multiplier.

\begin{lemma}[Bounds of the change of variables]
\label{lem:cutoff-diffeo-bounds}
\lean{Rd.locallyLipschitzOn_inter, Rd.exists_lipschitzOnWith_of_contDiffOn_inter,
  Rd.iteratedFDerivWithin_eq_zero_of_eqOn, Rd.CutoffDiffeo.exists_lipschitzOnWith_Psi,
  Rd.CutoffDiffeo.exists_lipschitzOnWith_Phi, Rd.CutoffDiffeo.exists_iteratedFDerivWithin_chi_bound,
  Rd.CutoffDiffeo.exists_iteratedFDerivWithin_Phi_bound}
\leanok
\uses{def:cutoff-diffeo, lem:bounded-derivatives-compact}
A map of class $C^1$ within a relatively open part $\Omega \cap U$ of a convex set is Lipschitz on
the compact subsets of $\Omega \cap U$. Hence $\Psi$ is Lipschitz on $K_0$ and $\Phi$ on
$\Psi(K_0)$. For every $K$ the derivatives of $\chi$ of order at most $K$ are bounded on
$\Omega \cap U'$, and those of $\Phi$ on $\Psi(K_0)$.
\end{lemma}

\begin{proof}
\leanok
The mean value inequality on the convex sets $\Omega \cap B$ for small balls $B$ gives local
Lipschitz bounds within $\Omega$; on a compact set finitely many of them suffice. The derivatives
are continuous within the sets, hence bounded on the compact sets $K_0$ and $\Psi(K_0)$
(\Cref{lem:bounded-derivatives-compact}); the
derivatives of $\chi$ vanish on $\Omega \cap U' \setminus K_0$, where $\chi$ vanishes near
every point.
\end{proof}

\begin{lemma}[Overlap of transported centres]
\label{lem:transported-centres}
\lean{Rd.card_le_of_lipschitz_centres}
\leanok
\uses{lem:cells-partition}
Let $\Phi$ be $L$-Lipschitz on $B$, and let $z_m \in B$ with
$\|\Phi(z_m) - c_{j,m}\| \leq d 2^{-j}$ for the centres $c_{j,m}$ of level $j$. Then at most
$\lceil (2(2L\rho + d) + 2)^n \rceil$ of the points $z_m$ lie in a ball of radius $\rho 2^{-j}$.
\end{lemma}

\begin{proof}
\leanok
If $z_{m_0}$ and $z_m$ lie in such a ball, then
$\|\Phi(z_{m_0}) - c_{j,m}\| \leq (2L\rho + d) 2^{-j}$, and \Cref{lem:cells-partition} counts the
centres in the ball around $\Phi(z_{m_0})$.
\end{proof}

\begin{theorem}[Smooth atoms under changes of variables]
\label{thm:smooth-atoms-change-of-variables}
\lean{Rd.CutoffDiffeo.norm_iteratedFDerivWithin_mul_atom_le, Rd.CutoffDiffeo.apply_eq_zero_of_far,
  Rd.CutoffDiffeo.isSmoothAtom_apply}
\leanok
\uses{def:smooth-atoms, def:cutoff-diffeo, lem:cutoff-diffeo-bounds}
Let $\Omega$ be measurable and convex with nonempty interior, $T$ a smooth change of variables
with a cut-off on $\Omega$, and $K \in \N$. There are $L, A' \geq 0$ such that for every smooth
atom $a$ on $V$ of order $K$, level $j$, radius factor $d$, constant $A$ and centre $c$, and every
$y_0 \in K_0$ with $\|y_0 - c\| \leq d 2^{-j}$, $T a$ is a smooth atom on $\Omega$ of order $K$
and level $j$ with centre $\Psi(y_0)$, radius factor $2Ld$ and constant $A' A$.
\end{theorem}

\begin{proof}
\leanok
If $T a(x) \neq 0$, then $\Phi(x) \in K_0$ and $\|\Phi(x) - c\| \leq d 2^{-j}$, so
$\|x - \Psi(y_0)\| \leq L \|\Phi(x) - y_0\| \leq 2 L d 2^{-j}$ with a Lipschitz constant $L$ of
$\Psi$ on $K_0$. By Leibniz's rule within $\Omega \cap U'$,
$\|D^i (\chi a)\| \leq 2^i M A 2^{ji}$ (Mathlib's \texttt{norm\_iteratedFDerivWithin\_mul\_le}),
and by the chain rule within sets (Mathlib's \texttt{norm\_iteratedFDerivWithin\_comp\_le}),
$\|D^l (T a)(x)\| \leq l!\, 2^l M A D^l 2^{jl}$ on $\Psi(K_0)$, with the bounds $M$ of $\chi$ and
$D \geq 1$ of $\Phi$ from \Cref{lem:cutoff-diffeo-bounds}. Outside the compact $\Psi(K_0)$ the
function $T a$ vanishes near every point of $\Omega$, so its derivatives within $\Omega$ vanish.
Hence $A' = K!\, 2^K M D^K$.
\end{proof}

\section{Invariance under changes of variables}

A smooth change of variables with a cut-off maps the Besov spaces of $V$ into the intrinsic Besov
spaces of a convex set, and, on sets with a bounded extension, the intrinsic spaces into
themselves. The proof transports smooth atoms (\Cref{thm:smooth-atoms-change-of-variables}).

\begin{lemma}[The change of variables on $L_p$]
\label{lem:cutoff-diffeo-lp}
\lean{Rd.CutoffDiffeo.apply_add, Rd.CutoffDiffeo.apply_smul, Rd.CutoffDiffeo.apply_tsum,
  Rd.CutoffDiffeo.apply_ae_eq, Rd.CutoffDiffeo.aestronglyMeasurable_apply,
  Rd.CutoffDiffeo.lintegral_apply_eq, Rd.CutoffDiffeo.exists_det_bound,
  Rd.CutoffDiffeo.eLpNorm_apply_le, Rd.CutoffDiffeo.memLp_apply, Rd.CutoffDiffeo.applyLp,
  Rd.CutoffDiffeo.applyLp_ae}
\leanok
\uses{def:cutoff-diffeo}
Let $\Omega$ be measurable and convex with nonempty interior, and $T$ a smooth change of variables
with a cut-off on $\Omega$. Then $T$ is linear, $T g_1 = T g_2$ almost everywhere if $g_1 = g_2$
almost everywhere on $\Omega$, and for $e > 0$
\[
  \int |T g|^e = \int_{K_0} |\det D\Psi(y)|\, |\chi(y)|^e |g(y)|^e \, dy .
\]
Hence $T$ is a bounded linear operator on $L_p$, $1 \leq p < \infty$. If for every $m$ either
$\lambda'_m = \lambda_m$ or $T a_m = 0$, then $T(\sum_m \lambda_m a_m) = \sum_m \lambda'_m T a_m$.
\end{lemma}

\begin{proof}
\leanok
If $T g_1(x) \neq T g_2(x)$, then $x = \Psi(y)$ with $y \in K_0$ and $g_1(y) \neq g_2(y)$; the
differentiable map $\Psi$ maps the null set of these $y$ to a null set (Mathlib's
\texttt{addHaar\_image\_eq\_zero\_of\_differentiableOn\_of\_addHaar\_eq\_zero}). $T g$ vanishes
outside $\Psi(K_0)$, and the change of variables $x = \Psi(y)$ on the compact $K_0$ (Mathlib's
\texttt{lintegral\_image\_eq\_lintegral\_abs\_det\_fderiv\_mul}) gives the integral; the Jacobian
and $\chi$ are bounded on $K_0$.
\end{proof}

\begin{lemma}[Transport of smooth atoms]
\label{lem:transported-atoms}
\lean{Rd.CutoffDiffeo.exists_transported_atoms}
\leanok
\uses{thm:smooth-atoms-change-of-variables, lem:smooth-atoms-regular, lem:transported-centres}
Let $\Omega$ be measurable and convex with nonempty interior, $r \geq 1$, and $d, A \geq 0$. There
are $d'$, $C' < \infty$ and $N$ such that for smooth atoms $a_{j,m}$ on $V$ of order $r$, radius
factor $d$ and constant $A$ centred at the cell centres $c_{j,m}$, and all coefficients
$\lambda_{j,m}$, the functions $T a_{j,m}$ are regular atoms on $\Omega$ of order and exponent
$r$ with radius factor $d'$ and constant $C'$, with coefficients $\lambda'_{j,m}$ such that
$|\lambda'_{j,m}| \leq |\lambda_{j,m}|$, $\lambda'_{j,m} = \lambda_{j,m}$ unless $T a_{j,m} = 0$,
and at most $N$ centres with $\lambda'_{j,m} \neq 0$ lie in a ball of radius $(d' + r) 2^{-j}$.
\end{lemma}

\begin{proof}
\leanok
If a point $y \in K_0$ lies within $d 2^{-j}$ of $c_{j,m}$, the new centre is $\Psi(y)$ and
$\lambda'_{j,m} = \lambda_{j,m}$ (\Cref{thm:smooth-atoms-change-of-variables},
\Cref{lem:smooth-atoms-regular}); otherwise $T a_{j,m} = 0$ and $\lambda'_{j,m} = 0$. Since
$\Phi(\Psi(y)) = y$ and $\Phi$ is Lipschitz on $\Psi(K_0)$, \Cref{lem:transported-centres} counts
the centres.
\end{proof}

\begin{theorem}[Invariance under changes of variables]
\label{thm:diffeo-invariance}
\lean{Rd.CutoffDiffeo.modNormOn_applyLp_le}
\leanok
\uses{lem:cutoff-diffeo-lp, lem:transported-atoms, thm:smooth-atomic-decomposition,
  thm:rd-atomic-synthesis, thm:local-means}
Let $\Omega$ be measurable and convex with nonempty interior, $T$ a smooth change of variables
with a cut-off on $\Omega$, $1 \leq p < \infty$, $0 < s < r$ and $1 \leq q \leq \infty$. There is
$C < \infty$ such that for all $F \in L_p(V)$
\[
  \|T F\|_{L_p(\Omega)} + \|(2^{ks} \omega_r(T F, 2^{-k}, \Omega)_p)_k\|_{\ell_q} \leq
  C \bigl(\|F\|_p + \|(2^{ks} \omega_r(F, 2^{-k})_p)_k\|_{\ell_q}\bigr) .
\]
\end{theorem}

\begin{proof}
\leanok
Decompose $F$ into smooth cell atoms of order $r$ (\Cref{thm:smooth-atomic-decomposition}) for the
standard bump. $T$ is continuous on $L_p$, so $T F = \sum_j T f_j$ in $L_p$, and
$T f_j = \sum_m \lambda'_{j,m} T a_{j,m}$ almost everywhere by \Cref{lem:cutoff-diffeo-lp}. The
transported atoms (\Cref{lem:transported-atoms}) are regular atoms on $\Omega$ of exponent
$r > s$, and the synthesis on $\Omega$ (\Cref{thm:rd-atomic-synthesis}) bounds the left side by
$K \|\lambda'\|_{b^s_{p,q}} \leq K \|\lambda\|_{b^s_{p,q}}$, which is the local means quasi-norm of
$F$, bounded by the right side (\Cref{thm:local-means}).
\end{proof}

\begin{definition}[Besov domains]
\label{def:besov-domain}
\lean{Rd.BesovDomain, Rd.modNormOn_univ, Rd.BesovDomain.univ, Rd.BesovDomain.measurableSet}
\leanok
\uses{def:rd-modulus-on}
A Besov domain is a closed convex set $\Omega \subseteq V$ with nonempty interior such that for
$1 \leq p < \infty$, $0 < s < r$ and $1 \leq q \leq \infty$ there is $C < \infty$: every
$f \in L_p(V)$ with
$N_\Omega(f) = \|f\|_{L_p(\Omega)} + \|(2^{ks} \omega_r(f, 2^{-k}, \Omega)_p)_k\|_{\ell_q} < \infty$
agrees on $\Omega$ with some $F \in L_p(V)$ with
$\|F\|_p + \|(2^{ks} \omega_r(F, 2^{-k})_p)_k\|_{\ell_q} \leq C N_\Omega(f)$. The space $V$ is a
Besov domain, with $F = f$.
\end{definition}

The extension is part of the definition, so a Besov domain is a hypothesis. Its instances are
$V$, the closed half-space and the models of manifolds with and without boundary
(\Cref{thm:half-space-besov-domain}), the closures of Lipschitz graph domains with convex graph
function (\Cref{thm:graph-besov-domain}), the quadrant (the model of manifolds with corners,
\Cref{thm:quadrant-besov-domain}), and the closures of nonempty convex bounded Lipschitz domains
(\Cref{thm:bounded-lipschitz-besov-domain}) and of nonempty bounded open convex sets, among them
the closed balls of positive radius and the closed unit cube (\Cref{cor:convex-besov-domain}).
The class consists of closed convex sets; every Lipschitz graph domain and every bounded
Lipschitz domain, convex or not, satisfies the extension condition of
\Cref{def:besov-extension-property} (\Cref{thm:graph-besov-domain} and
\Cref{thm:bounded-lipschitz-besov-domain}).

\begin{theorem}[Invariance on Besov domains]
\label{thm:besov-domain-invariance}
\lean{Rd.CutoffDiffeo.applyLp_eq_of_indicatorLp_eq, Rd.CutoffDiffeo.modNormOn_applyLp_le_modNormOn,
  Rd.BesovDomain.modNormOn_order_le_of_besovDomain}
\leanok
\uses{def:besov-domain, thm:diffeo-invariance, thm:rd-besov}
Let $\Omega$ be a Besov domain, $T$ a smooth change of variables with a cut-off on $\Omega$,
$1 \leq p < \infty$, $0 < s < r$ and $1 \leq q \leq \infty$. There is $C < \infty$ with
$N_\Omega(T f) \leq C N_\Omega(f)$ for all $f$ with $N_\Omega(f) < \infty$. For $s < r'$ there is
$C < \infty$ with $N_\Omega(f; r') \leq C N_\Omega(f; r)$, where $N_\Omega(f; r)$ uses the
differences of order $r$; so the order is a free parameter.
\end{theorem}

\begin{proof}
\leanok
$T f$ depends only on $f$ almost everywhere on $\Omega$ (\Cref{lem:cutoff-diffeo-lp}). Extend $f$
to $F$ with $\|F\|_p + \|(2^{ks} \omega_r(F, 2^{-k})_p)_k\|_{\ell_q} \leq C N_\Omega(f)$; then
$T f = T F$, and \Cref{thm:diffeo-invariance} applies. For the order, change it on $V$
(\Cref{thm:rd-besov}) and restrict, $N_\Omega(f; r') \leq
\|F\|_p + \|(2^{ks} \omega_{r'}(F, 2^{-k})_p)_k\|_{\ell_q}$ (\Cref{lem:rd-modulus-on}).
\end{proof}

\begin{definition}[The extension property]
\label{def:besov-extension-property}
\lean{Rd.HasBesovExtension, Rd.BesovDomain.hasBesovExtension, Rd.BesovDomain.of_hasBesovExtension}
\leanok
\uses{def:besov-domain}
A measurable set $\Omega \subseteq V$ has the extension property if for every $1 \leq p < \infty$,
all $r \in \N$ and $s \in \R$ with $0 < s < r$, and every $1 \leq q \leq \infty$ there is
$C < \infty$ such that every $f \in L_p(V)$ with $N_\Omega(f) < \infty$ agrees on $\Omega$ with
some $F \in L_p(V)$ with $\|F\|_p + \|(2^{ks} \omega_r(F, 2^{-k})_p)_k\|_{\ell_q} \leq
C N_\Omega(f)$. For every closed convex set $\Omega \subseteq V$ with nonempty interior, $\Omega$ is
a Besov domain if and only if $\Omega$ has the extension property.
\end{definition}

The extension property is the extension condition of \Cref{def:besov-domain} without convexity.
It is the same for sets that agree up to a null set, and it makes the order of the differences a
free parameter, as on a Besov domain (\Cref{thm:besov-domain-invariance}). DeVore and Sharpley
define Besov spaces on domains by the same moduli $\omega_r(f, t, \Omega)_p$ and construct
extension operators to $\R^d$ on Lipschitz graph domains \cite[Section~4]{DeVoreSharpley} and on
$(\varepsilon, \delta)$-domains \cite[Section~5]{DeVoreSharpley}.

\begin{lemma}[The extension property up to null sets]
\label{lem:besov-extension-null-sets}
\lean{Rd.HasBesovExtension.congr_set, Rd.BesovDomain.of_ae_eq, Rd.BesovDomain.range_of_ae_eq}
\leanok
\uses{def:besov-extension-property, lem:null-set-domains}
For all measurable sets $U, U' \subseteq V$ that agree up to a null set, if $U$ has the extension
property, then so does $U'$. For every closed convex set $\Omega \subseteq V$ with nonempty
interior and every measurable set $U \subseteq V$ with the extension property that agrees with
$\Omega$ up to a null set, $\Omega$ is a Besov domain. For
every real model with corners $I$ with model vector space $V$ and every measurable set
$U \subseteq V$ with the extension property that agrees with $\operatorname{range} I$ up to a null
set, $\operatorname{range} I$ is a Besov domain.
\end{lemma}

\begin{proof}
\leanok
The cut-offs $1_U f$, $1_{U'} f$ and the moduli on $U$ and $U'$ agree (\Cref{lem:null-set-domains}),
so $N_U = N_{U'}$, and $F$ agrees with $f$ on $U$ if and only if it does on $U'$; the proof keeps
the constant $C$. The range of a real model with corners is closed, convex and has nonempty
interior (Mathlib's
\texttt{ModelWithCorners.\allowbreak isClosed\_\allowbreak range},
\texttt{ModelWithCorners.\allowbreak convex\_\allowbreak range} and
\texttt{ModelWithCorners.\allowbreak nonempty\_\allowbreak interior}).
\end{proof}

\begin{corollary}[The order of the differences]
\label{cor:besov-extension-order}
\lean{Rd.HasBesovExtension.modNormOn_order_le_of_hasBesovExtension}
\leanok
\uses{def:besov-extension-property, thm:rd-besov, lem:rd-modulus-on}
For every measurable set $\Omega \subseteq V$ with the extension property, every
$1 \leq p < \infty$, all $r, r' \in \N$ and $s \in \R$ with $0 < s < r$ and $s < r'$, and every
$1 \leq q \leq \infty$, there is $C < \infty$ with $N_\Omega(f; r') \leq C N_\Omega(f; r)$ for every
$f \in L_p(V)$ with $N_\Omega(f; r) < \infty$.
\end{corollary}

\begin{proof}
\leanok
Extend $f$ to $F$ with $\|F\|_p + \|(2^{ks} \omega_r(F, 2^{-k})_p)_k\|_{\ell_q} \leq
C_1 N_\Omega(f; r)$, change the order on $V$ (\Cref{thm:rd-besov}) and restrict,
$N_\Omega(f; r') \leq \|F\|_p + \|(2^{ks} \omega_{r'}(F, 2^{-k})_p)_k\|_{\ell_q}$
(\Cref{lem:rd-modulus-on}).
\end{proof}

\begin{theorem}[The half-space is a Besov domain]
\label{thm:half-space-besov-domain}
\lean{Rd.orthogonal_singleton_ne_top, Rd.volume_orthogonal_singleton_eq_zero,
  Rd.closedHalfSpace, Rd.isClosed_closedHalfSpace, Rd.convex_closedHalfSpace,
  Rd.mem_interior_closedHalfSpace, Rd.LipschitzGraph.halfSpaceGraph,
  Rd.LipschitzGraph.halfSpaceGraph_domain, Rd.closedHalfSpace_ae_eq, Rd.BesovDomain.closedHalfSpace,
  Rd.firstUnitVector, Rd.norm_firstUnitVector, Rd.coord_zero_eq_inner,
  Manifold.range_eq_closedHalfSpace, Rd.BesovDomain.range_euclideanHalfSpace,
  Rd.BesovDomain.range_of_boundaryless}
\leanok
\uses{def:besov-domain, lem:null-set-domains, thm:graph-besov}
For every nonzero $e \in V$, the hyperplane $e^\perp$ is a null set. For a unit vector $e$ the
closed half-space $\{x : \langle x, e \rangle \geq 0\}$ is a Besov domain. Hence the models of
manifolds are Besov domains: the range $\{y : y_0 \geq 0\} = \{y : \langle y, e_0 \rangle \geq 0\}$
of the model with corners of the Euclidean half-space, for the unit vector
$e_0 = (1, 0, \dots, 0)$, and the range $V$ of a model without boundary.
\end{theorem}

\begin{proof}
\leanok
The hyperplane $e^\perp$ is a proper subspace (Mathlib's \texttt{Submodule.\allowbreak
orthogonal\_\allowbreak eq\_\allowbreak top\_\allowbreak iff}), hence a null set (Mathlib's
\texttt{MeasureTheory.\allowbreak Measure.\allowbreak addHaar\_\allowbreak submodule}). The closed
half-space differs from the Lipschitz graph domain $\{x : \langle x, e \rangle > 0\}$ with graph
function $0$ by the null set $e^\perp$ (\Cref{lem:null-set-domains}), and Rychkov's
extension is bounded on the latter (\Cref{thm:graph-besov}).
\end{proof}

\section{Besov spaces as interpolation spaces}
\label{sec:rd-besov-interpolation}

The band-limited approximation spaces interpolate as the approximation spaces of every scheme
(\Cref{thm:approximation-spaces-reiteration}), and the Besov spaces by differences as the
classical spaces of every compatible modulus (\Cref{cor:modulus-spaces-interpolate}). This is
the interpolation theorem for Besov spaces~\cite[Theorem~6.4.5]{BerghLofstrom}
and~\cite[Section~2.4]{TriebelInterpolation}, here with the quasi-norms by differences, for
$1 \leq p < \infty$ and $0 < s_0 < s_1 < r$.

\begin{corollary}[Band-limited approximation spaces interpolate]
\label{cor:band-approximation-interpolate}
\lean{Rd.memKMethod_bandApproxCouple_iff, Rd.memKMethod_bandLeftApproxCouple_iff}
\leanok
\uses{def:band-space, def:approximation-space-couples, thm:approximation-spaces-reiteration}
Let $1 \leq p < \infty$, $0 < \eta < 1$, $q_0, q_1, q \in [1, \infty]$ and $f \in L_p$, and let
$A^s_q$ be the approximation spaces of the band-limited scheme.
\begin{enumerate}
  \item For every $0 < s_0 < s_1$, $f \in (A^{s_0}_{q_0}, A^{s_1}_{q_1})_{\eta, q}$ if and only
    if $f \in A^{(1-\eta)s_0 + \eta s_1}_q$.
  \item For every $s_1 > 0$, $f \in (L_p, A^{s_1}_{q_1})_{\eta, q}$ if and only if
    $f \in A^{\eta s_1}_q$.
\end{enumerate}
\end{corollary}

\begin{proof}
\leanok
This is \Cref{thm:approximation-spaces-reiteration} for the band-limited scheme.
\end{proof}

\begin{theorem}[Besov spaces by differences interpolate]
\label{thm:rd-besov-interpolate}
\lean{Rd.memKMethod_bandModCouple_iff, Rd.memKMethod_modCouple_iff_littlewoodPaley,
  Rd.memKMethod_bandLeftModCouple_iff}
\leanok
\uses{thm:rd-besov, def:modulus-couples, cor:modulus-spaces-interpolate, def:lp-blocks}
Let $1 \leq p < \infty$, let $r \geq 1$ be an integer, let $0 < \eta < 1$,
$q_0, q_1, q \in [1, \infty]$ and $f \in L_p$, and write
$N_{s, q}(f) = \|f\|_p + \|(2^{ks} \omega_r(f, 2^{-k})_p)_k\|_{\ell_q}$.
\begin{enumerate}
  \item For every $0 < s_0 < s_1 < r$, with $s = (1 - \eta) s_0 + \eta s_1$, the following are
    equivalent: $f$ lies in the space of the continuous $K$-method with $(\eta, q)$ of the
    couple $(N_{s_0, q_0}, N_{s_1, q_1})$; $N_{s, q}(f) < \infty$;
    $\|f\|^{\mathrm{LP}}_{s, q} < \infty$.
  \item For every $0 < s_1 < r$, $f$ lies in the space of the continuous $K$-method with
    $(\eta, q)$ of the couple $(\|\cdot\|_p, N_{s_1, q_1})$ if and only if
    $N_{\eta s_1, q}(f) < \infty$.
\end{enumerate}
\end{theorem}

\begin{proof}
\leanok
$N_{s, q}$ is the classical quasi-norm of the modulus $\omega_r(f, 2^{-k})_p$, which is
compatible with the band-limited scheme with inverse order $r$ (\Cref{thm:rd-besov}). Apply
\Cref{cor:modulus-spaces-interpolate}; the Littlewood--Paley form is part (2) of
\Cref{thm:rd-besov}.
\end{proof}

% !TEX root = print.tex
\chapter{Daubechies wavelets}
\label{ch:daubechies}

This chapter constructs the compactly supported orthonormal wavelets of
Daubechies~\cite{Daubechies} on $\R$ and describes the Besov spaces $\Besov{s}{p}{q}(\R)$ by
their coefficients. A mask of order $N$ is a real polynomial $L$ of degree less than $N$ with
$L(1) = 1$ and $|L(e^{-2\pi i \xi})|^2 = P_N(\sin^2 \pi\xi)$, where
$P_N(y) = \sum_{k < N} \binom{N-1+k}{k} y^k$. The scaling function $\varphi$ is the inverse
Fourier transform of the infinite product of the symbol of the mask. Everything is proved with
the Fourier transform of integrable functions: the refinement equation, the compact support
(through the cascade algorithm), the orthonormality of the integer translates, and the
reproduction of the polynomials of degree less than $N$. The regularity of $\varphi$ comes from
the transfer operator of Eirola~\cite{Eirola} and Villemoes~\cite{Villemoes}, which gives
$\varphi \in H^s$ for $s < 1$ if $N = 2$ and for $s < 3 - \log_2 3$ if $N = 3$, and, for the
first wavelet $N = 2$, from the matrices of Daubechies and Lagarias~\cite{DaubechiesLagarias},
which give the H\"older exponent $2 - \log_2(1 + \sqrt 3)$. With Jackson's and Bernstein's
inequalities the Daubechies coefficients then characterise $\Besov{s}{p}{q}(\R)$.

\section{The Daubechies polynomial}

This section proves the B\'ezout identity of the Daubechies polynomial $P_N$ and the bounds on
$P_N$ that the later sections use. At $y = \sin^2 \pi\xi$ the identity becomes the trigonometric
identity behind the quadrature mirror filter identity of the masks. Throughout,
$\binom{2N-1}{N}$ stands for $\binom{-1}{0} = 1$ if $N = 0$.

The polynomial $P_N$ consists of the first $N$ terms of the binomial series
$(1 - y)^{-N} = \sum_{k \geq 0} \binom{N-1+k}{k} y^k$.

\begin{definition}[The Daubechies polynomial]
\label{def:daub-poly}
\lean{Daubechies.daubPoly, Daubechies.eval_daubPoly, Daubechies.daubPoly_zero,
  Daubechies.daubPoly_one, Daubechies.daubPoly_two, Daubechies.daubPoly_three,
  Daubechies.natDegree_daubPoly_le, Daubechies.eval_daubPoly_succ_eq_one_add}
\leanok
For $N \in \N$ the \emph{Daubechies polynomial} of order $N$ is the real polynomial
\[
  P_N(y) = \sum_{k=0}^{N-1} \binom{N-1+k}{k} y^k
\]
of degree at most $N - 1$. Thus $P_0 = 0$, $P_1 = 1$, $P_2(y) = 1 + 2y$ and
$P_3(y) = 1 + 3y + 6y^2$, and for every $M \in \N$ the constant term of $P_{M+1}$ is $1$:
\[
  P_{M+1}(y) = 1 + \sum_{k=0}^{M-1} \binom{M+k+1}{k+1} y^{k+1} .
\]
\end{definition}

\begin{lemma}[Values at $0$ and $1$]
\label{lem:daub-poly-values}
\lean{Daubechies.eval_zero_daubPoly, Daubechies.eval_one_daubPoly}
\leanok
\uses{def:daub-poly}
For every $N \geq 1$, $P_N(0) = 1$ and $P_N(1) = \binom{2N-1}{N}$.
\end{lemma}

\begin{proof}
\leanok
Only the constant term $\binom{N-1}{0} = 1$ contributes to $P_N(0)$. For $P_N(1)$, write
$\binom{N-1+k}{k} = \binom{N-1+k}{N-1}$ and apply the hockey-stick identity
$\sum_{i=0}^{n} \binom{i+m}{m} = \binom{n+m+1}{m+1}$ with $n = m = N - 1$.
\end{proof}

The coefficients of $P_N$ and $P_{N+1}$ are linked by Pascal's rule.

\begin{lemma}[The recurrence]
\label{lem:daub-poly-recurrence}
\lean{Daubechies.one_sub_mul_eval_daubPoly_succ, Daubechies.one_sub_X_mul_daubPoly_succ,
  Daubechies.choose_two_mul_eq_two_mul_choose}
\leanok
\uses{def:daub-poly}
For every $N \in \N$ and every $y \in \R$,
\[
  (1 - y) P_{N+1}(y) = P_N(y) + \binom{2N-1}{N} y^N - \binom{2N}{N} y^{N+1} ,
\]
and the same identity holds for polynomials. For every $N \geq 1$,
$\binom{2N}{N} = 2 \binom{2N-1}{N}$.
\end{lemma}

\begin{proof}
\leanok
On the left the constant term is $1$, the coefficient of $y^{k+1}$ for $k < N$ is
$\binom{N+k+1}{k+1} - \binom{N+k}{k} = \binom{N+k}{k+1}$ by Pascal's rule, and the top term is
$-\binom{2N}{N} y^{N+1}$. On the right,
$P_N(y) + \binom{2N-1}{N} y^N = \sum_{k=0}^{N} \binom{N-1+k}{k} y^k$ has the same
coefficients. Pascal's rule $\binom{2N}{N} = \binom{2N-1}{N-1} + \binom{2N-1}{N}$ and the
symmetry $\binom{2N-1}{N-1} = \binom{2N-1}{N}$ give the second identity.
\end{proof}

\begin{theorem}[The B\'ezout identity]
\label{thm:daub-bezout}
\lean{Daubechies.eval_daubPoly_bezout_succ, Daubechies.eval_daubPoly_bezout,
  Daubechies.daubPoly_bezout}
\leanok
\uses{def:daub-poly}
For every $N \geq 1$ and every $y \in \R$,
\[
  (1 - y)^N P_N(y) + y^N P_N(1 - y) = 1 ,
\]
and the same identity holds for polynomials.
\end{theorem}

\begin{proof}
\leanok
\uses{lem:daub-poly-recurrence}
Induction on $N$; for $N = 1$ the identity reads $(1 - y) + y = 1$. By
\Cref{lem:daub-poly-recurrence}, from $N$ to $N + 1$ the first summand changes by
$y^N (1 - y)^N \bigl(\binom{2N-1}{N} - \binom{2N}{N} y\bigr)$ and the second by
$y^N (1 - y)^N \bigl(\binom{2N-1}{N} - \binom{2N}{N} (1 - y)\bigr)$. The changes add up to
$y^N (1 - y)^N \bigl(2 \binom{2N-1}{N} - \binom{2N}{N}\bigr) = 0$. Two real polynomials with the
same value at every real number are equal.
\end{proof}

\begin{lemma}[Bounds for $P_N$]
\label{lem:daub-poly-bounds}
\lean{Daubechies.eval_daubPoly_nonneg, Daubechies.one_le_eval_daubPoly,
  Daubechies.eval_daubPoly_le_eval_one, Daubechies.eval_daubPoly_le}
\leanok
\uses{def:daub-poly}
Let $N \in \N$. For every $y \geq 0$, $P_N(y) \geq 0$, and $P_N(y) \geq 1$ if $N \geq 1$. For
every $y \in [0, 1]$,
\[
  P_N(y) \leq P_N(1) \quad\text{and}\quad P_N(y) \leq 1 + \binom{2N-1}{N} y .
\]
\end{lemma}

\begin{proof}
\leanok
\uses{lem:daub-poly-values}
The coefficients of $P_N$ are nonnegative, and for $N \geq 1$ the constant term is $1$. For
$y \in [0, 1]$, $y^k \leq 1$ for every $k$ and $y^k \leq y$ for $k \geq 1$, and the coefficients
of the terms with $k \geq 1$ add up to $P_N(1) - 1 = \binom{2N-1}{N} - 1$
(\Cref{lem:daub-poly-values}). For $N = 0$, $P_N = 0$.
\end{proof}

\begin{lemma}[A binomial bound]
\label{lem:daub-choose-bound}
\lean{Daubechies.choose_lt_four_pow}
\leanok
For every $N \geq 2$, $\binom{2N-1}{N} < 4^{N-1}$.
\end{lemma}

\begin{proof}
\leanok
Let $N = M + 1$. The coefficients $\binom{2M+1}{i}$ with $i \leq M$ form one half of the row
$2M + 1$ of Pascal's triangle, so $\sum_{i=0}^{M} \binom{2M+1}{i} = 2^{2M} = 4^M$. This sum
contains the term $\binom{2M+1}{M} = \binom{2M+1}{M+1}$ and, since $M \geq 1$, the further term
$\binom{2M+1}{0} = 1$, so $\binom{2N-1}{N} = \binom{2M+1}{M+1} \leq 4^M - 1$.
\end{proof}

Equivalently, $N - \log_4 \bigl(2 \binom{2N-1}{N}\bigr) > \frac12$ for every $N \geq 2$.

\paragraph{The trigonometric form.}
On the circle the Daubechies polynomial is evaluated at $\sin^2 \pi\xi$: a mask $L$ of order $N$
has $|L(e^{-2\pi i \xi})|^2 = P_N(\sin^2 \pi\xi)$.

\begin{definition}[The trigonometric Daubechies polynomial]
\label{def:daub-fun}
\lean{Daubechies.daubFun, Daubechies.daubFun_two, Daubechies.daubFun_three}
\leanok
\uses{def:daub-poly}
For $N \in \N$ and $\xi \in \R$ let $P(\xi) = P_N(\sin^2 \pi\xi)$; the order $N$ is fixed and
omitted from the notation. For $N = 2$, $P(\xi) = 1 + 2 \sin^2 \pi\xi$, and for $N = 3$,
$P(\xi) = 1 + 3 \sin^2 \pi\xi + 6 \sin^4 \pi\xi$.
\end{definition}

\begin{lemma}[Symmetries of $P$]
\label{lem:daub-fun-symmetry}
\lean{Daubechies.daubFun_periodic, Daubechies.daubFun_neg, Daubechies.daubFun_add_half,
  Daubechies.continuous_daubFun}
\leanok
\uses{def:daub-fun}
For every $N \in \N$ the function $P$ is continuous, and for every $\xi \in \R$
\[
  P(\xi + 1) = P(\xi), \qquad P(-\xi) = P(\xi), \qquad
  P\bigl(\xi + \tfrac12\bigr) = P_N(\cos^2 \pi\xi) .
\]
\end{lemma}

\begin{proof}
\leanok
$\sin(\pi(\xi + 1)) = -\sin \pi\xi$, $\sin(-\pi\xi) = -\sin \pi\xi$ and
$\sin(\pi(\xi + \frac12)) = \cos \pi\xi$, and $P(\xi)$ is a polynomial in $\sin \pi\xi$.
\end{proof}

\begin{proposition}[Bounds for $P$]
\label{prop:daub-fun-bounds}
\lean{Daubechies.daubFun_nonneg, Daubechies.daubFun_le_choose, Daubechies.daubFun_le,
  Daubechies.daubFun_zero, Daubechies.one_le_daubFun}
\leanok
\uses{def:daub-fun}
For every $N \in \N$ and every $\xi \in \R$,
\[
  0 \leq P(\xi) \leq \binom{2N-1}{N} \quad\text{and}\quad
  P(\xi) \leq 1 + \binom{2N-1}{N} \pi^2 \xi^2 .
\]
If $N \geq 1$, then moreover $P(0) = 1$ and $P(\xi) \geq 1$ for every $\xi \in \R$.
\end{proposition}

\begin{proof}
\leanok
\uses{lem:daub-poly-bounds, lem:daub-poly-values}
\Cref{lem:daub-poly-bounds} at $y = \sin^2 \pi\xi \in [0, 1]$, with $P_N(1) = \binom{2N-1}{N}$
and $P_N(0) = 1$ (\Cref{lem:daub-poly-values}), and $\sin^2 \pi\xi \leq \pi^2 \xi^2$.
\end{proof}

\begin{theorem}[The trigonometric B\'ezout identity]
\label{thm:daub-fun-bezout}
\lean{Daubechies.daubFun_bezout}
\leanok
\uses{def:daub-fun}
For every $N \geq 1$ and every $\xi \in \R$,
\[
  \cos(\pi\xi)^{2N} P(\xi) + \sin(\pi\xi)^{2N} P\bigl(\xi + \tfrac12\bigr) = 1 .
\]
\end{theorem}

\begin{proof}
\leanok
\uses{thm:daub-bezout, lem:daub-fun-symmetry}
\Cref{thm:daub-bezout} at $y = \sin^2 \pi\xi$, where $1 - y = \cos^2 \pi\xi$ and
$P(\xi + \frac12) = P_N(1 - y)$ (\Cref{lem:daub-fun-symmetry}).
\end{proof}

The symbol $m_0$ of a mask of order $N$ has $|m_0(\xi)|^2 = \cos(\pi\xi)^{2N} P(\xi)$, and the
identity becomes the quadrature mirror filter identity
\[
  |m_0(\xi)|^2 + \bigl|m_0\bigl(\xi + \tfrac12\bigr)\bigr|^2 = 1 .
\]

\section{Masks and filters}

This section defines the Daubechies masks with their filter coefficients and symbols, and proves
the identities between the filter coefficients on which the scaling function and the wavelet
rest: the quadrature mirror filter identity, the orthogonality of the even shifts, the sum
rules, and their counterparts for the dual filter. The first Daubechies mask, D4, is explicit.

\paragraph{The unit circle.}
A symbol is a polynomial evaluated at the point $e^{-2\pi i \xi}$ of the unit circle; the sign
of the exponent is the one of the Fourier transform.

\begin{lemma}[The exponentials $e^{-2\pi i \xi}$]
\label{lem:daub-exp}
\lean{Daubechies.expFreq, Daubechies.norm_expFreq, Daubechies.expFreq_ne_zero,
  Daubechies.expFreq_eq_exp_mul, Daubechies.expFreq_re, Daubechies.expFreq_im,
  Daubechies.expFreq_zero, Daubechies.expFreq_add, Daubechies.expFreq_pow,
  Daubechies.expFreq_intCast, Daubechies.expFreq_one, Daubechies.expFreq_periodic,
  Daubechies.expFreq_half_add_int, Daubechies.conj_expFreq, Daubechies.contDiff_expFreq,
  Daubechies.continuous_expFreq, Daubechies.norm_expFreq_sub_one_le,
  Daubechies.norm_expFreq_pow_sub_one_le}
\leanok
The function $\xi \mapsto e^{-2\pi i \xi}$ on $\R$ is smooth and has period $1$. For all
$\xi, \eta \in \R$, every $k \in \N$ and every $n \in \Z$:
\begin{enumerate}
  \item $|e^{-2\pi i \xi}| = 1$, $\operatorname{Re} e^{-2\pi i \xi} = \cos 2\pi\xi$ and
    $\operatorname{Im} e^{-2\pi i \xi} = -\sin 2\pi\xi$;
  \item $\overline{e^{-2\pi i \xi}} = e^{-2\pi i (-\xi)}$ and
    $e^{-2\pi i (\xi + \eta)} = e^{-2\pi i \xi} e^{-2\pi i \eta}$;
  \item $(e^{-2\pi i \xi})^k = e^{-2\pi i k \xi}$, $e^{-2\pi i n} = 1$ and
    $e^{-2\pi i (1/2 + n)} = -1$;
  \item $|e^{-2\pi i k \xi} - 1| \leq 2\pi k |\xi|$.
\end{enumerate}
\end{lemma}

\begin{proof}
\leanok
The exponential turns sums into products and commutes with conjugation, $e^{-2\pi i n} = 1$
for $n \in \Z$, and $e^{-\pi i} = -1$. The bound is $|e^{ix} - 1| \leq |x|$ for real $x$, at
$x = -2\pi k \xi$.
\end{proof}

\begin{lemma}[The factor $(1 + z)/2$ on the unit circle]
\label{lem:daub-half-factor}
\lean{Daubechies.one_add_expFreq_div_two, Daubechies.norm_one_add_expFreq_div_two,
  Daubechies.normSq_one_add_expFreq_div_two}
\leanok
For every $\xi \in \R$,
\[
  \frac{1 + e^{-2\pi i \xi}}{2} = e^{-\pi i \xi} \cos \pi\xi ;
\]
in particular $|1 + e^{-2\pi i \xi}|/2 = |\cos \pi\xi|$ and
$|1 + e^{-2\pi i \xi}|^2/4 = \cos^2 \pi\xi$.
\end{lemma}

\begin{proof}
\leanok
\uses{lem:daub-exp}
$2 \cos \pi\xi = e^{\pi i \xi} + e^{-\pi i \xi}$; multiply by $e^{-\pi i \xi}$, which has
modulus $1$ (\Cref{lem:daub-exp}).
\end{proof}

\paragraph{Masks.}
The factor $((1 + z)/2)^N$ of the mask polynomial produces the zeros of the symbol at the
half-integers and the sum rules of the filter coefficients.

\begin{definition}[Daubechies masks]
\label{def:daub-mask}
\lean{Daubechies.Mask, Daubechies.Mask.one_le}
\leanok
\uses{def:daub-fun}
Let $N \in \N$. A \emph{Daubechies mask of order $N$} is a real polynomial $L$ of degree less
than $N$ with $L(1) = 1$ and
\[
  |L(e^{-2\pi i \xi})|^2 = P(\xi) = P_N(\sin^2 \pi\xi) \quad\text{for every } \xi \in \R .
\]
Every Daubechies mask has order $N \geq 1$.
\end{definition}

\begin{definition}[Mask polynomial, filter coefficients and symbol]
\label{def:daub-symbol}
\lean{Daubechies.maskPoly, Daubechies.coeff, Daubechies.coeff_natCast, Daubechies.symbol}
\leanok
\uses{def:daub-mask}
Let $L$ be a Daubechies mask of order $N$. Its \emph{mask polynomial} is
\[
  m(z) = \Bigl(\frac{1 + z}{2}\Bigr)^N L(z) .
\]
Its \emph{filter coefficients} are $c_k = 2 [z^k] m$ for $k \in \N$, where $[z^k] m$ is the
coefficient of $z^k$ in $m$, and $c_k = 0$ for the integers $k < 0$. Its \emph{symbol} is the
function $m_0(\xi) = m(e^{-2\pi i \xi})$ of $\xi \in \R$.
\end{definition}

The filter coefficients, the symbol and everything built from them below belong to a Daubechies
mask $L$ of order $N$, which the notation omits.

\begin{lemma}[The mask polynomial]
\label{lem:daub-mask-poly}
\lean{Daubechies.X_add_one_pow_dvd_maskPoly, Daubechies.natDegree_maskPoly_le,
  Daubechies.natDegree_maskPoly_lt, Daubechies.eval_one_maskPoly, Daubechies.coeff_eq_zero,
  Daubechies.sum_coeff}
\leanok
\uses{def:daub-symbol}
For every Daubechies mask of order $N$ the mask polynomial $m$ is divisible by $(z + 1)^N$, has
degree at most $2N - 1$, and satisfies $m(1) = 1$. The filter coefficients satisfy $c_k = 0$ for
every integer $k < 0$ and every integer $k \geq 2N$, and
\[
  \sum_{k=0}^{2N-1} c_k = 2 .
\]
\end{lemma}

\begin{proof}
\leanok
$m(z) = (z + 1)^N 2^{-N} L(z)$ is the product of $N$ factors of degree $1$ and of $L$, of degree
at most $N - 1$, and $m(1) = L(1) = 1$. The coefficients of $m$ add up to $m(1) = 1$, and
$c_k = 2 [z^k] m$.
\end{proof}

\begin{lemma}[The symbol as a trigonometric polynomial]
\label{lem:daub-symbol-sum}
\lean{Daubechies.symbol_eq_mul, Daubechies.symbol_eq_sum, Daubechies.symbol_zero}
\leanok
\uses{def:daub-symbol}
For every Daubechies mask of order $N$ and every $\xi \in \R$,
\[
  m_0(\xi) = \Bigl(\frac{1 + e^{-2\pi i \xi}}{2}\Bigr)^N L(e^{-2\pi i \xi})
  = \frac12 \sum_{k=0}^{2N-1} c_k e^{-2\pi i k \xi} ,
\]
and $m_0(0) = 1$.
\end{lemma}

\begin{proof}
\leanok
\uses{lem:daub-mask-poly, lem:daub-exp}
By \Cref{lem:daub-mask-poly}, $m(z) = \frac12 \sum_{k < 2N} c_k z^k$. Evaluate at
$z = e^{-2\pi i \xi}$, where $z^k = e^{-2\pi i k \xi}$ (\Cref{lem:daub-exp}). At $\xi = 0$,
$z = 1$ and $m(1) = 1$.
\end{proof}

\begin{proposition}[The modulus of the symbol]
\label{prop:daub-symbol-modulus}
\lean{Daubechies.normSq_symbol}
\leanok
\uses{def:daub-symbol}
For every Daubechies mask of order $N$ and every $\xi \in \R$,
\[
  |m_0(\xi)|^2 = \cos(\pi\xi)^{2N} P(\xi) .
\]
\end{proposition}

\begin{proof}
\leanok
\uses{lem:daub-symbol-sum, lem:daub-half-factor, def:daub-mask}
By \Cref{lem:daub-symbol-sum}, $|m_0(\xi)|^2$ is the product of
$\bigl(|1 + e^{-2\pi i \xi}|^2/4\bigr)^N = \cos(\pi\xi)^{2N}$ (\Cref{lem:daub-half-factor}) and
$|L(e^{-2\pi i \xi})|^2 = P(\xi)$ (\Cref{def:daub-mask}).
\end{proof}

\begin{theorem}[The quadrature mirror filter identity]
\label{thm:daub-qmf}
\lean{Daubechies.qmf, Daubechies.norm_symbol_le_one}
\leanok
\uses{def:daub-symbol}
For every Daubechies mask and every $\xi \in \R$,
\[
  |m_0(\xi)|^2 + \bigl|m_0\bigl(\xi + \tfrac12\bigr)\bigr|^2 = 1 ;
\]
in particular $|m_0(\xi)| \leq 1$.
\end{theorem}

\begin{proof}
\leanok
\uses{prop:daub-symbol-modulus, thm:daub-fun-bezout}
By \Cref{prop:daub-symbol-modulus} at $\xi$ and at $\xi + \frac12$, where
$\cos(\pi\xi + \frac{\pi}{2})^{2N} = \sin(\pi\xi)^{2N}$, the left side is
$\cos(\pi\xi)^{2N} P(\xi) + \sin(\pi\xi)^{2N} P(\xi + \frac12)$. This is $1$ by the
trigonometric B\'ezout identity (\Cref{thm:daub-fun-bezout}), since $N \geq 1$.
\end{proof}

\begin{lemma}[Symmetries and zeros of the symbol]
\label{lem:daub-symbol-props}
\lean{Daubechies.symbol_periodic, Daubechies.symbol_neg, Daubechies.symbol_ne_zero,
  Daubechies.symbol_half_add_int, Daubechies.symbol_half, Daubechies.contDiff_symbol,
  Daubechies.continuous_symbol}
\leanok
\uses{def:daub-symbol}
For every Daubechies mask the symbol $m_0$ is smooth and has period $1$, and
$m_0(-\xi) = \overline{m_0(\xi)}$ for every $\xi \in \R$. Moreover $m_0(\xi) \neq 0$ for every
$\xi \in \R$ with $|\xi| < \frac12$, and $m_0(\frac12 + n) = 0$ for every $n \in \Z$.
\end{lemma}

\begin{proof}
\leanok
\uses{lem:daub-exp, lem:daub-symbol-sum, prop:daub-symbol-modulus, prop:daub-fun-bounds}
The symbol is a polynomial with real coefficients in $e^{-2\pi i \xi}$, which is smooth, has
period $1$ and satisfies $\overline{e^{-2\pi i \xi}} = e^{-2\pi i (-\xi)}$ (\Cref{lem:daub-exp}).
For $|\xi| < \frac12$, $\cos \pi\xi > 0$ and $P(\xi) \geq 1$ (\Cref{prop:daub-fun-bounds}), so
$|m_0(\xi)|^2 > 0$ by \Cref{prop:daub-symbol-modulus}. At $\xi = \frac12 + n$,
$z = e^{-2\pi i \xi} = -1$, where the factor $((1 + z)/2)^N$ of \Cref{lem:daub-symbol-sum}
vanishes since $N \geq 1$.
\end{proof}

\begin{lemma}[The symbol near $0$]
\label{lem:daub-symbol-lipschitz}
\lean{Daubechies.symbolLip, Daubechies.symbolLip_nonneg, Daubechies.norm_symbol_sub_one_le}
\leanok
\uses{def:daub-symbol}
For a Daubechies mask of order $N$ let $C_0 = \pi \sum_{k=0}^{2N-1} k |c_k|$. Then $C_0 \geq 0$,
and for every $\xi \in \R$
\[
  |m_0(\xi) - 1| \leq C_0 |\xi| .
\]
\end{lemma}

\begin{proof}
\leanok
\uses{lem:daub-mask-poly, lem:daub-symbol-sum, lem:daub-exp}
Since $\sum_k c_k = 2$ (\Cref{lem:daub-mask-poly}), \Cref{lem:daub-symbol-sum} gives
$m_0(\xi) - 1 = \frac12 \sum_k c_k (e^{-2\pi i k \xi} - 1)$, and
$|e^{-2\pi i k \xi} - 1| \leq 2\pi k |\xi|$ (\Cref{lem:daub-exp}).
\end{proof}

\paragraph{Sums over $\Z$.}
The identities below are sums over $\Z$ with finitely many nonzero terms (Mathlib's
\texttt{finsum}). Throughout, $\delta_{ab} = 1$ if $a = b$ and $\delta_{ab} = 0$ otherwise.

\begin{lemma}[Sums over $\Z$ by parity]
\label{lem:daub-parity-split}
\lean{Daubechies.finsum_eq_sum_range_of_eq_zero,
  Daubechies.finsum_two_mul_add_finsum_two_mul_add_one, Daubechies.finsum_two_mul_add_add_finsum,
  Daubechies.finsum_two_mul_add_add_two_mul, Daubechies.neg_one_zpow_eq_of_even_sub,
  Daubechies.neg_one_zpow_eq_neg_of_odd_sub, Daubechies.finsum_two_mul_add_eq_half}
\leanok
Let $g$ be a function on $\Z$ with values in a commutative monoid, written additively, and let
$n \in \N$ and $r, q \in \Z$.
\begin{enumerate}
  \item If $g(k) = 0$ for every $k < 0$ and every $k \geq n$, then
    $\sum_{k \in \Z} g(k) = \sum_{k=0}^{n-1} g(k)$.
  \item $\sum_{i \in \Z} g(2i + r + 2q) = \sum_{i \in \Z} g(2i + r)$, and if $g$ has finite
    support, then
    \[
      \sum_{i \in \Z} g(2i + r) + \sum_{i \in \Z} g(2i + r + 1) = \sum_{k \in \Z} g(k) .
    \]
  \item If $g$ is real valued with finite support and $\sum_{k \in \Z} (-1)^k g(k) = 0$, then
    $\sum_{i \in \Z} g(2i + r) = \frac12 \sum_{k \in \Z} g(k)$.
\end{enumerate}
For all integers $a$ and $b$, $(-1)^a = (-1)^b$ if $a - b$ is even, and $(-1)^a = -(-1)^b$ if
$a - b$ is odd.
\end{lemma}

\begin{proof}
\leanok
The even and the odd integers partition $\Z$, and $k \mapsto k + r$ and $i \mapsto i + q$ are
bijections of $\Z$. In item~3 the alternating sum is
$(-1)^r \bigl(\sum_i g(2i + r) - \sum_i g(2i + r + 1)\bigr)$, so the two parity classes carry the
same sum, and by item~2 the two sums add up to $\sum_k g(k)$.
\end{proof}

\begin{lemma}[The filter coefficients as a sequence on $\Z$]
\label{lem:daub-coeff-finsum}
\lean{Daubechies.finite_support_coeff, Daubechies.finsum_eq_sum_range_of_coeff,
  Daubechies.finsum_coeff, Daubechies.symbol_eq_finsum}
\leanok
\uses{def:daub-symbol}
For every Daubechies mask of order $N$ the sequence $(c_k)_{k \in \Z}$ has finite support,
$\sum_{k \in \Z} f(k) = \sum_{k=0}^{2N-1} f(k)$ for every function $f$ on $\Z$ that vanishes
wherever $c_k = 0$, and $\sum_{k \in \Z} c_k = 2$. For every $\xi \in \R$,
\[
  m_0(\xi) = \frac12 \sum_{k \in \Z} c_k e^{-2\pi i k \xi} .
\]
\end{lemma}

\begin{proof}
\leanok
\uses{lem:daub-mask-poly, lem:daub-parity-split, lem:daub-symbol-sum}
$c_k = 0$ for $k < 0$ and for $k \geq 2N$ (\Cref{lem:daub-mask-poly}); apply item~1 of
\Cref{lem:daub-parity-split}, and \Cref{lem:daub-symbol-sum}.
\end{proof}

\paragraph{Orthogonality of the even shifts.}
The coefficients of a trigonometric polynomial are its Fourier coefficients on $[0, 1]$, so
integrating the quadrature mirror filter identity against exponentials gives identities for the
filter coefficients.

\begin{lemma}[Orthonormality of the exponentials]
\label{lem:daub-exp-orthonormal}
\lean{Daubechies.continuous_expFreq_mul, Daubechies.integral_expFreq_intCast_mul,
  Daubechies.integral_sum_mul_conj_sum, Daubechies.integral_finsum_mul_conj_finsum,
  Daubechies.integral_expFreq_mul_normSq_finsum}
\leanok
For every $c \in \R$ the function $\xi \mapsto e^{-2\pi i c \xi}$ is continuous, and for every
$n \in \Z$, $\int_0^1 e^{-2\pi i n \xi} \, d\xi = \delta_{n0}$. For every finite set
$u \subseteq \Z$ and every $a : \Z \to \C$ let
$T_a(\xi) = \sum_{k \in u} a_k e^{-2\pi i k \xi}$. Then for all $a, b : \Z \to \C$
\[
  \int_0^1 T_a(\xi) \, \overline{T_b(\xi)} \, d\xi = \sum_{k \in u} a_k \overline{b_k} ,
\]
and the same holds for the sums over $k \in \Z$ if $a$ and $b$ have finite support. For every
$a : \Z \to \C$ with finite support and every $n \in \Z$,
\[
  \int_0^1 e^{-2\pi i n \xi} \Bigl|\sum_{k \in \Z} a_k e^{-2\pi i k \xi}\Bigr|^2 \, d\xi
  = \sum_{k \in \Z} a_{k-n} \overline{a_k} .
\]
\end{lemma}

\begin{proof}
\leanok
\uses{lem:daub-exp}
For $n \neq 0$, $\int_0^1 e^{-2\pi i n \xi} \, d\xi = (e^{-2\pi i n} - 1)/(-2\pi i n) = 0$.
Expanding the product gives the double sum of the terms
$a_k \overline{b_l} e^{-2\pi i (k - l) \xi}$ (\Cref{lem:daub-exp}), of which only those with
$k = l$ have nonzero integrals. For the last identity, multiplication by $e^{-2\pi i n \xi}$
shifts the coefficients to $a_{k-n}$, and $|z|^2 = z \bar z$.
\end{proof}

\begin{lemma}[Fourier coefficients of $|m_0|^2$]
\label{lem:daub-symbol-fourier}
\lean{Daubechies.norm_half_mul_sq, Daubechies.integral_expFreq_mul_normSq_symbol,
  Daubechies.integral_expFreq_mul_normSq_symbol_add_half}
\leanok
\uses{def:daub-symbol}
For every Daubechies mask and every $n \in \Z$,
\[
  \int_0^1 e^{-2\pi i n \xi} |m_0(\xi)|^2 \, d\xi = \frac14 \sum_{k \in \Z} c_{k-n} c_k
\]
and
\[
  \int_0^1 e^{-2\pi i n \xi} \bigl|m_0\bigl(\xi + \tfrac12\bigr)\bigr|^2 \, d\xi
  = \frac14 e^{\pi i n} \sum_{k \in \Z} c_{k-n} c_k .
\]
\end{lemma}

\begin{proof}
\leanok
\uses{lem:daub-coeff-finsum, lem:daub-exp-orthonormal}
By \Cref{lem:daub-coeff-finsum}, $m_0(\xi) = \frac12 \sum_k c_k e^{-2\pi i k \xi}$ and
$m_0(\xi + \frac12) = \frac12 \sum_k c_k e^{-\pi i k} e^{-2\pi i k \xi}$; apply the last identity
of \Cref{lem:daub-exp-orthonormal}, with $|z/2|^2 = |z|^2/4$. In the second case the factors
combine to $e^{-\pi i (k - n)} \, \overline{e^{-\pi i k}} = e^{\pi i n}$.
\end{proof}

\begin{theorem}[Orthogonality of the even shifts]
\label{thm:daub-even-shifts}
\lean{Daubechies.sum_coeff_mul_coeff_add}
\leanok
\uses{def:daub-symbol}
For every Daubechies mask and every $j \in \Z$,
\[
  \sum_{k \in \Z} c_k c_{k+2j} = 2 \delta_{j0} .
\]
\end{theorem}

\begin{proof}
\leanok
\uses{thm:daub-qmf, lem:daub-symbol-fourier, lem:daub-exp-orthonormal}
Multiply the quadrature mirror filter identity (\Cref{thm:daub-qmf}) by $e^{2\pi i \cdot 2j \xi}$
and integrate over $[0, 1]$. By \Cref{lem:daub-symbol-fourier} with $n = -2j$, where
$e^{\pi i n} = 1$, the left side becomes $\frac12 \sum_k c_k c_{k+2j}$, and the right side is
$\delta_{j0}$ by \Cref{lem:daub-exp-orthonormal}.
\end{proof}

\paragraph{Sum rules.}
The factor $(z + 1)^N$ of the mask polynomial makes the alternating moments of order less than
$N$ of the filter coefficients vanish. For a real polynomial $p$ let $E p = z p'$ be the Euler
operator.

\begin{lemma}[The Euler operator]
\label{lem:daub-euler-operator}
\lean{Daubechies.coeff_X_mul_derivative, Daubechies.coeff_iterate_X_mul_derivative,
  Daubechies.pow_dvd_iterate_X_mul_derivative,
  Daubechies.sum_range_neg_one_pow_mul_pow_mul_coeff}
\leanok
Let $p$ be a real polynomial and $j, k, l, n \in \N$. The coefficient of $z^k$ in $E^l p$ is
$k^l$ times the coefficient of $z^k$ in $p$. If $(z + 1)^j$ divides $p$ and $l \leq j$, then
$(z + 1)^{j-l}$ divides $E^l p$. If $(z + 1)^j$ divides $p$, $l < j$ and $p$ has degree less
than $n$, then
\[
  \sum_{k=0}^{n-1} (-1)^k k^l [z^k] p = 0 .
\]
\end{lemma}

\begin{proof}
\leanok
$E z^k = k z^k$. Differentiation lowers the power of $z + 1$ dividing a polynomial by at most
one, and multiplication by $z$ keeps it. The sum is the value of $E^l p$ at $z = -1$, a zero of
$E^l p$ since $j - l \geq 1$.
\end{proof}

\begin{theorem}[The sum rules]
\label{thm:daub-sum-rules}
\lean{Daubechies.sum_range_rule, Daubechies.sum_rule, Daubechies.sum_rule_poly}
\leanok
\uses{def:daub-symbol}
For every Daubechies mask of order $N$ and every real polynomial $Q$ of degree less than $N$,
\[
  \sum_{k \in \Z} (-1)^k Q(k) \, c_k = 0 .
\]
In particular $\sum_{k \in \Z} (-1)^k k^l c_k = \sum_{k=0}^{2N-1} (-1)^k k^l c_k = 0$ for every
$l < N$.
\end{theorem}

\begin{proof}
\leanok
\uses{lem:daub-euler-operator, lem:daub-mask-poly, lem:daub-coeff-finsum}
\Cref{lem:daub-euler-operator} for $p = m$, $j = N$ and $n = 2N$ (\Cref{lem:daub-mask-poly}),
with $c_k = 2 [z^k] m$, gives the case $Q(k) = k^l$. The sum over $\Z$ is the sum over
$0 \leq k < 2N$ (\Cref{lem:daub-coeff-finsum}), and $Q$ is a linear combination of the monomials
of degree less than $N$.
\end{proof}

\begin{corollary}[Parity sums]
\label{cor:daub-parity-sums}
\lean{Daubechies.finsum_neg_one_zpow_mul_coeff, Daubechies.sum_coeff_two_mul_add,
  Daubechies.sum_coeff_two_mul, Daubechies.sum_coeff_two_mul_add_one}
\leanok
\uses{def:daub-symbol}
For every Daubechies mask and every integer $r$,
\[
  \sum_{k \in \Z} (-1)^k c_k = 0 \quad\text{and}\quad \sum_{i \in \Z} c_{2i+r} = 1 ;
\]
in particular $\sum_{i \in \Z} c_{2i} = \sum_{i \in \Z} c_{2i+1} = 1$.
\end{corollary}

\begin{proof}
\leanok
\uses{thm:daub-sum-rules, lem:daub-parity-split, lem:daub-coeff-finsum}
The first identity is \Cref{thm:daub-sum-rules} for $l = 0 < N$. Item~3 of
\Cref{lem:daub-parity-split} and $\sum_k c_k = 2$ (\Cref{lem:daub-coeff-finsum}) give the second.
\end{proof}

The first identity says $m_0(\frac12) = 0$.

\begin{proposition}[The first moment]
\label{prop:daub-first-moments}
\lean{Daubechies.firstMoment, Daubechies.sum_two_mul_add_mul_coeff}
\leanok
\uses{def:daub-symbol}
For a Daubechies mask let $\mu = \frac12 \sum_{k \in \Z} k c_k$ be its \emph{first moment}. If
the mask has order $N \geq 2$, then for every $r \in \Z$
\[
  \sum_{i \in \Z} (2i + r) \, c_{2i+r} = \mu .
\]
\end{proposition}

\begin{proof}
\leanok
\uses{thm:daub-sum-rules, lem:daub-parity-split, lem:daub-coeff-finsum}
Item~3 of \Cref{lem:daub-parity-split} for $g(k) = k c_k$, which has finite support
(\Cref{lem:daub-coeff-finsum}) and whose alternating sum vanishes by \Cref{thm:daub-sum-rules}
with $l = 1 < N$.
\end{proof}

\begin{lemma}[Derivatives of the symbol]
\label{lem:daub-symbol-derivatives}
\lean{Daubechies.iteratedDeriv_sum_mul_cexp, Daubechies.symbol_eq_sum_cexp,
  Daubechies.iteratedDeriv_symbol}
\leanok
\uses{def:daub-symbol}
For every finite family of pairs $a_i, b_i \in \C$ and every $l \in \N$, the $l$-th derivative
of $t \mapsto \sum_i a_i e^{b_i t}$ on $\R$ is $t \mapsto \sum_i a_i b_i^l e^{b_i t}$. For every
Daubechies mask of order $N$, every $l \in \N$ and every $\xi \in \R$,
\[
  m_0^{(l)}(\xi) = \frac12 \sum_{k=0}^{2N-1} c_k (-2\pi i k)^l e^{-2\pi i k \xi} .
\]
\end{lemma}

\begin{proof}
\leanok
\uses{lem:daub-symbol-sum}
Induction on $l$, with $\frac{d}{dt} e^{bt} = b e^{bt}$, and
$m_0(\xi) = \frac12 \sum_k c_k e^{-2\pi i k \xi}$ (\Cref{lem:daub-symbol-sum}).
\end{proof}

\begin{proposition}[Zeros of the symbol at the half-integers]
\label{prop:daub-symbol-zeros}
\lean{Daubechies.iteratedDeriv_symbol_half_add_int}
\leanok
\uses{def:daub-symbol}
For every Daubechies mask of order $N$, every $l < N$ and every $n \in \Z$,
$m_0^{(l)}(\frac12 + n) = 0$.
\end{proposition}

\begin{proof}
\leanok
\uses{lem:daub-symbol-derivatives, lem:daub-exp, thm:daub-sum-rules}
At $\xi = \frac12 + n$, $e^{-2\pi i k \xi} = (-1)^k$ (\Cref{lem:daub-exp}), so
$m_0^{(l)}(\frac12 + n) = \frac12 (-2\pi i)^l \sum_k (-1)^k k^l c_k$ by
\Cref{lem:daub-symbol-derivatives}, which vanishes by \Cref{thm:daub-sum-rules}.
\end{proof}

\paragraph{The dual filter.}
The dual filter carries the coefficients of the wavelet. Its identities follow from those of
the filter by reflections $k \mapsto a - k$ of the index and a split by parity.

\begin{definition}[The dual filter]
\label{def:daub-dual-filter}
\lean{Daubechies.dualCoeff, Daubechies.dualCoeff_eq_zero, Daubechies.finite_support_dualCoeff,
  Daubechies.finsum_eq_sum_Icc_of_dualCoeff}
\leanok
\uses{def:daub-symbol}
The \emph{dual filter} of a Daubechies mask of order $N$ is $g_k = (-1)^k c_{1-k}$, $k \in \Z$.
It vanishes outside $2 - 2N \leq k \leq 1$, so it has finite support, and
$\sum_{k \in \Z} f(k) = \sum_{k=2-2N}^{1} f(k)$ for every function $f$ on $\Z$ that vanishes
wherever $g_k = 0$.
\end{definition}

\begin{proposition}[Orthogonality of the dual filter]
\label{prop:daub-dual-orthogonality}
\lean{Daubechies.sum_dualCoeff_mul_dualCoeff_add, Daubechies.sum_coeff_mul_dualCoeff_add}
\leanok
\uses{def:daub-dual-filter}
For every Daubechies mask and every $j \in \Z$,
\[
  \sum_{k \in \Z} g_k g_{k+2j} = 2 \delta_{j0} \quad\text{and}\quad
  \sum_{k \in \Z} c_k g_{k+2j} = 0 .
\]
\end{proposition}

\begin{proof}
\leanok
\uses{thm:daub-even-shifts, lem:daub-parity-split}
The substitution $k \mapsto 1 - 2j - k$ turns the first sum into $\sum_k c_k c_{k+2j}$, since
$1 - 2j - k$ and $1 - k$ have the same parity (\Cref{lem:daub-parity-split}); apply
\Cref{thm:daub-even-shifts}. It turns the second sum into its negative, since $1 - k$ and
$k + 2j$ have different parity.
\end{proof}

\begin{lemma}[Reindexing over one parity class]
\label{lem:daub-dual-reindex}
\lean{Daubechies.finsum_coeff_sub_two_mul_mul, Daubechies.finsum_dualCoeff_sub_two_mul_mul}
\leanok
\uses{def:daub-dual-filter}
For every Daubechies mask and all $a, b \in \Z$,
\[
  \sum_{j \in \Z} c_{a-2j} c_{b-2j} = \sum_{i \in \Z} c_{2i+b+(a-b)} c_{2i+b}
\]
and
\[
  \sum_{j \in \Z} g_{a-2j} g_{b-2j}
  = (-1)^{a+b} \sum_{i \in \Z} c_{2i+(1-a)+(a-b)} c_{2i+(1-a)} .
\]
\end{lemma}

\begin{proof}
\leanok
\uses{lem:daub-parity-split}
In the first sum substitute $j = -i$. In the second,
$(-1)^{a-2j} (-1)^{b-2j} = (-1)^{a+b}$ (\Cref{lem:daub-parity-split}),
$1 - (a - 2j) = 2j + (1 - a)$ and $1 - (b - 2j) = 2j + (1 - a) + (a - b)$.
\end{proof}

\begin{theorem}[The filter and its dual resolve the identity]
\label{thm:daub-filter-identity}
\lean{Daubechies.sum_coeff_mul_add_dual}
\leanok
\uses{def:daub-dual-filter}
For every Daubechies mask and all $a, b \in \Z$,
\[
  \sum_{j \in \Z} \bigl(c_{a-2j} c_{b-2j} + g_{a-2j} g_{b-2j}\bigr) = 2 \delta_{ab} .
\]
\end{theorem}

\begin{proof}
\leanok
\uses{lem:daub-dual-reindex, lem:daub-parity-split, thm:daub-even-shifts}
By \Cref{lem:daub-dual-reindex} both sums run over one parity class of
$h(k) = c_{k+(a-b)} c_k$, the first over the class of $b$ and the second over the class of
$1 - a$. If $a - b$ is even, then $(-1)^{a+b} = 1$ and the two classes are complementary, so the
sum is $\sum_k c_k c_{k+(a-b)} = 2 \delta_{ab}$ (\Cref{lem:daub-parity-split},
\Cref{thm:daub-even-shifts}). If $a - b$ is odd, the two classes agree and $(-1)^{a+b} = -1$, so
the two sums cancel.
\end{proof}

\begin{proposition}[Vanishing moments of the dual filter]
\label{prop:daub-dual-moments}
\lean{Daubechies.sum_dualCoeff_mul_pow}
\leanok
\uses{def:daub-dual-filter}
For every Daubechies mask of order $N$ and every $l < N$, $\sum_{k \in \Z} g_k k^l = 0$.
\end{proposition}

\begin{proof}
\leanok
\uses{thm:daub-sum-rules}
The substitution $k \mapsto 1 - k$ turns the sum into $-\sum_k (-1)^k (1 - k)^l c_k$, which
vanishes by \Cref{thm:daub-sum-rules} for $Q(z) = (1 - z)^l$.
\end{proof}

\paragraph{The mask D4.}
For $N = 2$ the polynomial $L$ is linear, with its root $2 + \sqrt 3$ outside the unit disk.

\begin{lemma}[Linear polynomials on the unit circle]
\label{lem:daub-linear-circle}
\lean{Daubechies.normSq_add_mul_expFreq, Daubechies.daubFun_two_eq_two_sub_cos}
\leanok
\uses{def:daub-fun}
For all $a, b, \xi \in \R$,
\[
  |a + b e^{-2\pi i \xi}|^2 = a^2 + b^2 + 2ab \cos 2\pi\xi .
\]
For $N = 2$ and every $\xi \in \R$, $P(\xi) = 2 - \cos 2\pi\xi$.
\end{lemma}

\begin{proof}
\leanok
\uses{lem:daub-exp}
$|a + bz|^2 = a^2 + b^2 |z|^2 + 2ab \operatorname{Re} z$ for real $a$ and $b$, with $|z| = 1$ and
$\operatorname{Re} z = \cos 2\pi\xi$ at $z = e^{-2\pi i \xi}$ (\Cref{lem:daub-exp}). For $N = 2$,
$P(\xi) = 1 + 2 \sin^2 \pi\xi$ (\Cref{def:daub-fun}) and $\cos 2\pi\xi = 1 - 2 \sin^2 \pi\xi$.
\end{proof}

\begin{proposition}[The mask D4]
\label{prop:daub-d4}
\lean{Daubechies.d4, Daubechies.d4_L}
\leanok
\uses{def:daub-mask}
The polynomial
\[
  L(z) = \frac{1 + \sqrt 3}{2} + \frac{1 - \sqrt 3}{2} \, z
\]
is a Daubechies mask of order $2$, the mask \emph{D4}.
\end{proposition}

\begin{proof}
\leanok
\uses{lem:daub-linear-circle}
$L$ has degree $1 < 2$ and $L(1) = 1$. With $a = (1 + \sqrt 3)/2$ and $b = (1 - \sqrt 3)/2$,
$a^2 + b^2 = 2$ and $2ab = -1$, so $|L(e^{-2\pi i \xi})|^2 = 2 - \cos 2\pi\xi = P(\xi)$ by
\Cref{lem:daub-linear-circle}.
\end{proof}

\begin{proposition}[The filter coefficients of D4]
\label{prop:daub-d4-coeff}
\lean{Daubechies.maskPoly_d4, Daubechies.coeff_d4_zero, Daubechies.coeff_d4_one,
  Daubechies.coeff_d4_two, Daubechies.coeff_d4_three, Daubechies.coeff_d4,
  Daubechies.coeff_d4_eq_zero}
\leanok
\uses{prop:daub-d4, def:daub-symbol}
The mask polynomial of D4 is
\[
  m(z) = \frac18 \bigl((1 + \sqrt 3) + (3 + \sqrt 3) z + (3 - \sqrt 3) z^2
    + (1 - \sqrt 3) z^3\bigr) ,
\]
its filter coefficients are
\[
  c_0 = \frac{1 + \sqrt 3}{4} , \quad c_1 = \frac{3 + \sqrt 3}{4} , \quad
  c_2 = \frac{3 - \sqrt 3}{4} , \quad c_3 = \frac{1 - \sqrt 3}{4} ,
\]
and $c_k = 0$ for every other $k \in \Z$.
\end{proposition}

\begin{proof}
\leanok
\uses{lem:daub-mask-poly}
Expand $m(z) = \frac14 (1 + 2z + z^2) L(z)$; two real polynomials agree once they agree at every
real point. Then $c_k = 2 [z^k] m$, and $c_k = 0$ for $k < 0$ and for $k \geq 4$
(\Cref{lem:daub-mask-poly}).
\end{proof}

The orthonormal filter $c_k/\sqrt 2$ is the form in which these coefficients are usually
tabulated.

\section{The symbol product}

The scaling function of a mask is the inverse Fourier transform of the infinite product
$\Phi(\xi) = \prod_{j \geq 0} m_0(\xi/2^{j+1})$. This section proves that the product converges
locally uniformly, satisfies $\Phi(2\xi) = m_0(\xi) \Phi(\xi)$, and has the modulus
$|\Phi(\xi)|^2 = \operatorname{sinc}(\pi\xi)^{2N} \prod_{j \geq 0} P(\xi/2^{j+1})$.

\paragraph{The cosine product.}
Infinite products are understood as in Mathlib: $\prod_{j \geq 0} a_j$ converges if the products
$\prod_{j \in F} a_j$ over the finite sets $F \subseteq \N$ converge as $F$ increases, and then
the partial products $\prod_{j<n} a_j$ converge to the same value. Here
$\operatorname{sinc} x = \sin x / x$ for $x \neq 0$ and $\operatorname{sinc} 0 = 1$. The doubling
formula $\sin 2y = 2 \sin y \cos y$ telescopes into the cosine product, which for $x = \pi/2$ is
Vi\`ete's product $2/\pi = \prod_{j \geq 0} \cos(\pi/2^{j+2})$.

\begin{lemma}[Dyadic scalings]
\label{lem:daub-dyadic}
\lean{Daubechies.div_two_pow_add_succ, Daubechies.sq_div_two_pow_succ,
  Daubechies.hasSum_sq_div_two_pow}
\leanok
For every $x \in \R$ and all $n, j \in \N$, $x/2^{n+j+1} = (x/2^n)/2^{j+1}$ and
$(x/2^{j+1})^2 = \frac{x^2}{4} \, 4^{-j}$, and
\[
  \sum_{j \geq 0} \Bigl(\frac{x}{2^{j+1}}\Bigr)^2 = \frac{x^2}{3} .
\]
\end{lemma}

\begin{proof}
\leanok
A geometric series: $\frac{x^2}{4} \sum_{j \geq 0} 4^{-j} = \frac{x^2}{4} \cdot \frac43$.
\end{proof}

\begin{lemma}[The finite cosine product]
\label{lem:daub-finite-cos-product}
\lean{Daubechies.sinc_mul_cos, Daubechies.sinc_div_two_pow_mul_prod_cos,
  Daubechies.norm_cos_sub_one_le, Daubechies.multipliable_cos_div_two_pow}
\leanok
For every $x \in \R$ and every $n \in \N$,
\[
  \operatorname{sinc}\Bigl(\frac{x}{2^n}\Bigr) \prod_{j<n} \cos\Bigl(\frac{x}{2^{j+1}}\Bigr)
  = \operatorname{sinc} x ,
\]
and the product $\prod_{j \geq 0} \cos(x/2^{j+1})$ converges. For every $y \in \R$,
$\operatorname{sinc}(y) \cos(y) = \operatorname{sinc}(2y)$ and $|\cos y - 1| \leq y^2/2$.
\end{lemma}

\begin{proof}
\leanok
\uses{lem:daub-dyadic}
$\sin 2y = 2 \sin y \cos y$, and $1 - y^2/2 \leq \cos y \leq 1$. Induction on $n$ with
$y = x/2^{n+1}$ gives the identity. The factors are $1 + a_j$ with
$|a_j| \leq (x/2^{j+1})^2/2$, a summable bound (\Cref{lem:daub-dyadic}).
\end{proof}

\begin{theorem}[The cosine product]
\label{thm:daub-cos-product}
\lean{Daubechies.hasProd_cos_div_two_pow, Daubechies.tprod_cos}
\leanok
For every $x \in \R$,
\[
  \prod_{j \geq 0} \cos\Bigl(\frac{x}{2^{j+1}}\Bigr) = \operatorname{sinc} x .
\]
\end{theorem}

\begin{proof}
\leanok
\uses{lem:daub-finite-cos-product}
The partial products converge (\Cref{lem:daub-finite-cos-product}), and
$\operatorname{sinc}(x/2^n) \to \operatorname{sinc} 0 = 1$, so by the identity of
\Cref{lem:daub-finite-cos-product} their limit is $\operatorname{sinc} x$.
\end{proof}

\paragraph{The product of $P$.}
The modulus of $\Phi$ contains the real product $\prod_{j \geq 0} P(\xi/2^{j+1})$; its tail after
the first $n$ factors is bounded uniformly for $|\xi| \leq 2^n/2$.

\begin{lemma}[Convergence of the product of $P$]
\label{lem:daub-fun-product}
\lean{Daubechies.multipliable_daubFun, Daubechies.multipliable_daubFun_shift,
  Daubechies.tprod_daubFun_nonneg, Daubechies.tprod_daubFun_eq_prod_mul}
\leanok
\uses{def:daub-fun}
For every $N \in \N$, every $n \in \N$ and every $\xi \in \R$ the products
$\prod_{j \geq 0} P(\xi/2^{j+1})$ and $\prod_{j \geq 0} P(\xi/2^{n+j+1})$ converge, the first
one is nonnegative, and
\[
  \prod_{j \geq 0} P\Bigl(\frac{\xi}{2^{j+1}}\Bigr)
  = \prod_{j<n} P\Bigl(\frac{\xi}{2^{j+1}}\Bigr) \cdot
    \prod_{j \geq 0} P\Bigl(\frac{\xi}{2^{n+j+1}}\Bigr) .
\]
\end{lemma}

\begin{proof}
\leanok
\uses{prop:daub-fun-bounds, lem:daub-dyadic}
For $N \geq 1$ the factors are $1 + a_j$ with
$0 \leq a_j \leq \binom{2N-1}{N} \pi^2 (\xi/2^{j+1})^2$ (\Cref{prop:daub-fun-bounds}), a
summable bound (\Cref{lem:daub-dyadic}); for $N = 0$ every factor is $0$. The second product is
the first one at the point $\xi/2^n$ (\Cref{lem:daub-dyadic}), and $P \geq 0$.
\end{proof}

\begin{proposition}[Bounds for the product of $P$]
\label{prop:daub-fun-product-bound}
\lean{Daubechies.tprod_daubFun_le, Daubechies.tprod_daubFun_shift_le}
\leanok
\uses{def:daub-fun}
For every $N \in \N$ and every $x \in \R$,
\[
  \prod_{j \geq 0} P\Bigl(\frac{x}{2^{j+1}}\Bigr)
  \leq \exp\Bigl(\binom{2N-1}{N} \frac{\pi^2 x^2}{3}\Bigr) .
\]
For every $n \in \N$ and every $\xi \in \R$ with $|\xi| \leq 2^n/2$,
\[
  \prod_{j \geq 0} P\Bigl(\frac{\xi}{2^{n+j+1}}\Bigr)
  \leq \exp\Bigl(\binom{2N-1}{N} \frac{\pi^2}{12}\Bigr) .
\]
\end{proposition}

\begin{proof}
\leanok
\uses{prop:daub-fun-bounds, lem:daub-dyadic, lem:daub-fun-product}
With $t_j = \binom{2N-1}{N} \pi^2 (x/2^{j+1})^2$, each factor is at most $1 + t_j \leq e^{t_j}$
(\Cref{prop:daub-fun-bounds}), and $\sum_j (x/2^{j+1})^2 = x^2/3$ (\Cref{lem:daub-dyadic}). The
tail is the product at $x = \xi/2^n$ (\Cref{lem:daub-fun-product}), where $|x| \leq \frac12$
gives $x^2/3 \leq 1/12$.
\end{proof}

\paragraph{Convergence.}
Near $0$ the symbol is $1 + O(|\eta|)$, so the factors of $\Phi$ approach $1$ geometrically.

\begin{definition}[The symbol product]
\label{def:daub-symbol-product}
\lean{Daubechies.symbolProduct}
\leanok
\uses{def:daub-symbol}
The \emph{symbol product} of a Daubechies mask is the function
\[
  \Phi(\xi) = \prod_{j \geq 0} m_0\Bigl(\frac{\xi}{2^{j+1}}\Bigr) , \qquad \xi \in \R .
\]
\end{definition}

\begin{lemma}[The factors of the symbol product]
\label{lem:daub-symbol-product-factors}
\lean{Daubechies.norm_symbol_div_two_pow_sub_one_le, Daubechies.summable_norm_symbol_sub_one,
  Daubechies.multipliable_symbolProduct}
\leanok
\uses{def:daub-symbol-product, lem:daub-symbol-lipschitz}
For every Daubechies mask, every $\xi \in \R$ and every $j \in \N$,
\[
  \Bigl|m_0\Bigl(\frac{\xi}{2^{j+1}}\Bigr) - 1\Bigr| \leq \frac{C_0 |\xi|}{2^{j+1}}
\]
with the constant $C_0$ of \Cref{lem:daub-symbol-lipschitz}. Moreover
$\sum_{j \geq 0} |m_0(\xi/2^{j+1}) - 1| < \infty$, and the product $\Phi(\xi)$ converges.
\end{lemma}

\begin{proof}
\leanok
\uses{lem:daub-symbol-lipschitz}
\Cref{lem:daub-symbol-lipschitz} at the point $\xi/2^{j+1}$. The bounds form a geometric
series, and a product $\prod_j (1 + a_j)$ of complex numbers with $\sum_j |a_j| < \infty$
converges.
\end{proof}

\begin{theorem}[Locally uniform convergence]
\label{thm:daub-symbol-product-continuous}
\lean{Daubechies.hasProdUniformlyOn_symbolProduct,
  Daubechies.hasProdLocallyUniformly_symbolProduct,
  Daubechies.tendstoLocallyUniformly_prod_symbol, Daubechies.tendsto_prod_symbol,
  Daubechies.continuous_symbolProduct}
\leanok
\uses{def:daub-symbol-product}
For every Daubechies mask the product $\Phi$ converges uniformly on every compact set
$K \subseteq \R$. The partial products $\prod_{j<n} m_0(\xi/2^{j+1})$ converge to $\Phi(\xi)$ as
$n \to \infty$, for every $\xi \in \R$ and locally uniformly on $\R$, and $\Phi$ is continuous.
\end{theorem}

\begin{proof}
\leanok
\uses{lem:daub-symbol-product-factors, lem:daub-symbol-props}
On a compact set $K \subseteq [-R, R]$ the bound of \Cref{lem:daub-symbol-product-factors} is at
most $C_0 R/2^{j+1}$, uniformly in $\xi$, and the factors are continuous
(\Cref{lem:daub-symbol-props}); Mathlib's uniform version of the convergence criterion for
products $\prod_j (1 + a_j)$ applies. Locally uniform limits of continuous functions are
continuous.
\end{proof}

\paragraph{Algebraic properties.}
The product inherits the bounds and the symmetry of the symbol, and splitting off its first
factor gives the refinement identity.

\begin{proposition}[Values, bounds and symmetry of $\Phi$]
\label{prop:daub-symbol-product-basic}
\lean{Daubechies.symbolProduct_zero, Daubechies.norm_prod_symbol_le_one,
  Daubechies.norm_symbolProduct_le_one, Daubechies.symbolProduct_neg}
\leanok
\uses{def:daub-symbol-product}
For every Daubechies mask, $\Phi(0) = 1$, and for every $\xi \in \R$ and every $n \in \N$,
\[
  \Bigl|\prod_{j<n} m_0\Bigl(\frac{\xi}{2^{j+1}}\Bigr)\Bigr| \leq 1 , \qquad
  |\Phi(\xi)| \leq 1 , \qquad \Phi(-\xi) = \overline{\Phi(\xi)} .
\]
\end{proposition}

\begin{proof}
\leanok
\uses{lem:daub-symbol-sum, thm:daub-qmf, lem:daub-symbol-props,
  lem:daub-symbol-product-factors}
$m_0(0) = 1$ (\Cref{lem:daub-symbol-sum}) and $|m_0| \leq 1$ (\Cref{thm:daub-qmf}), and the
modulus of a convergent product is the product of the moduli
(\Cref{lem:daub-symbol-product-factors}). Conjugation is continuous and multiplicative, and
$m_0(-\eta) = \overline{m_0(\eta)}$ (\Cref{lem:daub-symbol-props}).
\end{proof}

\begin{lemma}[Splitting off factors]
\label{lem:daub-symbol-product-split}
\lean{Daubechies.prod_range_succ_symbol, Daubechies.symbolProduct_eq_prod_mul}
\leanok
\uses{def:daub-symbol-product}
For every Daubechies mask, every $\xi \in \R$ and every $n \in \N$,
\[
  \prod_{j<n+1} m_0\Bigl(\frac{\xi}{2^{j+1}}\Bigr)
  = m_0\Bigl(\frac{\xi}{2}\Bigr) \prod_{j<n} m_0\Bigl(\frac{\xi/2}{2^{j+1}}\Bigr)
\]
and
\[
  \Phi(\xi) = \prod_{j<n} m_0\Bigl(\frac{\xi}{2^{j+1}}\Bigr) \cdot \Phi\Bigl(\frac{\xi}{2^n}\Bigr) .
\]
\end{lemma}

\begin{proof}
\leanok
\uses{lem:daub-dyadic, lem:daub-symbol-product-factors}
$\xi/2^{j+2} = (\xi/2)/2^{j+1}$ and $\xi/2^{n+j+1} = (\xi/2^n)/2^{j+1}$
(\Cref{lem:daub-dyadic}), and a convergent product splits into its first $n$ factors and its
tail, which converges (\Cref{lem:daub-symbol-product-factors}).
\end{proof}

\begin{theorem}[The refinement identity]
\label{thm:daub-refinement-identity}
\lean{Daubechies.symbolProduct_two_mul}
\leanok
\uses{def:daub-symbol-product}
For every Daubechies mask and every $\xi \in \R$, $\Phi(2\xi) = m_0(\xi) \Phi(\xi)$.
\end{theorem}

\begin{proof}
\leanok
\uses{lem:daub-symbol-product-split}
\Cref{lem:daub-symbol-product-split} with $n = 1$ at the point $2\xi$.
\end{proof}

\begin{proposition}[$\Phi$ near the origin]
\label{prop:daub-symbol-product-nonzero}
\lean{Daubechies.symbolProduct_ne_zero, Daubechies.exists_pos_le_norm_symbolProduct}
\leanok
\uses{def:daub-symbol-product}
For every Daubechies mask there is $\kappa > 0$ such that $\Phi(\xi) \neq 0$ and
$|\Phi(\xi)| \geq \kappa$ for every $\xi \in \R$ with $|\xi| \leq \frac12$.
\end{proposition}

\begin{proof}
\leanok
\uses{lem:daub-symbol-props, lem:daub-symbol-product-factors,
  thm:daub-symbol-product-continuous}
For $|\xi| \leq \frac12$ every argument $\xi/2^{j+1}$ lies in $[-\frac14, \frac14]$, where
$m_0 \neq 0$ (\Cref{lem:daub-symbol-props}), and a convergent product of nonzero factors
$1 + a_j$ with $\sum_j |a_j| < \infty$ is nonzero (\Cref{lem:daub-symbol-product-factors}). The
continuous function $|\Phi|$ (\Cref{thm:daub-symbol-product-continuous}) attains its minimum on
the compact interval $[-\frac12, \frac12]$.
\end{proof}

\paragraph{The modulus.}
The squared modulus of each factor is $\cos(\pi\eta)^{2N} P(\eta)$, and the cosines multiply to
the function $\operatorname{sinc}$.

\begin{theorem}[The modulus of the symbol product]
\label{thm:daub-symbol-product-modulus}
\lean{Daubechies.normSq_symbolProduct}
\leanok
\uses{def:daub-symbol-product, def:daub-fun}
For every Daubechies mask of order $N$ and every $\xi \in \R$,
\[
  |\Phi(\xi)|^2
  = \operatorname{sinc}(\pi\xi)^{2N} \prod_{j \geq 0} P\Bigl(\frac{\xi}{2^{j+1}}\Bigr) .
\]
\end{theorem}

\begin{proof}
\leanok
\uses{prop:daub-symbol-modulus, thm:daub-cos-product, lem:daub-fun-product,
  lem:daub-symbol-product-factors}
The factors satisfy $|m_0(\eta)|^2 = \cos(\pi\eta)^{2N} P(\eta)$
(\Cref{prop:daub-symbol-modulus}), the cosines multiply to
$\prod_{j \geq 0} \cos(\pi\xi/2^{j+1}) = \operatorname{sinc}(\pi\xi)$
(\Cref{thm:daub-cos-product}), and the products of $P$ and of the factors of $\Phi$ converge
(\Cref{lem:daub-fun-product,lem:daub-symbol-product-factors}). The factorwise product of two
convergent products converges to the product of their values.
\end{proof}

\begin{proposition}[A bound by finite products]
\label{prop:daub-symbol-product-bound}
\lean{Daubechies.normSq_symbolProduct_le}
\leanok
\uses{def:daub-symbol-product, def:daub-fun}
For every Daubechies mask of order $N$, every $n \in \N$ and every $\xi \in \R$ with
$|\xi| \leq 2^n/2$,
\[
  |\Phi(\xi)|^2 \leq \exp\Bigl(\binom{2N-1}{N} \frac{\pi^2}{12}\Bigr)
  \operatorname{sinc}(\pi\xi)^{2N} \prod_{j<n} P\Bigl(\frac{\xi}{2^{j+1}}\Bigr) .
\]
\end{proposition}

\begin{proof}
\leanok
\uses{thm:daub-symbol-product-modulus, lem:daub-fun-product, prop:daub-fun-product-bound}
Split the product of $P$ in \Cref{thm:daub-symbol-product-modulus} into its first $n$ factors
and its tail (\Cref{lem:daub-fun-product}), and bound the tail by
\Cref{prop:daub-fun-product-bound}.
\end{proof}

The next section bounds the integrals of $\prod_{j<n} P(\xi/2^{j+1})$ through the transfer
operator.

\section{The transfer operator}

The transfer operator of Eirola~\cite{Eirola} and Villemoes~\cite{Villemoes} turns the decay of
the symbol product into the growth of its iterates: a bound $T^n 1 \leq C \lambda^n$ gives
$\int_\R (1 + |\xi|)^{2s} |\Phi(\xi)|^2 \, d\xi < \infty$ for every $s < N - \log_4 \lambda$. For
$N = 2$ and $N = 3$ explicit eigenfunctions give $\lambda = 4$ and $\lambda = 9$, hence the
exponents $1$ and $3 - \log_2 3$.

\paragraph{The operator.}
The points $\xi/2$ and $(\xi + 1)/2$ are the two preimages of $\xi$ under the doubling map
$\eta \mapsto 2\eta$ modulo $1$.

\begin{definition}[The transfer operator]
\label{def:daub-transfer}
\lean{Daubechies.transfer, Daubechies.transfer_add, Daubechies.transfer_smul}
\leanok
\uses{def:daub-fun}
The \emph{transfer operator} $T$ of order $N \in \N$ maps a function $g : \R \to \R$ to
\[
  (T g)(\xi) = P\Bigl(\frac{\xi}{2}\Bigr) g\Bigl(\frac{\xi}{2}\Bigr)
    + P\Bigl(\frac{\xi + 1}{2}\Bigr) g\Bigl(\frac{\xi + 1}{2}\Bigr) , \qquad \xi \in \R .
\]
It is linear.
\end{definition}

\begin{lemma}[Elementary properties of $T$]
\label{lem:daub-transfer-basic}
\lean{Daubechies.transfer_periodic, Daubechies.continuous_transfer, Daubechies.transfer_nonneg,
  Daubechies.transfer_mono, Daubechies.abs_transfer_le}
\leanok
\uses{def:daub-transfer}
Let $N \in \N$, and let $g$, $g_1$ and $g_2$ be real functions on $\R$. If $g$ is continuous,
so is $T g$; if $g$ has period $1$, so has $T g$; if $g \geq 0$, then $T g \geq 0$; and if
$g_1 \leq g_2$, then $T g_1 \leq T g_2$. If $|g| \leq B$ on $\R$, then
$|T g| \leq 2 \binom{2N-1}{N} B$ on $\R$.
\end{lemma}

\begin{proof}
\leanok
\uses{lem:daub-fun-symmetry, prop:daub-fun-bounds}
$P$ is continuous, has period $1$ (\Cref{lem:daub-fun-symmetry}) and satisfies
$0 \leq P \leq \binom{2N-1}{N}$ (\Cref{prop:daub-fun-bounds}). From $\xi$ to $\xi + 1$ the points
$\xi/2$ and $(\xi + 1)/2$ become $(\xi + 1)/2$ and $\xi/2 + 1$, so for $1$-periodic $g$ the two
summands swap.
\end{proof}

\begin{lemma}[The iterates of $T$]
\label{lem:daub-transfer-iterates}
\lean{Daubechies.continuous_iterate_transfer, Daubechies.periodic_iterate_transfer,
  Daubechies.iterate_transfer_nonneg, Daubechies.abs_iterate_transfer_le,
  Daubechies.continuous_iterate_transfer_one, Daubechies.periodic_iterate_transfer_one,
  Daubechies.iterate_transfer_one_nonneg}
\leanok
\uses{def:daub-transfer}
Let $N, n \in \N$ and $g : \R \to \R$. If $g$ is continuous, has period $1$, or is nonnegative,
then so is $T^n g$. If $|g| \leq B$ on $\R$, then $|T^n g| \leq (2 \binom{2N-1}{N})^n B$ on $\R$.
In particular $T^n 1$ is continuous, nonnegative and has period $1$.
\end{lemma}

\begin{proof}
\leanok
\uses{lem:daub-transfer-basic}
Induction on $n$ with \Cref{lem:daub-transfer-basic}.
\end{proof}

\paragraph{The adjoint identity.}
Integrated against a $1$-periodic function $h$, the operator $T$ becomes the multiplication by
$P$ together with the dilation of $h$ by $2$. Iterated, this identifies the integrals of the
finite products of $P$ with the integrals of $T^n 1$.

\begin{lemma}[The two branches of halving]
\label{lem:daub-halving}
\lean{Daubechies.periodic_comp_two_mul,
  Daubechies.integral_comp_div_two_add_comp_add_one_div_two}
\leanok
If $h : \R \to \R$ has period $1$, then so has $\eta \mapsto h(2\eta)$. For every continuous
$f : \R \to \R$ with period $1$,
\[
  \int_{-1/2}^{1/2} \Bigl(f\Bigl(\frac{\eta}{2}\Bigr) + f\Bigl(\frac{\eta + 1}{2}\Bigr)\Bigr)
  \, d\eta = 2 \int_{-1/2}^{1/2} f(\eta) \, d\eta .
\]
\end{lemma}

\begin{proof}
\leanok
The substitutions $u = \eta/2$ and $u = (\eta + 1)/2$ give $2 \int_{-1/4}^{1/4} f$ and
$2 \int_{1/4}^{3/4} f$, together $2 \int_{-1/4}^{3/4} f$, and the integral of a $1$-periodic
function over an interval of length $1$ does not depend on the interval.
\end{proof}

\begin{theorem}[The adjoint identity]
\label{thm:daub-transfer-adjoint}
\lean{Daubechies.integral_transfer}
\leanok
\uses{def:daub-transfer}
For every $N \in \N$ and all continuous $h, g : \R \to \R$ with period $1$,
\[
  \int_{-1/2}^{1/2} h(\eta) \, (T g)(\eta) \, d\eta
  = 2 \int_{-1/2}^{1/2} h(2\eta) P(\eta) g(\eta) \, d\eta .
\]
\end{theorem}

\begin{proof}
\leanok
\uses{lem:daub-halving, lem:daub-fun-symmetry}
The function $f(u) = h(2u) P(u) g(u)$ is continuous and has period $1$
(\Cref{lem:daub-halving,lem:daub-fun-symmetry}), and the integrand on the left is
$f(\eta/2) + f((\eta + 1)/2)$, since $h(\eta + 1) = h(\eta)$. Apply \Cref{lem:daub-halving}.
\end{proof}

\begin{proposition}[The adjoint identity for the iterates]
\label{prop:daub-transfer-iterates-adjoint}
\lean{Daubechies.integral_iterate_transfer, Daubechies.integral_iterate_transfer_one}
\leanok
\uses{def:daub-transfer}
For all $N, n \in \N$ and all continuous $h, g : \R \to \R$ with period $1$,
\[
  \int_{-1/2}^{1/2} h(\eta) \, (T^n g)(\eta) \, d\eta
  = 2^n \int_{-1/2}^{1/2} h(2^n \eta) \prod_{i<n} P(2^i \eta) \, g(\eta) \, d\eta .
\]
In particular
\[
  \int_{-1/2}^{1/2} (T^n 1)(\eta) \, d\eta
  = 2^n \int_{-1/2}^{1/2} \prod_{i<n} P(2^i \eta) \, d\eta .
\]
\end{proposition}

\begin{proof}
\leanok
\uses{thm:daub-transfer-adjoint, lem:daub-transfer-iterates, lem:daub-halving}
Induction on $n$: $T^{n+1} g = T (T^n g)$, \Cref{thm:daub-transfer-adjoint} for the continuous
$1$-periodic function $T^n g$ (\Cref{lem:daub-transfer-iterates}), and the induction hypothesis
for the continuous $1$-periodic function $\eta \mapsto h(2\eta) P(\eta)$ in place of $h$
(\Cref{lem:daub-halving}).
\end{proof}

\begin{theorem}[The finite products and the transfer operator]
\label{thm:daub-transfer-products}
\lean{Daubechies.prod_daubFun_two_pow_mul_div, Daubechies.integral_prod_daubFun}
\leanok
\uses{def:daub-transfer}
For all $N, n \in \N$,
\[
  \int_{-2^n/2}^{2^n/2} \prod_{j<n} P\Bigl(\frac{\xi}{2^{j+1}}\Bigr) \, d\xi
  = \int_{-1/2}^{1/2} (T^n 1)(\eta) \, d\eta .
\]
\end{theorem}

\begin{proof}
\leanok
\uses{prop:daub-transfer-iterates-adjoint}
The substitution $\xi = 2^n \eta$ turns the left side into
\[
  2^n \int_{-1/2}^{1/2} \prod_{j<n} P\Bigl(\frac{2^n \eta}{2^{j+1}}\Bigr) \, d\eta ,
\]
and $\prod_{j<n} P(2^n \eta/2^{j+1}) = \prod_{i<n} P(2^i \eta)$ for every $\eta \in \R$, by the
reindexing $i = n - 1 - j$. Apply \Cref{prop:daub-transfer-iterates-adjoint}.
\end{proof}

\paragraph{Growth bounds.}
Bounds on $T^n 1$ thereby bound the integrals of the finite products of $P$.

\begin{definition}[Growth bounds]
\label{def:daub-transfer-growth}
\lean{Daubechies.TransferGrowth, Daubechies.TransferGrowth.one_le}
\leanok
\uses{def:daub-transfer}
Let $N \in \N$. A \emph{growth bound} for the transfer operator of order $N$ is a pair of real
numbers $\lambda > 0$ and $C$ with
\[
  (T^n 1)(\xi) \leq C \lambda^n \quad\text{for every } n \in \N \text{ and every } \xi \in \R .
\]
Every growth bound has $C \geq 1$.
\end{definition}

\begin{lemma}[Integrals under a growth bound]
\label{lem:daub-growth-integrals}
\lean{Daubechies.prod_daubFun_nonneg, Daubechies.continuous_prod_daubFun,
  Daubechies.TransferGrowth.integral_iterate_transfer_one_le,
  Daubechies.TransferGrowth.integral_prod_daubFun_le,
  Daubechies.TransferGrowth.setIntegral_prod_daubFun_le}
\leanok
\uses{def:daub-transfer-growth}
Let $N, n \in \N$. The function $\xi \mapsto \prod_{j<n} P(\xi/2^{j+1})$ is continuous and
nonnegative. If $\lambda, C$ is a growth bound for the transfer operator of order $N$, then for
every set $S \subseteq [-2^n/2, 2^n/2]$
\[
  \int_{-1/2}^{1/2} (T^n 1)(\eta) \, d\eta \leq C \lambda^n
  \quad\text{and}\quad
  \int_S \prod_{j<n} P\Bigl(\frac{\xi}{2^{j+1}}\Bigr) \, d\xi \leq C \lambda^n ,
\]
in particular for $S = [-2^n/2, 2^n/2]$.
\end{lemma}

\begin{proof}
\leanok
\uses{thm:daub-transfer-products, lem:daub-fun-symmetry, prop:daub-fun-bounds}
$P$ is continuous and nonnegative (\Cref{lem:daub-fun-symmetry,prop:daub-fun-bounds}). The
first bound integrates $T^n 1 \leq C \lambda^n$ over an interval of length $1$. The integrand of
the second is nonnegative, so the integral over $S$ is at most the integral over
$[-2^n/2, 2^n/2]$, which equals the first integral by \Cref{thm:daub-transfer-products}.
\end{proof}

\begin{proposition}[The crude growth bound]
\label{prop:daub-crude-growth}
\lean{Daubechies.transferGrowth_crude, Daubechies.logb_four_two_mul_choose_lt}
\leanok
\uses{def:daub-transfer-growth}
For every $N \in \N$, $\lambda = 2 \binom{2N-1}{N}$ and $C = 1$ form a growth bound for the
transfer operator of order $N$. For every $N \geq 2$,
$\log_4 \bigl(2 \binom{2N-1}{N}\bigr) < N - \frac12$.
\end{proposition}

\begin{proof}
\leanok
\uses{lem:daub-transfer-iterates, lem:daub-choose-bound}
\Cref{lem:daub-transfer-iterates} with $g = 1$ and $B = 1$. For $N \geq 2$,
$2 \binom{2N-1}{N} < 2 \cdot 4^{N-1} = 4^{N - 1/2}$ by \Cref{lem:daub-choose-bound}.
\end{proof}

\paragraph{Decay of the symbol product.}
Under a growth bound the integrals of $|\Phi|^2$ over the dyadic annuli decay geometrically,
since $\operatorname{sinc}(\pi\xi)^{2N}$ is at most $(\pi |\xi|)^{-2N}$.

\begin{definition}[Dyadic annuli]
\label{def:daub-annulus}
\lean{Daubechies.annulus, Daubechies.measurableSet_annulus, Daubechies.annulus_subset_Icc,
  Daubechies.exists_mem_annulus, Daubechies.univ_subset_Icc_union_annulus}
\leanok
For $n \in \N$ the \emph{dyadic annulus} is
\[
  A_n = \{\xi \in \R : 2^n/4 < |\xi| \leq 2^n/2\} .
\]
It is measurable and contained in $[-2^n/2, 2^n/2]$. Every $\xi \in \R$ with $|\xi| > \frac14$
lies in some $A_n$, so $\R = [-\frac14, \frac14] \cup \bigcup_{n \in \N} A_n$.
\end{definition}

\begin{lemma}[Integrability from the annuli]
\label{lem:daub-annulus-integrable}
\lean{Daubechies.integrable_of_integral_annulus_le}
\leanok
\uses{def:daub-annulus}
Let $f$ be a continuous function from $\R$ to a normed abelian group, and let $A \in \R$ and
$0 \leq r < 1$ with $\int_{A_n} |f| \leq A r^n$ for every $n \in \N$. Then $f$ is integrable on
$\R$.
\end{lemma}

\begin{proof}
\leanok
On $[-\frac14, \frac14]$ the function $f$ is integrable as a continuous function on a compact
interval, and on $\bigcup_n A_n$ since the integrals over the annuli are summable; these two
sets cover $\R$ (\Cref{def:daub-annulus}).
\end{proof}

\begin{lemma}[$\Phi$ on the annuli]
\label{lem:daub-sinc-annulus}
\lean{Daubechies.abs_sinc_le_inv_abs, Daubechies.sinc_pow_le_of_mem_annulus,
  Daubechies.normSq_symbolProduct_le_of_mem_annulus}
\leanok
\uses{def:daub-annulus, def:daub-symbol-product}
For every $x \in \R$ with $x \neq 0$, $|\operatorname{sinc} x| \leq 1/|x|$. For all
$N, n \in \N$ and every $\xi \in A_n$,
$\operatorname{sinc}(\pi\xi)^{2N} \leq (4/(\pi 2^n))^{2N}$, and for every Daubechies mask of
order $N$
\[
  |\Phi(\xi)|^2 \leq \exp\Bigl(\binom{2N-1}{N} \frac{\pi^2}{12}\Bigr)
  \Bigl(\frac{4}{\pi 2^n}\Bigr)^{2N} \prod_{j<n} P\Bigl(\frac{\xi}{2^{j+1}}\Bigr) .
\]
\end{lemma}

\begin{proof}
\leanok
\uses{prop:daub-symbol-product-bound, lem:daub-growth-integrals}
$|\sin x| \leq 1$, and $\pi |\xi| > \pi 2^n/4$ on $A_n$. Since $|\xi| \leq 2^n/2$ on $A_n$,
\Cref{prop:daub-symbol-product-bound} applies, and the finite product is nonnegative
(\Cref{lem:daub-growth-integrals}).
\end{proof}

\begin{proposition}[$|\Phi|^2$ on the annuli]
\label{prop:daub-annulus-integral}
\lean{Daubechies.integral_annulus_normSq_le}
\leanok
\uses{def:daub-annulus, def:daub-symbol-product, def:daub-transfer-growth}
Let $\lambda, C$ be a growth bound for the transfer operator of order $N$. Then for every
Daubechies mask of order $N$ and every $n \in \N$,
\[
  \int_{A_n} |\Phi(\xi)|^2 \, d\xi
  \leq \exp\Bigl(\binom{2N-1}{N} \frac{\pi^2}{12}\Bigr)
  \Bigl(\frac{4}{\pi 2^n}\Bigr)^{2N} C \lambda^n .
\]
\end{proposition}

\begin{proof}
\leanok
\uses{lem:daub-sinc-annulus, lem:daub-growth-integrals, thm:daub-symbol-product-continuous}
Both sides of the bound of \Cref{lem:daub-sinc-annulus} are continuous
(\Cref{thm:daub-symbol-product-continuous,lem:daub-growth-integrals}); integrate it over
$A_n \subseteq [-2^n/2, 2^n/2]$ and apply \Cref{lem:daub-growth-integrals}.
\end{proof}

Since $(4/(\pi 2^n))^{2N} = (4/\pi)^{2N} 4^{-Nn}$, these bounds decay geometrically in $n$, with
the ratio $\lambda/4^N$.

\begin{lemma}[Weights on the annuli]
\label{lem:daub-weights}
\lean{Daubechies.rpow_le_four_rpow_abs_mul, Daubechies.one_add_abs_rpow_le_of_mem_annulus,
  Daubechies.four_div_pi_mul_two_pow_pow, Daubechies.two_rpow_two_mul_mul_lt}
\leanok
\uses{def:daub-annulus}
The following hold.
\begin{enumerate}
  \item If $x, y \in \R$ satisfy $y > 0$ and $y/4 \leq x \leq 4y$, then $x^t \leq 4^{|t|} y^t$ for
    every $t \in \R$.
  \item For every $t \in \R$, every $n \in \N$ and every $\xi \in A_n$,
    $(1 + |\xi|)^t \leq 4^{|t|} (2^t)^n$.
  \item For all $N, n \in \N$, $(4/(\pi 2^n))^{2N} = (4/\pi)^{2N}/(4^N)^n$.
  \item If $N \in \N$ and $\lambda, s \in \R$ satisfy $\lambda > 0$ and
    $s < N - \log_4 \lambda$, then $2^{2s} \lambda < 4^N$.
\end{enumerate}
\end{lemma}

\begin{proof}
\leanok
Item~1 follows from $x \leq 4y$ for $t \geq 0$ and from $y/4 \leq x$ for $t \leq 0$. Item~2 is
item~1 with $y = 2^n$, since $2^n/4 \leq 1 + |\xi| \leq 4 \cdot 2^n$ on $A_n$. Item~3 is
$(2^n)^{2N} = (4^N)^n$, and item~4 says $\lambda < 4^{N-s} = 4^N/2^{2s}$.
\end{proof}

\begin{lemma}[Splitting a product of powers]
\label{lem:daub-weight-split}
\lean{Daubechies.rpow_mul_le_rpow_add_rpow_mul_sq}
\leanok
For every $x > 0$, every $c \geq 0$ and all $a, b \in \R$,
$x^a c \leq x^{2(a-b)} + x^{2b} c^2$.
\end{lemma}

\begin{proof}
\leanok
With $u = x^{a-b}$ and $v = x^b c$, $x^a c = uv \leq 2uv \leq u^2 + v^2$.
\end{proof}

\begin{theorem}[Weighted $L_2$ decay of the symbol product]
\label{thm:daub-decay-l2}
\lean{Daubechies.integrable_rpow_mul_normSq}
\leanok
\uses{def:daub-symbol-product, def:daub-transfer-growth}
Let $\lambda, C$ be a growth bound for the transfer operator of order $N$. Then for every
Daubechies mask of order $N$ and every $s \in \R$ with $s < N - \log_4 \lambda$, the function
$\xi \mapsto (1 + |\xi|)^{2s} |\Phi(\xi)|^2$ is integrable on $\R$.
\end{theorem}

\begin{proof}
\leanok
\uses{lem:daub-weights, prop:daub-annulus-integral, lem:daub-annulus-integrable,
  thm:daub-symbol-product-continuous}
On $A_n$ the weight is at most $4^{|2s|} (2^{2s})^n$ (\Cref{lem:daub-weights}), so by
\Cref{prop:daub-annulus-integral} the integral over $A_n$ is at most a constant times
$(2^{2s} \lambda/4^N)^n$, and $2^{2s} \lambda/4^N < 1$ (\Cref{lem:daub-weights}). The function is
continuous (\Cref{thm:daub-symbol-product-continuous}); apply \Cref{lem:daub-annulus-integrable}.
\end{proof}

\begin{theorem}[Weighted $L_1$ decay of the symbol product]
\label{thm:daub-decay-l1}
\lean{Daubechies.integrable_rpow_mul_norm}
\leanok
\uses{def:daub-symbol-product, def:daub-transfer-growth}
Let $\lambda, C$ be a growth bound for the transfer operator of order $N$. Then for every
Daubechies mask of order $N$ and every $a \in \R$ with $a < N - \log_4 \lambda - \frac12$, the
function $\xi \mapsto (1 + |\xi|)^a |\Phi(\xi)|$ is integrable on $\R$.
\end{theorem}

\begin{proof}
\leanok
\uses{lem:daub-weight-split, thm:daub-decay-l2, thm:daub-symbol-product-continuous}
Choose $b$ with $a + \frac12 < b < N - \log_4 \lambda$. By \Cref{lem:daub-weight-split},
\[
  (1 + |\xi|)^a |\Phi(\xi)| \leq (1 + |\xi|)^{-2(b - a)} + (1 + |\xi|)^{2b} |\Phi(\xi)|^2 ,
\]
where the first summand is integrable since $2(b - a) > 1$, and the second by
\Cref{thm:daub-decay-l2}. The function on the left is continuous
(\Cref{thm:daub-symbol-product-continuous}).
\end{proof}

\begin{corollary}[The symbol product is integrable and square integrable]
\label{cor:daub-symbol-product-integrable}
\lean{Daubechies.integrable_symbolProduct, Daubechies.memLp_symbolProduct_two}
\leanok
\uses{def:daub-symbol-product}
For every Daubechies mask of order $N \geq 2$ the symbol product $\Phi$ is integrable and square
integrable on $\R$.
\end{corollary}

\begin{proof}
\leanok
\uses{prop:daub-crude-growth, thm:daub-decay-l1, thm:daub-decay-l2}
The crude growth bound has $\log_4 \lambda < N - \frac12$ (\Cref{prop:daub-crude-growth}); take
$a = 0$ in \Cref{thm:daub-decay-l1} and $s = 0$ in \Cref{thm:daub-decay-l2}.
\end{proof}

\paragraph{Orders two and three.}
For $N = 2$ and $N = 3$ the transfer operator has eigenfunctions among the trigonometric
polynomials, and $T^n 1$ is known exactly. In the variable $y = \sin^2(\pi\xi/2)$ every
eigenvalue equation becomes a polynomial identity in $y$.

\begin{lemma}[Linearity of the iterates and eigenfunctions]
\label{lem:daub-transfer-eigen}
\lean{Daubechies.iterate_transfer_add, Daubechies.iterate_transfer_smul,
  Daubechies.iterate_transfer_eq_pow_smul}
\leanok
\uses{def:daub-transfer}
For all $N, n \in \N$, all $g, g_1, g_2 : \R \to \R$ and every $c \in \R$,
$T^n (g_1 + g_2) = T^n g_1 + T^n g_2$ and $T^n (c g) = c \, T^n g$. If $e : \R \to \R$ and
$\mu \in \R$ satisfy $T e = \mu e$, then $T^n e = \mu^n e$ for every $n \in \N$.
\end{lemma}

\begin{proof}
\leanok
Induction on $n$, since $T$ is linear (\Cref{def:daub-transfer}).
\end{proof}

\begin{lemma}[The variable $y$]
\label{lem:daub-transfer-sin-sq}
\lean{Daubechies.transfer_eq_eval_sin_sq, Daubechies.cos_two_pi_mul_div_two,
  Daubechies.cos_two_pi_mul_add_one_div_two, Daubechies.cos_two_pi_mul_eq_sin_sq}
\leanok
\uses{def:daub-transfer, def:daub-poly}
Let $\xi \in \R$ and $y = \sin^2(\pi\xi/2)$. For every $N \in \N$ and every $g : \R \to \R$,
\[
  (T g)(\xi) = P_N(y) \, g\Bigl(\frac{\xi}{2}\Bigr)
    + P_N(1 - y) \, g\Bigl(\frac{\xi + 1}{2}\Bigr) ,
\]
and $\cos \pi\xi = 1 - 2y$, $\cos(\pi(\xi + 1)) = 2y - 1$ and $\cos 2\pi\xi = 1 - 8y(1 - y)$.
\end{lemma}

\begin{proof}
\leanok
\uses{lem:daub-fun-symmetry}
$P(\xi/2) = P_N(y)$ and
$P((\xi + 1)/2) = P(\xi/2 + \frac12) = P_N(\cos^2(\pi\xi/2)) = P_N(1 - y)$
(\Cref{lem:daub-fun-symmetry}). The double angle formulas give $\cos \pi\xi = 1 - 2y$ and
$\cos 2\pi\xi = 1 - 2 \sin^2 \pi\xi = 1 - 8y(1 - y)$, and $\cos(\pi\xi + \pi) = -\cos \pi\xi$.
\end{proof}

\begin{proposition}[Order two]
\label{prop:daub-transfer-two}
\lean{Daubechies.transfer_two_one, Daubechies.iterate_transfer_two_one,
  Daubechies.transferGrowth_two, Daubechies.two_sub_logb_four_four}
\leanok
\uses{def:daub-transfer-growth}
For $N = 2$, $T 1 = 4$ and $T^n 1 = 4^n$ for every $n \in \N$. The numbers $\lambda = 4$ and
$C = 1$ form a growth bound for the transfer operator of order $2$, and $2 - \log_4 4 = 1$.
\end{proposition}

\begin{proof}
\leanok
\uses{lem:daub-transfer-sin-sq, lem:daub-transfer-eigen, def:daub-poly}
$P_2(y) = 1 + 2y$ (\Cref{def:daub-poly}), so $(T 1)(\xi) = (1 + 2y) + (1 + 2(1 - y)) = 4$ by
\Cref{lem:daub-transfer-sin-sq}. Then \Cref{lem:daub-transfer-eigen} with $e = 1$ and $\mu = 4$.
\end{proof}

\begin{proposition}[Eigenfunctions of order three]
\label{prop:daub-transfer-three-eigen}
\lean{Daubechies.transfer_three_eigen_one, Daubechies.transfer_three_eigen_two}
\leanok
\uses{def:daub-transfer}
Let $N = 3$ and $c(\xi) = \cos 2\pi\xi$. Then
\[
  T (9 + c) = 9 \, (9 + c) \quad\text{and}\quad T (1 + 3c) = -4 \, (1 + 3c) .
\]
\end{proposition}

\begin{proof}
\leanok
\uses{lem:daub-transfer-sin-sq, def:daub-poly}
$P_3(y) = 1 + 3y + 6y^2$ (\Cref{def:daub-poly}). By \Cref{lem:daub-transfer-sin-sq}, at $\xi$
both sides of the first identity equal $90 - 72 y(1 - y)$ and both sides of the second equal
$-16 + 96 y(1 - y)$, where $y = \sin^2(\pi\xi/2)$.
\end{proof}

\begin{proposition}[Order three]
\label{prop:daub-transfer-three}
\lean{Daubechies.iterate_transfer_three_one, Daubechies.iterate_transfer_three_one_le,
  Daubechies.transferGrowth_three, Daubechies.three_sub_logb_four_nine}
\leanok
\uses{def:daub-transfer-growth}
For $N = 3$, every $n \in \N$ and every $\xi \in \R$,
\[
  (T^n 1)(\xi) = \frac{3}{26} \, 9^n (9 + \cos 2\pi\xi)
    - \frac{1}{26} \, (-4)^n (1 + 3 \cos 2\pi\xi) ,
\]
and $(T^n 1)(\xi) \leq \frac{17}{13} \, 9^n$. The numbers $\lambda = 9$ and $C = 2$ form a
growth bound for the transfer operator of order $3$, and $3 - \log_4 9 = 3 - \log_2 3$.
\end{proposition}

\begin{proof}
\leanok
\uses{prop:daub-transfer-three-eigen, lem:daub-transfer-eigen}
$1 = \frac{3}{26} (9 + c) - \frac{1}{26} (1 + 3c)$; apply \Cref{lem:daub-transfer-eigen} with the
eigenfunctions of \Cref{prop:daub-transfer-three-eigen}. Then $9 + \cos 2\pi\xi \leq 10$ and
$|(-4)^n (1 + 3 \cos 2\pi\xi)| \leq 4 \cdot 9^n$ give
$T^n 1 \leq (\frac{30}{26} + \frac{4}{26}) \, 9^n$. Finally
$\log_4 9 = \log 3^2 / \log 2^2 = \log_2 3$.
\end{proof}

By \Cref{thm:daub-decay-l2}, $\int_\R (1 + |\xi|)^{2s} |\Phi(\xi)|^2 \, d\xi < \infty$ for every
Daubechies mask of order $2$ and every $s < 1$, and for every Daubechies mask of order $3$ and
every $s < 3 - \log_2 3 \approx 1.415$.

\section{Fourier transforms of integrable functions}

This section is general analysis on $\R$, used for the scaling function
$\varphi = \mathcal F^{-1} \Phi$ in the following sections. The Fourier transform of
$f : \R \to \C$ is
\[
  \widehat f(\xi) = \mathcal F f(\xi) = \int_\R f(x) e^{-2\pi i x \xi} \, dx ,
\]
and its inverse transform is $\mathcal F^{-1} f(x) = \int_\R f(\xi) e^{2\pi i x \xi} \, d\xi$.
Lean reads the integral of a function that is not integrable as $0$; with this convention the
statements below without integrability hypotheses hold for every $f$.

\begin{lemma}[Continuity and the uniform bound]
\label{lem:daub-fourier-continuous}
\lean{Real.continuous_fourier_of_integrable, Real.continuous_fourierInv_of_integrable,
  Real.norm_fourier_le_integral_norm, Real.norm_fourierInv_le_integral_norm}
\leanok
For every integrable $f : \R \to \C$ the functions $\widehat f$ and $\mathcal F^{-1} f$ are
continuous. For every $f : \R \to \C$ and all $\xi, x \in \R$,
\[
  |\widehat f(\xi)| \leq \int_\R |f| \quad\text{and}\quad |\mathcal F^{-1} f(x)| \leq \int_\R |f| .
\]
\end{lemma}

\begin{proof}
\leanok
$|e^{-2\pi i x \xi}| = 1$, and dominated convergence with the dominating function $|f|$; the
inverse transform is $\mathcal F^{-1} f(x) = \widehat f(-x)$.
\end{proof}

\begin{lemma}[Linearity]
\label{lem:daub-fourier-linear}
\lean{Real.fourier_const_mul_apply, Real.fourierInv_const_mul_apply, Real.fourier_sum_apply,
  Real.fourierInv_sum_apply}
\leanok
For every $f : \R \to \C$ and every $c \in \C$, $\mathcal F(c f) = c \widehat f$ and
$\mathcal F^{-1}(c f) = c \mathcal F^{-1} f$. For every finite family $(F_i)_{i \in S}$ of
integrable functions $F_i : \R \to \C$,
\[
  \mathcal F\Bigl(\sum_{i \in S} F_i\Bigr) = \sum_{i \in S} \widehat{F_i}
  \quad\text{and}\quad
  \mathcal F^{-1}\Bigl(\sum_{i \in S} F_i\Bigr) = \sum_{i \in S} \mathcal F^{-1} F_i .
\]
\end{lemma}

\begin{proof}
\leanok
Linearity of the integral.
\end{proof}

\begin{lemma}[Translation and modulation]
\label{lem:daub-fourier-translate}
\lean{Real.fourier_comp_add_right_apply, Real.fourier_comp_sub_right_apply,
  Real.fourierInv_comp_add_right_apply, Real.fourier_exp_mul_apply,
  Real.fourierInv_exp_mul_apply}
\leanok
For $b \in \R$ let $e_b(x) = e^{2\pi i b x}$. For every $f : \R \to \C$ and all
$a, b, \xi, x \in \R$:
\begin{enumerate}
  \item $\mathcal F[f(\cdot + a)](\xi) = e^{2\pi i a \xi} \widehat f(\xi)$ and
    $\mathcal F[f(\cdot - a)](\xi) = e^{-2\pi i a \xi} \widehat f(\xi)$;
  \item $\mathcal F^{-1}[f(\cdot + a)](x) = e^{-2\pi i a x} \mathcal F^{-1} f(x)$;
  \item $\mathcal F(e_b f)(\xi) = \widehat f(\xi - b)$ and
    $\mathcal F^{-1}(e_b f)(x) = \mathcal F^{-1} f(x + b)$.
\end{enumerate}
\end{lemma}

\begin{proof}
\leanok
The substitution $y = x + a$, and $e^{-2\pi i x \xi} e^{2\pi i b x} = e^{-2\pi i x (\xi - b)}$.
\end{proof}

\begin{lemma}[Dilation]
\label{lem:daub-fourier-dilate}
\lean{Real.fourier_comp_mul_left_apply_of_ne_zero, Real.fourier_comp_mul_left_apply,
  Real.fourierInv_comp_mul_left_apply, Real.fourier_comp_mul_sub_apply}
\leanok
For every $f : \R \to \C$ and all $b, \xi, x \in \R$:
\begin{enumerate}
  \item $\mathcal F[f(a\,\cdot)](\xi) = |a|^{-1} \widehat f(\xi / a)$ for every $a \neq 0$;
  \item $\mathcal F^{-1}[f(a\,\cdot)](x) = a^{-1} \mathcal F^{-1} f(x / a)$ for every $a > 0$;
  \item $\widehat g(\xi) = a^{-1} e^{-2\pi i b \xi / a} \widehat f(\xi / a)$ for every $a > 0$,
    where $g(x) = f(ax - b)$.
\end{enumerate}
\end{lemma}

\begin{proof}
\leanok
\uses{lem:daub-fourier-translate}
The substitution $y = ax$, with $dx = |a|^{-1} \, dy$. The third identity is the first one for
$f(\cdot - b)$, together with \Cref{lem:daub-fourier-translate}.
\end{proof}

The third identity is the form met by the cascade algorithm
$\eta_{n+1}(x) = \sum_k c_k \eta_n(2x - k)$.

\begin{lemma}[Conjugation]
\label{lem:daub-fourier-conj}
\lean{Real.conj_fourier, Real.conj_fourierInv}
\leanok
For every $f : \R \to \C$ and all $\xi, x \in \R$, with $\bar f(x) = \overline{f(x)}$,
\[
  \overline{\widehat f(\xi)} = \mathcal F^{-1} \bar f(\xi) \quad\text{and}\quad
  \overline{\mathcal F^{-1} f(x)} = \mathcal F \bar f(x) .
\]
\end{lemma}

\begin{proof}
\leanok
$\overline{e^{-2\pi i x \xi}} = e^{2\pi i x \xi}$, and conjugation commutes with the integral.
\end{proof}

\paragraph{The multiplication formula and Parseval's identity.}
The transform moves from one factor of a product to the other. Together with Fourier inversion
this gives Parseval's identity for integrable functions with integrable transforms.

\begin{theorem}[The multiplication formula]
\label{thm:daub-multiplication-formula}
\lean{Real.integral_fourier_mul, Real.integral_fourierInv_mul}
\leanok
For all integrable $f, g : \R \to \C$,
\[
  \int_\R \widehat f(\xi) \, g(\xi) \, d\xi = \int_\R f(x) \, \widehat g(x) \, dx ,
\]
and the same holds with $\mathcal F^{-1}$ in place of $\mathcal F$:
\[
  \int_\R \mathcal F^{-1} f(\xi) \, g(\xi) \, d\xi = \int_\R f(x) \, \mathcal F^{-1} g(x) \, dx .
\]
\end{theorem}

\begin{proof}
\leanok
By Fubini's theorem both sides of the first identity equal the absolutely convergent double
integral $\int_\R \int_\R e^{-2\pi i x \xi} f(x) g(\xi) \, dx \, d\xi$. The second identity is
the same with $e^{2\pi i x \xi}$.
\end{proof}

\begin{theorem}[Parseval's identity]
\label{thm:daub-parseval}
\lean{Real.integral_mul_conj_eq_integral_fourier, Real.integral_norm_sq_fourier}
\leanok
Let $f : \R \to \C$ be integrable, and let $g : \R \to \C$ be continuous and integrable with
integrable $\widehat g$. Then
\[
  \int_\R f(x) \, \overline{g(x)} \, dx = \int_\R \widehat f(\xi) \, \overline{\widehat g(\xi)}
  \, d\xi .
\]
In particular $\int_\R |\widehat f|^2 = \int_\R |f|^2$ for every continuous integrable
$f : \R \to \C$ with integrable $\widehat f$.
\end{theorem}

\begin{proof}
\leanok
\uses{thm:daub-multiplication-formula, lem:daub-fourier-conj}
By \Cref{lem:daub-fourier-conj} and \Cref{thm:daub-multiplication-formula} for $\mathcal F^{-1}$,
\[
  \int_\R \widehat g \, \overline{\widehat f} = \int_\R \widehat g \, \mathcal F^{-1} \bar f
  = \int_\R \mathcal F^{-1} \widehat g \, \bar f ,
\]
and $\mathcal F^{-1} \widehat g = g$ by the Fourier inversion theorem, which holds for
continuous integrable $g$ with integrable $\widehat g$. Conjugation gives the identity. For
$g = f$ it reads $\int |f|^2 = \int |\widehat f|^2$, since $z \bar z = |z|^2$.
\end{proof}

\begin{lemma}[Testing against Schwartz functions]
\label{lem:daub-fourier-schwartz}
\lean{Real.integral_mul_schwartz_eq}
\leanok
For every integrable $f : \R \to \C$ and every Schwartz function $\psi : \R \to \C$,
\[
  \int_\R f(x) \, \psi(x) \, dx = \int_\R \widehat f(\xi) \, \mathcal F^{-1} \psi(\xi) \, d\xi .
\]
\end{lemma}

\begin{proof}
\leanok
\uses{thm:daub-multiplication-formula}
$\mathcal F^{-1} \psi$ is a Schwartz function, hence integrable, and
$\mathcal F \mathcal F^{-1} \psi = \psi$ by Fourier inversion on the Schwartz space. Apply
\Cref{thm:daub-multiplication-formula} to $f$ and $\mathcal F^{-1} \psi$.
\end{proof}

\paragraph{The fundamental lemma.}
A function whose integrals against all test functions supported in an open set vanish is zero
almost everywhere on that set, and everywhere on it if it is continuous.

\begin{proposition}[The fundamental lemma of the calculus of variations]
\label{prop:daub-fundamental-lemma}
\lean{Real.ae_eq_zero_of_integral_mul_schwartz_eq_zero,
  Real.eqOn_zero_of_integral_mul_schwartz_eq_zero}
\leanok
Let $U \subseteq \R$ be open, and let $f : \R \to \C$ satisfy $\int_\R \psi f = 0$ for every
Schwartz function $\psi : \R \to \C$ whose support is compact and contained in $U$.
\begin{enumerate}
  \item If $f$ is locally integrable on $U$, then $f(x) = 0$ for almost every $x \in U$.
  \item If $f$ is continuous, then $f(x) = 0$ for every $x \in U$.
\end{enumerate}
\end{proposition}

\begin{proof}
\leanok
Mathlib's fundamental lemma tests with real smooth functions $g$ with compact support in $U$,
and $x \mapsto g(x)$, read as a complex function, is a Schwartz function. Two continuous
functions that agree almost everywhere on an open set agree on it, since nonempty open sets
have positive Lebesgue measure.
\end{proof}

\section{Masks of every order}

The masks of every order come from the Fej\'er--Riesz theorem, as in~\cite{Daubechies}: a real
polynomial $Q$ with $Q > 0$ on $[-1, 1]$ satisfies $Q(\cos\theta) = |L(e^{i\theta})|^2$ with a
real polynomial $L$. For $Q(t) = P_N((1 - t)/2)$, where $Q(\cos 2\pi\xi) = P_N(\sin^2 \pi\xi)$,
this gives the polynomial $L$ of a mask of order $N$. The proof factors $Q$ over $\C$ and
splits each factor $\cos\theta - t$ with the inverse of the Joukowsky map.

\paragraph{The inverse of the Joukowsky map.}
The Joukowsky map $a \mapsto (a + a^{-1})/2$ takes $e^{i\varphi}$ to $\cos\varphi$, and for
$a \neq 0$ the equation $(a + a^{-1})/2 = t$ reads $a^2 - 2ta + 1 = 0$.

\begin{lemma}[Roots of $a^2 - 2ta + 1$]
\label{lem:daub-joukowsky-equation}
\lean{Complex.joukowsky_eq_re_of_norm_eq_one, Complex.joukowsky_root_ne_zero,
  Complex.joukowsky_root_conj}
\leanok
Let $a, t \in \C$ with $a^2 - 2ta + 1 = 0$. Then $a \neq 0$ and
$\bar a^2 - 2 \bar t \bar a + 1 = 0$. If $|a| = 1$, then $t = \operatorname{Re} a$; in
particular $t$ is a real number in $[-1, 1]$.
\end{lemma}

\begin{proof}
\leanok
The constant term of the equation is $1$, and conjugation is a ring homomorphism of $\C$. If
$|a| = 1$, then $a^{-1} = \bar a$, and the equation divided by $a$ says
$2t = a + a^{-1} = a + \bar a = 2 \operatorname{Re} a$.
\end{proof}

\begin{lemma}[The root in the disk]
\label{lem:daub-joukowsky-root}
\lean{Complex.exists_joukowsky_root, Complex.joukowsky_root_unique}
\leanok
For every $t \in \C \setminus [-1, 1]$ the equation $a^2 - 2ta + 1 = 0$ has a root $a$ with
$|a| < 1$. For every $t \in \C$ it has at most one root $a$ with $|a| < 1$.
\end{lemma}

\begin{proof}
\leanok
\uses{lem:daub-joukowsky-equation}
The two roots $t \pm s$, $s^2 = t^2 - 1$, have product $1$, and for $t \notin [-1, 1]$ neither
lies on the unit circle (\Cref{lem:daub-joukowsky-equation}), so one of them lies in the open
unit disk. Two different roots are the two roots of the equation and have product $1$, so they
cannot both lie in the open unit disk.
\end{proof}

\begin{definition}[The inverse of the Joukowsky map]
\label{def:daub-joukowsky-inv}
\lean{Complex.joukowskyInv, Complex.joukowskyInv_spec, Complex.joukowskyInv_eq}
\leanok
\uses{lem:daub-joukowsky-root}
For $t \in \C$ let $a(t)$ be the root $a$ of $a^2 - 2ta + 1 = 0$ with $|a| < 1$ if there is
one, and $a(t) = 0$ otherwise. If the equation has a root $a$ with $|a| < 1$, then $a = a(t)$;
in particular $|a(t)| < 1$ and $a(t)^2 - 2t \, a(t) + 1 = 0$. This is the case for every
$t \in \C \setminus [-1, 1]$; on this set $t \mapsto a(t)$ is the inverse of the Joukowsky map
on the punctured open unit disk.
\end{definition}

For $t \in [-1, 1]$ the two roots $t \pm i \sqrt{1 - t^2}$ lie on the unit circle, and
$a(t) = 0$.

\begin{lemma}[The inverse of the Joukowsky map and conjugation]
\label{lem:daub-joukowsky-conj}
\lean{Complex.joukowskyInv_conj}
\leanok
\uses{def:daub-joukowsky-inv}
For every $t \in \C$, $a(\bar t) = \overline{a(t)}$.
\end{lemma}

\begin{proof}
\leanok
\uses{lem:daub-joukowsky-equation, lem:daub-joukowsky-root}
Conjugation maps the roots of $a^2 - 2ta + 1$ in the open unit disk to the roots of
$a^2 - 2\bar t a + 1$ in the open unit disk (\Cref{lem:daub-joukowsky-equation}), and such a root
is unique (\Cref{lem:daub-joukowsky-root}). If there is no such root for $t$, there is none
for $\bar t$, and both sides are $0$.
\end{proof}

\begin{lemma}[The factorisation on the unit circle]
\label{lem:daub-joukowsky-factor}
\lean{Complex.joukowsky_two_mul_mul_re_sub, Complex.joukowsky_two_mul_norm_mul_norm_re_sub}
\leanok
Let $a, t, z \in \C$ with $a^2 - 2ta + 1 = 0$. If $z \bar z = 1$, then
\[
  2a (\operatorname{Re} z - t) = -(1 - az)(1 - a \bar z) ,
\]
and hence $2 |a| \, |\operatorname{Re} z - t| = |1 - az| \, |1 - \bar a z|$ if $|z| = 1$.
\end{lemma}

\begin{proof}
\leanok
Expand the right side with $z \bar z = 1$, $z + \bar z = 2 \operatorname{Re} z$ and
$a^2 + 1 = 2ta$. For the moduli, $|1 - a \bar z| = |\overline{1 - a \bar z}| = |1 - \bar a z|$.
\end{proof}

For $z = e^{i\theta}$ the left side is $2 |a| \, |\cos\theta - t|$: on the unit circle the factor
$\cos\theta - t$ splits into two linear factors in $e^{i\theta}$.

\paragraph{The complex roots of a real polynomial.}
For a real polynomial $Q$, products $\prod_t$ run over the complex roots $t$ of $Q$, counted
with multiplicity.

\begin{lemma}[Real polynomials over $\C$]
\label{lem:daub-real-poly-roots}
\lean{Multiset.norm_prod, Complex.aeval_ofReal, Polynomial.map_conj_aroots,
  Polynomial.map_comp_conj_aroots, Polynomial.ofReal_ne_of_mem_aroots,
  Polynomial.aeval_eq_mul_prod_aroots, Polynomial.eval_eq_abs_leadingCoeff_mul_prod_norm}
\leanok
Let $Q$ be a real polynomial with leading coefficient $c$.
\begin{enumerate}
  \item The conjugates $\bar t$ of the complex roots $t$ of $Q$ are again the complex roots of
    $Q$, with the same multiplicities; so for every function $g$ on $\C$ the values $g(\bar t)$
    and $g(t)$ over the roots agree as multisets.
  \item $Q(z) = c \prod_t (z - t)$ for every $z \in \C$.
  \item $Q(x) = |c| \prod_t |x - t|$ for every $x \in \R$ with $Q(x) \geq 0$.
  \item No $r \in \R$ with $Q(r) \neq 0$ is a complex root of $Q$.
\end{enumerate}
\end{lemma}

\begin{proof}
\leanok
$\C$ is algebraically closed, so $Q$ has $\deg Q$ complex roots, counted with multiplicity, and
item~2 holds; the coefficients of $Q$ are real, so $\bar t$ is a root of the same multiplicity
as $t$. Item~3 is the modulus of item~2, since $Q(x) = |Q(x)|$: the modulus of a product is the
product of the moduli, and a real polynomial takes real values at real points. Item~4 holds
since $Q(t) = 0$ at every root $t$.
\end{proof}

\paragraph{The Fej\'er--Riesz factor.}
For $Q > 0$ on $[-1, 1]$ each root $t$ of $Q$ contributes the factor $1 - a(t) z$, whose zero
$1/a(t)$ lies outside the closed unit disk.

\begin{definition}[The Fej\'er--Riesz factor]
\label{def:daub-fejer-riesz-factor}
\lean{Polynomial.fejerRieszFactor, Polynomial.eval_fejerRieszFactor,
  Polynomial.natDegree_fejerRieszFactor_le}
\leanok
\uses{def:daub-joukowsky-inv}
The \emph{Fej\'er--Riesz factor} of a real polynomial $Q$ is the complex polynomial
\[
  L_0(z) = \prod_t \bigl(1 - a(t) z\bigr) ,
\]
the product over the complex roots $t$ of $Q$, counted with multiplicity. Its degree is at most
$\deg Q$.
\end{definition}

\begin{lemma}[The Fej\'er--Riesz factor is real]
\label{lem:daub-fejer-riesz-factor-real}
\lean{Polynomial.map_conj_fejerRieszFactor, Polynomial.exists_map_eq_of_map_conj_eq}
\leanok
\uses{def:daub-fejer-riesz-factor}
A complex polynomial whose coefficients are fixed by conjugation has real coefficients. For
every real polynomial $Q$ the coefficients of $L_0$ are fixed by conjugation, so $L_0$ is a real
polynomial.
\end{lemma}

\begin{proof}
\leanok
\uses{lem:daub-joukowsky-conj, lem:daub-real-poly-roots}
A complex number fixed by conjugation equals its real part. Conjugating the coefficients of
$L_0$ replaces $a(t)$ by $\overline{a(t)} = a(\bar t)$ (\Cref{lem:daub-joukowsky-conj}), and the
roots of $Q$ are invariant under conjugation (\Cref{lem:daub-real-poly-roots}).
\end{proof}

\begin{lemma}[Zeros of the Fej\'er--Riesz factor]
\label{lem:daub-fejer-riesz-factor-zeros}
\lean{Polynomial.joukowskyInv_spec_of_mem_aroots, Polynomial.eval_fejerRieszFactor_ne_zero}
\leanok
\uses{def:daub-fejer-riesz-factor}
Let $Q$ be a real polynomial with $Q(r) > 0$ for every $r \in [-1, 1]$. Then every complex root
$t$ of $Q$ satisfies $|a(t)| < 1$ and $a(t)^2 - 2t \, a(t) + 1 = 0$, and $L_0(z) \neq 0$ for
every $z \in \C$ with $|z| \leq 1$.
\end{lemma}

\begin{proof}
\leanok
\uses{lem:daub-real-poly-roots, lem:daub-joukowsky-root}
No root of $Q$ lies in $[-1, 1]$ (\Cref{lem:daub-real-poly-roots}), so the equation for $a(t)$
has a root in the open unit disk (\Cref{lem:daub-joukowsky-root}). For $|z| \leq 1$ each factor
$1 - a(t) z$ of $L_0(z)$ is nonzero, since $|a(t) z| < 1$.
\end{proof}

\begin{lemma}[The Fej\'er--Riesz factor on the unit circle]
\label{lem:daub-fejer-riesz-factor-modulus}
\lean{Polynomial.prod_two_mul_norm_mul_prod_norm_eq}
\leanok
\uses{def:daub-fejer-riesz-factor}
Let $Q$ be a real polynomial with $Q(r) > 0$ for every $r \in [-1, 1]$. Then for every
$\theta \in \R$
\[
  \prod_t 2 |a(t)| \cdot \prod_t |\cos\theta - t| = |L_0(e^{i\theta})|^2 .
\]
\end{lemma}

\begin{proof}
\leanok
\uses{lem:daub-joukowsky-factor, lem:daub-fejer-riesz-factor-zeros, lem:daub-joukowsky-conj,
  lem:daub-real-poly-roots}
Let $z = e^{i\theta}$, so $|z| = 1$ and $\operatorname{Re} z = \cos\theta$. For each root $t$,
\Cref{lem:daub-joukowsky-factor,lem:daub-fejer-riesz-factor-zeros} give
\[
  2 |a(t)| \, |\cos\theta - t| = |1 - a(t) z| \, \bigl|1 - \overline{a(t)} z\bigr| .
\]
Since $\overline{a(t)} = a(\bar t)$ (\Cref{lem:daub-joukowsky-conj}) and the roots are invariant
under conjugation (\Cref{lem:daub-real-poly-roots}),
$\prod_t |1 - \overline{a(t)} z| = \prod_t |1 - a(t) z| = |L_0(z)|$.
\end{proof}

\paragraph{The Fej\'er--Riesz theorem.}
On the unit circle $|L_0|^2$ is $Q(\cos\theta)$ up to a positive constant; normalising the
constant and the sign gives the theorem.

\begin{lemma}[Fej\'er--Riesz up to a constant]
\label{lem:daub-fejer-riesz-constant}
\lean{Polynomial.exists_eval_cos_eq_mul_norm_sq}
\leanok
Let $Q$ be a real polynomial with $Q(r) > 0$ for every $r \in [-1, 1]$. Then there are $K > 0$
and a real polynomial $L$ with $\deg L \leq \deg Q$ and $L(z) \neq 0$ for every $z \in \C$ with
$|z| \leq 1$ such that $Q(\cos\theta) = K \, |L(e^{i\theta})|^2$ for every $\theta \in \R$.
\end{lemma}

\begin{proof}
\leanok
\uses{def:daub-fejer-riesz-factor, lem:daub-fejer-riesz-factor-real,
  lem:daub-fejer-riesz-factor-zeros, lem:daub-fejer-riesz-factor-modulus,
  lem:daub-real-poly-roots, lem:daub-joukowsky-equation}
Let $L$ be $L_0$ read as a real polynomial (\Cref{lem:daub-fejer-riesz-factor-real}); its degree
is at most $\deg Q$ (\Cref{def:daub-fejer-riesz-factor}), and it has no zeros in the closed unit
disk (\Cref{lem:daub-fejer-riesz-factor-zeros}). Since $Q(0) > 0$, $Q \neq 0$; let $c \neq 0$ be its
leading coefficient and $K = |c| / \prod_t 2 |a(t)|$, which is positive since $a(t) \neq 0$
(\Cref{lem:daub-joukowsky-equation}). At $x = \cos\theta \in [-1, 1]$, $Q(x) > 0$, so
\[
  Q(x) = |c| \prod_t |x - t| = K \, |L_0(e^{i\theta})|^2
\]
by \Cref{lem:daub-real-poly-roots} and \Cref{lem:daub-fejer-riesz-factor-modulus}.
\end{proof}

\begin{lemma}[Fixing the sign]
\label{lem:daub-fejer-riesz-sign}
\lean{Polynomial.exists_eval_one_pos_of_aeval_ne_zero}
\leanok
Let $L$ be a real polynomial with $L(z) \neq 0$ for every $z \in \C$ with $|z| \leq 1$. Then
there is a real polynomial $L'$ with $\deg L' = \deg L$ and $L'(1) > 0$, with $L'(z) \neq 0$ for
every $z \in \C$ with $|z| \leq 1$, and with $|L'(z)| = |L(z)|$ for every $z \in \C$.
\end{lemma}

\begin{proof}
\leanok
$L(1) \neq 0$; take $L' = L$ if $L(1) > 0$, and $L' = -L$ otherwise.
\end{proof}

\begin{theorem}[The Fej\'er--Riesz theorem]
\label{thm:daub-fejer-riesz}
\lean{Polynomial.fejer_riesz_outer, Polynomial.fejer_riesz}
\leanok
Let $Q$ be a real polynomial with $Q(t) > 0$ for every $t \in [-1, 1]$. Then there is a real
polynomial $L$ with $\deg L \leq \deg Q$, $L(1) > 0$ and $L(z) \neq 0$ for every $z \in \C$ with
$|z| \leq 1$, such that
\[
  |L(e^{i\theta})|^2 = Q(\cos\theta) \quad\text{for every } \theta \in \R .
\]
\end{theorem}

\begin{proof}
\leanok
\uses{lem:daub-fejer-riesz-constant, lem:daub-fejer-riesz-sign}
Take $K$ and $L$ from \Cref{lem:daub-fejer-riesz-constant}, and apply
\Cref{lem:daub-fejer-riesz-sign} to $\sqrt K L$; then
$|\sqrt K L(e^{i\theta})|^2 = K \, |L(e^{i\theta})|^2 = Q(\cos\theta)$.
\end{proof}

\begin{lemma}[Two readings of the identity]
\label{lem:daub-fejer-riesz-readings}
\lean{Polynomial.forall_norm_aeval_exp_neg_sq_iff, Polynomial.eval_one_sq_of_norm_aeval_exp_sq}
\leanok
Let $Q$ and $L$ be real polynomials. Then $|L(e^{-i\theta})|^2 = Q(\cos\theta)$ for every
$\theta \in \R$ if and only if $|L(e^{i\theta})|^2 = Q(\cos\theta)$ for every $\theta \in \R$,
and then $L(1)^2 = Q(1)$.
\end{lemma}

\begin{proof}
\leanok
Replace $\theta$ by $-\theta$, and use $\cos(-\theta) = \cos\theta$. At $\theta = 0$ the identity
reads $|L(1)|^2 = Q(1)$, and $L(1)$ is real.
\end{proof}

\begin{corollary}[The Fej\'er--Riesz theorem with $e^{-i\theta}$]
\label{cor:daub-fejer-riesz-exp-neg}
\lean{Polynomial.fejer_riesz_exp_neg}
\leanok
Let $Q$ be a real polynomial with $Q(t) > 0$ for every $t \in [-1, 1]$. Then there is a real
polynomial $L$ with $\deg L \leq \deg Q$ and $L(1) > 0$ such that
\[
  |L(e^{-i\theta})|^2 = Q(\cos\theta) \quad\text{for every } \theta \in \R .
\]
\end{corollary}

\begin{proof}
\leanok
\uses{thm:daub-fejer-riesz, lem:daub-fejer-riesz-readings}
\Cref{thm:daub-fejer-riesz} and \Cref{lem:daub-fejer-riesz-readings}.
\end{proof}

With $\theta = 2\pi\xi$ this is the form $|L(e^{-2\pi i \xi})|^2$ of the masks. For example,
$Q(t) = 2 - t = P_2((1 - t)/2)$ has the single root $t = 2$, with $a(2) = 2 - \sqrt 3$, and the
construction gives $L(z) = \frac12 \bigl((1 + \sqrt 3) + (1 - \sqrt 3) z\bigr)$, the polynomial of
the first Daubechies wavelet, $N = 2$.


\chapter{Dominating mixed smoothness on \texorpdfstring{$\R^d$}{R\^{}d}}
\label{ch:mixed-smoothness}

This chapter works on $\R^d$ with Lebesgue measure, for functions of $d$ real variables
$x = (x_1, \dots, x_d)$. The spaces of dominating mixed smoothness measure smoothness in each
coordinate separately: their quasi-norms are $\ell_q$ norms over $k \in \N^d$ with the weights
$2^{s|k|_1}$, $|k|_1 = k_1 + \dots + k_d$. All kernels are tensor products of the one-dimensional
kernels of \Cref{ch:fourier-rd} and \Cref{ch:rd} with $V = \R$, and every convolution is a
$d$-dimensional convolution, so Young's inequality and the approximate identities apply to
them. Three
descriptions are shown to agree, mixed differences after T.~Ullrich \cite{UllrichT}, mixed local
means, and tensor Littlewood--Paley blocks after Vyb\'{\i}ral \cite{Vybiral}; the last gives the
rate of approximation on hyperbolic crosses.

\section{Tensor kernels}

\begin{definition}[Tensor kernels]
\label{def:tensor-kernel}
\lean{Rd.tensorKernel, Rd.tensorKernel_apply, Rd.integrable_tensorKernel,
  Rd.continuous_tensorKernel, Rd.aestronglyMeasurable_tensorKernel}
\leanok
For one-dimensional kernels $K_1, \dots, K_d$ the tensor product is
$(\otimes K)(x) = \prod_i K_i(x_i)$. Tensor products of integrable (continuous, measurable)
kernels are integrable (continuous, measurable).
\end{definition}

\begin{lemma}[Integrals and convolutions of tensor kernels]
\label{lem:tensor-kernel-fubini}
\lean{Rd.integral_tensorKernel, Rd.lintegral_enorm_tensorKernel, Rd.eLpNorm_tensorKernel_one,
  Rd.lintegral_fin_nat_prod_volume_eq_prod, Rd.eLpNorm_tensorKernel,
  Rd.conv_tensorKernel, Rd.dilate_tensorKernel, Rd.sum_tensorKernel_piFinset}
\leanok
\uses{def:tensor-kernel, thm:young}
$\int \otimes K = \prod_i \int K_i$, $\|\otimes K\|_1 = \prod_i \|K_i\|_1$,
$(\otimes K) * (\otimes K') = \otimes (K_i * K'_i)$, $(\otimes K)_c = \otimes (K_i)_c$, and
$\sum_{k \in T_1 \times \dots \times T_d} \otimes (K_i(k_i)) = \otimes \bigl(\sum_{j \in T_i} K_i(j)\bigr)$
for finite sets $T_i$. For measurable $f_1, \dots, f_d : \R \to [0, \infty]$,
$\int_{\R^d} \prod_i f_i(x_i) \, dx = \prod_i \int_\R f_i$, and for $0 < p < \infty$ and
measurable $K_1, \dots, K_d$, $\|\otimes K\|_p = \prod_i \|K_i\|_p$.
\end{lemma}

\begin{proof}
\leanok
Fubini's theorem for product measures (Tonelli's theorem for nonnegative functions, by
induction on $d$), applied to $|K_i|^p$ for the $L_p$ norms, and $\prod_i \sum_j = \sum \prod$
for finite sums.
\end{proof}

\begin{lemma}[Diagonal changes of variables]
\label{lem:diagonal-substitution}
\lean{Rd.diagMap, Rd.diagMap_apply, Rd.det_diagMap, Rd.diagEquiv, Rd.diagEquiv_apply,
  Rd.integral_comp_diagMap, Rd.integral_tensorKernel_mul_comp_eq_conv, Rd.integrable_comp_diag}
\leanok
\uses{lem:tensor-kernel-fubini}
For $a_1, \dots, a_d > 0$, $\int F((a_i z_i)_i) \, dz = (\prod_i a_i)^{-1} \int F$, the map
$F \mapsto F((a_i \cdot_i)_i)$ keeps integrability, and
\[
  \int (\otimes\varphi)(z) f(x - (a_i z_i)_i) \, dz = ((\otimes \varphi_i^{(1/a_i)}) * f)(x),
  \qquad \varphi^{(c)}(t) = c \, \varphi(ct).
\]
\end{lemma}

\begin{proof}
\leanok
The diagonal map scales Lebesgue measure by its determinant $\prod_i a_i$.
\end{proof}

\section{Sequences on \texorpdfstring{$\N^d$}{N\^{}d}}

\begin{definition}[$\ell_q$ norms on $\N^d$]
\label{def:mixed-lq}
\lean{Rd.mixedWeight, Rd.mixedWeight_ne_top, Rd.mixedWeight_neg_mul_self}
\leanok
\uses{def:lq-norm}
For $a : \N^d \to [0, \infty]$, $\|a\|_{\ell_q} = (\sum_k a_k^q)^{1/q}$ for $q < \infty$ and
$\|a\|_{\ell_\infty} = \sup_k a_k$, the norm of \Cref{def:lq-norm} on the countable set $\N^d$.
The mixed weights are $2^{s|k|_1} = \prod_i 2^{s k_i}$.
\end{definition}

\begin{lemma}[The Schur test]
\label{lem:schur-test}
\lean{Rd.lqNorm_schur_le, Rd.tsum_prod_le_prod_tsum}
\leanok
\uses{def:mixed-lq, lem:weighted-jensen}
Let $w \geq 0$ on $\N^d \times \N^d$ with $\sum_k w(m, k) \leq A$ for every $m$ and
$\sum_m w(m, k) \leq A$ for every $k$. Then for $1 \leq q \leq \infty$ and every $b \geq 0$,
$\|(\sum_k w(m, k) b_k)_m\|_{\ell_q} \leq A \|b\|_{\ell_q}$. Moreover
$\sum_{k \in \N^d} \prod_i f_i(k_i) \leq \prod_i \sum_j f_i(j)$ for $f_i \geq 0$.
\end{lemma}

\begin{proof}
\leanok
By the weighted Jensen inequality (\Cref{lem:weighted-jensen}) for the counting measure,
$(\sum_k w(m, k) b_k)^q \leq A^{q-1} \sum_k w(m, k) b_k^q$; summing over $m$ and exchanging the sums
gives at most $A^q \sum_k b_k^q$. The product bound holds on the finite boxes.
\end{proof}

\begin{lemma}[The almost diagonal lemma on $\N^d$]
\label{lem:mixed-almost-diagonal}
\lean{Rd.almostDiagCoeff, Rd.almostDiagProfile, Rd.dyadicWeight_mul_almostDiagCoeff_mul,
  Rd.tsum_almostDiagProfile_le, Rd.lqNorm_almostDiagonal_pi_le,
  Rd.lqNorm_le_of_le_almostDiagonal_pi}
\leanok
\uses{lem:schur-test, lem:almost-diagonal}
Let $c(i, j) = 2^{-(i - j)L}$ for $j \leq i$ and $c(i, j) = 1$ for $j > i$, let $s \geq 0$ and
$1 \leq q \leq \infty$, and let $A_{s,L}$ be the constant of \Cref{lem:almost-diagonal}. Then
for every $a \geq 0$
\[
  \bigl\|\bigl(2^{s|m|_1} \textstyle\sum_k \prod_i c(m_i, k_i) \, a_k\bigr)_m\bigr\|_{\ell_q}
  \leq A_{s,L}^d \, \|(2^{s|k|_1} a_k)_k\|_{\ell_q} .
\]
Consequently, for every $C_0 \in [0, \infty]$ and every $a' \geq 0$ on $\N^d$ with
$a'_m \leq C_0^d \sum_k \prod_i c(m_i, k_i) \, a_k$ for every $m$,
\[
  \|(2^{s|m|_1} a'_m)_m\|_{\ell_q} \leq (C_0 A_{s,L})^d \, \|(2^{s|k|_1} a_k)_k\|_{\ell_q} .
\]
\end{lemma}

\begin{proof}
\leanok
$2^{s|m|_1} \prod_i c(m_i, k_i) \, 2^{-s|k|_1} = \prod_i u(m_i - k_i)$ with $u(l) = 2^{-(L-s)l}$
for $l \geq 0$ and $u(l) = 2^{sl}$ for $l < 0$, and $\sum_{l \in \Z} u(l) \leq A_{s,L}$. All row
and column sums are therefore at most $A_{s,L}^d$ (\Cref{lem:schur-test}). This replaces the
one-dimensional Hardy inequalities of \Cref{lem:almost-diagonal}, which do not iterate over the
coordinates.
\end{proof}

\section{Mixed differences}

\begin{definition}[Mixed differences]
\label{def:mixed-differences}
\lean{Rd.mixedDiff, Rd.diffAct_coordElem_pow_eq_higherDiff, Rd.mixedDiff_empty, Rd.mixedDiff_insert,
  Rd.mixedDiff_add, Rd.mixedDiff_sub}
\leanok
\uses{def:translation-algebra}
For $e \subseteq \{1, \dots, d\}$ and steps $h \in \R^d$ the mixed difference of order $L$ is
$\Delta^{L,e}_h f = (\prod_{i \in e} \Delta^L_{h_i e_i}) f$, the action of
$\prod_{i \in e} (\delta_{h_i e_i} - 1)^L$ in the commutative algebra of translations; a further
coordinate $i \notin e$ applies $\Delta^L_{h_i e_i}$.
\end{definition}

\begin{lemma}[Properties of mixed differences]
\label{lem:mixed-differences}
\lean{Rd.aestronglyMeasurable_mixedDiff, Rd.integrable_mixedDiff, Rd.mixedDiff_congr_ae,
  Rd.eLpNorm_mixedDiff_le, Rd.mixedDiff_conv_ae, Rd.fwdDiff_single_tensorKernel,
  Rd.higherDiff_single_tensorKernel, Rd.mixedDiff_tensorKernel}
\leanok
\uses{def:mixed-differences, lem:rd-differences, def:tensor-kernel}
Mixed differences keep measurability and integrability, depend only on the class of $f$, satisfy
$\|\Delta^{L,e}_h f\|_p \leq 2^{L|e|} \|f\|_p$ and
$\Delta^{L,e}_h (K * f) = (\Delta^{L,e}_h K) * f$ almost everywhere, and act on tensor kernels
factorwise: $\Delta^{L,e}_h (\otimes K) = \otimes K'$ with $K'_i = \Delta^L_{h_i} K_i$ for
$i \in e$ and $K'_i = K_i$ otherwise.
\end{lemma}

\begin{proof}
\leanok
Induction over $e$ with the one-dimensional properties (\Cref{lem:rd-differences}).
\end{proof}

\begin{definition}[Mixed moduli and the mixed Besov quasi-norm]
\label{def:mixed-moduli}
\lean{Rd.mixedModulus, Rd.eLpNorm_mixedDiff_le_mixedModulus, Rd.mixedModulus_mono,
  Rd.mixedModulus_empty, Rd.mixedModulus_add_le, Rd.mixedModulus_le, Rd.mixedModulus_congr_ae,
  Rd.mixedSupp, Rd.mixedSupp_zero, Rd.lpMixedModulus, Rd.lpMixedModulus_zero,
  Rd.lpMixedModulus_add_le, Rd.lpMixedModulus_le, Rd.lpMixedModulus_zero_left, Rd.mixedModNorm,
  Rd.enorm_le_mixedModNorm}
\leanok
\uses{lem:mixed-differences, def:mixed-lq}
The mixed modulus of smoothness is
$\Omega^e_L(f, t)_p = \sup \{\|\Delta^{L,e}_h f\|_p : |h_i| \leq t_i \text{ for } i \in e\}$. For
$k \in \N^d$ let $e(k) = \{i : k_i \neq 0\}$. The Besov quasi-norm of dominating mixed smoothness
by differences is $\|(2^{s|k|_1} \Omega^{e(k)}_L(f, 2^{-k})_p)_{k \in \N^d}\|_{\ell_q}$; its term
$k = 0$ is $\|f\|_p$. The mixed modulus increases in $t$, is subadditive, is at most
$2^{L|e|} \|f\|_p$, and depends only on the class of $f$.
\end{definition}

\section{Averaged mixed differences}

\begin{lemma}[Bounded polynomials in translations]
\label{lem:acts-bounded}
\lean{Rd.ActsBounded, Rd.ActsBounded.one, Rd.ActsBounded.single, Rd.ActsBounded.mono,
  Rd.ActsBounded.add, Rd.ActsBounded.mul, Rd.ActsBounded.smul, Rd.ActsBounded.neg,
  Rd.ActsBounded.sub, Rd.ActsBounded.pow, Rd.ActsBounded.prod, Rd.actsBounded_momentPoly,
  Rd.actsBounded_addPoly}
\leanok
\uses{def:mixed-differences}
An element $a$ of the algebra of translations acts boundedly with constant $C$ if $a G$ stays
measurable and $\|a G\|_p \leq C \|G\|_p$. Translations $c \, \delta_v$ act with $|c|$, and sums,
products and powers add and multiply the constants. In particular $1 - (-1)^L (Y - 1)^L$ and
$(-1)^L ((Y + 1)^L - 1)$ act with $2^L + 1$ for every translation $Y$.
\end{lemma}

\begin{proof}
\leanok
Translations preserve $L_p$ norms; the triangle inequality and composition.
\end{proof}

\begin{lemma}[The moment kernel polynomial]
\label{lem:moment-polynomial}
\lean{Rd.sum_jacksonCoeff_smul_pow}
\leanok
\uses{def:moment-kernel}
In every commutative $\C$-algebra,
$\sum_{l=1}^L (-1)^{l+1} \binom{L}{l} Y^l = 1 - (-1)^L (Y - 1)^L$.
\end{lemma}

\begin{proof}
\leanok
The binomial theorem.
\end{proof}

\begin{theorem}[Tensor products of dilated bumps average translates]
\label{thm:tensor-averaged-differences}
\lean{Rd.diffAct_prod_sum_pow_apply, Rd.conv_tensorKernel_sum_dilate_ae}
\leanok
\uses{lem:diagonal-substitution, lem:tensor-kernel-fubini, lem:acts-bounded}
Let $\kappa_i = \sum_{t \in T} q_i(t) \, g_{c_i/\nu(t)}$ with integers $\nu(t) \geq 1$, scales
$c_i > 0$ and an integrable $g$ on $\R$, and let $Y_i = \delta_{-z_i e_i / c_i}$. For $f \in L_p$
and almost every $x$,
\[
  ((\otimes\kappa) * f)(x) = \int (\otimes g)(z)
    \Bigl(\prod_i \textstyle\sum_t q_i(t) Y_i^{\nu(t)}\Bigr) f(x) \, dz .
\]
\end{theorem}

\begin{proof}
\leanok
The product expands into the translates $f(x - (\nu(n_i) z_i / c_i)_i)$ with the coefficients
$\prod_i q_i(n_i)$, and each translate integrates to a convolution with an anisotropic dilation
of $\otimes g$ (\Cref{lem:diagonal-substitution}).
\end{proof}

\section{Mixed local means}

\begin{definition}[Mixed local means]
\label{def:mixed-local-means}
\lean{Rd.LocalMeansBump.mixedMeansKernel, Rd.LocalMeansBump.mixedDualKernel,
  Rd.LocalMeansBump.integrable_mixedMeansKernel, Rd.LocalMeansBump.integrable_mixedDualKernel,
  Rd.LocalMeansBump.mixedLocalMean, Rd.LocalMeansBump.mixedLocalMeansNorm}
\leanok
\uses{def:local-means, def:tensor-kernel}
For a bump $g$ on $\R$ with the kernels $k_j$ and $\eta_j$ of \Cref{def:local-means}, the mixed
kernels are $k_k = \otimes_i k_{k_i}$ and $\eta_k = \otimes_i \eta_{k_i}$, $k \in \N^d$, the mixed
local means of $f$ are $k_k * f$, and the mixed local means quasi-norm is
$\|(2^{s|k|_1} \|k_k * f\|_p)_k\|_{\ell_q}$.
\end{definition}

\begin{lemma}[Local means kernels as sums of dilated bumps]
\label{lem:kernel-dilation-sum}
\lean{Rd.LocalMeansBump.lmExp, Rd.LocalMeansBump.lmCoeff, Rd.LocalMeansBump.one_le_lmExp,
  Rd.LocalMeansBump.meansKernel_eq_sum_dilate, Rd.LocalMeansBump.sum_lmCoeff_smul_pow}
\leanok
\uses{def:moment-kernel, lem:moment-polynomial}
With $a_l = (-1)^l \binom{L}{l+1}$, $k_0 = \sum_{l < L} a_l g_{1/(l+1)}$ and
$k_j = \sum_{l < L} a_l (g_{2^j/(l+1)} - g_{2^j/(2(l+1))})$ for $j \geq 1$. The corresponding
polynomials are $P(Y) = 1 - (-1)^L (Y - 1)^L$ and
$P(Y) - P(Y^2) = (Y - 1)^L (-1)^L ((Y + 1)^L - 1)$.
\end{lemma}

\begin{proof}
\leanok
$m_c = \sum_l a_l g_{c/(l+1)}$ and $g_{2^j/(2(l+1))} = g_{2^{j-1}/(l+1)}$; the factorisation from
$Y^2 - 1 = (Y - 1)(Y + 1)$ and \Cref{lem:moment-polynomial}.
\end{proof}

\begin{theorem}[Mixed local means are bounded by the mixed modulus]
\label{thm:mixed-local-means-upper}
\lean{Rd.LocalMeansBump.lmRemainder, Rd.LocalMeansBump.prod_sum_lmCoeff_eq,
  Rd.LocalMeansBump.actsBounded_lmRemainder, Rd.LocalMeansBump.mixedLocalMean_ae_eq,
  Rd.LocalMeansBump.enorm_mixedLocalMean_le, Rd.LocalMeansBump.mixedLocalMeansNorm_le}
\leanok
\uses{def:mixed-local-means, lem:kernel-dilation-sum, thm:tensor-averaged-differences,
  lem:acts-bounded, lem:weighted-jensen, def:mixed-moduli}
If the bump vanishes outside $[-1, 1]$, then
$\|k_k * f\|_p \leq (\|g\|_1 (2^L + 1))^d \, \Omega^{e(k)}_L(f, 2^{-k})_p$. Hence the mixed local
means quasi-norm is at most $(\|g\|_1 (2^L + 1))^d$ times the quasi-norm by differences.
\end{theorem}

\begin{proof}
\leanok
By \Cref{lem:kernel-dilation-sum} and \Cref{thm:tensor-averaged-differences},
\[
  (k_k * f)(x) = \int (\otimes g)(z) \, B(z) \, \Delta^{L, e(k)}_{h(z)} f(x) \, dz
\]
with $h(z)_i = -z_i 2^{-k_i}$, so $|h(z)_i| \leq 2^{-k_i}$ on the support of $\otimes g$, and with
polynomials $B(z)$ in the translations that act with constant $(2^L + 1)^d$
(\Cref{lem:acts-bounded}). Minkowski's inequality (\Cref{lem:weighted-jensen}) concludes.
\end{proof}

\begin{lemma}[The mixed reproducing formula]
\label{lem:mixed-reproducing}
\lean{Rd.conv_dilate_dilate, Rd.convLp_congr, Rd.convLp_finsetSum_kernel, Rd.conv_add_conv_sub,
  Rd.LocalMeansBump.sum_conv_dualKernel_meansKernel,
  Rd.LocalMeansBump.sum_conv_mixedDualKernel_mixedMeansKernel,
  Rd.LocalMeansBump.integral_conv_momentKernel,
  Rd.LocalMeansBump.tendsto_sum_mixedDualKernel_mixedLocalMean}
\leanok
\uses{lem:reproducing-formula, lem:reproducing-pointwise, lem:tensor-kernel-fubini,
  thm:approx-identity}
In one variable $\sum_{j \leq J} \eta_j * k_j = m_{2^J} * m_{2^J}$ at every point. Over the boxes
$\{k : k_i \leq J\}$, $\sum_k \eta_k * k_k = (\otimes (m * m))_{2^J}$, and
$\sum_k \eta_k * k_k * f \to f$ in $L_p$.
\end{lemma}

\begin{proof}
\leanok
$(a + b) * (a - b) = a * a - b * b$ for the commuting continuous kernels with compact support;
tensor products and dilations commute with convolution (\Cref{lem:tensor-kernel-fubini}), and
$\otimes (m * m)$ has integral $1$, so \Cref{thm:approx-identity} applies.
\end{proof}

\begin{lemma}[Mixed moduli of the reproducing terms]
\label{lem:mixed-dual-moduli}
\lean{Rd.LocalMeansBump.lpMixedModulus_convLp_mixedDualKernel_le}
\leanok
\uses{lem:dual-kernel-moduli, lem:mixed-differences, lem:tensor-kernel-fubini}
$\Omega^{e(m)}_L(\eta_k * h, 2^{-m})_p \leq C_\eta^d \prod_i c(m_i, k_i) \, \|h\|_p$ with the
coefficients $c$ of \Cref{lem:mixed-almost-diagonal}.
\end{lemma}

\begin{proof}
\leanok
$\Delta^{L, e(m)}_v (\eta_k * h) = (\Delta^{L, e(m)}_v \eta_k) * h$, a convolution with a tensor
product whose factors have $L_1$ norms at most $C_\eta c(m_i, k_i)$
(\Cref{lem:dual-kernel-moduli}).
\end{proof}

\begin{theorem}[Mixed local means and mixed differences]
\label{thm:mixed-local-means}
\lean{ENNReal.le_tsum_of_tendsto_finsetSum, Rd.LocalMeansBump.mixedModNorm_le_of_tendsto,
  Rd.LocalMeansBump.mixedModNorm_le_mixedLocalMeansNorm,
  Rd.LocalMeansBump.mixedBesov_iff_mixedLocalMeans}
\leanok
\uses{thm:mixed-local-means-upper, lem:mixed-reproducing, lem:mixed-dual-moduli,
  lem:mixed-almost-diagonal}
Let $1 \leq p < \infty$, $0 < s < L$, $1 \leq q \leq \infty$, and let the bump vanish outside
$[-1, 1]$. Then $\|(2^{s|k|_1} \Omega^{e(k)}_L(f, 2^{-k})_p)_k\|_{\ell_q} < \infty$ if and only if
$\|(2^{s|k|_1} \|k_k * f\|_p)_k\|_{\ell_q} < \infty$, with equivalent quasi-norms.
\end{theorem}

\begin{proof}
\leanok
One direction is \Cref{thm:mixed-local-means-upper}. Conversely
$f = \sum_k \eta_k * (k_k * f)$ (\Cref{lem:mixed-reproducing}), and the mixed modulus is
subadditive and bounded on $L_p$, so
$\Omega^{e(m)}_L(f, 2^{-m})_p \leq C_\eta^d \sum_k \prod_i c(m_i, k_i) \|k_k * f\|_p$
(\Cref{lem:mixed-dual-moduli}); \Cref{lem:mixed-almost-diagonal} sums over $k$.
\end{proof}

\section{Tensor Littlewood--Paley blocks}
\label{sec:mixed-lp-blocks}

\begin{definition}[Tensor Littlewood--Paley blocks]
\label{def:mixed-lp}
\lean{Rd.mixedLpBlockKernel, Rd.mixedDilatedBandKernel, Rd.integrable_mixedLpBlockKernel,
  Rd.integrable_mixedDilatedBandKernel, Rd.conv_mixedDilatedBandKernel_mixedLpBlockKernel,
  Rd.sum_mixedLpBlockKernel, Rd.mixedBlock, Rd.mixedLPNorm, Rd.tendsto_sum_mixedBlock}
\leanok
\uses{def:lp-kernels, thm:kernel-identities, lem:tensor-kernel-fubini, thm:approx-identity}
The tensor blocks are $\delta_k = \otimes_i \psi_{k_i}$, so that
$\delta_k * f = \mathcal{F}^{-1}(\varphi_k \mathcal{F} f)$ for the tensor decomposition of unity
$\varphi_k = \otimes_i \mathcal{F}\psi_{k_i}$. Vyb\'{\i}ral's quasi-norm is
$\|(2^{s|k|_1} \|\delta_k * f\|_p)_k\|_{\ell_q}$. The blocks satisfy
$(\otimes_i \Theta_{k_i}) * \delta_k = \delta_k$, sum over the boxes $\{k : k_i \leq J\}$ to
$(\otimes \Phi)_{2^J}$, and $\sum_k \delta_k * f = f$ in $L_p$.
\end{definition}

\begin{lemma}[One-dimensional estimates in $L_1$]
\label{lem:mixed-lp-kernels}
\lean{Rd.eLpNorm_conv_eq_enorm_convLp, Rd.one_le_dyadicWeight_mul_almostDiagCoeff,
  Rd.LocalMeansBump.exists_eLpNorm_conv_lpBlockKernel_dualKernel_le,
  Rd.LocalMeansBump.exists_eLpNorm_conv_meansKernel_dilatedBandKernel_le}
\leanok
\uses{thm:band-jackson, thm:kernel-identities, lem:local-means-upper, thm:band-inverse,
  lem:dual-kernel-moduli}
There is $C < \infty$ with $\|\psi_m * \eta_k\|_1 \leq C c(m, k)$ and
$\|k_m * \Theta_k\|_1 \leq C c(m, k)$ for all $m, k \in \N$.
\end{lemma}

\begin{proof}
\leanok
For $k + 3 \leq m$, $\psi_m$ annihilates the band-limited Jackson approximant $J_{m-3} \eta_k$
(\Cref{thm:kernel-identities}), so
$\|\psi_m * \eta_k\|_1 \leq \|\psi_m\|_1 \|\eta_k - J_{m-3} \eta_k\|_1
\leq C \omega_L(\eta_k, 2^{-(m-3)})_1$ (\Cref{thm:band-jackson}), and
\Cref{lem:dual-kernel-moduli} bounds the modulus. For $k < m$,
$\|k_m * \Theta_k\|_1 \leq C \omega_L(\Theta_k, 2^{1-m})_1$ (\Cref{lem:local-means-upper}), and the
differences of $\Theta_k$ give the decay (\Cref{thm:band-inverse}). In the remaining cases Young's
inequality suffices.
\end{proof}

\begin{theorem}[Littlewood--Paley blocks, local means and differences]
\label{thm:mixed-lp}
\lean{Rd.LocalMeansBump.enorm_mixedBlock_convLp_mixedDualKernel_le,
  Rd.LocalMeansBump.enorm_mixedLocalMean_mixedBlock_le, Rd.LocalMeansBump.exists_mixedLPNorm_le,
  Rd.LocalMeansBump.exists_mixedLocalMeansNorm_le, Rd.mixedBesov_iff_mixedLP}
\leanok
\uses{def:mixed-lp, lem:mixed-lp-kernels, lem:mixed-reproducing, thm:mixed-local-means,
  lem:mixed-almost-diagonal, lem:bumps-every-radius, lem:bump-of-smooth-bump}
Let $1 \leq p < \infty$, $0 < s < L$ and $1 \leq q \leq \infty$. Then
$\|(2^{s|k|_1} \Omega^{e(k)}_L(f, 2^{-k})_p)_k\|_{\ell_q} < \infty$ if and only if
$\|(2^{s|k|_1} \|\delta_k * f\|_p)_k\|_{\ell_q} < \infty$.
\end{theorem}

\begin{proof}
\leanok
$\delta_m * f = \sum_k (\delta_m * \eta_k) * (k_k * f)$ and
$k_m * f = \sum_k (k_m * \otimes_i \Theta_{k_i}) * (\delta_k * f)$, and the tensor kernels have
$L_1$ norms at most $C^d \prod_i c(m_i, k_i)$ (\Cref{lem:mixed-lp-kernels}), so
\Cref{lem:mixed-almost-diagonal} compares Vyb\'{\i}ral's quasi-norm with the mixed local means.
\Cref{thm:mixed-local-means} for the centred bump of radius $1$ (\Cref{lem:bump-of-smooth-bump})
concludes.
\end{proof}

\section{Hyperbolic crosses}

\begin{definition}[Levels and hyperbolic crosses]
\label{def:hyperbolic-cross}
\lean{Rd.gridLevel, Rd.le_gridLevel, Rd.gridLevelSet, Rd.mem_gridLevelSet,
  Rd.mixedWeight_eq_dyadicWeight_gridLevel, Rd.hyperbolicCross, Rd.mem_hyperbolicCross,
  Rd.hyperbolicSum}
\leanok
\uses{def:mixed-lp}
The level of $k \in \N^d$ is $|k|_1 = \sum_i k_i$, the hyperbolic cross of level $N$ is
$\{k : |k|_1 \leq N\}$, and $S_N f = \sum_{|k|_1 \leq N} \delta_k * f$ is band-limited to the union
of the dyadic rectangles of these blocks.
\end{definition}

\begin{lemma}[Counting and Hölder's inequality on levels]
\label{lem:level-count}
\lean{Rd.card_gridLevelSet_le, Rd.tsum_le_tsum_gridLevelSet, Rd.sum_le_card_rpow_mul_lqNorm}
\leanok
\uses{def:hyperbolic-cross, def:mixed-lq}
$\#\{k \in \N^d : |k|_1 = n\} \leq (n + 1)^{d-1}$; sums over $\N^d$ are at most the sums over the
levels of the sums on each level; and $\sum_{k \in S} x_k \leq (\# S)^{1 - 1/q} \|x\|_{\ell_q}$ for
finite $S$.
\end{lemma}

\begin{proof}
\leanok
An index of level $n$ is determined by its first $d - 1$ coordinates; Hölder's inequality on $S$.
\end{proof}

\begin{lemma}[The rate series]
\label{lem:rate-series}
\lean{Rd.tsum_dyadicWeight_mul_rpow_ne_top, Rd.tsum_tail_le}
\leanok
For $s > 0$ and $a \geq 0$, $C = \sum_l 2^{-sl} (l + 1)^a < \infty$ and
$\sum_{n > N} 2^{-sn} (n + 1)^a \leq C \, 2^{-sN} (N + 1)^a$.
\end{lemma}

\begin{proof}
\leanok
$(l + 1)^a \leq (l + 1)^{\lceil a \rceil}$ and $\sum_l l^K r^l < \infty$ for $r < 1$;
$n + 1 \leq (n - N + 1)(N + 1)$ for $n \geq N$.
\end{proof}

\begin{theorem}[Hyperbolic cross approximation]
\label{thm:hyperbolic-cross}
\lean{Rd.enorm_sub_hyperbolicSum_le, Rd.enorm_sub_hyperbolicSum_le_mixedLPNorm,
  Rd.hyperbolicCrossConst, Rd.hyperbolicCrossConst_ne_top, Rd.enorm_sub_hyperbolicSum_le_of_pos}
\leanok
\uses{def:hyperbolic-cross, lem:level-count, lem:rate-series, def:mixed-lp}
For every $1 \leq p < \infty$, every $f \in L_p(\R^d)$, every $s > 0$, every
$1 \leq q \leq \infty$ and every $N \in \N$, the constant
$C = \sum_l 2^{-sl} (l + 1)^{(d-1)(1 - 1/q)}$ is finite and
\[
  \|f - S_N f\|_p \leq C \, 2^{-sN} (N + 1)^{(d-1)(1 - 1/q)} \,
    \|(2^{s|k|_1} \|\delta_k * f\|_p)_k\|_{\ell_q} .
\]
\end{theorem}

The inequality itself holds for every real $s$, with $C \in [0, \infty]$; for $s \leq 0$ the
series $C$ is $\infty$.

\begin{proof}
\leanok
$\|f - S_N f\|_p \leq \sum_{|k|_1 > N} \|\delta_k * f\|_p$, since $\sum_k \delta_k * f = f$. On the
level $n$, Hölder's inequality costs $(n + 1)^{(d-1)(1 - 1/q)}$ (\Cref{lem:level-count}), and
\Cref{lem:rate-series} sums over $n > N$.
\end{proof}
